%----------------------------------------------------------------------------------------------------%
%----------------------------------------------------------------------------------------------------%
%%%%%%%%%%%                                              HEADER                                                %%%%%%%%%%%
%----------------------------------------------------------------------------------------------------%
%----------------------------------------------------------------------------------------------------%
%
% Universale header.tex – gilt für Arbeitsblätter, Leistungskontrollen, Handreichungen, Nachschlagewerke.
% Neue Pakete und TikZ-Bibliotheken werden HIER ergänzt (nicht im Dokument) und unten protokolliert.
% Quelle: https://github.com/mq90/latex-skills-schule  (CC BY 4.0, Martin Quandt)
%
% Änderungsprotokoll
% 2026-11-09  zusammengeführt aus allen Handreichungs-Headern: TikZ-Bibliotheken fit, shapes.multipart,
%             arrows.meta, shapes.callouts, circuits.logic.US, circuits.logic.IEC; Pakete array, booktabs,
%             longtable; xskak, chessboard, sudoku, tikzpeople, tipa (safe); MusiXTeX als auskommentierte
%             Option (definiert \cup als Achtelnote um – bricht $A \cup B$)
% 2026-11-11  Naturwissenschaften: circuitikz [european] (Schaltpläne), mhchem [version=4] (\ce),
%             chemfig (Strukturformeln), pgf-PeriodicTable (\pgfPT), genealogytree [all] (Stammbäume).
%             NICHT siunitx: kollidiert mit units (\unit[10]{cm}); siunitx nur in eigenständigen Dokumenten
% 2026-11-11  TikZ-Bibliothek decorations.pathmorphing (Federn, Wellen: decoration=coil/snake)
% 2026-11-23  twemojis: farbige Emojis mit pdflatex, z.B. \twemoji[height=1cm]{man standing} (Skill latex-sachskizzen)
%






%----------------------------------------------------------------------------------------------------%
%%%%%%%%%%%%%                            Gestaltung des Dokuments                                %%%%%%%%%%
%----------------------------------------------------------------------------------------------------%



%Definition der Dokumentenklasse
\documentclass[12pt]{article} %book, article, report, standalone



%Pakete für eine saubere Kodierung
\usepackage[utf8]{inputenc} %UTF-8 als Zeichenkodierung  ... ABER nicht immer notwendig, je nach Compiler
\usepackage[T1]{fontenc} %Verwendet westeuropäische Schriftkodierung
\usepackage[ngerman]{babel} %Ermöglicht deutsche Übersetzung und Silbentrennung
\usepackage{lmodern} %Verwendet moderne Schriftarten



%Pakete für die Gestaltung des Dokuments
\usepackage{geometry} %Ermöglicht die Ausrichtung des Blattes mit \geometry{...}
\geometry{a4paper, top=15mm, left=15mm, right=20mm, bottom=15mm, headsep=10mm, footskip=10mm} %gewählte Ausrichtung des Blattes
\pagestyle{empty} %Art der Seitennummerierung, z.B.: empty , plain
\usepackage{fancyhdr} %Ermöglicht die Gestaltung einer Kopfzeile für LK's oder KA's 
\usepackage{lscape} %Ermöglicht eventuelles Querformat des Blattes mit \begin{landscape} ... \end{landscape}
\usepackage{pdfpages} %Ermöglicht einbinden PDF-Dokumente bzw. einzelene Seiten mit \includepdf[pages={1,3,5}]{Dateiname} bzw. pages=-
\usepackage[german]{layout} %Ermöglicht es mit den Befehl \layout sich die Default-Layout beschreiben zu lassen (Bild vom Layout)
%----------------------------------------------------------------------------------------------------%








%----------------------------------------------------------------------------------------------------%
%%%%%%%%%%%%                                        Packages                                                  %%%%%%%%%%%
%----------------------------------------------------------------------------------------------------%



%Pakete für die Strukturierung des Dokuments
\usepackage[shortlabels]{enumitem} %Ermöglicht die Nummerierung selbst festzulegen mit \begin{enumerate}[Art der Nummerierung] ...   \end{enumerate}
\parindent=0px %Setzt die Absatzeinrückung auf 0 Pixel und ernfernt damit unnötige Leerzeichen vor dem ersten Wort eines neuen Abschnitts
\usepackage{setspace} %Ermöglicht verschiedene Funktion zur Abstandsregelung, z.B.: \hspace{2.5cm} ; \onehalfspacing
\onehalfspacing %Setzt den Zeilenabstand auf 1,5
\usepackage[style=german]{csquotes}%Ermöglicht den deutschen Zitierstil darstellen zu lassen mt \enquote{Zitat}
\usepackage[breaklinks=true,%Zeilenumbruch der Verlinkung ist bündig
			colorlinks=true,%Links können farbig sein
			linkcolor=blue,%Linkfarbe wird auf blau gesetzt
			urlcolor=darkgray,%Farb der URL wird auf dunkelgrau gesetzt
			plainpages=false,%
			hyperfootnotes=false,%
			pdfpagelabels=false%
		]{hyperref} %Ermöglicht die Verlinkung innerhalb des Dokuments 
\usepackage{multicol} %Ermöglicht gleichgroße Aufteilung des Dokuments in n Spalten mit \begin{multicols}{n} ... \end{multicols}
\usepackage{xcolor,colortbl} %Ermöglicht Zeilen, Spalten und Zellen einer Tabelle zu färben
\usepackage{caption} %Ermöglicht die Bezeichnung und Strukturierung von Abbildungen oder Quellcode
\usepackage{fancybox} % Ermöglicht die Umrandung einzelner Abschnitte
\usepackage[most]{tcolorbox} %Ermöglicht die Nutzung von Colorboxen für Merksätze, Lösungen, Definitionen, usw mit \newtcolorbox
\tcbuselibrary{breakable} %Aufteilen der Colorboxen, welche über mehrere Seiten gehen
\usepackage{ifthen} % Ermöglicht \ifthenelse Befehle
\allowdisplaybreaks %Erlaubt Seitenumbrüche innerhalb von align-Umgebung



%Nützliche Pakete für Symbole und mathematischen Formelsatz
\usepackage{amsmath, amssymb, mathrsfs, amsfonts, amsthm, dsfont} %sämtliche Pakete für mathematische Zeichen, Symbole und Formelsatz
\usepackage{units} %Ermöglicht die Angabe von Einheiten mit \unit[Wert]{Einheit}
\usepackage{gensymb} %Ermöglicht Nutzung von \degree für das macro des Grad-Symbols °
\usepackage[official]{eurosym} %Ermöglicht die Nutzung von \euro{}
\usepackage{pifont} %Ermöglicht die Nutzung verschiedener Fonts, z.B.: \ding{42} , \ding{52} , \ding{56}
\usepackage{tikzsymbols} %Ermöglicht die Nutzung von u.a. Smileys und anderen Symbolen, z.B.: \Winkey , \Smiley[1.5]



%Pakete zur Darstellung von Grafiken
\usepackage{tikz} %\begin{tikzpicture} .... \end{tikzpicture}
\usepackage{pgf, pgfplots} %Ermöglicht 2D und 3D Plof-Funktion
\pgfplotsset{compat=1.17}
\usetikzlibrary{arrows,%Pfeile
			calc,%Rechnen in Koordinaten, z.B. mit ($(1,2)+5*(0,1)$) auf Punkt (1,7) kommen
			angles, plotmarks, quotes, babel, intersections, shapes.geometric, decorations.pathreplacing, decorations.pathmorphing, automata, positioning, shadows, patterns, patterns.meta,%
			fit,%Rahmen um mehrere Knoten (UML, Automaten)
			shapes.multipart,%geteilte Knoten (UML-Klassen)
			arrows.meta,%moderne Pfeilspitzen, z.B. -{Triangle[open]}
			shapes.callouts,%Sprechblasen
			circuits.logic.US, circuits.logic.IEC%Logikgatter (digitale Schaltungen)
			} %nützliche Bibliotheken für TikZ
\usepackage{animate} %animierte Darstellungen
\usepackage{graphicx} %\includegraphics[width=..px,height=..px]{Dateiname}
\usepackage{multirow} %Zellen über mehrere Zeilen mit \multirow{n}{*}{Text}
\usepackage{array} %erweiterte Spaltentypen, z.B. >{\centering\arraybackslash}p{2cm}
\usepackage{booktabs} %schöne Tabellenlinien \toprule, \midrule, \bottomrule
\usepackage{longtable} %Tabellen über mehrere Seiten



%gelegentlich hilfreiche Zusatzpakete
\usepackage{hologo} %\hologo{LaTeX}
\usepackage{suetterl} %Sütterlin als Schriftart mit \mathfrak{Text}
\usepackage{qrcode} %Einbindung eines QR-Codes 
\usepackage{url} %Einbindung eines URL-Links
\usepackage{ccicons} %Creative Commons Symbole
\usepackage{listings} %\begin{lstlisting} ... \end{lstlisting}
\usepackage{struktex} %Struktogrammdarstellung
\usepackage{ulem} %Unterstreichung mit \uline{} ggf. mit \hfill im Inneren für Unterstrich bis zum Ende der Zeile
\usepackage{xskak} %Schach: Partien und Züge mit \newchessgame, \mainline{1. e4 e5} (lädt skak)
\usepackage{chessboard} %Schachbretter mit \chessboard[setfen=...]
\usepackage{sudoku} %Sudoku-Gitter mit \begin{sudoku} |5|3| | ... |. \end{sudoku}
\usepackage{tikzpeople} %Personen-Figuren als TikZ-Knoten, z.B. \node[alice, minimum size=2cm]{};
\usepackage{twemojis} %farbige Emojis (Twemoji, CC BY 4.0) als Vektorgrafik: \twemoji[height=1cm]{man standing}, im Text \texttwemoji{light bulb}
\usepackage[safe]{tipa} %Lautschrift mit \textipa{...}; [safe], weil tipa sonst \; \: \! \* \s umdefiniert (bricht z.B. \(\;\)) – Langformen siehe latex-lautschrift
\usepackage[european]{circuitikz} %Schaltpläne mit deutschen Symbolen: \begin{circuitikz} \draw (0,0) to[lamp] (2,0); \end{circuitikz}
\usepackage[version=4]{mhchem} %Reaktionsgleichungen mit \ce{2 H2 + O2 -> 2 H2O}
\usepackage{chemfig} %Strukturformeln mit \chemfig{H-C(-[2]H)(-[6]H)-H}, freie Elektronenpaare mit \charge
\usepackage{pgf-PeriodicTable} %Periodensystem mit \pgfPT[...], deutsche Namen mit \pgfPTsetLanguage{de}
\usepackage[all]{genealogytree} %Stammbäume mit \begin{genealogypicture} ... \end{genealogypicture}
%\usepackage{musixtex} %Notensatz – NUR für Musik-Dokumente einkommentieren: definiert u.a. \cup als Note um (bricht $A \cup B$)
%\usepackage[list=true, font=large, labelfont=bf, labelformat=brace, position=top]{subcaption}



%abgeänderte Befehle
\renewcommand{\arraystretch}{1.25} %
%\renewcommand{\familydefault}{\sfdefault} %serifenlose Schriftart
\renewcommand{\lstlistingname}{Quellcode} %Umbenennung
\renewcommand{\lstlistlistingname}{Quellcodeverzeichnis} %Umbenennung
%----------------------------------------------------------------------------------------------------%
%----------------------------------------------------------------------------------------------------%
%----------------------------------------------------------------------------------------------------%
%%%% %%%%%%                                     NEWCOMMANDS                                         %%%%%%%%%%%%
%----------------------------------------------------------------------------------------------------%
%----------------------------------------------------------------------------------------------------%
%
% Universale newcommands.tex – alle eigenen Befehle, Farben, Boxen, Kopfzeilen, TikZ- und Listing-Stile.
% Neue Befehle werden HIER ergänzt (nicht im Dokument) und unten protokolliert.
% Quelle: https://github.com/mq90/latex-skills-schule  (CC BY 4.0, Martin Quandt)
%
% Änderungsprotokoll
% 2026-11-09  zusammengeführt aus allen Handreichungen: \TikZ, Box »ausgabe«, Stile für
%             Programmablaufpläne (input, operation, condition, output) und UML (umlclass, inherit,
%             compose, aggregate, associate, depend); Listing-Stil »latex« wieder wie in den
%             Handreichungen ohne Zeilennummern, nummerierte Variante als »latexnummeriert«;
%             Umlaute und · € ° → ≤ ≥ ≠ – — „ “ … × ² ³ π in allen Listings erlaubt (\lstset literate); \lehrerblase, \schuelerblase; PAP-Stile mit align=center (Zeilenumbruch mit \\);
%             \ipa{...} für Lautschrift; \koordinatensystem[Schritt]{xmin}{xmax}{ymin}{ymax};
%             TikZ-Stile punkt, punktoffen, baumdiagramm, wkeit
% 2026-11-09  TikZ-Stil baumdiagrammvertikal (Baumdiagramm von oben nach unten)
% 2026-11-09  Formularfelder für ausfüllbare PDFs: \feld, \zeichen, \kreuz, \antwort, einheitliches
%             Aussehen über \DefaultOptionsofText/CheckBox/Radio/PopdownBox (nur wirksam in einer Form-Umgebung)
% 2026-11-11  Akzentfarben akzentblau, akzentrot, akzentgruen, akzentorange, akzentlila (Karten,
%             Naturwissenschaften, Schulorganisation); Box »hinweis« mit frei wählbarem Titel
% 2026-11-11  Umgebung »passend«: verkleinert den Inhalt auf \linewidth, aber nur wenn er breiter ist
% 2026-11-11  \oz{+II}{Cu}: Oxidationszahl über ein Elementsymbol (Chemie, auch innerhalb \ce{... $\oz{0}{Zn}$ ...})
% 2026-11-22  Farbige Schachbretter: \schachbrettfarben{braun|blau|gruen|grau|klassisch} (Standard: braun),
%             Farben brett...hell/brett...dunkel und brettmarke (braucht LSBC5 aus der header.tex)
%








%----------------------------------------------------------------------------------------------------%
%%%%%%%%%%%                                   Mathematische Kürzel                                       %%%%%%%%%%
%----------------------------------------------------------------------------------------------------%


%vereinfachte Schreibweise für den Displaystyle
\newcommand{\dis}[1]{\displaystyle{#1}}


%Zahlenbereiche
\newcommand{\N}{\mathbb{N}} %natürliche Zahlen
\newcommand{\Z}{\mathbb{Z}} %ganze Zahlen
\newcommand{\Q}{\mathbb{Q}} %rationale Zahlen
\newcommand{\R}{\mathbb{R}} %reelle Zahlen
\newcommand{\C}{\mathbb{C}} %komplexe Zahlen
\newcommand{\cP}{\mathcal{P}} %Potenzmenge


%Abkürzungen für Analysis
\newcommand{\eps}{\varepsilon} %Kürzel für Epsilon
\newcommand{\limn}{\underset{n\longrightarrow \infty}{\lim}} %Grenzwert für n gegen Unendlich
\newcommand{\limh}{\underset{h\longrightarrow 0}{\lim}} %Grenzwert für h gegen Null
\newcommand{\limxto}[1]{\underset{x\longrightarrow #1}{\lim}} %beliebiger funktionaler Grenzwert x gegen ....
\newcommand{\dx}{\text{\;d} x} %dx Operator für Integrale
\renewcommand{\d}[1]{\text{\;d}#1} %Operator für beliebige Integrationsvariable
\newcommand{\norm}[1]{\left\lVert#1\right\rVert} %\norm{...} für die Normdarstellung ||...||


%Abkürzungen für Analytische Geometrie
\newcommand{\vektor}[3]{\begin{tiny}\begin{pmatrix} \text{\begin{scriptsize}$#1$\end{scriptsize}}\\ \text{\begin{scriptsize}$#2$\end{scriptsize}}\\ \text{\begin{scriptsize}$#3$\end{scriptsize}}\end{pmatrix}\end{tiny}} %dreidimensionaler Vektor
\newcommand{\punktvektor}[2]{\begin{tiny}\begin{pmatrix} \text{\begin{scriptsize}$#1$\end{scriptsize}}\\ \text{\begin{scriptsize}$#2$\end{scriptsize}}\end{pmatrix}\end{tiny}} %zweidimensionaler Vektor


%Abkürzungen für Stochastik
\newcommand{\E}{\mathbb{E}} %Erwartungswert
\newcommand{\V}{\mathbb{V}} %Varianz


%Umgebung für gegeben und gesucht
\newcommand{\gegges}[2]{%
	\begin{multicols}{2}
		\begin{enumerate}
			\item[geg:] #1
			\item[ges:] #2
		\end{enumerate}
	\end{multicols}
}
%----------------------------------------------------------------------------------------------------%









%----------------------------------------------------------------------------------------------------%
%%%%%%%%%%%%                             Verschiedene  hilfreiche Kürzel                                %%%%%%%%%
%----------------------------------------------------------------------------------------------------%


%römische Darstellung einer Zahl \uproman{2021} ergibt MMXXI
\newcommand{\uproman}[1]{\uppercase\expandafter{\romannumeral#1}}


%Automatisch erstellte Lücke wie die Länge des Textes mit \luecke{<Text>}
\newcommand{\luecke}[1]{\hspace{0,1cm}\underline{\hspace{0,25cm}\hphantom{#1}\hspace{0,25cm}}\hspace{0,1cm}}


%fett-kursive Schrift für neu eingeführte Begriffe
\newcommand{\fk}[1]{\textbf{\textit{#1}}}


%Schwierigkeitsstufen für Aufgaben mit \leicht , \mittel , \schwer
\newcommand{\leicht}{\smash{\raisebox{-1pt}{{\color{black}\rule{4pt}{8pt}}\,{\color{lightgray}\rule{4pt}{8pt}}\,{\color{lightgray}\rule{4pt}{8pt}}}}\,}
\newcommand{\mittel}{\smash{\raisebox{-1pt}{{\color{black}\rule{4pt}{8pt}}\,{\color{black}\rule{4pt}{8pt}}\,{\color{lightgray}\rule{4pt}{8pt}}}}\,}
\newcommand{\schwer}{\smash{\raisebox{-1pt}{{\color{black}\rule{4pt}{8pt}}\,{\color{black}\rule{4pt}{8pt}}\,{\color{black}\rule{4pt}{8pt}}}}\,}
%----------------------------------------------------------------------------------------------------%










%----------------------------------------------------------------------------------------------------%
%%%%%%%                                            Verschiedene Farben                                                      %%%%%%%
%----------------------------------------------------------------------------------------------------%


%Formeln in herkömmlichen Farben schreiben
\newcommand{\frot}[1]{\text{\textcolor{red}{$#1$}}} % rote Formel für Lösungen
\newcommand{\fgruen}[1]{\text{\textcolor{green}{$#1$}}} % grüne Formel für Lösungen
\newcommand{\fblau}[1]{\text{\textcolor{blue}{$#1$}}} % blaue Formel für Lösungen
\newcommand{\fmagenta}[1]{\text{\textcolor{magenta}{$#1$}}} % magenta-farbige Formel für Lösungen
\newcommand{\forange}[1]{\text{\textcolor{orange}{$#1$}}} % orangene Formel für Lösungen
\newcommand{\fcyan}[1]{\text{\textcolor{cyan}{$#1$}}} % cyan-farbige Formel für Lösungen
\newcommand{\folive}[1]{\text{\textcolor{olive}{$#1$}}} % olive-farbige Formel für Lösungen


%Custom made colors
\definecolor{plum}{rgb}{0.36078, 0.20784, 0.4}
\definecolor{chameleon}{rgb}{0.30588, 0.60392, 0.023529}
\definecolor{cornflower}{rgb}{0.12549, 0.29020, 0.52941}
\definecolor{scarlet}{rgb}{0.8, 0, 0}
\definecolor{brick}{rgb}{0.64314, 0, 0}
\definecolor{sunrise}{rgb}{0.80784, 0.36078, 0}
\definecolor{lightblue}{rgb}{0.15,0.35,0.75}
\definecolor{carolina}{rgb}{153, 186, 221}

%Akzentfarben für Grafiken, Karten und Diagramme (kräftig, gut unterscheidbar; aufgehellt z.B. akzentblau!18)
\definecolor{akzentblau}{HTML}{2563EB}
\definecolor{akzentrot}{HTML}{DC2626}
\definecolor{akzentgruen}{HTML}{16A34A}
\definecolor{akzentorange}{HTML}{EA580C}
\definecolor{akzentlila}{HTML}{7C3AED}
%----------------------------------------------------------------------------------------------------%









%----------------------------------------------------------------------------------------------------%
%%%%%%%%                            Kopfzeilen für Leistungsüberprüfungen                              %%%%%%%%%
%----------------------------------------------------------------------------------------------------%


%Definition Kopfzeile für Leistungskontrolle mit \lk{Klasse}{Thema}
\newcommand{\lk}[2]{%
	\pagestyle{fancy}
	\fancyhead[L]{%
	\begin{tabular}[b]{l}
		\\Name:\\Punkte:\\Note:
	\end{tabular} 
	}
	\fancyhead[R]{%
	\begin{tabular}[b]{r}
		\\Leistungskontrolle\\Klasse #1\\ \today
	\end{tabular} 
	}
	\fancyhead[C]{%
	\textbf{\large{#2}}\\
	}
	\fancyfoot[C]{}
	\renewcommand{\headrulewidth}{1pt}
	\fancyheadoffset{1cm}
}


%Definition Kopfzeile für Klassenarbeit mit \ka{Klasse}{Thema}
\newcommand{\ka}[2]{%
	\pagestyle{fancy}
	\fancyhead[L]{%
	\begin{tabular}[b]{l}
		\\Name:\\Punkte:\\Note:
	\end{tabular} 
	}
	\fancyhead[R]{%
	\begin{tabular}[b]{r}
		\\Klassenarbeit\\Klasse #1\\ \today
	\end{tabular} 
	}
	\fancyhead[C]{%
	\textbf{\large{#2}}\\
	}
	\fancyfoot[C]{}
	\renewcommand{\headrulewidth}{1pt}
	\fancyheadoffset{1cm}
}


%Definition Kopfzeile für Klausur mit \kl{Klasse}{Thema}
\newcommand{\kl}[2]{%
	\pagestyle{fancy}
	\fancyhead[L]{%
	\begin{tabular}[b]{l}
		\\Name:\\Punkte:\\Note:
	\end{tabular} 
	}
	\fancyhead[R]{%
	\begin{tabular}[b]{r}
		\\Klausur\\Klasse #1\\ \today
	\end{tabular} 
	}
	\fancyhead[C]{%
	\textbf{\large{#2}}\\
	}
	\fancyfoot[C]{}
	\renewcommand{\headrulewidth}{1pt}
	\fancyheadoffset{1cm}
}


%Entfernt den Header für Leistungskontrollen, Klassenarbeiten oder Klausuren
\newcommand{\clearfancyheader}{%
	\fancyhf{}
	\renewcommand{\headrulewidth}{0pt}
	\renewcommand{\headheight}{0pt}
}


%Definition Kopfzeile für Tägliche Übungen mit \tu{Klasse}{Thema}
\newcommand{\tu}[2]{%
	\pagestyle{fancy}
	\fancyhead[L]{%
	\begin{tabular}[b]{l}
		\\Name:\\Punkte:\\Note:
	\end{tabular} 
	}
	\fancyhead[R]{%
	\begin{tabular}[b]{r}
		\\Tägliche Übung\\Klasse #1\\ \today
	\end{tabular} 
	}
	\fancyhead[C]{%
	\textbf{\large{#2}}\\
	}
	\fancyfoot[C]{}
	\renewcommand{\headrulewidth}{1pt}
	\fancyheadoffset{1cm}
}
%----------------------------------------------------------------------------------------------------%








%----------------------------------------------------------------------------------------------------%
%%%%%%%%%                                     Colorboxes für Hervorhebungen                              %%%%%%%%%
%----------------------------------------------------------------------------------------------------%


 %Colorboxes für Lösungen mit \begin{solution}<Text>\end{solution}
\newtcolorbox{solution}{%
        	lower separated=false,%
	colback=gray!5!white,%
	colframe=gray!75!black,%
        	fonttitle=\bfseries,%
	title=Lösung:,%
      	enhanced,%
	breakable,%
        	attach boxed title to top left={xshift=0.5cm,yshift=-2mm}%
}


%Colorboxes für das chronologische Vorgehen 
\newtcolorbox{todo}{%
        	lower separated=false,%
	colback=yellow!5!white,%
	colframe=orange!75!black,%
        	enhanced,%
	breakable%
}


%Colorboxes für Merksätze
\newtcolorbox{merke}{%
        	lower separated=false,%
	colback=red!5!white,%
	colframe=red!75!black,%
        	fonttitle=\bfseries,%
	title=Merke:,%
      	enhanced,%
	breakable,%
        	attach boxed title to top left={xshift=0.5cm,yshift=-2mm}%
}


%Hinweis-Box mit frei wählbarem Titel: \begin{hinweis}[Lizenz] ... \end{hinweis} (Standardtitel: Hinweis)
\newtcolorbox{hinweis}[1][Hinweis]{%
	enhanced,%
	breakable,%
	sharp corners,%
	colback=akzentblau!5,%
	colframe=akzentblau,%
	boxrule=0pt,%
	leftrule=3pt,%
	fonttitle=\bfseries,%
	coltitle=akzentblau!85!black,%
	attach title to upper={\ },%
	title={#1:}%
}


%Colorbox2 für Merksätze 
\newtcolorbox{merke2}{%
        	lower separated=false,%
        	colback=red!40,%
       	title=Merke,%
       	colframe=black,%
        	fonttitle=\bfseries,%
        	colbacktitle=red,%
        	coltitle=black,%
        	enhanced,%
        	breakable,%
        	attach boxed title to top left={xshift=0.5cm,yshift=-2mm},%
        	title=Merke:%
}
%----------------------------------------------------------------------------------------------------%










%----------------------------------------------------------------------------------------------------%
%%%%%%%%%%%%%                                         MC-Fragen                                    %%%%%%%%%%%%%
%----------------------------------------------------------------------------------------------------%

%Definition des Kästchens zum ankreuzen
\newcommand{\casket}{\LARGE{$\square$}}


%Umgebung für horizontale MC-Fragen
\newenvironment{mcquestionh}[1]{%
	#1%
	\begin{center}\begin{tabular}{>{\centering\arraybackslash}p{2.5cm}  >{\centering\arraybackslash}p{2.5cm} >{\centering\arraybackslash}p{2.5cm} >{\centering\arraybackslash}p{2.5cm} >{\centering\arraybackslash}p{2.5cm}}
	\casket & \casket & \casket & \casket & \casket\medskip\\
	}
	%hier wird die fehlende Tabellenzeile angegeben. Aufruf erfolgt durch:
	%\begin{mcquestionh}{<Frage>}
	%   	\answerh{<Antwort>} & \answerh{<Antwort>} & \answerh{<Antwort>} & \answerh{<Antwort>} & \answerh{<Antwort>}
	%\end{mcquestionh}
	{
	\\
	\end{tabular}\end{center}
}


%Antwort für horizontale MC-Fragen mit \answerh{<Antwort>}
\newcommand{\answerh}[1]{%
	\parbox{2.5cm}{\centering{#1}}
}


%Umgebung für vertikale MC-Fragen
\newenvironment{mcquestionv}[1]{%
    	#1%
    	\begin{enumerate}[\raisebox{-3pt}{\casket}]
    	}
    	%hier wird das fehlende \item angegeben werden. Aufruf erfolgt durch:
    	%\begin{mcquestionv}{<Frage>}
    	%   	\answerv{<Antwort>}
    	%  	\answerv{<Antwort>}
    	%   	\answerv{<Antwort>}
    	%   	\answerv{<Antwort>}
    	%   	\answerv{<Antwort>}
    	%\end{mcquestionv}
    	{
    	\end{enumerate}
}


%Antwort für vertikale MC-Frage mit \answerv{<Antwort>}
\newcommand{\answerv}[1]{%
	\item \hspace{0.1cm} #1
}
%----------------------------------------------------------------------------------------------------%









%----------------------------------------------------------------------------------------------------%
%%%%%%%%%%%                                          Tikz Settings                                          %%%%%%%%%%%%
%----------------------------------------------------------------------------------------------------%


%Koordinatensystem im Standardstil der Mathe-Arbeitsblätter (Gitter gray!30, Achsen -latex, Beschriftung \tiny)
%\koordinatensystem[Schritt]{xmin}{xmax}{ymin}{ymax};
%             TikZ-Stile punkt, punktoffen, baumdiagramm, wkeit
% 2026-11-09  TikZ-Stil baumdiagrammvertikal (Baumdiagramm von oben nach unten) – ganze Zahlen, Bereich muss 0 enthalten,
%Striche und Zahlen bei allen Vielfachen von Schritt (Standard 2). Innerhalb von tikzpicture[scale=0.42] verwenden.
\newcommand{\koordinatensystem}[5][2]{%
	\foreach \kosyx in {#2,...,#3} \draw[gray!30, thin] (\kosyx,#4) -- (\kosyx,#5);
	\foreach \kosyy in {#4,...,#5} \draw[gray!30, thin] (#2,\kosyy) -- (#3,\kosyy);
	\draw[-latex, thick] ({#2-0.3},0) -- ({#3+0.5},0) node[right] {$x$};
	\draw[-latex, thick] (0,{#4-0.3}) -- (0,{#5+0.5}) node[above] {$y$};
	\foreach \kosyx in {#2,...,#3} {%
		\pgfmathtruncatemacro{\kosyrest}{mod(abs(\kosyx),#1)}%
		\ifnum\kosyrest=0 \ifnum\kosyx=0 \else
			\draw (\kosyx,0.1) -- (\kosyx,-0.1) node[below, font=\tiny] {$\kosyx$};
		\fi\fi}
	\foreach \kosyy in {#4,...,#5} {%
		\pgfmathtruncatemacro{\kosyrest}{mod(abs(\kosyy),#1)}%
		\ifnum\kosyrest=0 \ifnum\kosyy=0 \else
			\draw (0.1,\kosyy) -- (-0.1,\kosyy) node[left, font=\tiny] {$\kosyy$};
		\fi\fi}
}


%optionale Eigenschaft für ein  Tikzpicture, damit dieses bündig am oberen Rand einer enumerate-Auflistung steht (benötigt tikzlibrary calc!!!)
\tikzset{
  topline/.style={baseline={([yshift=#1]current bounding box.north)}},
  topline/.default=-\ht\strutbox
}


%Punkte in Graphen und Figuren: \node[punkt, label=above:{$P$}] at (2,3) {}; (offen: punktoffen)
%Als Knoten, damit die Größe nicht mit scale=0.42 mitschrumpft (circle (3pt) würde mitskaliert)
\tikzset{
  punkt/.style={circle, fill, inner sep=0pt, minimum size=4pt},
  punktoffen/.style={circle, draw, fill=white, inner sep=0pt, minimum size=4pt}
}


%Baumdiagramme (Stochastik) von links nach rechts: \begin{tikzpicture}[baumdiagramm] \node[punkt]{} child{...}
%Wahrscheinlichkeiten an den Ästen: edge from parent node[wkeit, above] {$0{,}6$} (oberer Ast above, unterer below)
\tikzset{
  baumdiagramm/.style={grow'=right, sloped, level distance=2.6cm,
    level 1/.style={sibling distance=3cm}, level 2/.style={sibling distance=1.5cm},
    level 3/.style={sibling distance=0.75cm}, edge from parent/.style={draw, thin}},
  wkeit/.style={font=\scriptsize, inner sep=2pt}
}


%Baumdiagramme von oben nach unten: \begin{tikzpicture}[baumdiagrammvertikal] \node[punkt]{} child{...}
%erstes child = linker Ast; Wahrscheinlichkeit links mit node[wkeit, left], rechts mit node[wkeit, right]
\tikzset{
  baumdiagrammvertikal/.style={grow=down, level distance=1.6cm,
    level 1/.style={sibling distance=4.8cm}, level 2/.style={sibling distance=2.4cm},
    level 3/.style={sibling distance=1.2cm}, edge from parent/.style={draw, thin}}
}
%----------------------------------------------------------------------------------------------------%









%----------------------------------------------------------------------------------------------------%
%%%%%%%%%%%                                     Handreichungen                                        %%%%%%%%%%%
%----------------------------------------------------------------------------------------------------%


%Schriftzug TikZ im Fließtext (Leerzeichen danach ist eingebaut)
\newcommand{\TikZ}{Ti\textit{k}Z }


%Oxidationszahl über ein Elementsymbol: $\oz{+II}{Cu}$, in \ce{...} zwischen Dollarzeichen: \ce{$\oz{0}{Zn}$ + ...}
\newcommand{\oz}[2]{\overset{\scriptscriptstyle\mathrm{#1}}{\mathrm{\ce{#2}}}}


%Umgebung passend: Inhalt (z.B. eine große tikzpicture) höchstens so breit wie die Zeile setzen.
%Ist er schmaler, bleibt er unverändert. \begin{passend} ... \end{passend} (Inhalt nicht als Argument, # erlaubt)
\newsavebox{\passendbox}
\newenvironment{passend}{\begin{lrbox}{\passendbox}}{\end{lrbox}%
	\ifdim\wd\passendbox>\linewidth\resizebox{\linewidth}{!}{\usebox{\passendbox}}\else\usebox{\passendbox}\fi}


%Box für die Ausgabe unter einem Code-Beispiel: erst lstlisting[style=latex], darunter \begin{ausgabe} ... \end{ausgabe}
\newtcolorbox{ausgabe}{%
	colback=blue!5,%
	colframe=blue!50!black,%
	title=Ausgabe,%
	fonttitle=\bfseries%
}
%----------------------------------------------------------------------------------------------------%









%----------------------------------------------------------------------------------------------------%
%%%%%%%%%%%                                        Lautschrift                                         %%%%%%%%%%%
%----------------------------------------------------------------------------------------------------%


%Lautschrift in eckigen Klammern mit \ipa{"a:b@nt} – schaltet das babel-Kürzel " vorher aus (sonst wird "a zu ä)
\makeatletter
\newcommand{\ipa}{\begingroup\shorthandoff{"}\ipa@aux}
\newcommand{\ipa@aux}[1]{[\textipa{#1}]\endgroup}
\makeatother
%----------------------------------------------------------------------------------------------------%









%----------------------------------------------------------------------------------------------------%
%%%%%%%%%%%                                       Sprechblasen                                         %%%%%%%%%%%
%----------------------------------------------------------------------------------------------------%


%Sprechblase mit Zipfel links unten (Lehrkraft) mit \lehrerblase[TikZ-Optionen]{Text} (benötigt shapes.callouts)
\newcommand{\lehrerblase}[2][]{%
  \tikz[baseline=(B.base)]{%
    \node[rectangle callout, rounded corners=6pt, callout relative pointer={(-13mm,-7mm)},
      draw, thick, fill=blue!6, align=left, inner sep=3.5mm, #1] (B) {#2};}%
}


%Sprechblase mit Zipfel rechts unten (Schülerin/Schüler) mit \schuelerblase[TikZ-Optionen]{Text}
\newcommand{\schuelerblase}[2][]{%
  \tikz[baseline=(B.base)]{%
    \node[rectangle callout, rounded corners=6pt, callout relative pointer={(13mm,-7mm)},
      draw, thick, fill=green!6, align=left, inner sep=3.5mm, #1] (B) {#2};}%
}
%----------------------------------------------------------------------------------------------------%









%----------------------------------------------------------------------------------------------------%
%%%%%%%%%%%                            TikZ-Stile für Programmablaufpläne                            %%%%%%%%%%%
%----------------------------------------------------------------------------------------------------%


%Ein-/Ausgabe als Parallelogramm, Operation als Rechteck, Verzweigung als Raute (benötigt shapes.geometric)
%align=center, damit \\ im Knotentext eine neue Zeile ergibt (ohne align wird \\ stillschweigend verschluckt)
\tikzstyle{input}=[trapezium, trapezium left angle=60, trapezium right angle=120, fill=green!20, draw=green, align=center]
\tikzstyle{operation}=[rectangle, fill=orange!20, draw=orange, align=center]
\tikzstyle{condition}=[diamond, fill=magenta!20, draw=magenta, align=center]
\tikzstyle{output}=[trapezium, trapezium left angle=60, trapezium right angle=120, fill=red!20, draw=red, align=center]
%----------------------------------------------------------------------------------------------------%









%----------------------------------------------------------------------------------------------------%
%%%%%%%%%%%                            TikZ-Stile für UML-Klassendiagramme                           %%%%%%%%%%%
%----------------------------------------------------------------------------------------------------%


%Klasse mit drei Teilen: Name, Attribute, Methoden (benötigt shapes.multipart)
\tikzstyle{umlclass} = [rectangle split, rectangle split parts=3, rectangle split part align={center,left,left}, draw, thick, text width=4cm, minimum height=2cm]


%Beziehungen (benötigt arrows.meta)
\tikzstyle{inherit} = [-{Triangle[open]}, thick]	%Vererbung
\tikzstyle{compose} = [-{Diamond}, thick]	%Komposition
\tikzstyle{aggregate} = [-{Diamond[open]}, thick]	%Aggregation
\tikzstyle{associate} = [-, thick]	%Assoziation
\tikzstyle{depend} = [-, thick, dashed]	%Abhängigkeit
%----------------------------------------------------------------------------------------------------%








%----------------------------------------------------------------------------------------------------%
%%%%%%%%%%%                                         Listing Settings                                       %%%%%%%%%%%%
%----------------------------------------------------------------------------------------------------%


%Umlaute und häufige Sonderzeichen in ALLEN Listings (Python, Java, ...), sonst Fehler »Invalid UTF-8 byte sequence«
\lstset{literate=%
	{ä}{{\"a}}1 {ö}{{\"o}}1 {ü}{{\"u}}1 {Ä}{{\"A}}1 {Ö}{{\"O}}1 {Ü}{{\"U}}1 {ß}{{\ss}}1%
	{·}{{$\cdot$}}1 {€}{{\euro}}1 {°}{{\degree}}1 {→}{{$\rightarrow$}}2 {≤}{{$\leq$}}1 {≥}{{$\geq$}}1 {≠}{{$\neq$}}1%
	{–}{{--}}1 {—}{{---}}1 {„}{{\glqq}}1 {“}{{\grqq}}1 {…}{{\dots}}3 {×}{{$\times$}}1 {²}{{$^2$}}1 {³}{{$^3$}}1 {π}{{$\pi$}}1%
}


%Einstellung der Codeumgebung für LaTeX-Code-Beispiele (Stil der Handreichungen: ohne Zeilennummern)
%Verwendung: \begin{lstlisting}[style=latex] ... \end{lstlisting}, darunter \begin{ausgabe} ... \end{ausgabe}
\lstdefinestyle{latex}{%
  	language=TeX,%
  	frame=single,%
  	backgroundcolor=\color{yellow!20},%
  	rulecolor=\color{lightgray!40},%
  	breaklines=true,%
  	basicstyle=\scriptsize\ttfamily,%
  	keywordstyle=\bfseries\color{purple!40!black},%
  	commentstyle=\itshape\color{green!40!black},%
  	identifierstyle=\color{blue},%
	stringstyle=\color{black},%
	captionpos=b,%
	tabsize=2,%
	showspaces=false,%
	showstringspaces=false,%
  	showtabs=false,%
	deletekeywords={document},%
	morekeywords={begin}%
}


%Variante mit Zeilennummern: \begin{lstlisting}[style=latexnummeriert]
\lstdefinestyle{latexnummeriert}{%
	style=latex,%
	numbers=left,%
	numberstyle=\tiny\color{gray!90!black},%
  	numbersep=-5pt%
}
%----------------------------------------------------------------------------------------------------%






%----------------------------------------------------------------------------------------------------%
%%%%%%%%%%%                          Formularfelder (ausfüllbare PDFs)                          %%%%%%%%%%%
%----------------------------------------------------------------------------------------------------%
% Nur innerhalb von \begin{Form} ... \end{Form} (hyperref) verwendbar. Für Optionsfelder (radio)
% im Dokument zusätzlich \hypersetup{unicode=false}, sonst bekommen sie kaputte Namen.
% Jedes Feld steckt in einer eigenen Box (\raisebox/\hbox): Ohne sie übernimmt pdfTeX die Höhe
% der umgebenden Zeile, in Tabellen werden die Felder dann so hoch wie die ganze Zeile.

%Aussehen aller Felder: blauer Rand, hellblauer Hintergrund, Schrift 11 pt, wird mitgedruckt
\renewcommand{\DefaultOptionsofText}{print, bordercolor={0 0.353 0.612}, backgroundcolor={0.94 0.96 1}, charsize=11pt}
\renewcommand{\DefaultOptionsofCheckBox}{print, bordercolor={0 0.353 0.612}, backgroundcolor={0.94 0.96 1}, charsize=11pt}
\renewcommand{\DefaultOptionsofRadio}{print, notoggletooff, bordercolor={0 0.353 0.612}, backgroundcolor={0.94 0.96 1}, charsize=11pt}
\renewcommand{\DefaultOptionsofPopdownBox}{print, bordercolor={0 0.353 0.612}, backgroundcolor={0.94 0.96 1}, charsize=11pt}
\renewcommand{\DefaultHeightofText}{0.6cm}

%\feld[Breite]{Name} – einzeiliges Textfeld, sitzt auf der Grundlinie (Standardbreite 2,4 cm)
\newcommand{\feld}[2][2.4cm]{\raisebox{-0.17cm}{\TextField[name=#2, width=#1]{}}}

%\zeichen{Name} – Feld für genau ein Zeichen, zentriert (z. B. E/V/A, Buchstabe einer Reihenfolge)
\newcommand{\zeichen}[1]{\raisebox{-0.12cm}{\TextField[name=#1, width=0.8cm, height=0.5cm, maxlen=1, align=1]{}}}

%\kreuz{Name} – Ankreuzkästchen
\newcommand{\kreuz}[1]{\raisebox{-0.08cm}{\CheckBox[name=#1, width=0.42cm, height=0.42cm]{}}}

%\antwort[Höhe]{Name} – mehrzeiliges Antwortfeld über die ganze Zeilenbreite (Standardhöhe 3,2 cm ≈ 4 Zeilen)
\newcommand{\antwort}[2][3.2cm]{\par\vspace{0.15cm}\noindent\hbox{\TextField[name=#2, multiline=true, width=\linewidth, height=#1]{}}}
%----------------------------------------------------------------------------------------------------%




%----------------------------------------------------------------------------------------------------%
%%%%%%%%%%%                                   Schach (farbige Bretter)                                   %%%%%%%%%%
%----------------------------------------------------------------------------------------------------%
% Die Standardschrift von skak schraffiert die dunklen Felder; Figur und Schraffur sind dort ein
% Zeichen und lassen sich nicht getrennt einfärben. Deshalb: Schrift skaknew in der Kodierung LSBC5
% (zeichnet nur die Figuren, weiße Figuren weiß gefüllt), die Felder malt chessboard als Flächen
% darunter. Eigene Markierungen mit colorbackfields liegen über den Feldern und unter den Figuren.

%Feldfarben der Brettschemata (hell/dunkel) und Markierungsfarbe für den letzten Zug
\definecolor{brettblauhell}{HTML}{EEF3FD} %Schema blau, helle Felder
\definecolor{brettblaudunkel}{HTML}{A9B4CC} %Schema blau, dunkle Felder
\definecolor{brettbraunhell}{HTML}{F0D9B5} %Schema braun (klassisches Holzbrett), helle Felder
\definecolor{brettbraundunkel}{HTML}{B58863} %Schema braun, dunkle Felder
\definecolor{brettgruenhell}{HTML}{EEEED2} %Schema gruen (Turnierbrett), helle Felder
\definecolor{brettgruendunkel}{HTML}{769656} %Schema gruen, dunkle Felder
\definecolor{brettgrauhell}{HTML}{F2F2F2} %Schema grau (für Schwarz-Weiß-Druck), helle Felder
\definecolor{brettgraudunkel}{HTML}{B4B4B4} %Schema grau, dunkle Felder
\definecolor{brettmarke}{HTML}{F6D55C} %Start- und Zielfeld eines Zuges hervorheben: pgfstyle=color, color=brettmarke, colorbackfields={e2,e4}

% Setzt die Voreinstellungen von \setchessboard zurück – eigene \setchessboard{...} deshalb DANACH schreiben.
%\schachbrettfarben{Schema} – Farben aller folgenden Schachbretter: braun (Standard), blau, gruen, grau oder klassisch (schraffiert wie in skak)
\makeatletter
\let\brett@grundeinstellung\board@val@savedkeylist
\newcommand{\schachbrettfarben}[1]{%
	\let\board@val@savedkeylist\brett@grundeinstellung
	\ifthenelse{\equal{#1}{klassisch}}{}%
	{\setchessboard{boardfontfamily=skaknew, boardfontencoding=LSBC5,
		piececolor=black, piecemaskcolor=white, setfontcolors,
		pgfstyle=color, opacity=1, color=brett#1hell, colorbackboard, color=brett#1dunkel, colorblackbackfields,
		color=black, labelfont=\sffamily\color{black!65}}}}
\makeatother
\schachbrettfarben{braun}
%----------------------------------------------------------------------------------------------------%


% Mathematisches Wörterbuch Deutsch – Japanisch mit Lesung und Umschrift
% Skills: latex-grundstil, latex-grundlagen (Formelsatz), latex-geometrie, latex-funktionsgraphen,
%         latex-venn-diagramme
% Kompilieren mit pdflatex (dreimal, wegen Inhaltsverzeichnis und Seitenzahlen im Register).
% Benötigt die Pakete CJK und ipaex-type1 (beide in TeX Live/MacTeX enthalten).

%----------------------------------------------------------------------------------------------------%
% Dokumentspezifische Befehle (nur für dieses Wörterbuch, deshalb lokal und nicht in newcommands.tex)
%----------------------------------------------------------------------------------------------------%

% Japanisch mit pdflatex: CJKutf8 setzt Kanji und Kana, sobald der Text in \begin{CJK}{UTF8}{ipxm} steht.
% Schriften: IPAex Mincho (ipxm) und IPAex Gothic (ipxg) aus dem Paket ipaex-type1 (in TeX Live/MacTeX
% enthalten). Das Paket steht bewusst hier und nicht in der universalen header.tex, weil es nur dieses
% Dokument braucht und die japanischen Schriften sonst für jedes Dokument installiert sein müssten.
\usepackage{CJKutf8}
\newcommand{\jfett}[1]{\leavevmode{\CJKfamily{ipxg}#1}}   % Hervorhebung: im Japanischen Gothic statt fett
% Absatzschutz: Beginnt ein Absatz mit einem japanischen Zeichen, öffnet CJK erst seine Gruppe und beginnt
% dann den Absatz. Bricht TeX genau dort die Seite um, läuft die Ausgaberoutine innerhalb dieser Gruppe:
% Kopf- und Fußzeile bekommen dann die Zeichentabelle (ToUnicode) der japanischen Schrift, und der
% lateinische Text lässt sich im PDF nicht mehr suchen oder kopieren. Deshalb beginnt der Absatz hier
% zuerst. Der Befehl steht direkt hinter \begin{CJK}, weil CJK seine Zeichenbefehle erst dort anlegt.
\makeatletter
\newcommand{\CJKabsatzschutz}{\global\let\CJK@altchar@ohne\CJK@altchar
	\gdef\CJK@altchar{\leavevmode\CJK@altchar@ohne}}
\makeatother
\newcommand{\jp}[2]{#1 (\textit{#2})}           % japanisches Wort mit Umschrift im deutschen Fließtext

% \eintrag{Marke}{Artikel}{Begriff}{Japanisch}{Lesung in Hiragana}{Umschrift}{Zeichen/Skizze}
% Die Lesung bleibt leer, wenn der Begriff schon in Kana geschrieben ist.
% Schaltet das babel-Kürzel " aus, BEVOR die Argumente gelesen werden (wie in den anderen Wörterbüchern).
\makeatletter
\newcommand\eintrag{\begingroup\shorthandoff{"}\eintrag@aux}
\newcommand\eintrag@aux[7]{\endgroup
	#2 \textbf{#3}\phantomsection\label{e:#1}
	& \jfett{#4}\newline{\small\color{black!70}\ifx\relax#5\relax\else#5\hspace{0.6em}\fi\mbox{[#6]}}
	& #7 \\ \midrule}
\makeatother

% Tabellenkopf für jedes Kapitel (wird auf jeder Folgeseite wiederholt)
\newenvironment{woerter}{%
	\arrayrulecolor{black!20}%
	\begin{longtable}{>{\raggedright\arraybackslash}m{5.0cm} >{\raggedright\arraybackslash}m{6.4cm} >{\centering\arraybackslash}m{4.6cm}}
		\rowcolor{cornflower!15}\textbf{Deutsch} & \textbf{Japanisch} & \textbf{Zeichen \,/\, Skizze} \\ \midrule
		\endhead}{%
	\end{longtable}\arrayrulecolor{black}}

% Kleine Skizzen (Gestaltung wie im Skill latex-geometrie: Figur blau very thick, Maße/Hervorhebung rot)
\tikzset{
	fig/.style={blue, very thick},
	mass/.style={red, thick},
	hilf/.style={dashed},
	achse/.style={-{Stealth[length=4pt]}, thin},
}
\newcommand{\skizze}[1]{\tikz[baseline=(current bounding box.center), scale=0.45, line cap=round,
	line join=round, every node/.style={font=\scriptsize}, outer sep=0pt]{#1}}

% Grundbegriffe
\newcommand{\skPunkt}{\skizze{\fill[blue] (0,0) circle (2.5pt); \node[above right] at (0,0) {$P$};}}
\newcommand{\skGerade}{\skizze{\draw[fig] (0,0) -- (3.6,0.8); \node[below right] at (3.2,0.7) {$g$};}}
\newcommand{\skStrecke}{\skizze{\draw[fig] (0,0) -- (3.2,0.7); \fill (0,0) circle (2pt) (3.2,0.7) circle (2pt);
	\node[below] at (0,0) {$A$}; \node[below] at (3.2,0.7) {$B$};}}
\newcommand{\skStrahl}{\skizze{\draw[fig, -{Stealth}] (0,0) -- (3.6,0.8); \fill (0,0) circle (2pt);
	\node[below] at (0,0) {$S$};}}
\newcommand{\skParallel}{\skizze{\draw[fig] (0,0) -- (3.5,0.6) node[right] {$g$};
	\draw[fig] (0,-0.9) -- (3.5,-0.3) node[right] {$h$};}}
\newcommand{\skSenkrecht}{\skizze{\draw[fig] (-1.6,0) -- (1.6,0) node[right] {$g$};
	\draw[fig] (0,-0.5) -- (0,1.8) node[right] {$h$};
	\draw[mass] (0.45,0) arc[start angle=0, end angle=90, radius=0.45]; \fill[red] (45:0.24) circle (1pt);}}
\newcommand{\skAbstand}{\skizze{\draw[fig] (-0.4,0) -- (3.2,0) node[right] {$g$}; \fill (1.4,1.6) circle (2pt) node[above] {$P$};
	\draw[mass, hilf] (1.4,1.6) -- node[right] {$d$} (1.4,0);
	\draw[mass] (1.75,0) arc[start angle=0, end angle=90, radius=0.35]; \fill[red] (1.4,0) ++(45:0.19) circle (0.8pt);}}

% Winkel
\newcommand{\skWinkel}[1]{\skizze{\draw[mass] (0.7,0) arc[start angle=0, end angle=#1, radius=0.7];
	\draw[fig] (2.3,0) -- (0,0) -- (#1:2.3); \node at ({#1/2}:1.15) {$\alpha$};}}
\newcommand{\skRechterWinkel}{\skizze{\draw[mass] (0.6,0) arc[start angle=0, end angle=90, radius=0.6];
	\fill[red] (45:0.32) circle (1pt); \draw[fig] (2.3,0) -- (0,0) -- (0,2.3);}}
\newcommand{\skVollwinkel}{\skizze{\draw[mass] (0.6,0) arc[start angle=0, end angle=355, radius=0.6];
	\draw[fig] (0,0) -- (2.3,0); \fill[blue] (0,0) circle (1.5pt); \node at (135:1.0) {$\alpha$};}}
\newcommand{\skNebenwinkel}{\skizze{\draw[mass] (0.7,0) arc[start angle=0, end angle=60, radius=0.7];
	\draw[mass] (60:0.55) arc[start angle=60, end angle=180, radius=0.55];
	\draw[fig] (-2.2,0) -- (2.2,0); \draw[fig] (0,0) -- (60:2.2);
	\node at (28:1.15) {$\alpha$}; \node at (125:1.0) {$\beta$};}}
\newcommand{\skScheitelwinkel}{\skizze{\draw[mass] (143.1:0.6) arc[start angle=143.1, end angle=201.8, radius=0.6];
	\draw[mass] (323.1:0.6) arc[start angle=323.1, end angle=381.8, radius=0.6];
	\draw[fig] (-2,-0.8) -- (2,0.8); \draw[fig] (-1.6,1.2) -- (1.6,-1.2);
	\node at (172:1.05) {$\alpha$}; \node at (352:1.1) {$\alpha'$};}}
\newcommand{\skStufenwinkel}{\skizze{\draw[mass] (0.5,0) arc[start angle=0, end angle=53.13, radius=0.5];
	\draw[mass] (1.7,1.6) arc[start angle=0, end angle=53.13, radius=0.5];
	\draw[fig] (-1.2,0) -- (2.6,0); \draw[fig] (-1.2,1.6) -- (2.6,1.6); \draw[fig] (-0.6,-0.8) -- (1.8,2.4);
	\node at (26:0.85) {$\alpha$}; \node at ($(1.2,1.6)+(26:0.85)$) {$\beta$};}}
\newcommand{\skWechselwinkel}{\skizze{\draw[mass] (0.5,0) arc[start angle=0, end angle=53.13, radius=0.5];
	\draw[mass] (0.7,1.6) arc[start angle=180, end angle=233.13, radius=0.5];
	\draw[fig] (-1.2,0) -- (2.6,0); \draw[fig] (-1.2,1.6) -- (2.6,1.6); \draw[fig] (-0.6,-0.8) -- (1.8,2.4);
	\node at (26:0.85) {$\alpha$}; \node at ($(1.2,1.6)+(206:0.85)$) {$\beta$};}}

% Abbildungen
\newcommand{\skSymmetrieachse}{\skizze{\draw[fig] (0,0) -- (2.6,0) -- (1.3,2.1) -- cycle;
	\draw[mass, hilf] (1.3,-0.5) -- (1.3,2.5);}}
\newcommand{\skSpiegelung}{\skizze{\draw[mass, hilf] (0,-0.4) -- (0,1.8);
	\draw[fig] (-2,0) -- (-0.5,0.3) -- (-1.5,1.4) -- cycle;
	\draw[fig, blue!45] (2,0) -- (0.5,0.3) -- (1.5,1.4) -- cycle;}}
\newcommand{\skVerschiebung}{\skizze{\draw[fig] (0,0) -- (1.3,0) -- (0.3,1.1) -- cycle;
	\draw[fig, blue!45] (2,0.6) -- (3.3,0.6) -- (2.3,1.7) -- cycle;
	\draw[mass, -{Stealth}] (0.3,1.1) -- (2.3,1.7);}}
\newcommand{\skDrehung}{\skizze{\draw[fig] (0.6,0) -- (1.9,0) -- (0.8,0.9) -- cycle;
	\draw[fig, blue!45] (0,0.6) -- (0,1.9) -- (-0.9,0.8) -- cycle;
	\fill (0,0) circle (2pt) node[below left] {$Z$};
	\draw[mass, -{Stealth}] (10:2.2) arc[start angle=10, end angle=80, radius=2.2];}}

% Dreiecke
\newcommand{\skDreieck}{\skizze{\draw[fig] (0,0) -- (3,0) -- (1,2) -- cycle;
	\node[below left] at (0,0) {$A$}; \node[below right] at (3,0) {$B$}; \node[above] at (1,2) {$C$};}}
\newcommand{\skRechtwinklig}{\skizze{\draw[mass] (0.5,0) arc[start angle=0, end angle=90, radius=0.5];
	\fill[red] (45:0.27) circle (1pt); \draw[fig] (0,0) -- (3,0) -- (0,2) -- cycle;}}
\newcommand{\skGleichschenklig}{\skizze{\draw[fig] (0,0) -- (2.6,0) -- (1.3,2.2) -- cycle;
	\draw[mass] (0.52,1.176) -- (0.78,1.024); \draw[mass] (1.82,1.024) -- (2.08,1.176);}}
\newcommand{\skGleichseitig}{\skizze{\draw[fig] (0,0) -- (2.6,0) -- (1.3,2.252) -- cycle;
	\draw[mass] (1.3,-0.15) -- (1.3,0.15); \draw[mass] (0.52,1.201) -- (0.78,1.051);
	\draw[mass] (1.82,1.051) -- (2.08,1.201);}}
\newcommand{\skHypotenuse}{\skizze{\draw[fig] (0,2) -- (0,0) -- (3,0); \draw[red, very thick] (3,0) -- (0,2);
	\draw[mass] (0.5,0) arc[start angle=0, end angle=90, radius=0.5]; \fill[red] (45:0.27) circle (1pt);
	\node[above right] at (1.5,1) {$c$};}}
\newcommand{\skKathete}{\skizze{\draw[fig] (3,0) -- (0,2); \draw[red, very thick] (0,2) -- (0,0) -- (3,0);
	\draw[mass] (0.5,0) arc[start angle=0, end angle=90, radius=0.5]; \fill[red] (45:0.27) circle (1pt);
	\node[left] at (0,1) {$b$}; \node[below] at (1.5,0) {$a$};}}
\newcommand{\skPythagoras}{\skizze{\draw[fig] (0,0) -- (2,0) -- (0,1.5) -- cycle;
	\draw[fig, fill=blue!10] (0,0) rectangle (2,-2); \draw[fig, fill=blue!10] (0,0) rectangle (-1.5,1.5);
	\draw[fig, fill=red!10] (2,0) -- (3.5,2) -- (1.5,3.5) -- (0,1.5) -- cycle;
	\node at (1,-1) {$a^2$}; \node at (-0.75,0.75) {$b^2$}; \node at (1.75,1.75) {$c^2$};}}
\newcommand{\skHoehe}{\skizze{\draw[fig] (0,0) -- (3.2,0) -- (1,2) -- cycle;
	\draw[mass, hilf] (1,2) -- node[right] {$h$} (1,0);
	\draw[mass] (1.35,0) arc[start angle=0, end angle=90, radius=0.35]; \fill[red] (1,0) ++(45:0.19) circle (0.8pt);}}

% Vierecke, Vielecke
\newcommand{\skViereck}{\skizze{\draw[fig] (0,0) -- (2.8,0.3) -- (2.3,1.9) -- (0.4,1.5) -- cycle;}}
\newcommand{\skQuadrat}{\skizze{\draw[fig] (0,0) rectangle (1.9,1.9);}}
\newcommand{\skRechteck}{\skizze{\draw[fig] (0,0) rectangle (3.2,1.6);}}
\newcommand{\skParallelogramm}{\skizze{\draw[fig] (0,0) -- (2.6,0) -- (3.4,1.5) -- (0.8,1.5) -- cycle;}}
\newcommand{\skRaute}{\skizze{\draw[fig] (0,0) -- (1.4,-1) -- (2.8,0) -- (1.4,1) -- cycle;}}
\newcommand{\skTrapez}{\skizze{\draw[fig] (0,0) -- (3.2,0) -- (2.4,1.5) -- (0.6,1.5) -- cycle;}}
\newcommand{\skDrachen}{\skizze{\draw[fig] (1,0) -- (2,1.6) -- (1,2.3) -- (0,1.6) -- cycle;
	\draw[mass, hilf] (1,-0.3) -- (1,2.6);}}
\newcommand{\skFuenfeck}{\skizze{\draw[fig] (90:1.2) -- (162:1.2) -- (234:1.2) -- (306:1.2) -- (18:1.2) -- cycle;}}
\newcommand{\skSechseck}{\skizze{\draw[fig] (0:1.2) -- (60:1.2) -- (120:1.2) -- (180:1.2) -- (240:1.2) -- (300:1.2) -- cycle;}}
\newcommand{\skFlaeche}{\skizze{\filldraw[gray, opacity=0.5] (0,0) rectangle (3,1.6);
	\draw[fig] (0,0) rectangle (3,1.6); \node at (1.5,0.8) {$A$};}}
\newcommand{\skUmfangViereck}{\skizze{\draw[red, very thick] (0,0) rectangle (3,1.6); \node at (1.5,0.8) {$u$};}}

% Kreis
\newcommand{\skKreis}{\skizze{\draw[fig] (0,0) circle (1.2); \fill (0,0) circle (1.5pt);}}
\newcommand{\skMittelpunkt}{\skizze{\draw[fig] (0,0) circle (1.2); \fill[red] (0,0) circle (2.5pt) node[below] {$M$};}}
\newcommand{\skRadius}{\skizze{\draw[fig] (0,0) circle (1.2); \fill (0,0) circle (1.5pt) node[below] {$M$};
	\draw[mass] (0,0) -- node[above] {$r$} (30:1.2);}}
\newcommand{\skDurchmesser}{\skizze{\draw[fig] (0,0) circle (1.2); \fill (0,0) circle (1.5pt);
	\draw[mass] (200:1.2) -- (20:1.2); \node[above] at (100:0.25) {$d$};}}
\newcommand{\skKreisumfang}{\skizze{\draw[red, very thick] (0,0) circle (1.2); \fill (0,0) circle (1.5pt);}}
\newcommand{\skSektor}{\skizze{\filldraw[gray, opacity=0.5] (0,0) -- (1.2,0) arc[start angle=0, end angle=70, radius=1.2] -- cycle;
	\draw[fig] (0,0) circle (1.2); \draw[mass] (1.2,0) -- (0,0) -- (70:1.2);}}
\newcommand{\skSehne}{\skizze{\draw[fig] (0,0) circle (1.2); \draw[mass] (160:1.2) -- (40:1.2);}}
\newcommand{\skTangente}{\skizze{\draw[fig] (0,0) circle (1.1); \draw[mass] ($(50:1.1)+(140:1.5)$) -- ($(50:1.1)+(-40:1.5)$);
	\fill[red] (50:1.1) circle (2.5pt);}}

% Körper (Schrägbild: Tiefe unter 45° und halb so lang, verdeckte Kanten gestrichelt)
\newcommand{\skWuerfel}{\skizze{\def\d{0.57}
	\draw[fig, hilf, thick] (0,0) -- (\d,\d) -- (1.6+\d,\d) (\d,\d) -- (\d,1.6+\d);
	\draw[fig] (0,0) rectangle (1.6,1.6); \draw[fig] (1.6,0) -- (1.6+\d,\d) -- (1.6+\d,1.6+\d) -- (\d,1.6+\d) -- (0,1.6)
		(1.6,1.6) -- (1.6+\d,1.6+\d);}}
\newcommand{\skQuader}{\skizze{\def\d{0.6}
	\draw[fig, hilf, thick] (0,0) -- (\d,\d) -- (2.6+\d,\d) (\d,\d) -- (\d,1.3+\d);
	\draw[fig] (0,0) rectangle (2.6,1.3); \draw[fig] (2.6,0) -- (2.6+\d,\d) -- (2.6+\d,1.3+\d) -- (\d,1.3+\d) -- (0,1.3)
		(2.6,1.3) -- (2.6+\d,1.3+\d);}}
\newcommand{\skPrisma}{\skizze{
	\draw[fig, hilf, thick] (0,0) -- (1.4,0.7) -- (3,0.7) (1.4,0.7) -- (2.2,2);
	\draw[fig] (0,0) -- (1.6,0) -- (0.8,1.3) -- cycle;
	\draw[fig] (1.6,0) -- (3,0.7) -- (2.2,2) -- (0.8,1.3);}}
\newcommand{\skZylinder}{\skizze{
	\draw[fig, hilf, thick] (1,0) arc[start angle=0, end angle=180, x radius=1, y radius=0.3];
	\draw[fig] (-1,0) arc[start angle=180, end angle=360, x radius=1, y radius=0.3];
	\draw[fig] (0,1.9) ellipse[x radius=1, y radius=0.3];
	\draw[fig] (-1,0) -- (-1,1.9) (1,0) -- (1,1.9);}}
\newcommand{\skKegel}{\skizze{
	\draw[fig, hilf, thick] (1,0) arc[start angle=0, end angle=180, x radius=1, y radius=0.3];
	\draw[fig] (-1,0) arc[start angle=180, end angle=360, x radius=1, y radius=0.3];
	\draw[fig] (-1,0) -- (0,2.2) -- (1,0);}}
\newcommand{\skPyramide}{\skizze{
	\draw[fig, hilf, thick] (0,0) -- (0.8,0.6) -- (2.8,0.6) (0.8,0.6) -- (1.4,2.3);
	\draw[fig] (0,0) -- (2,0) -- (2.8,0.6) -- (1.4,2.3) -- cycle (2,0) -- (1.4,2.3);}}
\newcommand{\skKugel}{\skizze{\draw[fig] (0,0) circle (1.1);
	\draw[fig, hilf, thick] (1.1,0) arc[start angle=0, end angle=180, x radius=1.1, y radius=0.3];
	\draw[fig] (-1.1,0) arc[start angle=180, end angle=360, x radius=1.1, y radius=0.3];}}
\newcommand{\skKante}{\skizze{\def\d{0.57}
	\draw[fig, hilf, thick] (0,0) -- (\d,\d) -- (1.6+\d,\d) (\d,\d) -- (\d,1.6+\d);
	\draw[fig] (0,0) rectangle (1.6,1.6); \draw[fig] (1.6,0) -- (1.6+\d,\d) -- (1.6+\d,1.6+\d) -- (\d,1.6+\d) -- (0,1.6);
	\draw[red, very thick] (1.6,1.6) -- (1.6+\d,1.6+\d);}}
\newcommand{\skEcke}{\skizze{\def\d{0.57}
	\draw[fig, hilf, thick] (0,0) -- (\d,\d) -- (1.6+\d,\d) (\d,\d) -- (\d,1.6+\d);
	\draw[fig] (0,0) rectangle (1.6,1.6); \draw[fig] (1.6,0) -- (1.6+\d,\d) -- (1.6+\d,1.6+\d) -- (\d,1.6+\d) -- (0,1.6)
		(1.6,1.6) -- (1.6+\d,1.6+\d); \fill[red] (1.6+\d,1.6+\d) circle (2.5pt);}}
\newcommand{\skSeitenflaeche}{\skizze{\def\d{0.57}
	\draw[fig, hilf, thick] (0,0) -- (\d,\d) -- (1.6+\d,\d) (\d,\d) -- (\d,1.6+\d);
	\fill[gray, opacity=0.5] (0,0) rectangle (1.6,1.6);
	\draw[fig] (0,0) rectangle (1.6,1.6); \draw[fig] (1.6,0) -- (1.6+\d,\d) -- (1.6+\d,1.6+\d) -- (\d,1.6+\d) -- (0,1.6)
		(1.6,1.6) -- (1.6+\d,1.6+\d);}}
\newcommand{\skNetz}{\skizze{\begin{scope}[scale=0.62]
	\foreach \x/\y in {1/0, 0/1, 1/1, 2/1, 3/1, 1/2} {\draw[fig, fill=blue!8] (\x,\y) rectangle ++(1,1);}
	\end{scope}}}

% Funktionen (Achsen dünn mit Pfeil, Graph blau, Hervorhebungen rot)
\newcommand{\achsen}[4]{\draw[achse] (#1,0) -- (#2,0) node[below left] {$x$}; \draw[achse] (0,#3) -- (0,#4) node[below left] {$y$};}
\newcommand{\skKoordinaten}{\skizze{\draw[black!25, very thin] (-1.5,-1.5) grid (2.5,2); \achsen{-1.6}{2.7}{-1.6}{2.2}
	\fill[red] (1,1.5) circle (2pt) node[right] {$P$};}}
\newcommand{\skGraph}{\skizze{\achsen{-1.8}{2.4}{-0.6}{2.4}
	\begin{scope}\clip (-1.8,-0.6) rectangle (2.2,2.2); \draw[fig, domain=-1.8:2.2, samples=40] plot (\x, {0.3*\x*\x*\x - 0.6*\x + 0.9});\end{scope}}}
\newcommand{\skXAchse}{\skizze{\draw[achse, red, thick] (-1.8,0) -- (2.4,0) node[below left] {$x$}; \draw[achse] (0,-1.2) -- (0,1.6) node[below left] {$y$};}}
\newcommand{\skParabel}{\skizze{\achsen{-1.9}{1.9}{-0.6}{2.4}
	\begin{scope}\clip (-1.9,-0.6) rectangle (1.9,2.2); \draw[fig, domain=-1.8:1.8, samples=40] plot (\x, {0.7*\x*\x});\end{scope}}}
\newcommand{\skScheitelpunkt}{\skizze{\achsen{-0.6}{3}{-0.6}{2.4}
	\begin{scope}\clip (-0.6,-0.6) rectangle (3,2.2); \draw[fig, domain=-0.3:3, samples=40] plot (\x, {0.7*(\x-1.4)*(\x-1.4)+0.4});\end{scope}
	\fill[red] (1.4,0.4) circle (2.5pt) node[below] {$S$};}}
\newcommand{\skNullstelle}{\skizze{\achsen{-1.9}{1.9}{-1.2}{1.8}
	\begin{scope}\clip (-1.9,-1.2) rectangle (1.9,1.6); \draw[fig, domain=-1.8:1.8, samples=40] plot (\x, {0.6*\x*\x-0.9});\end{scope}
	\fill[red] (-1.2247,0) circle (2.5pt) (1.2247,0) circle (2.5pt);}}
\newcommand{\skLinear}{\skizze{\achsen{-0.8}{3}{-0.6}{2.4}
	\begin{scope}\clip (-0.8,-0.6) rectangle (3,2.2); \draw[fig, domain=-0.8:3] plot (\x, {0.5*\x+0.5});\end{scope}}}
\newcommand{\skSteigung}{\skizze{\achsen{-0.8}{3}{-0.6}{2.4}
	\begin{scope}\clip (-0.8,-0.6) rectangle (3,2.2); \draw[fig, domain=-0.8:3] plot (\x, {0.5*\x+0.5});\end{scope}
	\draw[mass] (0.6,0.8) -- node[below] {$\Delta x$} (2.2,0.8) -- node[right] {$\Delta y$} (2.2,1.6);}}
\newcommand{\skAchsenabschnitt}{\skizze{\achsen{-0.8}{3}{-0.6}{2.4}
	\begin{scope}\clip (-0.8,-0.6) rectangle (3,2.2); \draw[fig, domain=-0.8:3] plot (\x, {0.5*\x+0.5});\end{scope}
	\fill[red] (0,0.5) circle (2.5pt) node[left] {$n$};}}
\newcommand{\skSchnittpunkt}{\skizze{\draw[fig] (-0.2,0) -- (3,1.6); \draw[fig] (-0.2,1.8) -- (3,0.2);
	\fill[red] (1.6,0.9) circle (2.5pt) node[above=2pt] {$S$};}}
\newcommand{\skExponential}{\skizze{\achsen{-1.9}{1.9}{-0.6}{2.4}
	\begin{scope}\clip (-1.9,-0.6) rectangle (1.9,2.2); \draw[fig, domain=-1.9:1.9, samples=40] plot (\x, {0.5*2^\x});\end{scope}}}
\newcommand{\skLogarithmus}{\skizze{\achsen{-0.6}{3.2}{-1.4}{1.6}
	\begin{scope}\clip (-0.6,-1.4) rectangle (3.2,1.4); \draw[fig, domain=0.2:3.2, samples=40] plot (\x, {ln(\x)});\end{scope}}}
\newcommand{\sinusgraph}{\achsen{-0.3}{3.5}{-1.2}{1.4}\draw[fig, domain=0:3.3, samples=50] plot (\x, {0.9*sin(2*\x r)});}
\newcommand{\skSinus}{\skizze{\sinusgraph}}
\newcommand{\skKosinus}{\skizze{\achsen{-0.3}{3.5}{-1.2}{1.4}\draw[fig, domain=0:3.3, samples=50] plot (\x, {0.9*cos(2*\x r)});}}
\newcommand{\skHochpunkt}{\skizze{\sinusgraph \fill[red] (0.785,0.9) circle (2.5pt) node[above] {$H$};}}
\newcommand{\skTiefpunkt}{\skizze{\sinusgraph \fill[red] (2.356,-0.9) circle (2.5pt) node[below] {$T$};}}
\newcommand{\skWendepunkt}{\skizze{\sinusgraph \fill[red] (1.571,0) circle (2.5pt) node[above right] {$W$};}}
\newcommand{\skMonotonSteigend}{\skizze{\achsen{-0.4}{3}{-0.4}{2.2}\draw[fig, domain=0:2.7, samples=40] plot (\x, {0.25*\x*\x+0.2});}}
\newcommand{\skMonotonFallend}{\skizze{\achsen{-0.4}{3}{-0.4}{2.2}\draw[fig, domain=0.25:2.8, samples=40] plot (\x, {0.5/\x+0.1});}}
\newcommand{\skAchsensymmetrisch}{\skizze{\draw[mass, hilf] (0,-0.6) -- (0,2.2); \achsen{-1.9}{1.9}{-0.6}{2.4}
	\begin{scope}\clip (-1.9,-0.6) rectangle (1.9,2.2); \draw[fig, domain=-1.8:1.8, samples=40] plot (\x, {0.3*\x*\x*\x*\x - 0.9*\x*\x + 1});\end{scope}}}
\newcommand{\skPunktsymmetrisch}{\skizze{\achsen{-1.9}{1.9}{-1.4}{1.6}
	\begin{scope}\clip (-1.9,-1.4) rectangle (1.9,1.4); \draw[fig, domain=-1.8:1.8, samples=40] plot (\x, {0.4*\x*\x*\x - 0.9*\x});\end{scope}
	\fill[red] (0,0) circle (2.5pt);}}
\newcommand{\skAsymptote}{\skizze{\achsen{-0.4}{3.2}{-0.4}{2.4}\draw[mass, hilf] (-0.3,0.4) -- (3.1,0.4);
	\draw[fig, domain=0.3:3.1, samples=40] plot (\x, {0.5/\x+0.4});}}
\newcommand{\skIntegral}{\skizze{\fill[gray, opacity=0.5] (0.5,0) -- plot[domain=0.5:2.5, samples=30] (\x, {0.3*(\x-1.5)*(\x-1.5)+0.6}) -- (2.5,0) -- cycle;
	\achsen{-0.4}{3.2}{-0.4}{2.2}\draw[fig, domain=0.1:2.9, samples=40] plot (\x, {0.3*(\x-1.5)*(\x-1.5)+0.6});
	\node[below] at (0.5,0) {$a$}; \node[below] at (2.5,0) {$b$};}}
\newcommand{\skTangenteGraph}{\skizze{\achsen{-0.4}{3.2}{-0.4}{2.4}
	\begin{scope}\clip (-0.4,-0.4) rectangle (3.2,2.2); \draw[fig, domain=0:3, samples=40] plot (\x, {0.25*\x*\x});
	\draw[mass] (0.4,-0.4) -- (3.2,2.4);\end{scope} \fill[red] (2,1) circle (2pt);}}
\newcommand{\skFolge}{\skizze{\achsen{-0.3}{3.4}{-0.3}{2.2}
	\foreach \n in {1,...,10} {\fill[blue] ({0.3*\n},{1.6-1.2/\n}) circle (1.6pt);}
	\draw[mass, hilf] (0,1.6) -- (3.2,1.6);}}

% Vektoren
\newcommand{\skVektor}{\skizze{\draw[mass, very thick, -{Stealth}] (0,0) -- node[above left] {$\vec{a}$} (2.6,1.3);}}
\newcommand{\skOrtsvektor}{\skizze{\achsen{-0.4}{3}{-0.4}{2.2}\draw[mass, very thick, -{Stealth}] (0,0) -- (2.2,1.5);
	\fill (2.2,1.5) circle (2pt) node[above] {$P$}; \node[below left] at (0,0) {$O$};}}
\newcommand{\skKreuzprodukt}{\skizze{\draw[fig, -{Stealth}] (0,0) -- (2.2,0) node[below] {$\vec{a}$};
	\draw[fig, -{Stealth}] (0,0) -- (1,-0.8) node[below] {$\vec{b}$};
	\draw[mass, very thick, -{Stealth}] (0,0) -- (0,2) node[right] {$\vec{a}\times\vec{b}$};}}
\newcommand{\skEbene}{\skizze{\filldraw[blue!10, draw=blue, very thick] (0,0) -- (2.8,0) -- (3.6,1.4) -- (0.8,1.4) -- cycle;
	\node at (1.8,0.7) {$E$};}}

% Mengen (wie im Skill latex-venn-diagramme: Schnitt über \clip)
\newcommand{\venn}{\draw[fig] (0,0) circle (1) (1.2,0) circle (1); \node at (-0.5,0) {$A$}; \node at (1.7,0) {$B$};}
\newcommand{\skSchnittmenge}{\skizze{\begin{scope}\clip (0,0) circle (1); \fill[gray, opacity=0.5] (1.2,0) circle (1);\end{scope}\venn}}
\newcommand{\skVereinigung}{\skizze{\fill[gray, opacity=0.5] (0,0) circle (1) (1.2,0) circle (1);\venn}}
\newcommand{\skTeilmenge}{\skizze{\draw[fig] (0,0) ellipse[x radius=1.6, y radius=1.1]; \draw[fig, fill=gray!40] (-0.5,-0.1) circle (0.55);
	\node at (-0.5,-0.1) {$A$}; \node at (0.8,0.3) {$B$};}}

% Stochastik
\newcommand{\skBaum}{\skizze{\draw[fig, thick] (0,0) -- (1.4,0.9) (0,0) -- (1.4,-0.9) (1.4,0.9) -- (2.8,1.35) (1.4,0.9) -- (2.8,0.45)
	(1.4,-0.9) -- (2.8,-0.45) (1.4,-0.9) -- (2.8,-1.35);
	\foreach \p in {(0,0),(1.4,0.9),(1.4,-0.9),(2.8,1.35),(2.8,0.45),(2.8,-0.45),(2.8,-1.35)} {\fill[blue] \p circle (1.6pt);}}}
\newcommand{\skVierfeldertafel}{{\scriptsize\begin{tabular}{c|cc|c}
	& $B$ & $\overline{B}$ & \\ \hline $A$ & & & \\ $\overline{A}$ & & & \\ \hline & & & $1$\end{tabular}}}
\newcommand{\skWertetabelle}{{\scriptsize\begin{tabular}{c|c|c|c|c} $x$ & $0$ & $1$ & $2$ & $3$ \\ \hline $y$ & $1$ & $3$ & $5$ & $7$\end{tabular}}}
\newcommand{\skSaeulen}{\skizze{\draw[achse] (-0.2,0) -- (3.4,0); \draw[achse] (0,-0.2) -- (0,2.2);
	\foreach \k/\h in {1/0.4, 2/1.2, 3/1.8, 4/1.2, 5/0.4} {\filldraw[fill=blue!25, draw=blue, thick] ({0.55*\k-0.2},0) rectangle ++(0.4,\h);}}}

% Ober- und Untersumme (Ergänzung aus Teil II)
\newcommand{\rechtecksumme}[2]{\skizze{\foreach \i in {0,...,3} {
		\pgfmathsetmacro{\xl}{0.4+0.6*\i} \pgfmathsetmacro{\xr}{1.0+0.6*\i}
		\pgfmathsetmacro{\h}{0.18*(#1)*(#1)+0.3}
		\filldraw[fill=#2!30, draw=#2] (\xl,0) rectangle (\xr,\h);}
	\achsen{-0.3}{3.2}{-0.3}{2.4}\draw[fig, domain=0:2.9, samples=30] plot (\x, {0.18*\x*\x+0.3});}}
\newcommand{\skUntersumme}{\rechtecksumme{\xl}{gray}}
\newcommand{\skObersumme}{\rechtecksumme{\xr}{red}}

%----------------------------------------------------------------------------------------------------%
% Ergänzungen: Geometrie, Körper, analytische Geometrie
%----------------------------------------------------------------------------------------------------%
\newcommand{\skScheitel}{\skizze{\draw[fig] (2.3,0) -- (0,0) -- (50:2.3); \fill[red] (0,0) circle (3pt) node[below] {$S$};}}
\newcommand{\skSchenkel}{\skizze{\draw[fig] (0,0) -- (50:2.3); \draw[red, very thick] (0,0) -- (2.3,0); \fill (0,0) circle (1.5pt);}}
\newcommand{\skKomplement}{\skizze{\draw[mass] (0.8,0) arc[start angle=0, end angle=35, radius=0.8];
	\draw[mass] (35:1.1) arc[start angle=35, end angle=90, radius=1.1];
	\draw[fig] (2.3,0) -- (0,0) -- (0,2.3); \draw[fig] (0,0) -- (35:2.3);
	\node at (17:1.25) {$\alpha$}; \node at (62:1.5) {$\beta$};}}
\newcommand{\skEntgegen}{\skizze{\draw[mass] (0.5,0) arc[start angle=0, end angle=53.13, radius=0.5];
	\draw[mass] ($(1.2,1.6)+(233.13:0.45)$) arc[start angle=233.13, end angle=360, radius=0.45];
	\draw[fig] (-1.2,0) -- (2.6,0); \draw[fig] (-1.2,1.6) -- (2.6,1.6); \draw[fig] (-0.6,-0.8) -- (1.8,2.4);
	\node at (26:0.85) {$\alpha$}; \node at ($(1.2,1.6)+(300:0.8)$) {$\beta$};}}
\newcommand{\skMittelsenkrechte}{\skizze{\draw[fig] (0,0) -- (3,0); \fill (0,0) circle (2pt) (3,0) circle (2pt);
	\draw[mass] (1.5,-0.8) -- (1.5,1.4); \draw[mass] (1.85,0) arc[start angle=0, end angle=90, radius=0.35];
	\fill[red] (1.5,0) ++(45:0.19) circle (0.8pt); \draw[mass] (0.75,-0.13) -- (0.75,0.13) (2.25,-0.13) -- (2.25,0.13);}}
\newcommand{\skWinkelhalbierende}{\skizze{\draw[mass] (0.8,0) arc[start angle=0, end angle=30, radius=0.8];
	\draw[mass] (30:0.95) arc[start angle=30, end angle=60, radius=0.95];
	\draw[fig] (2.3,0) -- (0,0) -- (60:2.3); \draw[mass, hilf] (0,0) -- (30:2.6);}}
\newcommand{\skInnenwinkel}{\skizze{\draw[mass] (0.55,0) arc[start angle=0, end angle=63.43, radius=0.55];
	\draw[mass] ($(3,0)+(135:0.55)$) arc[start angle=135, end angle=180, radius=0.55];
	\draw[mass] ($(1,2)+(243.43:0.5)$) arc[start angle=243.43, end angle=315, radius=0.5];
	\draw[fig] (0,0) -- (3,0) -- (1,2) -- cycle;}}
\newcommand{\skAussenwinkel}{\skizze{\draw[mass] (3.6,0) arc[start angle=0, end angle=135, radius=0.6];
	\draw[fig] (0,0) -- (3,0) -- (1,2) -- cycle; \draw[fig, hilf, thick] (3,0) -- (4.3,0);
	\node at ($(3,0)+(60:0.98)$) {$\delta$};}}
\newcommand{\skBasis}{\skizze{\draw[fig] (0,0) -- (1.3,2.2) -- (2.6,0); \draw[red, very thick] (0,0) -- (2.6,0);}}
\newcommand{\skBasiswinkel}{\skizze{\draw[mass] (0.6,0) arc[start angle=0, end angle=59.4, radius=0.6];
	\draw[mass] ($(2.6,0)+(120.6:0.6)$) arc[start angle=120.6, end angle=180, radius=0.6];
	\draw[fig] (0,0) -- (2.6,0) -- (1.3,2.2) -- cycle;}}
\newcommand{\rwdreieck}{\draw[mass] (0.85,0) arc[start angle=0, end angle=30.96, radius=0.85]; \node at (14:1.2) {$\alpha$};
	\draw[black!60, thick] (3,0.4) arc[start angle=90, end angle=180, radius=0.4]; \fill[black!60] ($(3,0)+(135:0.2)$) circle (0.8pt);}
\newcommand{\skAnkathete}{\skizze{\rwdreieck \draw[fig] (0,0) -- (3,1.8) -- (3,0); \draw[red, very thick] (0,0) -- (3,0);}}
\newcommand{\skGegenkathete}{\skizze{\rwdreieck \draw[fig] (3,0) -- (0,0) -- (3,1.8); \draw[red, very thick] (3,0) -- (3,1.8);}}
\newcommand{\skSeitenhalbierende}{\skizze{\draw[fig] (0,0) -- (3,0) -- (1,2) -- cycle; \draw[mass] (1,2) -- (1.5,0);
	\draw[mass] (0.75,-0.13) -- (0.75,0.13) (2.25,-0.13) -- (2.25,0.13);}}
\newcommand{\skInkreis}{\skizze{\draw[fig] (0,0) -- (3,0) -- (1,2) -- cycle; \draw[mass] (1.204,0.744) circle (0.744);}}
\newcommand{\skUmkreis}{\skizze{\draw[mass] (1.5,0.5) circle (1.581); \draw[fig] (0,0) -- (3,0) -- (1,2) -- cycle;}}
\newcommand{\skDiagonale}{\skizze{\draw[fig] (0,0) rectangle (3,1.6); \draw[mass] (0,0) -- (3,1.6);}}
\newcommand{\skNeck}[1]{\skizze{\draw[fig] (90:1.2) \foreach \k in {1,...,#1} { -- ({90+\k*360/#1}:1.2) } -- cycle;}}
\newcommand{\skSekante}{\skizze{\draw[fig] (0,0) circle (1.2); \draw[mass] ($(150:1.2)!-0.35!(30:1.2)$) -- ($(150:1.2)!1.35!(30:1.2)$);
	\fill[red] (150:1.2) circle (1.6pt) (30:1.2) circle (1.6pt);}}
\newcommand{\skBeruehrpunkt}{\skizze{\draw[fig] (0,0) circle (1.1); \draw[fig, thick] ($(50:1.1)+(140:1.5)$) -- ($(50:1.1)+(-40:1.5)$);
	\fill[red] (50:1.1) circle (3pt) node[above right] {$B$};}}
\newcommand{\skKreisbogen}{\skizze{\draw[fig] (0,0) circle (1.2); \draw[red, line width=1.8pt] (20:1.2) arc[start angle=20, end angle=110, radius=1.2];}}
\newcommand{\skKreisabschnitt}{\skizze{\fill[gray, opacity=0.5] (150:1.2) arc[start angle=150, end angle=30, radius=1.2] -- cycle;
	\draw[fig] (0,0) circle (1.2); \draw[mass] (150:1.2) -- (30:1.2);}}
\newcommand{\skKreisflaeche}{\skizze{\filldraw[fill=gray!50, draw=blue, very thick] (0,0) circle (1.2);}}

% Körper
\newcommand{\wuerfelbasis}{\def\d{0.57}\draw[fig, hilf, thick] (0,0) -- (\d,\d) -- (1.6+\d,\d) (\d,\d) -- (\d,1.6+\d);
	\draw[fig] (0,0) rectangle (1.6,1.6); \draw[fig] (1.6,0) -- (1.6+\d,\d) -- (1.6+\d,1.6+\d) -- (\d,1.6+\d) -- (0,1.6) (1.6,1.6) -- (1.6+\d,1.6+\d);}
\newcommand{\quaderbasis}{\def\d{0.6}\draw[fig, hilf, thick] (0,0) -- (\d,\d) -- (2.6+\d,\d) (\d,\d) -- (\d,1.3+\d);
	\draw[fig] (0,0) rectangle (2.6,1.3); \draw[fig] (2.6,0) -- (2.6+\d,\d) -- (2.6+\d,1.3+\d) -- (\d,1.3+\d) -- (0,1.3) (2.6,1.3) -- (2.6+\d,1.3+\d);}
\newcommand{\zylinderbasis}{\draw[fig, hilf, thick] (1,0) arc[start angle=0, end angle=180, x radius=1, y radius=0.3];
	\draw[fig] (-1,0) arc[start angle=180, end angle=360, x radius=1, y radius=0.3];
	\draw[fig] (0,1.9) ellipse[x radius=1, y radius=0.3]; \draw[fig] (-1,0) -- (-1,1.9) (1,0) -- (1,1.9);}
\newcommand{\skFlaechendiagonale}{\skizze{\wuerfelbasis \draw[red, very thick] (0,0) -- (1.6,1.6);}}
\newcommand{\skRaumdiagonale}{\skizze{\wuerfelbasis \draw[red, very thick] (0,0) -- (1.6+\d,1.6+\d);}}
\newcommand{\skGrundflaeche}{\skizze{\quaderbasis \fill[red, opacity=0.4] (0,0) -- (2.6,0) -- (2.6+\d,\d) -- (\d,\d) -- cycle;}}
\newcommand{\skDeckflaeche}{\skizze{\quaderbasis \fill[red, opacity=0.4] (0,1.3) -- (2.6,1.3) -- (2.6+\d,1.3+\d) -- (\d,1.3+\d) -- cycle;}}
\newcommand{\skMantel}{\skizze{\fill[red, opacity=0.35] (-1,0) arc[start angle=180, end angle=360, x radius=1, y radius=0.3] -- (1,1.9)
	arc[start angle=0, end angle=-180, x radius=1, y radius=0.3] -- cycle; \zylinderbasis}}
\newcommand{\skPyramidenstumpf}{\skizze{\draw[fig, hilf, thick] (0,0) -- (0.8,0.6) -- (2.8,0.6) (0.8,0.6) -- (1.1,1.45);
	\draw[fig] (0,0) -- (2,0) -- (2.8,0.6) -- (2.1,1.45) -- (1.1,1.45) -- (0.7,1.15) -- cycle (0.7,1.15) -- (1.7,1.15) -- (2.1,1.45) (1.7,1.15) -- (2,0);}}
\newcommand{\skKegelstumpf}{\skizze{\draw[fig, hilf, thick] (1.1,0) arc[start angle=0, end angle=180, x radius=1.1, y radius=0.3];
	\draw[fig] (-1.1,0) arc[start angle=180, end angle=360, x radius=1.1, y radius=0.3];
	\draw[fig] (0,1.5) ellipse[x radius=0.5, y radius=0.15]; \draw[fig] (-1.1,0) -- (-0.5,1.5) (1.1,0) -- (0.5,1.5);}}
\newcommand{\skTetraeder}{\skizze{\draw[fig, hilf, thick] (0,0) -- (2.6,0.3);
	\draw[fig] (0,0) -- (1.6,-0.4) -- (2.6,0.3) -- (1.3,1.9) -- cycle (1.6,-0.4) -- (1.3,1.9);}}
\newcommand{\skHohlzylinder}{\skizze{\zylinderbasis \draw[fig] (0,1.9) ellipse[x radius=0.5, y radius=0.15];
	\draw[fig, hilf, thin] (-0.5,0) -- (-0.5,1.9) (0.5,0) -- (0.5,1.9);}}
\newcommand{\skSchieferZylinder}{\skizze{\draw[fig, hilf, thick] (1,0) arc[start angle=0, end angle=180, x radius=1, y radius=0.3];
	\draw[fig] (-1,0) arc[start angle=180, end angle=360, x radius=1, y radius=0.3];
	\draw[fig] (0.8,1.8) ellipse[x radius=1, y radius=0.3]; \draw[fig] (-1,0) -- (-0.2,1.8) (1,0) -- (1.8,1.8);}}
\newcommand{\drehachse}[2]{\draw[dash dot] (0,#1) -- (0,#2); \draw[mass, -{Stealth[length=3pt]}] (-0.3,#2-0.15) arc[start angle=180, end angle=500, x radius=0.3, y radius=0.1];}
\newcommand{\skRotation}{\skizze{\filldraw[fill=gray!50, draw=blue, very thick] (0,-1.1) arc[start angle=-90, end angle=90, radius=1.1] -- cycle;
	\draw[fig, hilf, thick] (0,1.1) arc[start angle=90, end angle=270, radius=1.1]; \drehachse{-1.4}{1.6}}}
\newcommand{\skKugelsegment}{\skizze{\fill[gray, opacity=0.5] (27.04:1.1) arc[start angle=27.04, end angle=152.96, radius=1.1] -- cycle;
	\draw[fig] (0,0) circle (1.1); \draw[fig] (-0.98,0.5) arc[start angle=180, end angle=360, x radius=0.98, y radius=0.2];
	\draw[fig, hilf, thick] (0.98,0.5) arc[start angle=0, end angle=180, x radius=0.98, y radius=0.2];}}
\newcommand{\skKugelsektor}{\skizze{\fill[gray, opacity=0.5] (0,0) -- (55:1.1) arc[start angle=55, end angle=125, radius=1.1] -- cycle;
	\draw[fig] (0,0) circle (1.1); \draw[fig, thick] (55:1.1) -- (0,0) -- (125:1.1);
	\draw[fig, thick] (-0.63,0.9) arc[start angle=180, end angle=360, x radius=0.63, y radius=0.13];}}
\newcommand{\skEllipsoid}{\skizze{\draw[fig] (0,0) ellipse[x radius=1.5, y radius=0.9];
	\draw[fig, hilf, thick] (1.5,0) arc[start angle=0, end angle=180, x radius=1.5, y radius=0.3];
	\draw[fig] (-1.5,0) arc[start angle=180, end angle=360, x radius=1.5, y radius=0.3];}}
\newcommand{\skParaboloid}{\skizze{\draw[fig, domain=-1.3:1.3, samples=30] plot (\x, {0.9*\x*\x});
	\draw[fig] (0,1.521) ellipse[x radius=1.3, y radius=0.22];}}
\newcommand{\skHyperboloid}{\skizze{\draw[fig, domain=-1.2:1.2, samples=30, variable=\t] plot ({0.55*cosh(\t)}, {0.9*sinh(\t)});
	\draw[fig, domain=-1.2:1.2, samples=30, variable=\t] plot ({-0.55*cosh(\t)}, {0.9*sinh(\t)});
	\draw[fig] (0,1.358) ellipse[x radius=0.996, y radius=0.2];
	\draw[fig] (-0.996,-1.358) arc[start angle=180, end angle=360, x radius=0.996, y radius=0.2];
	\draw[fig, hilf, thick] (0.996,-1.358) arc[start angle=0, end angle=180, x radius=0.996, y radius=0.2];}}

% Trigonometrie und analytische Geometrie
\newcommand{\skEinheitskreis}{\skizze{\achsen{-1.6}{1.9}{-1.5}{1.8}\draw[fig] (0,0) circle (1.1); \node[below right] at (1.05,0) {$1$};}}
\newcommand{\skBogenmass}{\skizze{\draw[fig] (1.6,0) -- (0,0) -- (60:1.6); \draw[red, very thick] (1.6,0) arc[start angle=0, end angle=60, radius=1.6];
	\node at (30:2.0) {$b$}; \node[below] at (0.8,0) {$r$}; \node at (30:0.6) {$\alpha$};}}
\newcommand{\skPolar}{\skizze{\achsen{-0.4}{3}{-0.4}{2.4}\draw[mass] (0.9,0) arc[start angle=0, end angle=36, radius=0.9];
	\draw[mass, very thick] (0,0) -- node[above left, inner sep=1pt] {$r$} (2.2,1.6); \fill (2.2,1.6) circle (2pt) node[above] {$P$};
	\node at (17:1.35) {$\varphi$};}}
\newcommand{\skKreisgleichung}{\skizze{\achsen{-1.7}{1.9}{-1.6}{1.8}\draw[fig] (0,0) circle (1.2); \draw[mass] (0,0) -- node[above left, inner sep=1pt] {$r$} (35:1.2);}}
\newcommand{\ellipsebasis}{\achsen{-2}{2.2}{-1.4}{1.6}\draw[fig] (0,0) ellipse[x radius=1.6, y radius=1];}
\newcommand{\skEllipse}{\skizze{\ellipsebasis}}
\newcommand{\skBrennpunkt}{\skizze{\ellipsebasis \fill[red] (-1.249,0) circle (2.5pt) node[below] {$F_1$} (1.249,0) circle (2.5pt) node[below] {$F_2$};}}
\newcommand{\skHalbachse}{\skizze{\ellipsebasis \draw[red, very thick] (0,0) -- node[below] {$a$} (1.6,0) (0,0) -- node[left] {$b$} (0,1);}}
\newcommand{\skHyperbel}{\skizze{\achsen{-2.2}{2.4}{-1.6}{1.8}
	\draw[fig, domain=-1.6:1.6, samples=30, variable=\t] plot ({0.7*cosh(\t)}, {0.6*sinh(\t)});
	\draw[fig, domain=-1.6:1.6, samples=30, variable=\t] plot ({-0.7*cosh(\t)}, {0.6*sinh(\t)});}}
\newcommand{\skLeitlinie}{\skizze{\achsen{-1.2}{2.4}{-1.7}{1.9}
	\draw[fig, domain=-1.5:1.5, samples=30, variable=\t] plot ({\t*\t/1.6}, {\t});
	\draw[mass, hilf] (-0.4,-1.6) -- (-0.4,1.7) node[left] {$l$}; \fill (0.4,0) circle (2pt) node[below right] {$F$};}}

%----------------------------------------------------------------------------------------------------%
% Teil B: Lesen und Sprechen (dokumentspezifisch)
%----------------------------------------------------------------------------------------------------%

% \lies{Ausdruck}{Japanisch}{Umschrift}{deutsche Lesart}   (bauen.py erzeugt Japanisch und Umschrift
% aus der Kurzschreibweise in den Quelldateien)
\makeatletter
\newcommand\lies{\begingroup\shorthandoff{"}\lies@aux}
\newcommand\lies@aux[4]{\endgroup #1 & \leavevmode #2\newline{\small\color{black!70}[#3]}\ifx\relax#4\relax\else\newline{\color{black!60}#4}\fi \\[1.5pt]}
\makeatother

% Zeilenabstand in den Listen; bauen.py setzt ihn für einzelne Kapitel enger (siehe ENG dort)
\newcommand{\lesenabstand}{1.12}
% Liste ohne Linien und Rahmen: links der Ausdruck, rechts Japanisch (schwarz), Umschrift (klein) und Deutsch (grau)
\newenvironment{lesen}{%
	\setstretch{\lesenabstand}% drei Zeilen je Eintrag (Japanisch, Umschrift, Deutsch): enger als der Fließtext
	\begin{longtable}{@{}>{\raggedright\arraybackslash}p{4.6cm} >{\raggedright\arraybackslash}p{12.1cm}@{}}}{%
	\end{longtable}\vspace{-1.2em}}

% Kurzer Hinweis im Fließtext
\newcommand{\anmerkung}[1]{\par\smallskip\noindent{\small\sloppy\textbf{Hinweis:} #1\par}\smallskip}

% Zahlenstrahl für Intervalle: \intervall{g|o}{g|o} (g = Randpunkt gehört dazu, o = gehört nicht dazu)
\newcommand{\randpunkt}[2]{\if#1g\fill[blue] (#2,0) circle (3.2pt);\else\filldraw[fill=white, draw=blue, very thick] (#2,0) circle (3.2pt);\fi}
\newcommand{\intervall}[2]{\tikz[baseline=-0.6ex, scale=0.6]{\draw[achse] (-0.3,0) -- (5.6,0);
	\draw[blue, line width=2pt] (1,0) -- (4,0); \randpunkt{#1}{1} \randpunkt{#2}{4}
	\node[below=2pt, font=\scriptsize] at (1,0) {$a$}; \node[below=2pt, font=\scriptsize] at (4,0) {$b$};}}
\newcommand{\intervallR}{\tikz[baseline=-0.6ex, scale=0.6]{\draw[achse] (-0.3,0) -- (5.6,0);
	\draw[blue, line width=2pt, -{Stealth}] (1,0) -- (5.3,0); \randpunkt{g}{1}
	\node[below=2pt, font=\scriptsize] at (1,0) {$a$};}}

% Rechengröße mit Namen (japanisch/deutsch) unter einer Klammer
\newcommand{\ub}[3]{\underbrace{\text{\Large$#1$}}_{\text{\footnotesize\shortstack{#2\\\color{black!60}#3}}}}

% Verhindert, dass eine Zwischenüberschrift allein am Seitenende steht
\newcommand{\platz}{\par\ifdim\dimexpr\pagegoal-\pagetotal\relax<5.5cm \newpage\fi}


\title{Mathematisches Wörterbuch\\[0.3em] \large Deutsch -- Japanisch mit Lesung, Umschrift, Zeichen und Skizzen}
\author{Martin Quandt}
\date{\today}
\pagestyle{plain}
\setlength{\tabcolsep}{4pt}


% Zweisprachige Überschriften: \teil, \kapitel, \zwischen{Deutsch}{Japanisch}{Umschrift}
% Ins Inhaltsverzeichnis und in die PDF-Lesezeichen kommt nur „Deutsch -- Japanisch“ ohne Umschrift.
\newcommand\teil[3]{\part[#1 -- #2]{\raggedright #1\\[5pt]{\normalfont\Large\color{black!55}#2\hspace{0.5em}\penalty0 [#3]}}}
\newcommand\kapitel[3]{\section[#1 -- #2]{\raggedright #1\\[0pt]{\normalfont\small\color{black!55}#2\hspace{0.5em}\penalty0 [#3]}}}
\newcommand\zwischen[3]{\subsection*{\raggedright #1 {\normalfont\small\color{black!55}\,--\, #2\hspace{0.5em}\penalty0 [#3]}}}
% Die zweite Überschriftenzeile kostet Platz: Zeilen in den Listen und Tabellen etwas enger setzen
\setlength{\aboverulesep}{0.15ex}\setlength{\belowrulesep}{0.3ex}\renewcommand{\arraystretch}{0.93}

\begin{document}
\begin{CJK}{UTF8}{ipxm}\CJKabsatzschutz

\maketitle

Dieses Wörterbuch enthält 353 mathematische Fachbegriffe aus der Sekundarstufe I und II,
geordnet nach Sachgebieten. Zu jedem Begriff stehen der japanische Fachausdruck, seine Lesung, die
Umschrift in lateinischen Buchstaben und -- wo es hilft -- das Zeichen, eine Formel oder eine kleine
Skizze. Teil~II zeigt, wie man Zeichen, Formeln und Rechenwege auf Japanisch vorliest, mit vielen
Beispielen. Am Ende finden Sie ein Register Deutsch -- Japanisch und ein Register Japanisch -- Deutsch.

\zwischen{Hinweise zu Schrift und Umschrift}{表記とローマ字}{hyōki to rōmaji}

Japanische Fachbegriffe bestehen meist aus Kanji, den chinesischen Schriftzeichen. Lehnwörter wie
\jp{ベクトル}{bekutoru}, „Vektor“, stehen in der Silbenschrift Katakana. Unter jedem Begriff steht klein
seine Lesung in der Silbenschrift Hiragana und in eckigen Klammern die Umschrift nach dem
Hepburn-System. Ist der Begriff schon in Kana geschrieben, steht dort nur die Umschrift.

Die Umschrift liest man fast wie Deutsch. Das Japanische betont keine Silbe durch Lautstärke: Alle
Silben sind etwa gleich lang und gleich stark. Die Partikeln は, へ und を spricht man \textit{wa},
\textit{e} und \textit{o}; so stehen sie auch in der Umschrift.

\medskip
{\small
\begin{tabular}{@{}>{\raggedright\arraybackslash}p{2.4cm} >{\raggedright\arraybackslash}p{9.2cm} >{\raggedright\arraybackslash}p{5.0cm}@{}}
	\toprule
	\textbf{Umschrift} & \textbf{Aussprache} & \textbf{Beispiel} \\ \midrule
	ā \,ī \,ū \,ē \,ō & langer Vokal, etwa doppelt so lang wie der kurze & \leavevmode 数 [sū], 方程式 [hōteishiki] \\
	u & ungerundet; zwischen stimmlosen Lauten und am Wortende kaum hörbar & \leavevmode 確率 [kakuritsu] \\
	r & kurzer Zungenschlag zwischen r, l und d & \leavevmode 比例 [hirei] \\
	s, z & s stimmlos wie in „Tasse“, z stimmhaft wie das s in „Sonne“ & \leavevmode 整数 [seisū], 座標 [zahyō] \\
	sh, ch & wie „sch“ und „tsch“ & \leavevmode 小数 [shōsū], 頂点 [chōten] \\
	j, ts & wie „dsch“ und wie das deutsche „z“ & \leavevmode 二乗 [nijō], 通分 [tsūbun] \\
	y, w & y wie das deutsche „j“, w wie im Englischen & \leavevmode 要素 [yōso], 和 [wa] \\
	kk, pp, ss, tt & kurzes Innehalten vor dem Konsonanten & \leavevmode 括弧 [kakko], 直角 [chokkaku] \\
	n' & das n schließt die Silbe ab, danach setzt man neu an & \leavevmode 単位 [tan'i] \\
	\bottomrule
\end{tabular}}

\medskip
\textbf{Farben in den Skizzen:} Figuren sind blau, hervorgehobene Teile (der Begriff selbst, Winkel,
Maße) rot, verdeckte Kanten und Hilfslinien gestrichelt.

\newpage
{\small\setlength{\parskip}{0pt}\tableofcontents}
\newpage

\teil{Wörterbuch}{辞書}{jisho}

%%%%  ZAHLEN UND RECHENARTEN  %%%%
\kapitel{Zahlen und Rechenarten}{数と計算}{kazu to keisan}\renewcommand{\lesenabstand}{1.12}\renewcommand{\arraystretch}{0.93}

\begin{woerter}
\eintrag{1}{die}{Zahl}{数}{かず}{kazu}{$1;\ 2;\ 3;\ \ldots$}
\eintrag{2}{die}{Ziffer}{数字}{すうじ}{sūji}{$0;\ 1;\ \ldots;\ 9$}
\eintrag{3}{die}{Ordnungszahl}{序数}{じょすう}{josū}{1.\ 2.\ 3. \newline \footnotesize jap. 第1、第2、第3}
\eintrag{4}{der}{Index}{添え字}{そえじ}{soeji}{$x_1;\ x_2;\ \ldots;\ x_n$}
\eintrag{5}{die}{natürliche Zahl}{自然数}{しぜんすう}{shizensū}{$\N$}
\eintrag{6}{die}{ganze Zahl}{整数}{せいすう}{seisū}{$\Z$}
\eintrag{7}{die}{rationale Zahl}{有理数}{ゆうりすう}{yūrisū}{$\Q$}
\eintrag{8}{die}{reelle Zahl}{実数}{じっすう}{jissū}{$\R$}
\eintrag{9}{die}{irrationale Zahl}{無理数}{むりすう}{murisū}{$\sqrt{2};\ \pi$}
\eintrag{10}{die}{komplexe Zahl}{複素数}{ふくそすう}{fukusosū}{$\C$ \newline $z = a + b\,\mathrm{i}$}
\eintrag{11}{die}{imaginäre Einheit}{虚数単位}{きょすうたんい}{kyosū tan'i}{$\mathrm{i}^2 = -1$}
\eintrag{12}{die}{Primzahl}{素数}{そすう}{sosū}{$2;\ 3;\ 5;\ 7;\ 11;\ \ldots$}
\eintrag{13}{die}{Dezimalzahl}{小数}{しょうすう}{shōsū}{$3{,}14$}
\eintrag{14}{}{addieren}{足す}{たす}{tasu}{$a + b$}
\eintrag{15}{die}{Summe}{和}{わ}{wa}{$\underbrace{a + b}_{\text{Summe}}$}
\eintrag{16}{der}{Summand}{加数}{かすう}{kasū}{$\underbrace{a}_{\text{Summand}} + \underbrace{b}_{\text{Summand}}$}
\eintrag{17}{}{subtrahieren}{引く}{ひく}{hiku}{$a - b$}
\eintrag{18}{die}{Differenz}{差}{さ}{sa}{$\underbrace{a - b}_{\text{Differenz}}$}
\eintrag{19}{der}{Minuend}{被減数}{ひげんすう}{higensū}{$\underbrace{a}_{\text{Minuend}} - \underbrace{b}_{\text{Subtrahend}}$}
\eintrag{20}{der}{Subtrahend}{減数}{げんすう}{gensū}{$a - \frot{b}$}
\eintrag{21}{}{multiplizieren}{掛ける}{かける}{kakeru}{$a \cdot b$}
\eintrag{22}{das}{Produkt}{積}{せき}{seki}{$\underbrace{a \cdot b}_{\text{Produkt}}$}
\eintrag{23}{der}{Faktor}{因数}{いんすう}{insū}{$\underbrace{a}_{\text{Faktor}} \cdot \underbrace{b}_{\text{Faktor}}$}
\eintrag{24}{}{dividieren}{割る}{わる}{waru}{$a : b$}
\eintrag{25}{der}{Quotient}{商}{しょう}{shō}{$\underbrace{a : b}_{\text{Quotient}}$}
\eintrag{26}{der}{Dividend}{被除数}{ひじょすう}{hijosū}{$\underbrace{a}_{\text{Dividend}} : \underbrace{b}_{\text{Divisor}}$}
\eintrag{27}{der}{Divisor}{除数}{じょすう}{josū}{$a : \frot{b}$}
\eintrag{28}{der}{Bruch}{分数}{ぶんすう}{bunsū}{$\dfrac{a}{b}$}
\eintrag{29}{der}{Zähler}{分子}{ぶんし}{bunshi}{$\dfrac{\frot{a}}{b}$}
\eintrag{30}{der}{Nenner}{分母}{ぶんぼ}{bunbo}{$\dfrac{a}{\frot{b}}$}
\eintrag{31}{der}{echte Bruch}{真分数}{しんぶんすう}{shinbunsū}{$\dfrac{3}{4}$}
\eintrag{32}{der}{unechte Bruch}{仮分数}{かぶんすう}{kabunsū}{$\dfrac{7}{4}$}
\eintrag{33}{die}{gemischte Zahl}{帯分数}{たいぶんすう}{taibunsū}{$1\tfrac{3}{4}$}
\eintrag{34}{}{gleichnamige Brüche}{同分母分数}{どうぶんぼぶんすう}{dōbunbo bunsū}{$\dfrac{1}{5};\ \dfrac{3}{5}$}
\eintrag{35}{}{ungleichnamige Brüche}{異分母分数}{いぶんぼぶんすう}{ibunbo bunsū}{$\dfrac{1}{2};\ \dfrac{1}{3}$}
\eintrag{36}{der}{Hauptnenner}{最小公分母}{さいしょうこうぶんぼ}{saishō kōbunbo}{$\dfrac{1}{2} + \dfrac{1}{3} = \dfrac{3}{\frot{6}} + \dfrac{2}{\frot{6}}$}
\eintrag{37}{}{kürzen}{約分する}{やくぶんする}{yakubun suru}{$\dfrac{6}{8} = \dfrac{3}{4}$}
\eintrag{38}{}{erweitern}{倍分する}{ばいぶんする}{baibun suru}{$\dfrac{3}{4} = \dfrac{15}{20}$}
\eintrag{39}{der}{Kehrwert}{逆数}{ぎゃくすう}{gyakusū}{$\dfrac{a}{b} \;\rightarrow\; \dfrac{b}{a}$}
\eintrag{40}{das}{Verhältnis}{比}{ひ}{hi}{$a : b$}
\eintrag{41}{die}{Verhältnisgleichung}{比例式}{ひれいしき}{hireishiki}{$a : b = c : d$}
\eintrag{42}{die}{Potenz}{累乗}{るいじょう}{ruijō}{$a^n$}
\eintrag{43}{die}{Basis}{底}{てい}{tei}{$\frot{a}^{n}$}
\eintrag{44}{der}{Exponent}{指数}{しすう}{shisū}{$a^{\frot{n}}$}
\eintrag{45}{die}{Quadratwurzel}{平方根}{へいほうこん}{heihōkon}{$\sqrt{a}$}
\eintrag{46}{die}{Kubikwurzel}{立方根}{りっぽうこん}{rippōkon}{$\sqrt[3]{a}$}
\eintrag{47}{die}{$n$-te Wurzel}{$n$ 乗根}{$n$ じょうこん}{$n$ jōkon}{$\sqrt[n]{a}$}
\eintrag{48}{der}{Radikand}{被開方数}{ひかいほうすう}{hikaihōsū}{$\sqrt[n]{\frot{a}}$}
\eintrag{49}{der}{Wurzelexponent}{根指数}{こんしすう}{konshisū}{$\sqrt[\frot{n}]{a}$}
\eintrag{50}{}{potenzieren}{累乗する}{るいじょうする}{ruijō suru}{$2^5 = 32$}
\eintrag{51}{}{quadrieren}{二乗する}{にじょうする}{nijō suru}{$7^2 = 49$}
\eintrag{52}{}{radizieren, Wurzel ziehen}{累乗根を求める}{るいじょうこんをもとめる}{ruijōkon o motomeru}{$\sqrt{49} = 7$}
\eintrag{53}{der}{Betrag}{絶対値}{ぜったいち}{zettaichi}{$|a|$}
\eintrag{54}{der}{Teiler}{約数}{やくすう}{yakusū}{$3 \mid 12$}
\eintrag{55}{das}{Vielfache}{倍数}{ばいすう}{baisū}{$3;\ 6;\ 9;\ 12;\ \ldots$}
\eintrag{56}{der}{größte gemeinsame Teiler}{最大公約数}{さいだいこうやくすう}{saidai kōyakusū}{$\operatorname{ggT}(12;18) = 6$}
\eintrag{57}{das}{kleinste gemeinsame Vielfache}{最小公倍数}{さいしょうこうばいすう}{saishō kōbaisū}{$\operatorname{kgV}(4;6) = 12$}
\eintrag{58}{der}{Prozentsatz}{百分率}{ひゃくぶんりつ}{hyakubunritsu}{$p\,\%$}
\eintrag{59}{}{runden}{四捨五入する}{ししゃごにゅうする}{shishagonyū suru}{$3{,}14 \approx 3{,}1$}
\eintrag{60}{}{gleich}{等しい}{ひとしい}{hitoshii}{$=$}
\eintrag{61}{}{ungleich}{等しくない}{ひとしくない}{hitoshikunai}{$\neq$}
\eintrag{62}{}{identisch gleich}{恒等的に等しい}{こうとうてきにひとしい}{kōtōteki ni hitoshii}{$\equiv$}
\eintrag{63}{}{kleiner als}{より小さい}{よりちいさい}{yori chiisai}{$<$}
\eintrag{64}{}{größer als}{より大きい}{よりおおきい}{yori ōkii}{$>$}
\eintrag{65}{}{kleiner oder gleich}{以下}{いか}{ika}{$\le$ \newline \footnotesize jap. $\leqq$}
\eintrag{66}{}{größer oder gleich}{以上}{いじょう}{ijō}{$\ge$ \newline \footnotesize jap. $\geqq$}
\eintrag{67}{}{viel kleiner als}{はるかに小さい}{はるかにちいさい}{haruka ni chiisai}{$\ll$}
\eintrag{68}{}{proportional zu}{に比例する}{にひれいする}{ni hirei suru}{$y \sim x$ \newline \footnotesize jap. $y \propto x$}
\eintrag{69}{}{ungefähr}{ほぼ等しい}{ほぼひとしい}{hobo hitoshii}{$\approx$ \newline \footnotesize jap. $\fallingdotseq$}
\eintrag{70}{}{unendlich}{無限大}{むげんだい}{mugendai}{$\infty$}
\end{woerter}

\platz\zwischen{Beispiele}{例}{rei}

\begin{lesen}
\lies{$12 + 30 = 42$}{12と30の和は42である。}{jūni to sanjū no wa wa yonjūni de aru.}{Die Summe von zwölf und dreißig ist zweiundvierzig.}
\lies{$3 \mid 12$}{3は12の約数である。12は3の倍数である。12は3で割り切れる。}{san wa jūni no yakusū de aru. jūni wa san no baisū de aru. jūni wa san de warikireru.}{Drei ist ein Teiler von zwölf; zwölf ist ein Vielfaches von drei; zwölf ist durch drei teilbar.}
\lies{$\operatorname{ggT}(12;\,18) = 6$}{12と18の最大公約数は6である。}{jūni to jūhachi no saidai kōyakusū wa roku de aru.}{Der größte gemeinsame Teiler von zwölf und achtzehn ist sechs.}
\lies{$17$}{17は素数である。約数は1とその数自身の2つだけである。}{jūnana wa sosū de aru. yakusū wa ichi to sono kazu jishin no futatsu dake de aru.}{Siebzehn ist eine Primzahl: Sie hat genau zwei Teiler, die Eins und sich selbst.}
\lies{$3{,}14159 \approx 3{,}14$}{小数第3位を四捨五入して、小数第2位まで求めなさい。}{shōsū dai san i o shishagonyū shite, shōsū dai ni i made motomenasai.}{Runde die Zahl auf zwei Nachkommastellen.}
\lies{$|-7| = 7$}{マイナス7の絶対値は7である。}{mainasu nana no zettaichi wa nana de aru.}{Der Betrag von minus sieben ist sieben.}
\lies{$25\,\%$ von $80$ sind $20$}{80の25パーセントは20である。}{hachijū no nijūgo pāsento wa nijū de aru.}{Fünfundzwanzig Prozent von achtzig sind zwanzig.}
\end{lesen}


%%%%  ALGEBRA, GLEICHUNGEN UND MENGEN  %%%%
\newpage
\kapitel{Algebra, Gleichungen und Mengen}{代数、方程式、集合}{daisū, hōteishiki, shūgō}\renewcommand{\lesenabstand}{1.12}\renewcommand{\arraystretch}{0.93}

\begin{woerter}
\eintrag{71}{die}{Variable}{変数}{へんすう}{hensū}{$x;\ y;\ a$}
\eintrag{72}{der}{Term}{式}{しき}{shiki}{$3x + 2$}
\eintrag{73}{die}{Unbekannte}{未知数}{みちすう}{michisū}{$x$}
\eintrag{74}{der}{Koeffizient}{係数}{けいすう}{keisū}{$\frot{3}x^2$}
\eintrag{75}{der}{Grad}{次数}{じすう}{jisū}{$x^3 + x$: 3.\,Grad}
\eintrag{76}{die}{Gleichung}{方程式}{ほうていしき}{hōteishiki}{$2x + 3 = 7$}
\eintrag{77}{die}{Ungleichung}{不等式}{ふとうしき}{futōshiki}{$2x + 3 < 7$}
\eintrag{78}{die}{Lösung}{解}{かい}{kai}{$x = 2$}
\eintrag{79}{die}{Lösungsmenge}{解の集合}{かいのしゅうごう}{kai no shūgō}{$L = \{2\}$}
\eintrag{80}{das}{Gleichungssystem}{連立方程式}{れんりつほうていしき}{renritsu hōteishiki}{$\begin{aligned} x + y &= 5\\ x - y &= 1\end{aligned}$}
\eintrag{81}{die}{lineare Gleichung}{一次方程式}{いちじほうていしき}{ichiji hōteishiki}{$3x - 7 = 11$}
\eintrag{82}{die}{quadratische Gleichung}{二次方程式}{にじほうていしき}{niji hōteishiki}{$x^2 - 5x + 6 = 0$}
\eintrag{83}{die}{Wurzelgleichung}{無理方程式}{むりほうていしき}{muri hōteishiki}{$\sqrt{x + 3} = 5$}
\eintrag{84}{die}{Exponentialgleichung}{指数方程式}{しすうほうていしき}{shisū hōteishiki}{$2^x = 64$}
\eintrag{85}{die}{logarithmische Gleichung}{対数方程式}{たいすうほうていしき}{taisū hōteishiki}{$\log_3 x = 4$}
\eintrag{86}{die}{trigonometrische Gleichung}{三角方程式}{さんかくほうていしき}{sankaku hōteishiki}{$2\sin x = 1$}
\eintrag{87}{die}{Klammer}{括弧}{かっこ}{kakko}{$(\ \ )\quad [\ \ ]\quad \{\ \ \}$}
\eintrag{88}{die}{runden Klammern}{丸括弧}{まるかっこ}{marukakko}{$(\ \ )$}
\eintrag{89}{die}{eckigen Klammern}{角括弧}{かくかっこ}{kakukakko}{$[\ \ ]$}
\eintrag{90}{die}{geschweiften Klammern}{波括弧}{なみかっこ}{namikakko}{$\{\ \ \}$}
\eintrag{91}{}{ausmultiplizieren}{展開する}{てんかいする}{tenkai suru}{$a(b + c) = ab + ac$}
\eintrag{92}{}{ausklammern}{共通因数でくくる}{きょうつういんすうでくくる}{kyōtsū insū de kukuru}{$ab + ac = a(b + c)$}
\eintrag{93}{die}{binomische Formel}{乗法公式}{じょうほうこうしき}{jōhō kōshiki}{$(a + b)^2 = a^2 + 2ab + b^2$}
\eintrag{94}{die}{Proportionalität}{比例}{ひれい}{hirei}{$y = k \cdot x$}
\eintrag{95}{der}{Logarithmus}{対数}{たいすう}{taisū}{$\log_b a$}
\eintrag{96}{die}{Menge}{集合}{しゅうごう}{shūgō}{$M = \{1;\ 2;\ 3\}$}
\eintrag{97}{das}{Element}{要素}{ようそ}{yōso}{$2 \in M$}
\eintrag{98}{die}{Teilmenge}{部分集合}{ぶぶんしゅうごう}{bubun shūgō}{$A \subseteq B$ \newline \skTeilmenge}
\eintrag{99}{die}{echte Teilmenge}{真部分集合}{しんぶぶんしゅうごう}{shin bubun shūgō}{$A \subset B$ \newline \footnotesize jap. $A \subsetneqq B$}
\eintrag{100}{die}{Schnittmenge}{共通部分}{きょうつうぶぶん}{kyōtsū bubun}{$A \cap B$ \newline \skSchnittmenge}
\eintrag{101}{die}{Vereinigungsmenge}{和集合}{わしゅうごう}{washūgō}{$A \cup B$ \newline \skVereinigung}
\eintrag{102}{die}{Differenzmenge}{差集合}{さしゅうごう}{sashūgō}{$A \setminus B$}
\eintrag{103}{das}{kartesische Produkt}{直積}{ちょくせき}{chokuseki}{$A \times B$}
\eintrag{104}{die}{leere Menge}{空集合}{くうしゅうごう}{kūshūgō}{$\emptyset = \{\,\}$}
\end{woerter}

\platz\zwischen{Beispiele}{例}{rei}

\begin{lesen}
\lies{$3(x + 2) = 3x + 6$}{かっこをはずして展開しなさい。}{kakko o hazushite tenkai shinasai.}{Multipliziere die Klammer aus.}
\lies{$6x + 9 = 3(2x + 3)$}{共通因数3でくくりなさい。}{kyōtsū insū san de kukurinasai.}{Klammere aus: Der gemeinsame Faktor ist drei.}
\lies{$(x + 3)^2 = x^2 + 6x + 9$}{かっこ $x$ プラス3かっことじの2乗イコール $x$ 2乗プラス6$x$ プラス9。}{kakko $x$ purasu san kakkotoji no nijō ikōru $x$ nijō purasu roku $x$ purasu kyū.}{$x$ plus drei in Klammern zum Quadrat ist gleich $x$ Quadrat plus sechs $x$ plus neun.}
\lies{$2x + 3 = 11$ \newline $x = 4$}{方程式を $x$ について解きなさい。解は $x$ イコール4である。}{hōteishiki o $x$ ni tsuite tokinasai. kai wa $x$ ikōru yon de aru.}{Löse die Gleichung nach $x$ auf. Die Lösung ist $x$ gleich vier.}
\lies{$x + y = 5$ \newline $x - y = 1$}{連立方程式を解きなさい。解は $x$ イコール3、$y$ イコール2である。}{renritsu hōteishiki o tokinasai. kai wa $x$ ikōru san, $y$ ikōru ni de aru.}{Löse das Gleichungssystem: $x$ gleich drei und $y$ gleich zwei.}
\lies{$A = \{1;\,2;\,3\},\ B = \{3;\,4\}$ \newline $A \cap B = \{3\}$}{$A$ と $B$ の共通部分の要素は3だけである。}{$A$ to $B$ no kyōtsū bubun no yōso wa san dake de aru.}{Die Schnittmenge von $A$ und $B$ enthält nur das Element drei.}
\end{lesen}


%%%%  FUNKTIONEN  %%%%
\newpage
\kapitel{Funktionen}{関数}{kansū}\renewcommand{\lesenabstand}{1.12}\renewcommand{\arraystretch}{0.93}

\begin{woerter}
\eintrag{105}{die}{Funktion}{関数}{かんすう}{kansū}{$f(x)$}
\eintrag{106}{die}{unabhängige Variable}{独立変数}{どくりつへんすう}{dokuritsu hensū}{$x$}
\eintrag{107}{die}{abhängige Variable}{従属変数}{じゅうぞくへんすう}{jūzoku hensū}{$y = f(x)$}
\eintrag{108}{der}{Funktionswert}{関数の値}{かんすうのあたい}{kansū no atai}{$f(3) = 15$}
\eintrag{109}{der}{Definitionsbereich}{定義域}{ていぎいき}{teigiiki}{$D_f$}
\eintrag{110}{der}{Wertebereich}{値域}{ちいき}{chiiki}{$W_f$}
\eintrag{111}{der}{Funktionsgraph}{グラフ}{}{gurafu}{\skGraph}
\eintrag{112}{die}{Wertetabelle}{対応表}{たいおうひょう}{taiōhyō}{\skWertetabelle}
\eintrag{113}{das}{Koordinatensystem}{座標系}{ざひょうけい}{zahyōkei}{\skKoordinaten}
\eintrag{114}{die}{$x$-Achse}{$x$ 軸}{$x$ じく}{$x$ jiku}{\skXAchse}
\eintrag{115}{der}{Ursprung}{原点}{げんてん}{genten}{$O(0 \mid 0)$}
\eintrag{116}{die}{Koordinaten}{座標}{ざひょう}{zahyō}{$P(x_P \mid y_P)$ \newline \footnotesize jap. $P(x_P,\ y_P)$}
\eintrag{117}{die}{Abszisse}{$x$ 座標}{$x$ ざひょう}{$x$ zahyō}{$P(\frot{x_P} \mid y_P)$}
\eintrag{118}{die}{Ordinate}{$y$ 座標}{$y$ ざひょう}{$y$ zahyō}{$P(x_P \mid \frot{y_P})$}
\eintrag{119}{die}{lineare Funktion}{一次関数}{いちじかんすう}{ichiji kansū}{$f(x) = mx + n$ \newline \skLinear}
\eintrag{120}{die}{Steigung}{傾き}{かたむき}{katamuki}{$m = \dfrac{\Delta y}{\Delta x}$ \newline \skSteigung}
\eintrag{121}{der}{$y$-Achsenabschnitt}{$y$ 切片}{$y$ せっぺん}{$y$ seppen}{\skAchsenabschnitt}
\eintrag{122}{die}{Nullstelle}{零点}{れいてん}{reiten}{$f(x_0) = 0$ \newline \skNullstelle}
\eintrag{123}{der}{Schnittpunkt}{交点}{こうてん}{kōten}{\skSchnittpunkt}
\eintrag{124}{die}{quadratische Funktion}{二次関数}{にじかんすう}{niji kansū}{$f(x) = ax^2 + bx + c$}
\eintrag{125}{die}{Parabel}{放物線}{ほうぶつせん}{hōbutsusen}{\skParabel}
\eintrag{126}{die}{ganzrationale Funktion}{多項式関数}{たこうしきかんすう}{takōshiki kansū}{$f(x) = x^3 - 2x$}
\eintrag{127}{die}{gebrochenrationale Funktion}{分数関数}{ぶんすうかんすう}{bunsū kansū}{$f(x) = \dfrac{x + 1}{x - 2}$}
\eintrag{128}{die}{Wurzelfunktion}{無理関数}{むりかんすう}{muri kansū}{$f(x) = \sqrt{x}$}
\eintrag{129}{der}{Scheitelpunkt}{頂点}{ちょうてん}{chōten}{\skScheitelpunkt}
\eintrag{130}{die}{Exponentialfunktion}{指数関数}{しすうかんすう}{shisū kansū}{$f(x) = a \cdot b^x$ \newline \skExponential}
\eintrag{131}{die}{Logarithmusfunktion}{対数関数}{たいすうかんすう}{taisū kansū}{$f(x) = \log_b x$ \newline \skLogarithmus}
\eintrag{132}{der}{Sinus}{正弦、サイン}{せいげん、サイン}{seigen, sain}{$\sin x$ \newline \skSinus}
\eintrag{133}{der}{Kosinus}{余弦、コサイン}{よげん、コサイン}{yogen, kosain}{$\cos x$ \newline \skKosinus}
\eintrag{134}{der}{Tangens}{正接、タンジェント}{せいせつ、タンジェント}{seisetsu, tanjento}{$\tan x = \dfrac{\sin x}{\cos x}$}
\eintrag{135}{der}{Kotangens}{余接、コタンジェント}{よせつ、コタンジェント}{yosetsu, kotanjento}{$\cot x = \dfrac{\cos x}{\sin x}$}
\eintrag{136}{die}{Arkusfunktion}{逆三角関数}{ぎゃくさんかくかんすう}{gyaku sankaku kansū}{$\arcsin x$ \newline \footnotesize jap. meist $\sin^{-1} x$}
\eintrag{137}{der}{Einheitskreis}{単位円}{たんいえん}{tan'ien}{\skEinheitskreis}
\eintrag{138}{das}{Gradmaß}{度数法}{どすうほう}{dosūhō}{$360\degree$}
\eintrag{139}{das}{Bogenmaß}{弧度法}{こどほう}{kodohō}{$180\degree \mathrel{\widehat{=}} \pi$ \newline \skBogenmass}
\eintrag{140}{der}{Sinussatz}{正弦定理}{せいげんていり}{seigen teiri}{$\frac{a}{\sin\alpha} = \frac{b}{\sin\beta} = \frac{c}{\sin\gamma}$}
\eintrag{141}{der}{Kosinussatz}{余弦定理}{よげんていり}{yogen teiri}{\small$c^2 = a^2 + b^2 - 2ab\cos\gamma$}
\eintrag{142}{}{monoton steigend}{単調増加}{たんちょうぞうか}{tanchō zōka}{$x_1 < x_2 \Rightarrow f(x_1) \le f(x_2)$ \newline \skMonotonSteigend}
\eintrag{143}{}{monoton fallend}{単調減少}{たんちょうげんしょう}{tanchō genshō}{$x_1 < x_2 \Rightarrow f(x_1) \ge f(x_2)$ \newline \skMonotonFallend}
\eintrag{144}{}{achsensymmetrisch zur $y$-Achse}{$y$ 軸に関して対称}{$y$ じくにかんしてたいしょう}{$y$ jiku ni kanshite taishō}{$f(-x) = f(x)$ \newline \skAchsensymmetrisch}
\eintrag{145}{}{punktsymmetrisch zum Ursprung}{原点に関して対称}{げんてんにかんしてたいしょう}{genten ni kanshite taishō}{$f(-x) = -f(x)$ \newline \skPunktsymmetrisch}
\end{woerter}

\platz\zwischen{Beispiele}{例}{rei}

\begin{lesen}
\lies{$f(x) = 2x - 4$}{$f$ は傾き2、$y$ 切片マイナス4の一次関数である。零点は $x$ イコール2である。}{$f$ wa katamuki ni, $y$ seppen mainasu yon no ichiji kansū de aru. reiten wa $x$ ikōru ni de aru.}{$f$ ist eine lineare Funktion mit der Steigung zwei und dem $y$-Achsenabschnitt minus vier. Ihre Nullstelle ist $x$ gleich zwei.}
\lies{$f(x) = (x - 1)^2 - 4$}{グラフは頂点が $S(1,\,-4)$ の放物線である。零点はマイナス1と3である。}{gurafu wa chōten ga $S(1,\,-4)$ no hōbutsusen de aru. reiten wa mainasu ichi to san de aru.}{Der Graph ist eine Parabel mit dem Scheitelpunkt $S(1 \mid -4)$. Die Nullstellen sind minus eins und drei.}
\lies{$f(x) = x^2$}{グラフは $y$ 軸に関して対称である。関数は $-\infty < x \le 0$ で減少し、$0 \le x < \infty$ で増加する。}{gurafu wa $y$ jiku ni kanshite taishō de aru. kansū wa $-\infty < x \le 0$ de genshō shi, $0 \le x < \infty$ de zōka suru.}{Der Graph ist achsensymmetrisch zur $y$-Achse. Die Funktion ist fallend für $-\infty < x \le 0$ und steigend für $0 \le x < \infty$.}
\lies{$f(x) = 2^x$}{この指数関数のグラフは点 $(0,\,1)$ を通り、$x$ 軸が漸近線である。}{kono shisū kansū no gurafu wa ten $(0,\,1)$ o tōri, $x$ jiku ga zenkinsen de aru.}{Der Graph dieser Exponentialfunktion verläuft durch den Punkt $(0 \mid 1)$; die $x$-Achse ist Asymptote.}
\lies{$D_f = \R$ \newline $W_f = \{y \in \R \mid y \ge -4\}$}{定義域は実数全体、値域はマイナス4以上の実数全体である。}{teigiiki wa jissū zentai, chiiki wa mainasu yon ijō no jissū zentai de aru.}{Der Definitionsbereich ist die Menge aller reellen Zahlen; der Wertebereich ist die Menge aller reellen Zahlen, die größer oder gleich minus vier sind.}
\end{lesen}


%%%%  GEOMETRIE: GRUNDBEGRIFFE, WINKEL, ABBILDUNGEN  %%%%
\newpage
\kapitel{Geometrie: Grundbegriffe, Winkel, Abbildungen}{幾何：基本用語、角、移動}{kika: kihon yōgo, kaku, idō}\renewcommand{\lesenabstand}{1.12}\renewcommand{\arraystretch}{0.93}

\begin{woerter}
\eintrag{146}{der}{Punkt}{点}{てん}{ten}{\skPunkt}
\eintrag{147}{die}{Gerade}{直線}{ちょくせん}{chokusen}{\skGerade}
\eintrag{148}{die}{Strecke}{線分}{せんぶん}{senbun}{$\overline{AB}$ \newline \skStrecke}
\eintrag{149}{der}{Strahl}{半直線}{はんちょくせん}{hanchokusen}{\skStrahl}
\eintrag{150}{}{parallel}{平行}{へいこう}{heikō}{$g \parallel h$ \newline \skParallel}
\eintrag{151}{}{senkrecht}{垂直}{すいちょく}{suichoku}{$g \perp h$ \newline \skSenkrecht}
\eintrag{152}{der}{Abstand}{距離}{きょり}{kyori}{\skAbstand}
\eintrag{153}{der}{Winkel}{角}{かく}{kaku}{$\alpha;\ \beta;\ \gamma$ \newline \skWinkel{50}}
\eintrag{154}{der}{Scheitel}{頂点}{ちょうてん}{chōten}{\skScheitel}
\eintrag{155}{der}{Schenkel}{辺}{へん}{hen}{\skSchenkel}
\eintrag{156}{die}{Winkelhalbierende}{角の二等分線}{かくのにとうぶんせん}{kaku no nitōbunsen}{\skWinkelhalbierende}
\eintrag{157}{der}{spitze Winkel}{鋭角}{えいかく}{eikaku}{$0\degree < \alpha < 90\degree$ \newline \skWinkel{50}}
\eintrag{158}{der}{rechte Winkel}{直角}{ちょっかく}{chokkaku}{$\alpha = 90\degree$ \newline \skRechterWinkel}
\eintrag{159}{der}{stumpfe Winkel}{鈍角}{どんかく}{donkaku}{$90\degree < \alpha < 180\degree$ \newline \skWinkel{130}}
\eintrag{160}{der}{gestreckte Winkel}{平角}{へいかく}{heikaku}{$\alpha = 180\degree$ \newline \skWinkel{180}}
\eintrag{161}{der}{überstumpfe Winkel}{優角}{ゆうかく}{yūkaku}{$180\degree < \alpha < 360\degree$ \newline \skWinkel{250}}
\eintrag{162}{der}{Vollwinkel}{周角}{しゅうかく}{shūkaku}{$\alpha = 360\degree$ \newline \skVollwinkel}
\eintrag{163}{der}{Nebenwinkel}{隣接補角}{りんせつほかく}{rinsetsu hokaku}{$\alpha + \beta = 180\degree$ \newline \skNebenwinkel}
\eintrag{164}{der}{Komplementwinkel}{余角}{よかく}{yokaku}{$\alpha + \beta = 90\degree$ \newline \skKomplement}
\eintrag{165}{der}{Supplementwinkel}{補角}{ほかく}{hokaku}{$\alpha + \beta = 180\degree$}
\eintrag{166}{der}{Scheitelwinkel}{対頂角}{たいちょうかく}{taichōkaku}{$\alpha = \alpha'$ \newline \skScheitelwinkel}
\eintrag{167}{der}{Stufenwinkel}{同位角}{どういかく}{dōikaku}{$\alpha = \beta$ \newline \skStufenwinkel}
\eintrag{168}{der}{Wechselwinkel}{錯角}{さっかく}{sakkaku}{$\alpha = \beta$ \newline \skWechselwinkel}
\eintrag{169}{}{entgegengesetzt liegende Winkel}{同側内角}{どうそくないかく}{dōsoku naikaku}{$\alpha + \beta = 180\degree$ \newline \skEntgegen}
\eintrag{170}{die}{Mittelsenkrechte}{垂直二等分線}{すいちょくにとうぶんせん}{suichoku nitōbunsen}{\skMittelsenkrechte}
\eintrag{171}{die}{Symmetrieachse}{対称の軸}{たいしょうのじく}{taishō no jiku}{\skSymmetrieachse}
\eintrag{172}{die}{Spiegelung}{対称移動}{たいしょういどう}{taishō idō}{\skSpiegelung}
\eintrag{173}{die}{Verschiebung}{平行移動}{へいこういどう}{heikō idō}{\skVerschiebung}
\eintrag{174}{die}{Drehung}{回転移動}{かいてんいどう}{kaiten idō}{\skDrehung}
\eintrag{175}{}{kongruent}{合同}{ごうどう}{gōdō}{$\triangle ABC \cong \triangle A'B'C'$ \newline \footnotesize jap. $\triangle ABC \equiv \triangle A'B'C'$}
\eintrag{176}{}{ähnlich}{相似}{そうじ}{sōji}{$\triangle ABC \sim \triangle A'B'C'$ \newline \footnotesize jap. $\triangle ABC \backsim \triangle A'B'C'$}
\end{woerter}

\platz\zwischen{Beispiele}{例}{rei}

\begin{lesen}
\lies{$\overline{AB} = 5\,\mathrm{cm}$}{線分 $AB$ の長さは5センチメートルである。}{senbun $AB$ no nagasa wa go senchimētoru de aru.}{Die Strecke $A$ $B$ ist fünf Zentimeter lang.}
\lies{$g \parallel h$}{点 $P$ を通り、直線 $g$ に平行な直線 $h$ をひきなさい。}{ten $P$ o tōri, chokusen $g$ ni heikō na chokusen $h$ o hikinasai.}{Zeichne durch den Punkt $P$ die Parallele $h$ zu $g$.}
\lies{$\alpha = 35\degree$}{アルファは35度の鋭角である。}{arufa wa sanjūgo do no eikaku de aru.}{Alpha ist ein spitzer Winkel von fünfunddreißig Grad.}
\lies{$\beta = 180\degree - 70\degree = 110\degree$}{隣接補角の和は180度である。}{rinsetsu hokaku no wa wa hyakuhachijū do de aru.}{Nebenwinkel ergänzen sich zu hundertachtzig Grad.}
\lies{\skSpiegelung}{三角形を、直線 $g$ に関して対称移動しなさい。}{sankakukei o, chokusen $g$ ni kanshite taishō idō shinasai.}{Spiegele das Dreieck an der Geraden $g$.}
\lies{\skDrehung}{三角形を、点 $Z$ を中心に90度回転移動しなさい。}{sankakukei o, ten $Z$ o chūshin ni kyūjū do kaiten idō shinasai.}{Drehe das Dreieck um neunzig Grad um den Punkt $Z$.}
\lies{\skMittelsenkrechte}{線分 $AB$ の垂直二等分線を作図しなさい。}{senbun $AB$ no suichoku nitōbunsen o sakuzu shinasai.}{Konstruiere die Mittelsenkrechte der Strecke $A$ $B$.}
\end{lesen}


%%%%  EBENE FIGUREN  %%%%
\newpage
\kapitel{Ebene Figuren}{平面図形}{heimen zukei}\renewcommand{\lesenabstand}{1.12}\renewcommand{\arraystretch}{0.93}

\begin{woerter}
\eintrag{177}{das}{Dreieck}{三角形}{さんかくけい}{sankakukei}{$\triangle ABC$ \newline \skDreieck}
\eintrag{178}{der}{Innenwinkel}{内角}{ないかく}{naikaku}{$\alpha + \beta + \gamma = 180\degree$ \newline \skInnenwinkel}
\eintrag{179}{der}{Außenwinkel}{外角}{がいかく}{gaikaku}{\skAussenwinkel}
\eintrag{180}{das}{unregelmäßige Dreieck}{不等辺三角形}{ふとうへんさんかくけい}{futōhen sankakukei}{\skDreieck}
\eintrag{181}{das}{rechtwinklige Dreieck}{直角三角形}{ちょっかくさんかくけい}{chokkaku sankakukei}{\skRechtwinklig}
\eintrag{182}{das}{gleichschenklige Dreieck}{二等辺三角形}{にとうへんさんかくけい}{nitōhen sankakukei}{\skGleichschenklig}
\eintrag{183}{die}{Basis (im Dreieck)}{底辺}{ていへん}{teihen}{\skBasis}
\eintrag{184}{der}{Basiswinkel}{底角}{ていかく}{teikaku}{\skBasiswinkel}
\eintrag{185}{das}{gleichseitige Dreieck}{正三角形}{せいさんかくけい}{seisankakukei}{\skGleichseitig}
\eintrag{186}{die}{Hypotenuse}{斜辺}{しゃへん}{shahen}{\skHypotenuse}
\eintrag{187}{die}{Kathete}{直角をはさむ辺}{ちょっかくをはさむへん}{chokkaku o hasamu hen}{\skKathete}
\eintrag{188}{die}{Ankathete}{隣辺}{りんぺん}{rinpen}{\skAnkathete}
\eintrag{189}{die}{Gegenkathete}{対辺}{たいへん}{taihen}{\skGegenkathete}
\eintrag{190}{der}{Satz des Pythagoras}{三平方の定理}{さんへいほうのていり}{sanheihō no teiri}{$a^2 + b^2 = c^2$ \newline \skPythagoras}
\eintrag{191}{die}{Höhe}{高さ}{たかさ}{takasa}{\skHoehe}
\eintrag{192}{die}{Seitenhalbierende}{中線}{ちゅうせん}{chūsen}{\skSeitenhalbierende}
\eintrag{193}{der}{Inkreis}{内接円}{ないせつえん}{naisetsuen}{\skInkreis}
\eintrag{194}{der}{Umkreis}{外接円}{がいせつえん}{gaisetsuen}{\skUmkreis}
\eintrag{195}{}{flächengleich}{面積が等しい}{めんせきがひとしい}{menseki ga hitoshii}{$A_1 = A_2$}
\eintrag{196}{das}{Viereck}{四角形}{しかくけい}{shikakukei}{\skViereck}
\eintrag{197}{die}{Diagonale}{対角線}{たいかくせん}{taikakusen}{\skDiagonale}
\eintrag{198}{das}{Quadrat}{正方形}{せいほうけい}{seihōkei}{\skQuadrat}
\eintrag{199}{das}{Rechteck}{長方形}{ちょうほうけい}{chōhōkei}{\skRechteck}
\eintrag{200}{das}{Parallelogramm}{平行四辺形}{へいこうしへんけい}{heikō shihenkei}{\skParallelogramm}
\eintrag{201}{die}{Raute}{ひし形}{ひしがた}{hishigata}{\skRaute}
\eintrag{202}{das}{Trapez}{台形}{だいけい}{daikei}{\skTrapez}
\eintrag{203}{das}{Drachenviereck}{たこ形}{たこがた}{takogata}{\skDrachen}
\eintrag{204}{das}{Vieleck}{多角形}{たかくけい}{takakukei}{\skFuenfeck}
\eintrag{205}{das}{regelmäßige Sechseck}{正六角形}{せいろっかくけい}{seirokkakukei}{\skSechseck}
\eintrag{206}{das}{Fünfeck}{五角形}{ごかくけい}{gokakukei}{\skNeck{5}}
\eintrag{207}{das}{Achteck}{八角形}{はっかくけい}{hakkakukei}{\skNeck{8}}
\eintrag{208}{das}{Zwölfeck}{十二角形}{じゅうにかくけい}{jūnikakukei}{\skNeck{12}}
\eintrag{209}{der}{Umfang}{周の長さ}{しゅうのながさ}{shū no nagasa}{$u = 2a + 2b$ \newline \skUmfangViereck}
\eintrag{210}{der}{Flächeninhalt}{面積}{めんせき}{menseki}{$A = a \cdot b$ \newline \skFlaeche}
\eintrag{211}{der}{Kreis}{円}{えん}{en}{\skKreis}
\eintrag{212}{der}{Mittelpunkt}{中心}{ちゅうしん}{chūshin}{\skMittelpunkt}
\eintrag{213}{der}{Radius}{半径}{はんけい}{hankei}{\skRadius}
\eintrag{214}{der}{Durchmesser}{直径}{ちょっけい}{chokkei}{$d = 2r$ \newline \skDurchmesser}
\eintrag{215}{der}{Kreisumfang}{円周}{えんしゅう}{enshū}{$u = 2\pi r$ \newline \skKreisumfang}
\eintrag{216}{die}{Kreiszahl}{円周率}{えんしゅうりつ}{enshūritsu}{$\pi \approx 3{,}14159$}
\eintrag{217}{der}{Kreissektor}{扇形}{おうぎがた}{ōgigata}{\skSektor}
\eintrag{218}{die}{Sehne}{弦}{げん}{gen}{\skSehne}
\eintrag{219}{die}{Sekante}{割線}{かっせん}{kassen}{\skSekante}
\eintrag{220}{der}{Berührungspunkt}{接点}{せってん}{setten}{\skBeruehrpunkt}
\eintrag{221}{der}{Kreisbogen}{弧}{こ}{ko}{\skKreisbogen}
\eintrag{222}{der}{Kreisabschnitt}{弓形}{きゅうけい}{kyūkei}{\skKreisabschnitt}
\eintrag{223}{die}{Kreisfläche}{円の面積}{えんのめんせき}{en no menseki}{$A = \pi r^2$ \newline \skKreisflaeche}
\eintrag{224}{die}{Tangente}{接線}{せっせん}{sessen}{\skTangente}
\end{woerter}

\platz\zwischen{Beispiele}{例}{rei}

\begin{lesen}
\lies{$a = 6\,\mathrm{cm},\ b = 4\,\mathrm{cm}$ \newline $A = 24\,\mathrm{cm}^2$}{長方形の面積は、6センチメートル掛ける4センチメートルで、24平方センチメートルである。}{chōhōkei no menseki wa, roku senchimētoru kakeru yon senchimētoru de, nijūyon heihō senchimētoru de aru.}{Der Flächeninhalt des Rechtecks ist sechs Zentimeter mal vier Zentimeter, also vierundzwanzig Quadratzentimeter.}
\lies{$3^2 + 4^2 = 5^2$}{3辺の長さが3、4、5の三角形は、直角三角形である。}{sanpen no nagasa ga san, yon, go no sankakukei wa, chokkaku sankakukei de aru.}{Ein Dreieck mit den Seitenlängen drei, vier und fünf ist rechtwinklig.}
\lies{\skGleichschenklig}{二等辺三角形では、2辺が等しく、2つの底角も等しい。}{nitōhen sankakukei de wa, nihen ga hitoshiku, futatsu no teikaku mo hitoshii.}{Ein gleichschenkliges Dreieck hat zwei gleich lange Seiten und zwei gleich große Basiswinkel.}
\lies{\skDiagonale}{長方形の対角線は長さが等しく、それぞれの中点で交わる。}{chōhōkei no taikakusen wa nagasa ga hitoshiku, sorezore no chūten de majiwaru.}{Die Diagonalen eines Rechtecks sind gleich lang und halbieren einander.}
\lies{$d = 10\,\mathrm{cm}$ \newline $u = \pi \cdot d \approx 31{,}4\,\mathrm{cm}$}{直径10センチメートルの円の円周は、約31.4センチメートルである。}{chokkei jū senchimētoru no en no enshū wa, yaku sanjūitten yon senchimētoru de aru.}{Der Umfang eines Kreises mit zehn Zentimetern Durchmesser beträgt ungefähr einunddreißig Komma vier Zentimeter.}
\end{lesen}


%%%%  KÖRPER  %%%%
\newpage
\kapitel{Körper}{立体}{rittai}\renewcommand{\lesenabstand}{1.12}\renewcommand{\arraystretch}{0.93}

\begin{woerter}
\eintrag{225}{der}{Körper}{立体}{りったい}{rittai}{}
\eintrag{226}{der}{Würfel}{立方体}{りっぽうたい}{rippōtai}{\skWuerfel}
\eintrag{227}{der}{Quader}{直方体}{ちょくほうたい}{chokuhōtai}{\skQuader}
\eintrag{228}{das}{Prisma}{角柱}{かくちゅう}{kakuchū}{\skPrisma}
\eintrag{229}{der}{Zylinder}{円柱}{えんちゅう}{enchū}{\skZylinder}
\eintrag{230}{die}{Pyramide}{角錐}{かくすい}{kakusui}{\skPyramide}
\eintrag{231}{der}{Kegel}{円錐}{えんすい}{ensui}{\skKegel}
\eintrag{232}{die}{Kugel}{球}{きゅう}{kyū}{\skKugel}
\eintrag{233}{das}{Polyeder}{多面体}{ためんたい}{tamentai}{}
\eintrag{234}{das}{Tetraeder}{四面体}{しめんたい}{shimentai}{\skTetraeder}
\eintrag{235}{der}{Pyramidenstumpf}{角錐台}{かくすいだい}{kakusuidai}{\skPyramidenstumpf}
\eintrag{236}{der}{Kegelstumpf}{円錐台}{えんすいだい}{ensuidai}{\skKegelstumpf}
\eintrag{237}{der}{Hohlzylinder}{中空円柱}{ちゅうくうえんちゅう}{chūkū enchū}{\skHohlzylinder}
\eintrag{238}{der}{schiefe Zylinder}{斜円柱}{しゃえんちゅう}{shaenchū}{\skSchieferZylinder}
\eintrag{239}{der}{Rotationskörper}{回転体}{かいてんたい}{kaitentai}{\skRotation}
\eintrag{240}{das}{Kugelsegment}{球欠}{きゅうけつ}{kyūketsu}{\skKugelsegment}
\eintrag{241}{der}{Kugelsektor}{球扇形}{きゅうせんけい}{kyūsenkei}{\skKugelsektor}
\eintrag{242}{das}{Rotationsellipsoid}{回転楕円体}{かいてんだえんたい}{kaiten daentai}{\skEllipsoid}
\eintrag{243}{das}{Rotationsparaboloid}{回転放物面}{かいてんほうぶつめん}{kaiten hōbutsumen}{\skParaboloid}
\eintrag{244}{das}{Rotationshyperboloid}{回転双曲面}{かいてんそうきょくめん}{kaiten sōkyokumen}{\skHyperboloid}
\eintrag{245}{die}{Ecke}{頂点}{ちょうてん}{chōten}{\skEcke}
\eintrag{246}{die}{Kante}{辺}{へん}{hen}{\skKante}
\eintrag{247}{die}{Seitenfläche}{面}{めん}{men}{\skSeitenflaeche}
\eintrag{248}{die}{Grundfläche}{底面}{ていめん}{teimen}{\skGrundflaeche}
\eintrag{249}{die}{Deckfläche}{上面}{じょうめん}{jōmen}{\skDeckflaeche}
\eintrag{250}{die}{Mantelfläche}{側面}{そくめん}{sokumen}{$M = 2\pi r h$ \newline \skMantel}
\eintrag{251}{die}{Flächendiagonale}{面の対角線}{めんのたいかくせん}{men no taikakusen}{\skFlaechendiagonale}
\eintrag{252}{die}{Raumdiagonale}{立体の対角線}{りったいのたいかくせん}{rittai no taikakusen}{$d = \sqrt{a^2 + b^2 + c^2}$ \newline \skRaumdiagonale}
\eintrag{253}{das}{Netz}{展開図}{てんかいず}{tenkaizu}{\skNetz}
\eintrag{254}{das}{Schrägbild}{見取図}{みとりず}{mitorizu}{\skQuader}
\eintrag{255}{das}{Volumen}{体積}{たいせき}{taiseki}{$V = a \cdot b \cdot c$}
\eintrag{256}{der}{Oberflächeninhalt}{表面積}{ひょうめんせき}{hyōmenseki}{$O = 2(ab + ac + bc)$}
\end{woerter}

\platz\zwischen{Beispiele}{例}{rei}

\begin{lesen}
\lies{$a = 3\,\mathrm{cm}$ \newline $V = a^3 = 27\,\mathrm{cm}^3$}{1辺が3センチメートルの立方体の体積は、27立方センチメートルである。}{ippen ga san senchimētoru no rippōtai no taiseki wa, nijūnana rippō senchimētoru de aru.}{Das Volumen eines Würfels mit drei Zentimetern Kantenlänge beträgt siebenundzwanzig Kubikzentimeter.}
\lies{$O = 6a^2 = 54\,\mathrm{cm}^2$}{表面積は54平方センチメートルである。}{hyōmenseki wa gojūyon heihō senchimētoru de aru.}{Sein Oberflächeninhalt beträgt vierundfünfzig Quadratzentimeter.}
\lies{\skNetz}{立方体の展開図は、6つの正方形からなる。}{rippōtai no tenkaizu wa, muttsu no seihōkei kara naru.}{Das Netz eines Würfels besteht aus sechs Quadraten.}
\lies{\skPyramide}{正四角錐には、頂点が5つ、辺が8本、面が5つある。}{seishikakusui ni wa, chōten ga itsutsu, hen ga happon, men ga itsutsu aru.}{Eine quadratische Pyramide hat fünf Ecken, acht Kanten und fünf Flächen.}
\lies{$r = 2\,\mathrm{cm},\ h = 5\,\mathrm{cm}$ \newline $V = \pi r^2 h \approx 62{,}8\,\mathrm{cm}^3$}{円柱の体積は20パイで、約62.8立方センチメートルである。}{enchū no taiseki wa nijū pai de, yaku rokujūni ten hachi rippō senchimētoru de aru.}{Das Volumen des Zylinders ist zwanzig pi, ungefähr zweiundsechzig Komma acht Kubikzentimeter.}
\end{lesen}


%%%%  ANALYSIS  %%%%
\newpage
\kapitel{Analysis}{微分積分}{bibun sekibun}\renewcommand{\lesenabstand}{1.12}\renewcommand{\arraystretch}{0.93}

\begin{woerter}
\eintrag{257}{die}{Folge}{数列}{すうれつ}{sūretsu}{$(a_n)$ \newline \skFolge}
\eintrag{258}{die}{Nullfolge}{零に収束する数列}{れいにしゅうそくするすうれつ}{rei ni shūsoku suru sūretsu}{$\limn a_n = 0$}
\eintrag{259}{die}{arithmetische Folge}{等差数列}{とうさすうれつ}{tōsa sūretsu}{$a_{n+1} = a_n + d$}
\eintrag{260}{die}{geometrische Folge}{等比数列}{とうひすうれつ}{tōhi sūretsu}{$a_{n+1} = a_n \cdot q$}
\eintrag{261}{das}{Anfangsglied}{初項}{しょこう}{shokō}{$a_1$}
\eintrag{262}{}{konvergent}{収束する}{しゅうそくする}{shūsoku suru}{}
\eintrag{263}{}{divergent}{発散する}{はっさんする}{hassan suru}{}
\eintrag{264}{der}{Grenzwert}{極限値}{きょくげんち}{kyokugenchi}{$\limn a_n = g$}
\eintrag{265}{die}{Reihe}{級数}{きゅうすう}{kyūsū}{$\displaystyle\sum_{k=1}^{n} a_k$}
\eintrag{266}{die}{unendliche Reihe}{無限級数}{むげんきゅうすう}{mugen kyūsū}{$\displaystyle\sum_{k=1}^{\infty} a_k$}
\eintrag{267}{die}{Potenzreihe}{べき級数}{べききゅうすう}{beki kyūsū}{$\displaystyle\sum_{k=0}^{\infty} a_k x^k$}
\eintrag{268}{}{stetig}{連続}{れんぞく}{renzoku}{$\limxto{x_0} f(x) = f(x_0)$}
\eintrag{269}{}{unstetig}{不連続}{ふれんぞく}{furenzoku}{}
\eintrag{270}{das}{Intervall}{区間}{くかん}{kukan}{$[a;\,b]$}
\eintrag{271}{die}{Asymptote}{漸近線}{ぜんきんせん}{zenkinsen}{\skAsymptote}
\eintrag{272}{der}{Differenzenquotient}{平均変化率}{へいきんへんかりつ}{heikin henkaritsu}{$\dfrac{f(x_0 + h) - f(x_0)}{h}$}
\eintrag{273}{die}{Ableitung}{導関数}{どうかんすう}{dōkansū}{$f'(x)$}
\eintrag{274}{die}{Differenzialrechnung}{微分法}{びぶんほう}{bibunhō}{}
\eintrag{275}{der}{Differenzialquotient}{微分係数}{びぶんけいすう}{bibun keisū}{$\dfrac{\mathrm{d}y}{\mathrm{d}x}$}
\eintrag{276}{das}{Differenzial}{微分}{びぶん}{bibun}{$\mathrm{d}y$}
\eintrag{277}{die}{partielle Ableitung}{偏導関数}{へんどうかんすう}{hendōkansū}{$\dfrac{\partial z}{\partial x}$}
\eintrag{278}{}{ableiten}{微分する}{びぶんする}{bibun suru}{$\dfrac{\mathrm{d}f}{\mathrm{d}x}$}
\eintrag{279}{die}{Tangente}{接線}{せっせん}{sessen}{$m_t = f'(x_0)$ \newline \skTangenteGraph}
\eintrag{280}{die}{Funktionsuntersuchung}{関数の増減を調べる}{かんすうのぞうげんをしらべる}{kansū no zōgen o shiraberu}{}
\eintrag{281}{der}{Extrempunkt}{極値をとる点}{きょくちをとるてん}{kyokuchi o toru ten}{$f'(x_E) = 0$}
\eintrag{282}{der}{Hochpunkt}{極大点}{きょくだいてん}{kyokudaiten}{$f''(x_E) < 0$ \newline \skHochpunkt}
\eintrag{283}{der}{Tiefpunkt}{極小点}{きょくしょうてん}{kyokushōten}{$f''(x_E) > 0$ \newline \skTiefpunkt}
\eintrag{284}{der}{Wendepunkt}{変曲点}{へんきょくてん}{henkyokuten}{$f''(x_W) = 0$ \newline \skWendepunkt}
\eintrag{285}{die}{Stammfunktion}{原始関数}{げんしかんすう}{genshi kansū}{$F'(x) = f(x)$}
\eintrag{286}{die}{Integralrechnung}{積分法}{せきぶんほう}{sekibunhō}{}
\eintrag{287}{das}{unbestimmte Integral}{不定積分}{ふていせきぶん}{futei sekibun}{$\displaystyle\int f(x)\dx$}
\eintrag{288}{das}{bestimmte Integral}{定積分}{ていせきぶん}{teisekibun}{$\displaystyle\int_a^b f(x)\dx$}
\eintrag{289}{das}{Integralzeichen}{積分記号}{せきぶんきごう}{sekibun kigō}{$\displaystyle\int$}
\eintrag{290}{der}{Integrand}{被積分関数}{ひせきぶんかんすう}{hisekibun kansū}{$\displaystyle\int \frot{f(x)}\dx$}
\eintrag{291}{die}{Integrationsvariable}{積分変数}{せきぶんへんすう}{sekibun hensū}{$\displaystyle\int f(\frot{x})\,\mathrm{d}\frot{x}$}
\eintrag{292}{die}{Integrationskonstante}{積分定数}{せきぶんていすう}{sekibun teisū}{$F(x) + \frot{C}$}
\eintrag{293}{die}{Obersumme}{過剰和}{かじょうわ}{kajōwa}{\skObersumme}
\eintrag{294}{die}{Untersumme}{不足和}{ふそくわ}{fusokuwa}{\skUntersumme}
\eintrag{295}{das}{Integral}{積分}{せきぶん}{sekibun}{$\displaystyle\int_a^b f(x)\dx$ \newline \skIntegral}
\eintrag{296}{}{integrieren}{積分する}{せきぶんする}{sekibun suru}{$\displaystyle\int f(x)\dx = F(x) + C$}
\eintrag{297}{die}{Integrationsgrenze}{上端、下端}{じょうたん、かたん}{jōtan, katan}{$\displaystyle\int_{\frot{a}}^{\frot{b}} f(x)\dx$}
\eintrag{298}{die}{eulersche Zahl}{ネイピア数}{ネイピアすう}{neipiasū}{$e \approx 2{,}71828$}
\eintrag{299}{die}{natürliche Exponentialfunktion}{自然指数関数}{しぜんしすうかんすう}{shizen shisū kansū}{$f(x) = e^x$}
\eintrag{300}{der}{natürliche Logarithmus}{自然対数}{しぜんたいすう}{shizen taisū}{$\ln x$}
\end{woerter}

\platz\zwischen{Beispiele}{例}{rei}

\begin{lesen}
\lies{$f(x) = x^3 - 3x$ \newline $f'(x) = 3x^2 - 3$}{$f(x)$ を微分しなさい。導関数は3$x$ 2乗マイナス3である。}{$f(x)$ o bibun shinasai. dōkansū wa san $x$ nijō mainasu san de aru.}{Leite $f$ ab: Die Ableitung ist drei $x$ Quadrat minus drei.}
\lies{$f'(x) = 0 \ \Rightarrow\ x = \pm 1$}{導関数は、$x$ イコールプラスマイナス1で0になる。}{dōkansū wa, $x$ ikōru purasu mainasu ichi de zero ni naru.}{Die Ableitung ist null für $x$ gleich plus oder minus eins.}
\lies{$H(-1 \mid 2),\ T(1 \mid -2)$ \newline $W(0 \mid 0)$}{グラフは点 $(-1,\,2)$ で極大、点 $(1,\,-2)$ で極小となり、原点が変曲点である。}{gurafu wa ten $(-1,\,2)$ de kyokudai, ten $(1,\,-2)$ de kyokushō to nari, genten ga henkyokuten de aru.}{Der Graph hat den Hochpunkt $H$, den Tiefpunkt $T$ und den Wendepunkt $W$ im Ursprung.}
\lies{$F(x) = x^2$ \newline $f(x) = 2x$}{$F$ は $f$ の原始関数である。}{$F$ wa $f$ no genshi kansū de aru.}{Groß $F$ ist eine Stammfunktion von $f$.}
\end{lesen}


%%%%  VEKTOREN UND ANALYTISCHE GEOMETRIE  %%%%
\newpage
\kapitel{Vektoren und analytische Geometrie}{ベクトルと解析幾何}{bekutoru to kaiseki kika}\renewcommand{\lesenabstand}{1.0}\renewcommand{\arraystretch}{0.80}

\begin{woerter}
\eintrag{301}{der}{Vektor}{ベクトル}{}{bekutoru}{$\vec{a} = \vektor{a_1}{a_2}{a_3}$ \newline \skVektor}
\eintrag{302}{der}{Ortsvektor}{位置ベクトル}{いちベクトル}{ichi bekutoru}{$\overrightarrow{OP}$ \newline \skOrtsvektor}
\eintrag{303}{der}{Einheitsvektor}{単位ベクトル}{たんいベクトル}{tan'i bekutoru}{$|\vec{e}\,| = 1$}
\eintrag{304}{der}{Nullvektor}{零ベクトル}{れいベクトル}{rei bekutoru}{$\vec{0}$}
\eintrag{305}{der}{Betrag eines Vektors}{ベクトルの大きさ}{ベクトルのおおきさ}{bekutoru no ōkisa}{$|\vec{a}| = \sqrt{a_1^2 + a_2^2 + a_3^2}$}
\eintrag{306}{das}{Skalarprodukt}{内積}{ないせき}{naiseki}{$\vec{a} \cdot \vec{b}$}
\eintrag{307}{das}{Kreuzprodukt}{外積}{がいせき}{gaiseki}{$\vec{a} \times \vec{b}$ \newline \skKreuzprodukt}
\eintrag{308}{die}{Ebene}{平面}{へいめん}{heimen}{\skEbene}
\eintrag{309}{die}{Geradengleichung}{直線の方程式}{ちょくせんのほうていしき}{chokusen no hōteishiki}{$y = mx + n$}
\eintrag{310}{der}{Steigungswinkel}{傾斜角}{けいしゃかく}{keishakaku}{$m = \tan\alpha$}
\eintrag{311}{die}{Polarkoordinaten}{極座標}{きょくざひょう}{kyokuzahyō}{$P(r;\,\varphi)$ \newline \skPolar}
\eintrag{312}{der}{Kegelschnitt}{円錐曲線}{えんすいきょくせん}{ensui kyokusen}{}
\eintrag{313}{die}{Kreisgleichung}{円の方程式}{えんのほうていしき}{en no hōteishiki}{$x^2 + y^2 = r^2$ \newline \skKreisgleichung}
\eintrag{314}{die}{Ellipse}{楕円}{だえん}{daen}{$\frac{x^2}{a^2} + \frac{y^2}{b^2} = 1$ \newline \skEllipse}
\eintrag{315}{die}{Hyperbel}{双曲線}{そうきょくせん}{sōkyokusen}{$\frac{x^2}{a^2} - \frac{y^2}{b^2} = 1$ \newline \skHyperbel}
\eintrag{316}{der}{Brennpunkt}{焦点}{しょうてん}{shōten}{\skBrennpunkt}
\eintrag{317}{die}{Halbachse}{半軸}{はんじく}{hanjiku}{\skHalbachse}
\eintrag{318}{die}{Leitlinie}{準線}{じゅんせん}{junsen}{\skLeitlinie}
\eintrag{319}{}{linear abhängig}{一次従属}{いちじじゅうぞく}{ichiji jūzoku}{$\vec{b} = r \cdot \vec{a}$}
\eintrag{320}{die}{Matrix}{行列}{ぎょうれつ}{gyōretsu}{$\begin{pmatrix} a & b\\ c & d\end{pmatrix}$}
\eintrag{321}{die}{quadratische Matrix}{正方行列}{せいほうぎょうれつ}{seihō gyōretsu}{$n \times n$}
\eintrag{322}{die}{Zeile}{行}{ぎょう}{gyō}{}
\eintrag{323}{die}{Spalte}{列}{れつ}{retsu}{}
\eintrag{324}{die}{Hauptdiagonale}{主対角線}{しゅたいかくせん}{shutaikakusen}{$\begin{pmatrix} \frot{a} & b\\ c & \frot{d}\end{pmatrix}$}
\eintrag{325}{die}{Determinante}{行列式}{ぎょうれつしき}{gyōretsushiki}{$\begin{vmatrix} a & b\\ c & d\end{vmatrix} = ad - bc$}
\end{woerter}

\platz\zwischen{Beispiele}{例}{rei}

\begin{lesen}
\lies{$\vec{a} = \vektor{3}{4}{0}$ \newline $|\vec{a}\,| = \sqrt{9 + 16 + 0} = 5$}{ベクトル $a$ の大きさは5である。}{bekutoru $a$ no ōkisa wa go de aru.}{Der Betrag des Vektors $a$ ist fünf.}
\lies{$A(1 \mid 2 \mid 3),\ B(4 \mid 6 \mid 3)$ \newline $\overrightarrow{AB} = \vektor{3}{4}{0}$}{点 $A$ から点 $B$ へのベクトルの成分は、3、4、0である。}{ten $A$ kara ten $B$ e no bekutoru no seibun wa, san, yon, zero de aru.}{Der Vektor von $A$ nach $B$ ist drei, vier, null.}
\lies{$g\colon\ \vec{x} = \vektor{1}{2}{3} + t \cdot \vektor{3}{4}{0}$}{直線 $g$：位置ベクトルプラス $t$ 掛ける方向ベクトル}{chokusen $g$: ichi bekutoru purasu $t$ kakeru hōkō bekutoru}{die Gerade $g$: Stützvektor plus $t$ mal Richtungsvektor}
\end{lesen}


%%%%  STOCHASTIK  %%%%
\newpage
\kapitel{Stochastik}{確率と統計}{kakuritsu to tōkei}\renewcommand{\lesenabstand}{1.12}\renewcommand{\arraystretch}{0.93}

\begin{woerter}
\eintrag{326}{das}{Zufallsexperiment}{試行}{しこう}{shikō}{}
\eintrag{327}{das}{Ergebnis}{結果}{けっか}{kekka}{$\omega \in \Omega$}
\eintrag{328}{das}{Ereignis}{事象}{じしょう}{jishō}{$A \subseteq \Omega$}
\eintrag{329}{das}{Gegenereignis}{余事象}{よじしょう}{yojishō}{$\overline{A}$}
\eintrag{330}{die}{Wahrscheinlichkeit}{確率}{かくりつ}{kakuritsu}{$P(A)$}
\eintrag{331}{die}{relative Häufigkeit}{相対度数}{そうたいどすう}{sōtai dosū}{$h_n(A)$}
\eintrag{332}{das}{Baumdiagramm}{樹形図}{じゅけいず}{jukeizu}{\skBaum}
\eintrag{333}{die}{Vierfeldertafel}{分割表}{ぶんかつひょう}{bunkatsuhyō}{\skVierfeldertafel}
\eintrag{334}{die}{Zufallsgröße}{確率変数}{かくりつへんすう}{kakuritsu hensū}{$X$}
\eintrag{335}{der}{Erwartungswert}{期待値}{きたいち}{kitaichi}{$\E(X) = \mu$}
\eintrag{336}{die}{Varianz}{分散}{ぶんさん}{bunsan}{$\V(X)$}
\eintrag{337}{die}{Standardabweichung}{標準偏差}{ひょうじゅんへんさ}{hyōjun hensa}{$\sigma = \sqrt{\V(X)}$}
\eintrag{338}{der}{Mittelwert}{平均値}{へいきんち}{heikinchi}{$\bar{x}$}
\eintrag{339}{der}{Median}{中央値}{ちゅうおうち}{chūōchi}{$\tilde{x}$}
\eintrag{340}{die}{Fakultät}{階乗}{かいじょう}{kaijō}{$n! = 1 \cdot 2 \cdot \ldots \cdot n$}
\eintrag{341}{die}{Permutation}{順列}{じゅんれつ}{junretsu}{$n!$}
\eintrag{342}{die}{Variation}{$k$ 個を並べる順列}{$k$ こをならべるじゅんれつ}{$k$ ko o naraberu junretsu}{$\dfrac{n!}{(n-k)!}$ \newline \footnotesize jap. ${}_n\mathrm{P}_k$}
\eintrag{343}{die}{Kombination}{組合せ}{くみあわせ}{kumiawase}{$\dbinom{n}{k}$ \newline \footnotesize jap. ${}_n\mathrm{C}_k$}
\eintrag{344}{der}{Binomialkoeffizient}{二項係数}{にこうけいすう}{nikō keisū}{$\dbinom{n}{k}$}
\eintrag{345}{die}{Binomialverteilung}{二項分布}{にこうぶんぷ}{nikō bunpu}{$B_{n;p}(k)$ \newline \skSaeulen}
\end{woerter}

\platz\zwischen{Beispiele}{例}{rei}

\begin{lesen}
\lies{$P(\text{„6“}) = \dfrac16$}{正しく作られたさいころを投げて6の目が出る確率は、6分の1である。}{tadashiku tsukurareta saikoro o nagete roku no me ga deru kakuritsu wa, roku bun no ichi de aru.}{Die Wahrscheinlichkeit, mit einem fairen Würfel eine Sechs zu würfeln, beträgt ein Sechstel.}
\lies{$P(\text{„zweimal Kopf“}) = \dfrac12 \cdot \dfrac12 = \dfrac14$}{硬貨を2回投げる。2回とも表が出る確率は、4分の1である。}{kōka o nikai nageru. nikai tomo omote ga deru kakuritsu wa, yon bun no ichi de aru.}{Eine faire Münze wird zweimal geworfen. Die Wahrscheinlichkeit für zweimal Kopf beträgt ein Viertel.}
\lies{$P(\overline{A}) = 1 - P(A)$}{余事象の確率は、1引く $A$ の確率に等しい。}{yojishō no kakuritsu wa, ichi hiku $A$ no kakuritsu ni hitoshii.}{Die Wahrscheinlichkeit des Gegenereignisses ist eins minus die Wahrscheinlichkeit von $A$.}
\lies{$2;\ 4;\ 6;\ 8$ \newline $\bar{x} = 5$}{4つの値の平均値は5である。}{yottsu no atai no heikinchi wa go de aru.}{Der Mittelwert der vier Werte ist fünf.}
\lies{$\E(X) = 3{,}5$}{さいころの出る目の期待値は3.5である。}{saikoro no deru me no kitaichi wa san ten go de aru.}{Der Erwartungswert der Augenzahl eines fairen Würfels ist drei Komma fünf.}
\end{lesen}


%%%%  LOGIK UND BEWEIS  %%%%
\newpage
\kapitel{Logik und Beweis}{論理と証明}{ronri to shōmei}\renewcommand{\lesenabstand}{1.12}\renewcommand{\arraystretch}{0.93}

\begin{woerter}
\eintrag{346}{die}{Definition}{定義}{ていぎ}{teigi}{}
\eintrag{347}{der}{Satz}{定理}{ていり}{teiri}{}
\eintrag{348}{die}{Voraussetzung}{仮定}{かてい}{katei}{}
\eintrag{349}{der}{Beweis}{証明}{しょうめい}{shōmei}{$\square$}
\eintrag{350}{}{daraus folgt}{ならば}{}{naraba}{$A \Rightarrow B$}
\eintrag{351}{}{genau dann, wenn}{同値}{どうち}{dōchi}{$A \Leftrightarrow B$}
\eintrag{352}{}{für alle}{すべての}{}{subete no}{$\forall$}
\eintrag{353}{}{es gibt}{存在する}{そんざいする}{sonzai suru}{$\exists$}
\end{woerter}

\platz\zwischen{Beispiele}{例}{rei}

\begin{lesen}
\lies{$x > 2 \ \Rightarrow\ x^2 > 4$}{$x$ が2より大きいならば、$x$ 2乗は4より大きい。}{$x$ ga ni yori ōkii naraba, $x$ nijō wa yon yori ōkii.}{Wenn $x$ größer als zwei ist, dann ist $x$ Quadrat größer als vier.}
\lies{$n$ gerade $\Leftrightarrow$ $n^2$ gerade}{$n$ が偶数であることと、$n$ 2乗が偶数であることは同値である。}{$n$ ga gūsū de aru koto to, $n$ nijō ga gūsū de aru koto wa dōchi de aru.}{$n$ ist genau dann gerade, wenn $n$ Quadrat gerade ist.}
\lies{$\forall\, x \in \R\colon\ x^2 \ge 0$}{すべての実数 $x$ について、$x$ 2乗は0以上である。}{subete no jissū $x$ ni tsuite, $x$ nijō wa zero ijō de aru.}{Für alle reellen Zahlen $x$ gilt: $x$ Quadrat ist größer oder gleich null.}
\lies{$\exists\, x \in \R\colon\ x^2 = 2$}{$x$ 2乗イコール2となる実数 $x$ が存在する。}{$x$ nijō ikōru ni to naru jissū $x$ ga sonzai suru.}{Es gibt eine reelle Zahl $x$, für die $x$ Quadrat gleich zwei ist.}
\lies{Annahme}{ルート2が有理数であると仮定する。\ldots\ これは矛盾である。（背理法）}{rūto ni ga yūrisū de aru to katei suru. \ldots\ kore wa mujun de aru. (hairihō)}{Angenommen, die Wurzel aus zwei ist rational. \ldots\ Das ist ein Widerspruch. (Beweis durch Widerspruch)}
\end{lesen}



% Teil II: Mathematik lesen und sprechen.
% \lies{Ausdruck}{Japanisch in Kurzschreibweise}{Japanisch in Kurzschreibweise}{deutsche Lesart}
% Die Kurzschreibweise (Wort|Lesung, durch Leerzeichen getrennt) steht in umschrift.py; bauen.py macht
% daraus \lies{Ausdruck}{Japanisch}{Japanisch}{Umschrift}{deutsche Lesart}.
%%%%  TEIL B: MATHEMATIK LESEN UND SPRECHEN  %%%%
\newpage
\teil{Mathematik lesen und sprechen}{数学の読み方と話し方}{sūgaku no yomikata to hanashikata}

In diesem Teil steht, wie man Zeichen, Formeln und Rechenwege auf Japanisch (und zum Vergleich auf
Deutsch) vorliest. Links steht der Ausdruck, rechts die japanische Lesart, darunter klein ihre
Umschrift und grau die deutsche Lesart. Wo es mehrere übliche Lesarten gibt, sind sie durch
Schrägstriche getrennt. Die Beispiele sind bewusst einfach gehalten, damit sie sich direkt im
Unterricht verwenden lassen.

\smallskip
Variablen wie $x$ bleiben in der Umschrift stehen; gesprochen werden sie mit ihren englischen
Buchstabennamen (\textit{ekkusu}, \textit{wai}). Terme und Gleichungen liest man im Unterricht meist
mit den Lehnwörtern \jp{プラス}{purasu}, \jp{マイナス}{mainasu} und \jp{イコール}{ikōru}. Beim Rechnen
mit Zahlen sind auch die japanischen Wörter \jp{足す}{tasu}, \jp{引く}{hiku}, \jp{掛ける}{kakeru} und
\jp{割る}{waru} üblich; „gleich“ heißt dann einfach \jp{は}{wa}.

%--------------------------------------------------------------------------------------------------
\kapitel{Zahlen, Zeichen und Schreibweisen}{数、記号、表記}{kazu, kigō, hyōki}\renewcommand{\lesenabstand}{1.12}\renewcommand{\arraystretch}{0.93}

\platz\zwischen{Ordnungszahlen, Indizes und Aufzählungen}{序数、添え字、列挙}{josū, soeji, rekkyo}

\begin{lesen}
\lies{1.\quad 2.\quad 3.\quad 21.}{第1、第2、第3、第21 \;/\; 1番目、2番目、3番目}{dai ichi, dai ni, dai san, dai nijūichi \;/\; ichibanme, nibanme, sanbanme}{erstens, zweitens, drittens -- der Erste, der Zweite, \ldots}
\lies{1.\ \ldots, 2.\ \ldots}{第一に \ldots、第二に \ldots}{daiichi ni \ldots, daini ni \ldots}{erstens \ldots, zweitens \ldots}
\lies{$x_1,\ x_2$}{$x$ 1、$x$ 2}{$x$ ichi, $x$ ni}{$x$ eins, $x$ zwei \;/\; $x$ Index eins}
\lies{$a_n$}{$a$ $n$（添え字は $n$）}{$a$ $n$ (soeji wa $n$)}{$a$ Index $n$ \;/\; $a$ $n$}
\lies{$1;\ 2;\ 3;\ \ldots$}{1、2、3、\ldots\ と限りなく続く}{ichi, ni, san, \ldots\ to kagirinaku tsuzuku}{eins, zwei, drei und so weiter (unbegrenzt)}
\lies{$1;\ 2;\ \ldots;\ 25$}{1、2、\ldots、25まで}{ichi, ni, \ldots, nijūgo made}{eins, zwei und so weiter bis fünfundzwanzig}
\end{lesen}

\platz\zwischen{Dezimalzahlen, große Zahlen und Brüche}{小数、大きな数、分数}{shōsū, ōkina kazu, bunsū}

Im Japanischen trennt wie im Englischen ein \textbf{Punkt} die Nachkommastellen ab, das Komma gliedert
große Zahlen in Dreiergruppen -- im Deutschen ist es genau umgekehrt. Brüche liest man von unten nach
oben: erst den Nenner, dann \jp{分の}{bun no} und den Zähler.

\begin{lesen}
\lies{$3{,}14$ \newline \footnotesize jap. 3.14}{3.14}{san ten ichi yon}{drei Komma eins vier}
\lies{$0{,}3$ \newline \footnotesize jap. 0.3}{0.3}{rei ten san}{null Komma drei}
\lies{$30\,000$ \newline \footnotesize jap. 30,000}{3万}{sanman}{dreißigtausend}
\lies{$2{,}5 \cdot 10^6$ \newline \footnotesize jap. $2.5 \times 10^6$}{2.5掛ける10の6乗}{ni ten go kakeru jū no rokujō}{zwei Komma fünf mal zehn hoch sechs}
\lies{$\frac12;\ \frac13;\ \frac14$}{2分の1、3分の1、4分の1}{ni bun no ichi, san bun no ichi, yon bun no ichi}{ein Halb, ein Drittel, ein Viertel}
\lies{$\frac34;\ \frac58$}{4分の3、8分の5}{yon bun no san, hachi bun no go}{drei Viertel, fünf Achtel}
\lies{$\frac{1}{100}$}{100分の1}{hyaku bun no ichi}{ein Hundertstel}
\lies{$\frac{17}{25}$}{25分の17}{nijūgo bun no jūnana}{siebzehn Fünfundzwanzigstel}
\end{lesen}

\platz\zwischen{Japanische Schreibweisen}{日本の表記}{nihon no hyōki}

Japanische Schulbücher schreiben manches anders als deutsche. Die wichtigsten Unterschiede:

\begin{lesen}
\lies{$3 \times 5 \qquad 6 \div 2$}{3掛ける5、6割る2}{san kakeru go, roku waru ni}{Malzeichen ist das Kreuz, Geteiltzeichen ist $\div$; der Doppelpunkt steht nur für Verhältnisse}
\lies{$\pi \fallingdotseq 3.14$}{パイはほぼ3.14}{pai wa hobo san ten ichi yon}{$\fallingdotseq$ statt $\approx$; Dezimalpunkt statt Komma}
\lies{$a \leqq b \qquad a \geqq b$}{$a$ は $b$ 以下、$a$ は $b$ 以上}{$a$ wa $b$ ika, $a$ wa $b$ ijō}{$\leqq$ und $\geqq$ mit doppeltem Strich statt $\le$ und $\ge$}
\lies{$\triangle ABC \equiv \triangle DEF$}{合同}{gōdō}{kongruent: $\equiv$ statt $\cong$}
\lies{$\triangle ABC \backsim \triangle DEF$}{相似}{sōji}{ähnlich: $\backsim$ statt $\sim$}
\lies{$g \mathbin{/\!/} h$}{平行}{heikō}{parallel: schräge Striche statt $\parallel$}
\lies{$P(4,\ 3)$}{点 $P$ の座標は4、3}{ten $P$ no zahyō wa yon, san}{Koordinaten mit Komma statt mit senkrechtem Strich}
\lies{$y = ax + b$}{一次関数：傾き $a$、切片 $b$}{ichiji kansū: katamuki $a$, seppen $b$}{lineare Funktion mit $a$ und $b$ statt $m$ und $n$}
\lies{${}_n\mathrm{C}_r \qquad {}_n\mathrm{P}_r$}{組合せ、順列}{kumiawase, junretsu}{statt $\binom{n}{r}$ und $\frac{n!}{(n-r)!}$}
\lies{$10\,000 \qquad 100\,000\,000$}{1万、1億}{ichiman, ichioku}{große Zahlen zählt man in Viererschritten: zehntausend, hundert Millionen}
\lies{$a,\ b,\ c \qquad x,\ y,\ z$}{エー、ビー、シー、エックス、ワイ、ゼット}{ē, bī, shī, ekkusu, wai, zetto}{Buchstaben heißen wie im Englischen}
\end{lesen}

\anmerkung{Zwischen Variablen wird das Malzeichen weggelassen: $ab$. Der Doppelpunkt bedeutet in Japan
nie „geteilt durch“, sondern immer ein Verhältnis: $a : b$ liest man \jp{$a$ 対 $b$}{$a$ tai $b$}.}

%--------------------------------------------------------------------------------------------------

\platz\zwischen{Beispiele}{例}{rei}

\begin{lesen}
\lies{$7{,}25$}{7.25}{nana ten ni go}{sieben Komma zwei fünf}
\lies{$0{,}05$}{0.05}{rei ten rei go}{null Komma null fünf}
\lies{$1\,250\,000$ \newline \footnotesize jap. 1,250,000}{125万}{hyakunijūgoman}{eine Million zweihundertfünfzigtausend}
\end{lesen}

\newpage
\kapitel{Relationszeichen}{等号と不等号}{tōgō to futōgō}\renewcommand{\lesenabstand}{1.12}\renewcommand{\arraystretch}{0.93}

\begin{lesen}
\lies{$p = q$}{$p$ イコール $q$ \;/\; $p$ は $q$ に等しい}{$p$ ikōru $q$ \;/\; $p$ wa $q$ ni hitoshii}{$p$ (ist) gleich $q$}
\lies{$a \neq b$}{$a$ ノットイコール $b$ \;/\; $a$ は $b$ に等しくない}{$a$ notto ikōru $b$ \;/\; $a$ wa $b$ ni hitoshikunai}{$a$ (ist) ungleich $b$ \;/\; $a$ nicht gleich $b$}
\lies{$\pi \approx 3{,}14$ \newline \footnotesize jap. $\pi \fallingdotseq 3.14$}{パイニアリーイコール3.14 \;/\; パイは3.14にほぼ等しい}{pai niarī ikōru san ten ichi yon \;/\; pai wa san ten ichi yon ni hobo hitoshii}{pi (ist) ungefähr gleich drei Komma eins vier}
\lies{$f(x) \equiv 0$}{$f(x)$ は恒等的に0に等しい}{$f(x)$ wa kōtōteki ni zero ni hitoshii}{$f$ von $x$ ist identisch gleich null}
\lies{$a < b$}{$a$ 小なり $b$ \;/\; $a$ は $b$ より小さい}{$a$ shōnari $b$ \;/\; $a$ wa $b$ yori chiisai}{$a$ (ist) kleiner als $b$}
\lies{$a > b$}{$a$ 大なり $b$ \;/\; $a$ は $b$ より大きい}{$a$ dainari $b$ \;/\; $a$ wa $b$ yori ōkii}{$a$ (ist) größer als $b$}
\lies{$a \le b$ \newline \footnotesize jap. $a \leqq b$}{$a$ 小なりイコール $b$ \;/\; $a$ は $b$ 以下}{$a$ shōnari ikōru $b$ \;/\; $a$ wa $b$ ika}{$a$ (ist) kleiner oder gleich $b$ \;/\; höchstens gleich $b$}
\lies{$a \ge b$ \newline \footnotesize jap. $a \geqq b$}{$a$ 大なりイコール $b$ \;/\; $a$ は $b$ 以上}{$a$ dainari ikōru $b$ \;/\; $a$ wa $b$ ijō}{$a$ (ist) größer oder gleich $b$ \;/\; mindestens gleich $b$}
\lies{$a \ll b$}{$a$ は $b$ よりはるかに小さい}{$a$ wa $b$ yori haruka ni chiisai}{$a$ (ist) viel kleiner als $b$}
\lies{$a \gg b$}{$a$ は $b$ よりはるかに大きい}{$a$ wa $b$ yori haruka ni ōkii}{$a$ (ist) viel größer als $b$}
\lies{$y \sim x$ \newline \footnotesize jap. $y \propto x$}{$y$ は $x$ に比例する}{$y$ wa $x$ ni hirei suru}{$y$ ist (direkt) proportional zu $x$}
\lies{$a : b = c : d$}{$a$ 対 $b$ は $c$ 対 $d$ \;/\; $a$ 対 $b$ イコール $c$ 対 $d$}{$a$ tai $b$ wa $c$ tai $d$ \;/\; $a$ tai $b$ ikōru $c$ tai $d$}{$a$ verhält sich zu $b$ wie $c$ zu $d$}
\lies{$2 \le x \le 5$}{$x$ は2以上5以下}{$x$ wa ni ijō go ika}{$x$ ist mindestens zwei und höchstens fünf \;/\; $x$ liegt zwischen zwei und fünf (einschließlich)}
\end{lesen}

\anmerkung{\jp{以上}{ijō} und \jp{以下}{ika} schließen die Grenze ein ($\ge$, $\le$);
\jp{より大きい}{yori ōkii} und \jp{未満}{miman} schließen sie aus ($>$, $<$). Für „proportional“ schreibt
man in Japan $\propto$, das Zeichen $\backsim$ bedeutet „ähnlich“.}

%--------------------------------------------------------------------------------------------------

\platz\zwischen{Beispiele}{例}{rei}

\begin{lesen}
\lies{$5 < 8$}{5は8より小さい}{go wa hachi yori chiisai}{fünf ist kleiner als acht}
\lies{$x \ge 0$}{$x$ は0以上：$x$ は負でない}{$x$ wa zero ijō: $x$ wa fu de nai}{$x$ ist größer oder gleich null: $x$ ist nicht negativ}
\lies{$7 \neq 9$}{7は9に等しくない}{nana wa kyū ni hitoshikunai}{sieben ist ungleich neun}
\end{lesen}

\newpage
\kapitel{Grundrechenarten}{四則演算}{shisoku enzan}\renewcommand{\lesenabstand}{1.0}\renewcommand{\arraystretch}{0.93}

Die Namen der Rechengrößen auf einen Blick:
\[
\ub{a}{加数}{Summand} + \ub{b}{加数}{Summand} = \ub{s}{和}{Summe}
\qquad\qquad
\ub{a}{被減数}{Minuend} - \ub{b}{減数}{Subtrahend} = \ub{d}{差}{Differenz}
\]
\[
\ub{a}{因数}{Faktor} \cdot \ub{b}{因数}{Faktor} = \ub{p}{積}{Produkt}
\qquad\qquad
\ub{a}{被除数}{Dividend} : \ub{b}{除数}{Divisor} = \ub{q}{商}{Quotient}
\]

\begin{lesen}
\lies{$a_1 + a_2 = s$}{$a_1$ 足す $a_2$ は $s$ \;/\; $a_1$ プラス $a_2$ イコール $s$}{$a_1$ tasu $a_2$ wa $s$ \;/\; $a_1$ purasu $a_2$ ikōru $s$}{$a$ eins plus $a$ zwei gleich $s$}
\lies{$L - l = d$}{大文字の $L$ 引く小文字の $l$ は $d$}{ōmoji no $L$ hiku komoji no $l$ wa $d$}{groß $L$ minus klein $l$ gleich $d$}
\lies{$a \cdot b = c$ \newline \footnotesize jap. $a \times b$}{$a$ 掛ける $b$ は $c$}{$a$ kakeru $b$ wa $c$}{$a$ mal $b$ gleich $c$}
\lies{$ab = c$}{$a$ $b$ イコール $c$}{$a$ $b$ ikōru $c$}{$a$ $b$ gleich $c$}
\lies{$a : b = c$ \newline \footnotesize jap. $a \div b$}{$a$ 割る $b$ は $c$}{$a$ waru $b$ wa $c$}{$a$ geteilt durch $b$ gleich $c$ \;/\; $a$ durch $b$ gleich $c$}
\lies{$\dfrac{a}{b} = a/b$}{$b$ 分の $a$}{$b$ bun no $a$}{$a$ durch $b$ \;/\; $a$ $b$-tel}
\end{lesen}

\platz\zwischen{Beispiele}{例}{rei}

\begin{lesen}
\lies{$17 + 8 = 25$}{17足す8は25}{jūnana tasu hachi wa nijūgo}{siebzehn plus acht gleich fünfundzwanzig}
\lies{$42 - 15 = 27$}{42引く15は27}{yonjūni hiku jūgo wa nijūnana}{zweiundvierzig minus fünfzehn ist siebenundzwanzig}
\lies{$6 \cdot 7 = 42$}{6掛ける7は42 \;/\; 六七四十二（九九）}{roku kakeru nana wa yonjūni \;/\; rokushichi shijūni (kuku)}{sechs mal sieben ist zweiundvierzig}
\lies{$56 : 8 = 7$}{56割る8は7}{gojūroku waru hachi wa nana}{sechsundfünfzig geteilt durch acht ist sieben}
\lies{$(3 + 4) \cdot 5 = 35$}{かっこ3足す4かっことじ掛ける5は35}{kakko san tasu yon kakkotoji kakeru go wa sanjūgo}{Klammer auf drei plus vier Klammer zu mal fünf ist fünfunddreißig}
\end{lesen}

%--------------------------------------------------------------------------------------------------
\newpage
\kapitel{Rechnen mit Brüchen}{分数の計算}{bunsū no keisan}\renewcommand{\lesenabstand}{1.12}\renewcommand{\arraystretch}{0.93}

\begin{lesen}
\lies{$\dfrac{12}{18} = \dfrac{2}{3}$}{18分の12は3分の2。\newline 規則：約分するには、分子と分母を同じ数（ここでは6）で割る。}{jūhachi bun no jūni wa san bun no ni. \newline kisoku: yakubun suru ni wa, bunshi to bunbo o onaji kazu (koko de wa roku) de waru.}{zwölf Achtzehntel gleich zwei Drittel. \newline \textit{Regel:} Man kürzt einen Bruch, indem man Zähler und Nenner durch dieselbe Zahl dividiert (hier 6).}
\lies{$\dfrac{3}{4} = \dfrac{15}{20}$}{4分の3は20分の15。\newline 規則：倍分するには、分子と分母に同じ数（ここでは5）を掛ける。}{yon bun no san wa nijū bun no jūgo. \newline kisoku: baibun suru ni wa, bunshi to bunbo ni onaji kazu (koko de wa go) o kakeru.}{drei Viertel gleich fünfzehn Zwanzigstel. \newline \textit{Regel:} Man erweitert einen Bruch, indem man Zähler und Nenner mit derselben Zahl multipliziert (hier 5).}
\lies{$\dfrac{2}{3} + \dfrac{1}{4} = \dfrac{8}{12} + \dfrac{3}{12} = \dfrac{11}{12}$}{3分の2足す4分の1は、12分の8足す12分の3で、12分の11。\newline 規則：分母が異なる分数は、まず通分してから分子を足す（引く）。}{san bun no ni tasu yon bun no ichi wa, jūni bun no hachi tasu jūni bun no san de, jūni bun no jūichi. \newline kisoku: bunbo ga kotonaru bunsū wa, mazu tsūbun shite kara bunshi o tasu (hiku).}{zwei Drittel plus ein Viertel gleich acht Zwölftel plus drei Zwölftel gleich elf Zwölftel. \newline \textit{Regel:} Ungleichnamige Brüche macht man zuerst gleichnamig (Hauptnenner) und addiert bzw.\ subtrahiert dann die Zähler.}
\lies{$\dfrac{2}{3} \cdot \dfrac{9}{10} = \dfrac{18}{30} = \dfrac{3}{5}$}{3分の2掛ける10分の9は30分の18で、5分の3。\newline 規則：分子どうし、分母どうしを掛ける。}{san bun no ni kakeru jū bun no kyū wa sanjū bun no jūhachi de, go bun no san. \newline kisoku: bunshi dōshi, bunbo dōshi o kakeru.}{zwei Drittel mal neun Zehntel gleich achtzehn Dreißigstel gleich drei Fünftel. \newline \textit{Regel:} Man multipliziert Zähler mit Zähler und Nenner mit Nenner.}
\lies{$\dfrac{3}{4} : \dfrac{9}{8} = \dfrac{3}{4} \cdot \dfrac{8}{9} = \dfrac{2}{3}$}{4分の3割る8分の9は、4分の3掛ける9分の8で、3分の2。\newline 規則：分数で割るときは、その逆数を掛ける。}{yon bun no san waru hachi bun no kyū wa, yon bun no san kakeru kyū bun no hachi de, san bun no ni. \newline kisoku: bunsū de waru toki wa, sono gyakusū o kakeru.}{drei Viertel geteilt durch neun Achtel gleich drei Viertel mal acht Neuntel gleich zwei Drittel. \newline \textit{Regel:} Man dividiert durch einen Bruch, indem man mit seinem Kehrwert multipliziert.}
\lies{$\dfrac{11}{4} = 2\tfrac{3}{4} = 2 + \dfrac{3}{4}$}{4分の11は2と4分の3。\newline 規則：仮分数は、整数と真分数に分けて帯分数に直す。}{yon bun no jūichi wa ni to yon bun no san. \newline kisoku: kabunsū wa, seisū to shinbunsū ni wakete taibunsū ni naosu.}{elf Viertel gleich zwei drei Viertel. \newline \textit{Regel:} Einen unechten Bruch schreibt man als gemischte Zahl, indem man ihn in eine ganze Zahl und einen echten Bruch zerlegt.}
\end{lesen}

\anmerkung{Das Gleichnamigmachen heißt \jp{通分}{tsūbun}, das Kürzen \jp{約分}{yakubun}. Gemischte
Zahlen liest man mit \jp{と}{to}, „und“: $2\tfrac34$ -- \textit{ni to yon bun no san}.}

%--------------------------------------------------------------------------------------------------

\platz\zwischen{Beispiele}{例}{rei}

\begin{lesen}
\lies{$\dfrac56 - \dfrac14 = \dfrac{10}{12} - \dfrac{3}{12} = \dfrac{7}{12}$}{6分の5引く4分の1は、12分の10引く12分の3で、12分の7}{roku bun no go hiku yon bun no ichi wa, jūni bun no jū hiku jūni bun no san de, jūni bun no nana}{fünf Sechstel minus ein Viertel ist gleich zehn Zwölftel minus drei Zwölftel, also sieben Zwölftel}
\lies{$\dfrac45 \cdot 10 = 8$}{10の5分の4は8}{jū no go bun no yon wa hachi}{vier Fünftel von zehn sind acht}
\lies{$2\tfrac12 = \dfrac52$}{2と2分の1は、2分の5}{ni to ni bun no ichi wa, ni bun no go}{zweieinhalb ist gleich fünf Halbe}
\lies{$\dfrac38 = 0{,}375 = 37{,}5\,\%$}{8分の3は0.375、すなわち37.5パーセント}{hachi bun no san wa rei ten san nana go, sunawachi sanjūnana ten go pāsento}{drei Achtel ist gleich null Komma drei sieben fünf oder siebenunddreißig Komma fünf Prozent}
\end{lesen}

\newpage
\kapitel{Potenzen, Wurzeln und Logarithmen}{累乗、累乗根、対数}{ruijō, ruijōkon, taisū}\renewcommand{\lesenabstand}{1.12}\renewcommand{\arraystretch}{0.93}

Aus der Gleichung $b^n = c$ lassen sich drei Fragen stellen -- je nachdem, welche Größe gesucht ist:

\begin{center}
\begin{tikzpicture}[font=\small]
	\node[font=\huge] (f) at (0.6,0) {$b^{\,n} = c$};
	\coordinate (b) at ($(f.west)+(0.3,-0.15)$);
	\coordinate (n) at ($(f.north west)+(0.62,-0.2)$);
	\coordinate (c) at ($(f.east)+(-0.2,0.05)$);
	% Namen
	\node[align=center, blue!70!black] (nb) at (-2.4,1.3) {底\\[-2pt]\color{black!60}Basis};
	\node[align=center, blue!70!black] (nn) at (0.2,2.2) {指数\\[-2pt]\color{black!60}Exponent};
	\node[align=center, blue!70!black] (nc) at (3.9,1.3) {累乗\\[-2pt]\color{black!60}Potenz(wert)};
	\draw[-{Stealth}, blue!70!black, shorten >=4pt] (nb) -- (b);
	\draw[-{Stealth}, blue!70!black, shorten >=4pt] (nn) -- (n);
	\draw[-{Stealth}, blue!70!black, shorten >=4pt] (nc) -- (c);
	% drei Rechenarten
	\node[align=center] (p) at (-5,-2.2) {\textbf{$c$ gesucht:} $c = b^n$\\ 累乗する\\\color{black!60}Potenzieren};
	\node[align=center] (r) at (0.6,-2.6) {\textbf{$b$ gesucht:} $b = \sqrt[n]{c}$\\ 累乗根を求める\\\color{black!60}Radizieren};
	\node[align=center] (l) at (6,-2.2) {\textbf{$n$ gesucht:} $n = \log_b c$\\ 対数をとる\\\color{black!60}Logarithmieren};
	\draw[red, thick, -{Stealth}] (0.6,-0.5) -- (p.north east);
	\draw[red, thick, -{Stealth}] (0.6,-0.5) -- (r.north);
	\draw[red, thick, -{Stealth}] (0.6,-0.5) -- (l.north west);
\end{tikzpicture}
\end{center}

\platz\zwischen{Potenzen}{累乗}{ruijō}
\begin{lesen}
\lies{$b^n$}{$b$ の $n$ 乗}{$b$ no $n$ jō}{$b$ hoch $n$}
\lies{$5^2 = 25$}{5の2乗は25}{go no nijō wa nijūgo}{fünf Quadrat \;/\; fünf hoch zwei ist fünfundzwanzig}
\lies{$3^3 = 27$}{3の3乗は27}{san no sanjō wa nijūnana}{drei hoch drei ist siebenundzwanzig}
\lies{$2^5 = 32$}{2の5乗は32}{ni no gojō wa sanjūni}{zwei hoch fünf gleich zweiunddreißig}
\lies{$b^{-n} = \dfrac{1}{b^n}$}{$b$ のマイナス $n$ 乗は、$b$ の $n$ 乗分の1}{$b$ no mainasu $n$ jō wa, $b$ no $n$ jō bun no ichi}{$b$ hoch minus $n$ gleich eins durch $b$ hoch $n$}
\lies{$10^{-3} = 0{,}001$}{10のマイナス3乗は0.001}{jū no mainasu sanjō wa rei ten rei rei ichi}{zehn hoch minus drei gleich null Komma null null eins}
\end{lesen}

\platz\zwischen{Wurzeln}{累乗根}{ruijōkon}
\begin{lesen}
\lies{$\sqrt[n]{c} = b$}{$c$ の $n$ 乗根は $b$}{$c$ no $n$ jōkon wa $b$}{die $n$-te Wurzel aus $c$ ist gleich $b$}
\lies{$\sqrt{49} = 7$}{ルート49は7 \;/\; 49の正の平方根は7}{rūto yonjūkyū wa nana \;/\; yonjūkyū no sei no heihōkon wa nana}{die (Quadrat-)Wurzel aus neunundvierzig ist sieben}
\lies{$\sqrt[3]{8} = 2$}{8の3乗根は2 \;/\; 8の立方根は2}{hachi no sanjōkon wa ni \;/\; hachi no rippōkon wa ni}{die dritte Wurzel aus acht ist zwei}
\lies{$\sqrt[4]{81} = 3$}{81の4乗根は3}{hachijūichi no yonjōkon wa san}{die vierte Wurzel aus einundachtzig ist drei}
\lies{$\sqrt[n]{c^m} = c^{\frac{m}{n}}$}{$c$ の $m$ 乗の $n$ 乗根は、$c$ の $n$ 分の $m$ 乗}{$c$ no $m$ jō no $n$ jōkon wa, $c$ no $n$ bun no $m$ jō}{die $n$-te Wurzel aus $c$ hoch $m$ gleich $c$ hoch $m$ $n$-tel}
\lies{$\sqrt{a^2 + b^2}$}{ルート $a$ 2乗プラス $b$ 2乗}{rūto $a$ nijō purasu $b$ nijō}{die Wurzel aus $a$ Quadrat plus $b$ Quadrat}
\end{lesen}

\platz\zwischen{Logarithmen}{対数}{taisū}
\begin{lesen}
\lies{$\log_b c = n$}{$b$ を底とする $c$ の対数は $n$ \;/\; ログ $b$ の $c$ は $n$}{$b$ o tei to suru $c$ no taisū wa $n$ \;/\; rogu $b$ no $c$ wa $n$}{der Logarithmus von $c$ zur Basis $b$ ist $n$}
\lies{$\log_2 32 = 5$}{ログ2の32は5}{rogu ni no sanjūni wa go}{Logarithmus von zweiunddreißig zur Basis zwei gleich fünf}
\lies{$\lg 1000 = 3$}{1000の常用対数は3}{sen no jōyō taisū wa san}{der Zehnerlogarithmus (dekadische Logarithmus) von tausend ist drei}
\lies{$\ln \mathrm{e} = 1$}{$\mathrm{e}$ の自然対数は1}{$\mathrm{e}$ no shizen taisū wa ichi}{der natürliche Logarithmus von e ist eins}
\lies{$\log_b(x \cdot y) = \log_b x + \log_b y$}{積の対数は、対数の和に等しい}{seki no taisū wa, taisū no wa ni hitoshii}{der Logarithmus eines Produkts ist gleich der Summe der Logarithmen}
\end{lesen}

\anmerkung{Japanische Schulbücher schreiben den Zehnerlogarithmus $\log_{10} x$ und den natürlichen
Logarithmus $\log_{\mathrm{e}} x$ oder kurz $\log x$; $\lg x$ und $\ln x$ sind dort unüblich.
\jp{ルート}{rūto} kommt vom englischen \textit{root}, \jp{ログ}{rogu} von \textit{log}.}

%--------------------------------------------------------------------------------------------------

\platz\zwischen{Beispiele}{例}{rei}

\begin{lesen}
\lies{$2^{10} = 1024$}{2の10乗は1024}{ni no jūjō wa sennijūyon}{zwei hoch zehn ist eintausendvierundzwanzig}
\lies{$a^m \cdot a^n = a^{m+n}$}{$a$ の $m$ 乗掛ける $a$ の $n$ 乗は、$a$ の $m$ プラス $n$ 乗}{$a$ no $m$ jō kakeru $a$ no $n$ jō wa, $a$ no $m$ purasu $n$ jō}{$a$ hoch $m$ mal $a$ hoch $n$ ist gleich $a$ hoch $m$ plus $n$}
\lies{$9^{\frac12} = \sqrt{9} = 3$}{9の2分の1乗はルート9で、3}{kyū no ni bun no ichi jō wa rūto kyū de, san}{neun hoch ein Halb ist gleich der Wurzel aus neun, also drei}
\lies{$\sqrt[3]{27} = 3$}{27の3乗根は3}{nijūnana no sanjōkon wa san}{die dritte Wurzel aus siebenundzwanzig ist drei}
\end{lesen}

\newpage
\kapitel{Zahlbereiche}{数の集合}{kazu no shūgō}\renewcommand{\lesenabstand}{1.12}\renewcommand{\arraystretch}{0.93}

\begin{center}
\begin{tikzpicture}[font=\small]
	\foreach \k/\col in {4/blue!6, 3/blue!12, 2/blue!18, 1/blue!24, 0/blue!32} {
		\filldraw[fill=\col, draw=blue, thick] ({0.8*\k},0) ellipse[x radius={1.15+1.2*\k}, y radius={0.62+0.42*\k}];
	}
	\node[font=\large] at (-0.35,0.2) {$\N$};       \node at (0.4,-0.2) {$7$};
	\node[font=\large] at (2.15,0.4) {$\Z$};        \node at (2.15,-0.3) {$-5$};
	\node[font=\large] at (4.15,0.4) {$\Q$};        \node at (4.15,-0.3) {$\frac12;\ 0{,}\overline{3}$};
	\node[font=\large] at (6.15,0.4) {$\R$};        \node at (6.15,-0.3) {$\sqrt2;\ \pi$};
	\node[font=\large] at (8.1,0.4) {$\C$};         \node at (8.1,-0.3) {$\mathrm{i};\ 2 + 3\mathrm{i}$};
\end{tikzpicture}
\end{center}

\begin{lesen}
\lies{$\N = \{0;\ 1;\ 2;\ \ldots\}$}{自然数全体の集合}{shizensū zentai no shūgō}{die Menge der natürlichen Zahlen}
\lies{$\Z = \{\ldots;\ -1;\ 0;\ 1;\ \ldots\}$}{整数全体の集合：正の整数、負の整数と0}{seisū zentai no shūgō: sei no seisū, fu no seisū to zero}{die Menge der ganzen Zahlen: positive und negative ganze Zahlen und die Null}
\lies{$\Q = \left\{\dfrac{p}{q} \;\middle|\; p, q \in \Z;\ q \neq 0\right\}$}{$p$ と $q$ が整数で $q$ が0でないとき、$q$ 分の $p$ と表せる数全体の集合}{$p$ to $q$ ga seisū de $q$ ga zero de nai toki, $q$ bun no $p$ to arawaseru kazu zentai no shūgō}{die Menge aller Brüche $p$ durch $q$, wobei $p$ und $q$ ganze Zahlen sind und $q$ ungleich null ist}
\lies{$\R$}{実数全体の集合：有限小数と循環小数（有理数）、および循環しない無限小数（無理数）}{jissū zentai no shūgō: yūgen shōsū to junkan shōsū (yūrisū), oyobi junkan shinai mugen shōsū (murisū)}{die Menge der reellen Zahlen: alle abbrechenden und periodischen Dezimalbrüche (rationale Zahlen) und alle nichtabbrechenden, nichtperiodischen Dezimalbrüche (irrationale Zahlen)}
\lies{$\C = \{a + b\,\mathrm{i} \mid a, b \in \R\}$ \newline mit $\mathrm{i}^2 = -1$}{複素数全体の集合：$a$ プラス $b$ アイ。$a$ と $b$ は実数、アイは虚数単位}{fukusosū zentai no shūgō: $a$ purasu $b$ ai. $a$ to $b$ wa jissū, ai wa kyosū tan'i}{die Menge der komplexen Zahlen: $a$ plus $b$ i, wobei $a$ und $b$ reelle Zahlen sind und i die imaginäre Einheit ist}
\lies{$\N \subset \Z \subset \Q \subset \R \subset \C$}{それぞれの数の集合は、次の集合の真部分集合である}{sorezore no kazu no shūgō wa, tsugi no shūgō no shin bubun shūgō de aru}{jeder Zahlbereich ist eine echte Teilmenge des nächsten}
\end{lesen}

\anmerkung{In japanischen Schulbüchern beginnen die natürlichen Zahlen (\jp{自然数}{shizensū}) bei
$1$; die Null gehört nicht dazu. Periodische Dezimalbrüche schreibt man mit Punkten über den Ziffern:
$0.\dot{3}$ statt $0{,}\overline{3}$.}

%--------------------------------------------------------------------------------------------------

\platz\zwischen{Beispiele}{例}{rei}

\begin{lesen}
\lies{$-3 \in \Z,\ -3 \notin \N$}{マイナス3は整数であるが、自然数ではない}{mainasu san wa seisū de aru ga, shizensū de wa nai}{minus drei ist eine ganze Zahl, aber keine natürliche Zahl}
\lies{$\frac34 \in \Q$}{4分の3は有理数である}{yon bun no san wa yūrisū de aru}{drei Viertel ist eine rationale Zahl}
\lies{$\sqrt{2} \notin \Q$}{ルート2は無理数である}{rūto ni wa murisū de aru}{die Wurzel aus zwei ist irrational}
\lies{$0{,}\overline{3} = \frac13$ \newline \footnotesize jap. $0.\dot{3} = \frac13$}{循環小数0.333 \ldots\ は、3分の1に等しい}{junkan shōsū rei ten san san san \ldots\ wa, san bun no ichi ni hitoshii}{null Komma Periode drei ist gleich ein Drittel}
\end{lesen}

\newpage
\kapitel{Mengen}{集合}{shūgō}\renewcommand{\lesenabstand}{1.12}\renewcommand{\arraystretch}{0.93}

\begin{lesen}
\lies{$a \in M$}{$a$ は $M$ の要素である \;/\; $a$ は $M$ に属する}{$a$ wa $M$ no yōso de aru \;/\; $a$ wa $M$ ni zokusuru}{$a$ ist Element von $M$}
\lies{$a \notin M$}{$a$ は $M$ の要素でない}{$a$ wa $M$ no yōso de nai}{$a$ ist nicht Element von $M$}
\lies{$M = \{2;\ 4;\ 6\}$ \newline \footnotesize jap. $\{2,\ 4,\ 6\}$}{$M$ は2、4、6を要素とする集合}{$M$ wa ni, yon, roku o yōso to suru shūgō}{$M$ ist die Menge mit den Elementen zwei, vier und sechs}
\lies{$M = \emptyset$}{$M$ は空集合である}{$M$ wa kūshūgō de aru}{$M$ ist die leere Menge}
\lies{$A \subseteq B$}{$A$ は $B$ の部分集合である}{$A$ wa $B$ no bubun shūgō de aru}{$A$ ist (eine) Teilmenge von $B$}
\lies{$A \subset B$ \newline \footnotesize jap. $A \subsetneqq B$}{$A$ は $B$ の真部分集合である}{$A$ wa $B$ no shin bubun shūgō de aru}{$A$ ist (eine) echte Teilmenge von $B$}
\lies{$A \cup B$}{$A$ カップ $B$ \;/\; $A$ または $B$ \;/\; $A$ と $B$ の和集合}{$A$ kappu $B$ \;/\; $A$ mata wa $B$ \;/\; $A$ to $B$ no washūgō}{$A$ vereinigt $B$ \;/\; die Vereinigungsmenge von $A$ und $B$}
\lies{$A \cap B$}{$A$ キャップ $B$ \;/\; $A$ かつ $B$ \;/\; $A$ と $B$ の共通部分}{$A$ kyappu $B$ \;/\; $A$ katsu $B$ \;/\; $A$ to $B$ no kyōtsū bubun}{$A$ geschnitten $B$ \;/\; die Schnittmenge von $A$ und $B$}
\lies{$A \setminus B$}{$A$ と $B$ の差集合}{$A$ to $B$ no sashūgō}{$A$ ohne $B$ \;/\; die Differenzmenge von $A$ und $B$}
\lies{$A \times B$}{$A$ と $B$ の直積}{$A$ to $B$ no chokuseki}{$A$ kreuz $B$ \;/\; das kartesische Produkt von $A$ und $B$}
\lies{$\{x \in \R \mid x < 4\}$}{4より小さい実数 $x$ 全体の集合}{yon yori chiisai jissū $x$ zentai no shūgō}{die Menge aller reellen Zahlen $x$, für die gilt: $x$ ist kleiner als vier}
\end{lesen}

\anmerkung{In japanischen Schulbüchern bedeutet $A \subset B$ „Teilmenge“ einschließlich der Gleichheit.
Die Schnittmenge heißt dort \jp{共通部分}{kyōtsū bubun}, „gemeinsamer Teil“.}

%--------------------------------------------------------------------------------------------------

\platz\zwischen{Beispiele}{例}{rei}

\begin{lesen}
\lies{$A = \{1;\,2;\,3;\,4\}$ \newline $B = \{3;\,4;\,5\}$}{$A$ は1、2、3、4を要素とする集合、$B$ は3、4、5を要素とする集合}{$A$ wa ichi, ni, san, yon o yōso to suru shūgō, $B$ wa san, yon, go o yōso to suru shūgō}{$A$ ist die Menge mit den Elementen eins, zwei, drei und vier; $B$ ist die Menge mit den Elementen drei, vier und fünf}
\lies{$A \cap B = \{3;\,4\}$}{$A$ と $B$ の共通部分は、3と4を要素とする集合}{$A$ to $B$ no kyōtsū bubun wa, san to yon o yōso to suru shūgō}{die Schnittmenge von $A$ und $B$ ist die Menge mit den Elementen drei und vier}
\lies{$A \cup B = \{1;\,2;\,3;\,4;\,5\}$}{$A$ と $B$ の和集合は、1から5までの数を要素とする}{$A$ to $B$ no washūgō wa, ichi kara go made no kazu o yōso to suru}{die Vereinigungsmenge von $A$ und $B$ enthält die Zahlen von eins bis fünf}
\end{lesen}

\newpage
\kapitel{Gleichungen und Ungleichungen lösen}{方程式と不等式を解く}{hōteishiki to futōshiki o toku}\renewcommand{\lesenabstand}{1.12}\renewcommand{\arraystretch}{0.93}

\platz\zwischen{Beispiel: eine lineare Gleichung}{例：一次方程式}{rei: ichiji hōteishiki}

\begin{lesen}
\lies{$\dfrac{2(x - 1)}{3} + \dfrac{1}{2} = \dfrac{5}{2}$}{次の方程式を解きなさい。}{tsugi no hōteishiki o tokinasai.}{\textit{Löse die Gleichung.}}
\lies{$4(x - 1) + 3 = 15$}{両辺に分母の最小公倍数6を掛ける。}{ryōhen ni bunbo no saishō kōbaisū roku o kakeru.}{Multipliziere beide Seiten mit dem Hauptnenner sechs.}
\lies{$4x - 4 + 3 = 15$}{かっこをはずす。}{kakko o hazusu.}{Löse die Klammern auf (multipliziere aus).}
\lies{$4x - 1 = 15$}{同類項をまとめる。}{dōruikō o matomeru.}{Fasse zusammen.}
\lies{$4x = 16$}{両辺に1を足す。}{ryōhen ni ichi o tasu.}{Addiere eins auf beiden Seiten.}
\lies{$x = 4$}{両辺を $x$ の係数4で割る。}{ryōhen o $x$ no keisū yon de waru.}{Dividiere beide Seiten durch den Koeffizienten von $x$, also durch vier.}
\lies{$L = \{4\}$}{解は $x$ イコール4である。}{kai wa $x$ ikōru yon de aru.}{Die Lösungsmenge enthält nur das Element vier. \;/\; Die Lösung ist $x$ gleich vier.}
\lies{\textit{Probe:} \newline $\frac{2 \cdot 3}{3} + \frac12 = \frac52$ \;\checkmark}{確かめ：$x$ に4を代入すると、2足す2分の1は2分の5。\checkmark}{tashikame: $x$ ni yon o dainyū suru to, ni tasu ni bun no ichi wa ni bun no go. \checkmark}{\textit{Probe:} Setze vier für $x$ ein: zwei plus ein Halb gleich fünf Halbe. \checkmark}
\end{lesen}

\platz\zwischen{Beispiel: eine lineare Ungleichung}{例：一次不等式}{rei: ichiji futōshiki}

\begin{lesen}
\lies{$3 - 2x < 9$}{次の不等式を解きなさい。}{tsugi no futōshiki o tokinasai.}{\textit{Löse die Ungleichung.}}
\lies{$-2x < 6$}{両辺から3を引く。}{ryōhen kara san o hiku.}{Subtrahiere drei auf beiden Seiten.}
\lies{$x > -3$}{両辺をマイナス2で割る。注意：負の数を掛けたり割ったりすると、不等号の向きが変わる。}{ryōhen o mainasu ni de waru. chūi: fu no kazu o kaketari wattari suru to, futōgō no muki ga kawaru.}{Dividiere beide Seiten durch minus zwei. \textbf{Achtung:} Beim Multiplizieren oder Dividieren mit einer negativen Zahl dreht sich das Relationszeichen um.}
\lies{$L = \{x \in \R \mid x > -3\}$}{解は、マイナス3より大きいすべての実数である。}{kai wa, mainasu san yori ōkii subete no jissū de aru.}{Die Lösungsmenge enthält alle reellen Zahlen, die größer als minus drei sind.}
\end{lesen}

\begin{center}
\begin{tikzpicture}[scale=0.8, font=\small]
	\draw[achse] (-6.4,0) -- (3.6,0);
	\foreach \x in {-6,...,3} {\draw (\x,0.1) -- (\x,-0.1) node[below] {$\x$};}
	\draw[blue, very thick, -{Stealth}] (-3,0) -- (3.4,0);
	\filldraw[fill=white, draw=blue, very thick] (-3,0) circle (4pt);
	\node[above, blue] at (0.2,0.15) {$x > -3$};
\end{tikzpicture}
\end{center}

\platz\zwischen{Gleichungsarten}{方程式の種類}{hōteishiki no shurui}

\begin{lesen}
\lies{$3x - 7 = 11$}{一次方程式：3$x$ マイナス7イコール11}{ichiji hōteishiki: san $x$ mainasu nana ikōru jūichi}{eine lineare Gleichung (Gleichung ersten Grades): drei $x$ minus sieben gleich elf}
\lies{$x^2 - 5x + 6 = 0$}{二次方程式：$x$ 2乗マイナス5$x$ プラス6イコール0}{niji hōteishiki: $x$ nijō mainasu go $x$ purasu roku ikōru zero}{eine quadratische Gleichung (Gleichung zweiten Grades): $x$ Quadrat minus fünf $x$ plus sechs gleich null}
\lies{$\sqrt{x + 3} = 5$}{無理方程式：ルート $x$ プラス3イコール5}{muri hōteishiki: rūto $x$ purasu san ikōru go}{eine Wurzelgleichung: die Wurzel aus $x$ plus drei (aus der Summe) gleich fünf}
\lies{$2^x = 64$}{指数方程式：2の $x$ 乗イコール64}{shisū hōteishiki: ni no $x$ jō ikōru rokujūyon}{eine Exponentialgleichung: zwei hoch $x$ gleich vierundsechzig}
\lies{$\log_3 x = 4$}{対数方程式：ログ3の $x$ イコール4}{taisū hōteishiki: rogu san no $x$ ikōru yon}{eine logarithmische Gleichung: Logarithmus von $x$ zur Basis drei gleich vier}
\lies{$2 \sin x = 1$}{三角方程式：2サイン $x$ イコール1}{sankaku hōteishiki: ni sain $x$ ikōru ichi}{eine trigonometrische Gleichung: zwei mal Sinus $x$ gleich eins}
\lies{$\dfrac{x - 3}{2x + 1} = \dfrac16$}{2$x$ プラス1分の $x$ マイナス3イコール6分の1}{ni $x$ purasu ichi bun no $x$ mainasu san ikōru roku bun no ichi}{$x$ minus drei, alles geteilt durch zwei $x$ plus eins, gleich ein Sechstel}
\end{lesen}

\anmerkung{Auch bei längeren Brüchen nennt man zuerst den ganzen Nenner, dann \textit{bun no} und den
Zähler. Klammern liest man bei Bedarf mit: \jp{かっこ}{kakko} \ldots\ \jp{かっことじ}{kakko toji}.
Arbeitsaufträge enden im Schulbuch auf \jp{なさい}{nasai}: \jp{解きなさい}{tokinasai}, „löse“.}

%--------------------------------------------------------------------------------------------------

\platz\zwischen{Beispiele}{例}{rei}

\begin{lesen}
\lies{$x^2 - 5x + 6 = 0$ \newline $(x - 2)(x - 3) = 0$}{左辺を因数分解する。}{sahen o insū bunkai suru.}{Zerlege die linke Seite in Faktoren.}
\lies{$x_1 = 2,\ x_2 = 3$}{積が0ならば、因数のどれかが0である。解は2と3。}{seki ga zero naraba, insū no dore ka ga zero de aru. kai wa ni to san.}{Ein Produkt ist null, wenn einer der Faktoren null ist. Die Lösungen sind zwei und drei.}
\end{lesen}

\newpage
\kapitel{Funktionen}{関数}{kansū}\renewcommand{\lesenabstand}{1.12}\renewcommand{\arraystretch}{0.93}

\begin{lesen}
\lies{$y = f(x)$}{$y$ イコールエフエックス \;/\; $y$ は $x$ の関数である}{$y$ ikōru efu ekkusu \;/\; $y$ wa $x$ no kansū de aru}{$y$ gleich $f$ von $x$ \;/\; $y$ ist der Funktionswert zu $x$}
\lies{$x$}{独立変数}{dokuritsu hensū}{die unabhängige Variable}
\lies{$y$}{従属変数}{jūzoku hensū}{die abhängige Variable}
\lies{$f\colon x \mapsto x^2$}{$f$ は $x$ に $x$ の2乗を対応させる}{$f$ wa $x$ ni $x$ no nijō o taiō saseru}{$f$ ordnet $x$ das Quadrat von $x$ zu \;/\; $f$: $x$ wird abgebildet auf $x$ Quadrat}
\lies{$f(x) = 2x^2 - 3$}{エフエックスイコール2$x$ 2乗マイナス3}{efu ekkusu ikōru ni $x$ nijō mainasu san}{$f$ von $x$ gleich zwei $x$ Quadrat minus drei}
\lies{$f(3) = 15$}{エフ3イコール15 \;/\; $x$ イコール3のときの関数の値は15}{efu san ikōru jūgo \;/\; $x$ ikōru san no toki no kansū no atai wa jūgo}{$f$ von drei gleich fünfzehn \;/\; der Funktionswert an der Stelle drei ist fünfzehn}
\lies{$D_f,\ W_f$}{$f$ の定義域と値域}{$f$ no teigiiki to chiiki}{der Definitionsbereich und der Wertebereich von $f$}
\end{lesen}

\platz\zwischen{Der Graph einer Funktion}{関数のグラフ}{kansū no gurafu}

\begin{center}
\begin{tikzpicture}[font=\small, scale=1.05]
	\draw[achse] (-0.7,0) -- (8.6,0) node[below left] {$x$};
	\draw[achse] (0,-1.1) -- (0,5.4) node[below left] {$y$};
	\begin{scope}
		\clip (-0.7,-1.1) rectangle (8.3,5.2);
		\draw[blue, very thick, domain=-0.35:7.75, samples=80] plot (\x, {0.1*\x*\x*\x - 1.1*\x*\x + 3.2*\x + 0.3});
	\end{scope}
	\coordinate (P) at (6.5,2.0825);
	\draw[red, thick, dashed] (P) -- (6.5,0) (P) -- (0,2.0825);
	\fill[red] (P) circle (2.5pt) node[above left] {$P(x_P \mid y_P)$};
	\node[below] at (6.5,0) {$x_P$}; \node[left] at (0,2.0825) {$y_P$};
	\node[below left] at (0,0) {$O$};
	% Beschriftungen
	\node[align=left, anchor=west] (gr) at (8.0,4.6) {$f$ のグラフ\\\color{black!60}Graph von $f$};
	\draw[-{Stealth}, thin] (gr.west) -- (7.6,4.45);
	\node[align=center] (ab) at (5.1,-1.3) {$x$ 座標\\\color{black!60}Abszisse};
	\draw[-{Stealth}, thin] (ab.east) -- (6.35,-0.35);
	\node[align=center] (or) at (-2.0,3.7) {$y$ 座標\\\color{black!60}Ordinate};
	\draw[-{Stealth}, thin] (or.south) -- (-0.15,2.55);
	\node[align=center] (ur) at (1.6,-1.3) {原点\\\color{black!60}Ursprung};
	\draw[-{Stealth}, thin] (ur.north west) -- (-0.1,-0.2);
	\node[align=center] (xa) at (8.6,0.9) {$x$ 軸\\\color{black!60}$x$-Achse};
	\node[align=center] (ya) at (1.3,5.3) {$y$ 軸\\\color{black!60}$y$-Achse};
\end{tikzpicture}
\end{center}

\begin{lesen}
\lies{$P(x_P \mid y_P)$ \newline \footnotesize jap. $P(x_P,\ y_P)$}{直交座標系で、座標が $x_P$、$y_P$ の点 $P$}{chokkō zahyōkei de, zahyō ga $x_P$, $y_P$ no ten $P$}{der Punkt $P$ mit den Koordinaten $x_P$ und $y_P$ im rechtwinkligen (kartesischen) Koordinatensystem}
\lies{$x_P$}{点 $P$ の $x$ 座標}{ten $P$ no $x$ zahyō}{die Abszisse ($x$-Koordinate) von $P$}
\lies{$y_P$}{点 $P$ の $y$ 座標}{ten $P$ no $y$ zahyō}{die Ordinate ($y$-Koordinate) von $P$}
\end{lesen}

\platz\zwischen{Einteilung der elementaren Funktionen}{初等関数の分類}{shotō kansū no bunrui}

\begin{center}
\begin{tikzpicture}[font=\footnotesize, every node/.style={align=left, anchor=west}]
	\node (root) at (0,0)      {\jfett{初等関数}\\\color{black!60}elementare Funktionen};
	\node (alg) at (4.4,2.7)  {\jfett{代数関数}\\\color{black!60}algebraische Funktionen};
	\node (tra) at (4.4,-2.95){\jfett{超越関数}\\\color{black!60}transzendente Funktionen};
	\node (rat) at (9.2,3.7)  {有理関数\\\color{black!60}rationale Funktionen};
	\node (irr) at (9.2,1.7)  {無理関数\\\color{black!60}irrationale Funktionen\\ $y = \sqrt{x}$};
	\node (pol) at (13.0,4.75) {多項式関数\\\color{black!60}ganzrationale Funktionen\\ $y = x^3 - 2x$};
	\node (fra) at (13.0,2.75) {分数関数\\\color{black!60}gebrochenrationale F.\\ $y = \frac{x + 1}{x - 2}$};
	\node (exp) at (9.2,-0.2) {指数関数\\\color{black!60}Exponentialfunktionen\\ $y = \mathrm{e}^x$};
	\node (log) at (9.2,-2.05) {対数関数\\\color{black!60}Logarithmusfunktionen\\ $y = \ln x$};
	\node (tri) at (9.2,-3.9) {三角関数\\\color{black!60}Winkelfunktionen\\ $y = \sin x$};
	\node (arc) at (9.2,-5.75) {逆三角関数\\\color{black!60}Arkusfunktionen\\ $y = \arcsin x$};
	\foreach \a/\b in {root/alg, root/tra, alg/rat, alg/irr, rat/pol, rat/fra, tra/exp, tra/log, tra/tri, tra/arc}
		{\draw[blue, thick] (\a.east) -- ++(0.35,0) |- (\b.west);}
\end{tikzpicture}
\end{center}

\begin{lesen}
\lies{$y = \displaystyle\sum_{k=0}^{4} a_k x^k$}{$y$ イコールシグマ、$k$ が0から4までの $a_k$ 掛ける $x$ の $k$ 乗}{$y$ ikōru shiguma, $k$ ga zero kara yon made no $a_k$ kakeru $x$ no $k$ jō}{$y$ gleich Summe über $a$ Index $k$ mal $x$ hoch $k$ für $k$ von null bis vier}
\lies{$y = \dfrac{7 + x}{2 - x^2}$}{$y$ イコール2マイナス $x$ 2乗分の7プラス $x$}{$y$ ikōru ni mainasu $x$ nijō bun no nana purasu $x$}{$y$ ist gleich dem Bruch mit dem Zähler sieben plus $x$ und dem Nenner zwei minus $x$ Quadrat}
\lies{$y = -\sqrt{r^2 - x^2}$}{$y$ イコールマイナスルート $r$ 2乗マイナス $x$ 2乗}{$y$ ikōru mainasu rūto $r$ nijō mainasu $x$ nijō}{$y$ gleich minus die Wurzel aus $r$ Quadrat minus $x$ Quadrat}
\lies{$y = \mathrm{e}^x$}{$y$ イコール $\mathrm{e}$ の $x$ 乗}{$y$ ikōru $\mathrm{e}$ no $x$ jō}{$y$ gleich e hoch $x$}
\lies{$y = \ln x$}{$y$ イコールログ $x$ \;/\; $y$ は $x$ の自然対数}{$y$ ikōru rogu $x$ \;/\; $y$ wa $x$ no shizen taisū}{$y$ gleich ln $x$ \;/\; $y$ gleich natürlicher Logarithmus von $x$}
\end{lesen}

%--------------------------------------------------------------------------------------------------

\platz\zwischen{Beispiele}{例}{rei}

\begin{lesen}
\lies{$f(x) = 2x^2 - 3$ \newline $f(2) = 5$}{エフ2イコール5：点 $P(2,\,5)$ は $f$ のグラフ上にある}{efu ni ikōru go: ten $P(2,\,5)$ wa $f$ no gurafu jō ni aru}{$f$ von zwei ist fünf: Der Punkt $P(2 \mid 5)$ liegt auf dem Graphen von $f$}
\lies{$f(x) = \frac12 x + 1$}{傾き2分の1、$y$ 切片1の一次関数}{katamuki ni bun no ichi, $y$ seppen ichi no ichiji kansū}{eine lineare Funktion mit der Steigung ein Halb und dem $y$-Achsenabschnitt eins}
\lies{$f(x) = 0$}{$f$ の零点を求めなさい。$f(x)$ が0になる $x$ の値はいくつか。}{$f$ no reiten o motomenasai. $f(x)$ ga zero ni naru $x$ no atai wa ikutsu ka.}{Bestimme die Nullstellen von $f$: Für welche $x$ ist $f$ von $x$ gleich null?}
\lies{$f(x) = g(x)$}{2つのグラフの交点を求めなさい。}{futatsu no gurafu no kōten o motomenasai.}{Bestimme die Schnittpunkte der beiden Graphen.}
\end{lesen}

\newpage
\kapitel{Winkelfunktionen}{三角関数}{sankaku kansū}\renewcommand{\lesenabstand}{1.12}\renewcommand{\arraystretch}{0.93}

\begin{center}
\begin{tikzpicture}[scale=2.1, font=\small]
	\draw[achse] (-1.25,0) -- (1.5,0) node[below left] {$x$};
	\draw[achse] (0,-1.2) -- (0,1.35) node[below left] {$y$};
	\draw[blue, thick] (0,0) circle (1);
	\coordinate (P) at (35:1);
	\draw[mass] (0.3,0) arc[start angle=0, end angle=35, radius=0.3]; \node at (17:0.42) {$\alpha$};
	\draw[blue, very thick] (0,0) -- (P);
	\draw[red, very thick] (P) -- (P |- 0,0) node[midway, left, fill=white, inner sep=1pt] {$\sin\alpha$};
	\draw[green!50!black, very thick] (0,0) -- (P |- 0,0) node[midway, below] {$\cos\alpha$};
	\draw[violet, very thick] (1,0) -- (1,{tan(35)}) node[midway, right] {$\tan\alpha$};
	\draw[violet, dashed] (P) -- (1,{tan(35)});
	\fill[blue] (P) circle (1pt) node[above left] {$P$};
	\node[below] at (1,0) {$1$};
	\node[align=left, anchor=west, font=\footnotesize] at (1.55,0.9) {単位円\\\color{black!60}Einheitskreis};
\end{tikzpicture}
\end{center}

\begin{lesen}
\lies{$\sin x$}{サイン $x$ \;/\; $x$ の正弦}{sain $x$ \;/\; $x$ no seigen}{Sinus $x$}
\lies{$\cos x$}{コサイン $x$ \;/\; $x$ の余弦}{kosain $x$ \;/\; $x$ no yogen}{Kosinus $x$}
\lies{$\tan x$}{タンジェント $x$ \;/\; $x$ の正接}{tanjento $x$ \;/\; $x$ no seisetsu}{Tangens $x$}
\lies{$\cot x$}{コタンジェント $x$}{kotanjento $x$}{Kotangens $x$}
\lies{$\sec x,\ \csc x$}{セカント $x$、コセカント $x$}{sekanto $x$, kosekanto $x$}{Sekans und Kosekans von $x$ (im deutschen Unterricht kaum verwendet)}
\lies{$\arcsin x$ \newline \footnotesize jap. meist $\sin^{-1} x$}{アークサイン $x$ \;/\; サインインバース $x$}{ākusain $x$ \;/\; sain inbāsu $x$}{Arkussinus $x$ \;/\; der Winkel, dessen Sinus $x$ ist}
\lies{$\sin 30\degree = \dfrac12$}{サイン30度は2分の1}{sain sanjū do wa ni bun no ichi}{der Sinus von dreißig Grad ist ein Halb}
\lies{$\tan 45\degree = 1$}{タンジェント45度は1}{tanjento yonjūgo do wa ichi}{Tangens fünfundvierzig Grad gleich eins}
\lies{$\sin^2 x + \cos^2 x = 1$}{サイン2乗 $x$ プラスコサイン2乗 $x$ イコール1}{sain nijō $x$ purasu kosain nijō $x$ ikōru ichi}{Sinus Quadrat $x$ plus Kosinus Quadrat $x$ gleich eins}
\lies{$5\sin^2 x - \cos 2x = 0$}{5サイン2乗 $x$ マイナスコサイン2$x$ イコール0}{go sain nijō $x$ mainasu kosain ni $x$ ikōru zero}{fünf mal Sinus Quadrat $x$ minus Kosinus zwei $x$ ist gleich null}
\end{lesen}

\anmerkung{Die Winkelfunktionen heißen zusammen \jp{三角関数}{sankaku kansū}, im rechtwinkligen Dreieck
\jp{三角比}{sankakuhi}. Vorsicht bei $\sin^{-1} x$: Gemeint ist die Umkehrfunktion $\arcsin x$, nicht der
Kehrwert $\frac{1}{\sin x}$.}

%--------------------------------------------------------------------------------------------------

\platz\zwischen{Beispiele}{例}{rei}

\begin{lesen}
\lies{$\cos 60\degree = \dfrac12$}{コサイン60度は2分の1}{kosain rokujū do wa ni bun no ichi}{der Kosinus von sechzig Grad ist ein Halb}
\lies{$\sin x = 0{,}5$ \newline $x_1 = 30\degree,\ x_2 = 150\degree$}{0度から360度までの範囲で、この方程式の解は30度と150度の2つである。}{rei do kara sanbyakurokujū do made no han'i de, kono hōteishiki no kai wa sanjū do to hyakugojū do no futatsu de aru.}{Für Winkel zwischen null und dreihundertsechzig Grad hat die Gleichung zwei Lösungen: dreißig Grad und hundertfünfzig Grad.}
\lies{$\sin(x + 360\degree) = \sin x$}{サイン関数は、周期360度の周期関数である。}{sain kansū wa, shūki sanbyakurokujū do no shūki kansū de aru.}{Die Sinusfunktion ist periodisch mit der Periode dreihundertsechzig Grad.}
\lies{$-1 \le \sin x \le 1$}{サイン関数の値は、マイナス1以上1以下である。}{sain kansū no atai wa, mainasu ichi ijō ichi ika de aru.}{Die Werte der Sinusfunktion liegen zwischen minus eins und eins.}
\end{lesen}

\newpage
\kapitel{Intervalle und Grenzwerte}{区間と極限}{kukan to kyokugen}\renewcommand{\lesenabstand}{1.12}\renewcommand{\arraystretch}{0.93}

\begin{lesen}
\lies{\intervall{o}{o} \newline $]a;\,b[$ \quad jap. $(a, b)$}{$a$ から $b$ までの開区間：$a < x < b$}{$a$ kara $b$ made no kaikukan: $a < x < b$}{das offene Intervall von $a$ bis $b$: $a < x < b$}
\lies{\intervall{g}{g} \newline $[a;\,b]$ \quad jap. $[a, b]$}{$a$ から $b$ までの閉区間：$a \le x \le b$}{$a$ kara $b$ made no heikukan: $a \le x \le b$}{das abgeschlossene Intervall von $a$ bis $b$: $a \le x \le b$}
\lies{\intervall{o}{g} \newline $]a;\,b]$ \quad jap. $(a, b]$}{$a$ から $b$ までの半開区間：$a < x \le b$（$a$ は含まず、$b$ は含む）}{$a$ kara $b$ made no hankai kukan: $a < x \le b$ ($a$ wa fukumazu, $b$ wa fukumu)}{das halboffene Intervall von $a$ bis $b$, links offen und rechts abgeschlossen: $a < x \le b$}
\lies{\intervallR \newline $[a;\,\infty[$ \quad jap. $[a, \infty)$}{$a$ 以上の区間：$x \ge a$}{$a$ ijō no kukan: $x \ge a$}{das Intervall von $a$ bis unendlich: $x \ge a$}
\lies{$\R = \ ]{-\infty};\,\infty[$}{実数全体は、マイナス無限大からプラス無限大までの開区間である}{jissū zentai wa, mainasu mugendai kara purasu mugendai made no kaikukan de aru}{die Menge der reellen Zahlen ist das offene Intervall von minus unendlich bis plus unendlich}
\end{lesen}

\anmerkung{In Japan stehen runde Klammern für offene und eckige für abgeschlossene Intervalle, getrennt
durch ein Komma: $(a,\ b)$ und $[a,\ b]$. Für Monotonie-Angaben schreiben wir die Intervalle als
Ungleichungskette, z.\,B.\ \textit{fallend:} $-\infty < x \le 2$.}

\begin{lesen}
\lies{$x \to x_0$}{$x$ が $x_0$ に近づく}{$x$ ga $x_0$ ni chikazuku}{$x$ gegen $x$ null \;/\; $x$ strebt gegen $x$ null}
\lies{$\limxto{2} x^2 = 4$}{$x$ が2に近づくとき、$x$ 2乗の極限値は4 \;/\; リミット $x$ 矢印2の $x$ 2乗は4}{$x$ ga ni ni chikazuku toki, $x$ nijō no kyokugenchi wa yon \;/\; rimitto $x$ yajirushi ni no $x$ nijō wa yon}{der Grenzwert (Limes) von $x$ Quadrat für $x$ gegen zwei ist vier}
\lies{$\limxto{x_1} f(x) = L$}{$x$ が $x_1$ に限りなく近づくとき、$f(x)$ は $L$ に限りなく近づく}{$x$ ga $x_1$ ni kagirinaku chikazuku toki, $f(x)$ wa $L$ ni kagirinaku chikazuku}{Limes $f$ von $x$ für $x$ gegen $x$ eins ist groß $L$ -- $f(x)$ nähert sich $L$, wenn $x$ beliebig gegen $x_1$ strebt}
\lies{$\limxto{x_0} f(x) \neq f(x_0)$}{$x$ が $x_0$ に近づくときの $f(x)$ の極限値は $f(x_0)$ に等しくない：関数 $f$ は $x_0$ で不連続である}{$x$ ga $x_0$ ni chikazuku toki no $f(x)$ no kyokugenchi wa $f(x_0)$ ni hitoshikunai: kansū $f$ wa $x_0$ de furenzoku de aru}{der Grenzwert von $f$ für $x$ gegen $x$ null ist ungleich $f$ von $x$ null: $f$ ist an der Stelle $x_0$ unstetig}
\lies{$\limn \dfrac{1}{n} = 0$}{$n$ を限りなく大きくすると、$n$ 分の1は0に収束する}{$n$ o kagirinaku ōkiku suru to, $n$ bun no ichi wa zero ni shūsoku suru}{der Grenzwert von eins durch $n$ für $n$ gegen unendlich ist null: $\left(\frac1n\right)$ ist eine Nullfolge}
\lies{$\limn a_n = g$}{数列 $(a_n)$ は $g$ に収束する \;/\; 極限値は $g$ である}{sūretsu $(a_n)$ wa $g$ ni shūsoku suru \;/\; kyokugenchi wa $g$ de aru}{die Folge $(a_n)$ konvergiert gegen $g$ \;/\; sie ist konvergent mit dem Grenzwert $g$}
\end{lesen}

%--------------------------------------------------------------------------------------------------

\platz\zwischen{Beispiele}{例}{rei}

\begin{lesen}
\lies{$x \in [2;\,5]$}{$x$ は2以上5以下の閉区間に属する}{$x$ wa ni ijō go ika no heikukan ni zokusuru}{$x$ ist Element des abgeschlossenen Intervalls von zwei bis fünf}
\lies{$\lim\limits_{x \to \infty} \dfrac{1}{x} = 0$}{$x$ を限りなく大きくすると、$x$ 分の1の極限値は0}{$x$ o kagirinaku ōkiku suru to, $x$ bun no ichi no kyokugenchi wa zero}{der Grenzwert von eins durch $x$ für $x$ gegen unendlich ist null}
\lies{$\limn \dfrac{2n + 1}{n} = 2$}{$n$ を限りなく大きくすると、$n$ 分の2$n$ プラス1の極限値は2}{$n$ o kagirinaku ōkiku suru to, $n$ bun no ni $n$ purasu ichi no kyokugenchi wa ni}{der Grenzwert von zwei $n$ plus eins, geteilt durch $n$, für $n$ gegen unendlich ist zwei}
\lies{steigend: $2 \le x < \infty$}{関数は、$x$ が2以上の範囲で増加する}{kansū wa, $x$ ga ni ijō no han'i de zōka suru}{die Funktion ist steigend für alle $x$, die größer oder gleich zwei sind}
\end{lesen}

\newpage
\kapitel{Differenzialrechnung}{微分法}{bibunhō}\renewcommand{\lesenabstand}{1.12}\renewcommand{\arraystretch}{0.93}

\begin{center}
\begin{tikzpicture}[scale=1.25, font=\small]
	\draw[achse] (-0.4,0) -- (5,0) node[below left] {$x$};
	\draw[achse] (0,-0.4) -- (0,4.4) node[below left] {$y$};
	\begin{scope}
		\clip (-0.4,-0.4) rectangle (4.8,4.2);
		\draw[blue, very thick, domain=0:4.6, samples=60] plot (\x, {0.2*\x*\x + 0.4});
		% Sekante durch P(1|0,6) und Q(3,5|2,85), Steigung 0,9
		\draw[red, thick, dashed, domain=0:4.8] plot (\x, {0.9*(\x-1) + 0.6});
		% Tangente in P, Steigung 0,4
		\draw[red, very thick, domain=0:4.8] plot (\x, {0.4*(\x-1) + 0.6});
	\end{scope}
	\fill (1,0.6) circle (2pt) node[above left] {$P$};
	\fill (3.5,2.85) circle (2pt) node[above left] {$Q$};
	\draw[black!60] (1,0.6) -- node[below] {$h = \Delta x$} (3.5,0.6) -- node[right] {$\Delta y$} (3.5,2.85);
	\draw[black!40, dotted] (1,0.6) -- (1,0) node[below, black] {$x_0$};
	\draw[black!40, dotted] (3.5,0.6) -- (3.5,0) node[below, black] {$x_0 + h$};
	\node[red, align=left, anchor=west] at (4.85,3.8) {割線 -- Sekante};
	\node[red, align=left, anchor=west] at (4.85,2.2) {接線 -- Tangente};
\end{tikzpicture}
\end{center}

\begin{lesen}
\lies{$\dfrac{f(x_0 + h) - f(x_0)}{h}$}{平均変化率：2点 $P$、$Q$ を通る直線の傾き}{heikin henkaritsu: niten $P$, $Q$ o tōru chokusen no katamuki}{der Differenzenquotient: die Steigung der Sekante durch $P$ und $Q$}
\lies{$f'(x_0) = {}$ \newline $\displaystyle\lim_{h \to 0} \frac{f(x_0 + h) - f(x_0)}{h}$}{エフダッシュ $x_0$ は、$h$ が0に近づくときの平均変化率の極限値。図形的には、点 $P$ における接線の傾きである。}{efu dasshu $x_0$ wa, $h$ ga zero ni chikazuku toki no heikin henkaritsu no kyokugenchi. zukeiteki ni wa, ten $P$ ni okeru sessen no katamuki de aru.}{$f$ Strich von $x$ null ist der Grenzwert für $h$ gegen null von $f$ von $x$ null plus $h$ minus $f$ von $x$ null, alles durch $h$. Geometrisch ist das die Steigung der Tangente in $P$.}
\lies{$f'(x),\ f''(x),\ f^{(n)}(x)$}{第1次、第2次、第 $n$ 次導関数：エフダッシュエックス、エフツーダッシュエックス、エフエヌエックス}{dai ichi ji, dai ni ji, dai $n$ ji dōkansū: efu dasshu ekkusu, efu tsū dasshu ekkusu, efu enu ekkusu}{die erste, zweite, $n$-te Ableitung von $f$: $f$ Strich von $x$, $f$ zwei Strich von $x$, $f$ $n$ Strich von $x$}
\lies{$\dfrac{\mathrm{d}y}{\mathrm{d}x}$}{ディーワイディーエックス：$y$ を $x$ で微分したもの}{dī wai dī ekkusu: $y$ o $x$ de bibun shita mono}{d $y$ nach d $x$ -- der Differenzialquotient von $y$}
\lies{$\dfrac{\mathrm{d}^2 y}{\mathrm{d}x^2}$}{ディー2乗ワイディーエックス2乗}{dī nijō wai dī ekkusu nijō}{d zwei $y$ nach d $x$ Quadrat}
\lies{$\dfrac{\mathrm{d}^n y}{\mathrm{d}x^n}$}{ディー $n$ 乗ワイディーエックス $n$ 乗}{dī $n$ jō wai dī ekkusu $n$ jō}{d $n$ $y$ nach d $x$ hoch $n$}
\lies{$\dfrac{\partial z}{\partial x}$}{ラウンドディーゼットラウンドディーエックス：$z$ の $x$ に関する偏導関数}{raundo dī zetto raundo dī ekkusu: $z$ no $x$ ni kansuru hendōkansū}{d $z$ partiell nach d $x$}
\lies{$\dot{s} = \dfrac{\mathrm{d}s}{\mathrm{d}t}$}{エスドットイコールディーエスディーティー（時間に関する微分）}{esu dotto ikōru dī esu dī tī (jikan ni kansuru bibun)}{$s$ Punkt gleich d $s$ nach d $t$ (Ableitung nach der Zeit)}
\end{lesen}

\platz\zwischen{Beispiele}{例}{rei}

\begin{lesen}
\lies{$(x^3)' = 3x^2$}{$x$ 3乗の導関数は3$x$ 2乗}{$x$ sanjō no dōkansū wa san $x$ nijō}{die Ableitung von $x$ hoch drei ist drei $x$ Quadrat}
\lies{$(\sin x)' = \cos x$}{サイン $x$ を微分するとコサイン $x$}{sain $x$ o bibun suru to kosain $x$}{die (erste) Ableitung von Sinus $x$ ist Kosinus $x$}
\lies{$f(x) = x^2 - 4x$ \newline $f'(x) = 2x - 4$}{$f(x)$ を微分する：エフダッシュエックスイコール2$x$ マイナス4。}{$f(x)$ o bibun suru: efu dasshu ekkusu ikōru ni $x$ mainasu yon.}{Leite $f$ ab: $f$ Strich von $x$ gleich zwei $x$ minus vier.}
\lies{$2x - 4 = 0 \ \Rightarrow\ x_E = 2$}{導関数を0とおいて、$x$ について解く。}{dōkansū o zero to oite, $x$ ni tsuite toku.}{Setze die erste Ableitung gleich null und löse nach $x$ auf.}
\lies{$f''(2) = 2 > 0$}{$x$ イコール2での第2次導関数の値は正なので、グラフはそこで極小となる：$T(2,\,-4)$}{$x$ ikōru ni de no dai ni ji dōkansū no atai wa sei na no de, gurafu wa soko de kyokushō to naru: $T(2,\,-4)$}{Die zweite Ableitung an der Stelle zwei ist positiv, also hat der Graph dort einen Tiefpunkt: $T(2 \mid -4)$.}
\end{lesen}

%--------------------------------------------------------------------------------------------------
\newpage
\kapitel{Integralrechnung}{積分法}{sekibunhō}\renewcommand{\lesenabstand}{1.12}\renewcommand{\arraystretch}{0.93}

\begin{center}
\begin{tikzpicture}[scale=1.15, font=\small]
	% Untersumme (grau) und Obersumme (rot gestrichelt) für f(x) = 0,15x^2 + 0,6 auf [1; 5], n = 8
	\foreach \i in {0,...,7} {
		\pgfmathsetmacro{\xl}{1+0.5*\i} \pgfmathsetmacro{\xr}{1.5+0.5*\i}
		\pgfmathsetmacro{\yl}{0.15*\xl*\xl+0.6} \pgfmathsetmacro{\yr}{0.15*\xr*\xr+0.6}
		\filldraw[fill=gray!35, draw=gray] (\xl,0) rectangle (\xr,\yl);
		\draw[red, dashed] (\xl,\yl) -- (\xl,\yr) -- (\xr,\yr) -- (\xr,\yl);
	}
	\draw[achse] (-0.4,0) -- (6.2,0) node[below left] {$x$};
	\draw[achse] (0,-0.4) -- (0,5) node[below left] {$y$};
	\draw[blue, very thick, domain=0:5.6, samples=60] plot (\x, {0.15*\x*\x + 0.6}) node[right] {$f$};
	\draw[black!60] (1,0) -- (1,0.75); \draw[black!60] (5,0) -- (5,4.35);
	\node[below] at (1,0) {$a$}; \node[below] at (5,0) {$b$};
	\draw[|<->|, black!70] (3,-0.5) -- node[below] {$\Delta x$} (3.5,-0.5);
	\node[align=left, anchor=west] at (6.4,1.2) {\textcolor{gray}{\rule{8pt}{8pt}} 不足和 -- Untersumme\\[2pt] \textcolor{red}{\rule[2pt]{8pt}{0.8pt}} 過剰和 -- Obersumme};
\end{tikzpicture}
\end{center}

\begin{lesen}
\lies{$\displaystyle\int f(x)\dx = F(x) + C$, \newline wenn $F'(x) = f(x)$}{$F'(x)$ が $f(x)$ に等しいとき、インテグラルエフエックスディーエックスは、ラージエフエックスプラス $C$}{$F'(x)$ ga $f(x)$ ni hitoshii toki, integuraru efu ekkusu dī ekkusu wa, rāji efu ekkusu purasu $C$}{Integral $f$ von $x$ d $x$ ist gleich groß $F$ von $x$ plus groß $C$, wenn groß $F$ Strich von $x$ gleich $f$ von $x$ ist}
\lies{$\displaystyle\int f(x)\dx$}{$f$ の不定積分。$F$ は $f$ の原始関数、$C$ は積分定数}{$f$ no futei sekibun. $F$ wa $f$ no genshi kansū, $C$ wa sekibun teisū}{das unbestimmte Integral von $f$; $F$ ist eine Stammfunktion von $f$, $C$ die Integrationskonstante}
\lies{$\displaystyle\int_a^b f(x)\dx$}{インテグラル $a$ から $b$ までエフエックスディーエックス：下端 $a$、上端 $b$ の定積分}{integuraru $a$ kara $b$ made efu ekkusu dī ekkusu: katan $a$, jōtan $b$ no teisekibun}{Integral von $a$ bis $b$ über $f$ von $x$ d $x$ -- ein bestimmtes Integral mit der unteren Grenze $a$ und der oberen Grenze $b$}
\lies{$\big[F(x)\big]_a^b = F(b) - F(a)$}{$F(x)$ に $b$ を代入した値から、$a$ を代入した値を引く}{$F(x)$ ni $b$ o dainyū shita atai kara, $a$ o dainyū shita atai o hiku}{groß $F$ von $x$ in den Grenzen von $a$ bis $b$ ist gleich $F$ von $b$ minus $F$ von $a$}
\lies{$\displaystyle\lim_{\Delta x \to 0} \sum_{k} f(x_k)\,\Delta x$}{デルタ $x$ を0に近づけたときの、$f(x_k)$ 掛けるデルタ $x$ の和の極限値}{deruta $x$ o zero ni chikazuketa toki no, $f(x_k)$ kakeru deruta $x$ no wa no kyokugenchi}{der Grenzwert für Delta $x$ gegen null der Summe über $f$ von $x$ Index $k$ mal Delta $x$}
\lies{$\displaystyle\int_a^b f(x)\dx = A$}{$f(x) \ge 0$ のとき：$f$ のグラフ、$x$ 軸、直線 $x = a$ と $x = b$ で囲まれた部分の面積}{$f(x) \ge 0$ no toki: $f$ no gurafu, $x$ jiku, chokusen $x = a$ to $x = b$ de kakomareta bubun no menseki}{für $f(x) \ge 0$: der Inhalt der Fläche zwischen dem Graphen von $f$, der $x$-Achse und den Geraden $x = a$ und $x = b$}
\end{lesen}

\platz\zwischen{Beispiele}{例}{rei}

\begin{lesen}
\lies{$\displaystyle\int \cos x \dx = \sin x + C$}{インテグラルコサイン $x$ ディーエックスイコールサイン $x$ プラス $C$}{integuraru kosain $x$ dī ekkusu ikōru sain $x$ purasu $C$}{Integral Kosinus $x$ d $x$ gleich Sinus $x$ plus $C$}
\lies{$\displaystyle\int_0^2 x^2 \dx = \left[\tfrac13 x^3\right]_0^2 = \tfrac83$}{$x$ 2乗を0から2まで積分すると、3分の8}{$x$ nijō o zero kara ni made sekibun suru to, san bun no hachi}{Integral von null bis zwei über $x$ Quadrat d $x$ gleich ein Drittel $x$ hoch drei in den Grenzen von null bis zwei, also acht Drittel}
\lies{$\displaystyle\int_1^{\mathrm{e}} \frac{1}{x} \dx = \big[\ln x\big]_1^{\mathrm{e}} = 1$}{$x$ 分の1を1から $\mathrm{e}$ まで積分すると、1}{$x$ bun no ichi o ichi kara $\mathrm{e}$ made sekibun suru to, ichi}{Integral von eins bis e über eins durch $x$ d $x$ gleich ln $x$ in den Grenzen von eins bis e, also eins}
\end{lesen}

\anmerkung{Wichtige Anwendungen der Integralrechnung in der Geometrie: Flächeninhalte
(\jp{面積}{menseki}), Volumina (\jp{体積}{taiseki}) und Bogenlängen (\jp{弧の長さ}{ko no nagasa}).
Das Integralzeichen liest man \jp{インテグラル}{integuraru}, den Strich der Ableitung \jp{ダッシュ}{dasshu}.}

%--------------------------------------------------------------------------------------------------
\newpage
\kapitel{Geraden und Winkel}{直線と角}{chokusen to kaku}\renewcommand{\lesenabstand}{1.12}\renewcommand{\arraystretch}{0.93}

\begin{center}
\begin{tikzpicture}[font=\small]
	\draw[mass] (1.3,0) arc[start angle=0, end angle=40, radius=1.3]; \node at (20:0.9) {$\alpha$};
	\draw[blue, very thick] (4.6,0) -- (0,0) -- (40:4.6);
	\fill (0,0) circle (2pt) node[above left] {$S$}; \fill (3.2,0) circle (2pt) node[below] {$A$}; \fill (40:3.2) circle (2pt) node[above left] {$B$};
	\node[align=center] (v) at (-2.2,-0.9) {頂点\\\color{black!60}Scheitel};
	\draw[-{Stealth}, thin] (v.east) -- (-0.15,-0.1);
	\node[align=center] (l) at (6.6,0.9) {辺\\\color{black!60}Schenkel};
	\draw[-{Stealth}, thin] (l.west) -- (4.3,0.1); \draw[-{Stealth}, thin] (l.west) -- ($(40:4.3)+(0.15,-0.1)$);
\end{tikzpicture}
\end{center}

\begin{lesen}
\lies{$g$}{直線 $g$}{chokusen $g$}{die Gerade $g$}
\lies{$\overline{AB}$}{線分 $AB$ \;/\; $A$ から $B$ までの距離}{senbun $AB$ \;/\; $A$ kara $B$ made no kyori}{die Strecke $A$ $B$}
\lies{$g \parallel h$ \newline \footnotesize jap. $g \mathbin{/\!/} h$}{$g$ 平行 $h$ \;/\; $g$ は $h$ に平行である}{$g$ heikō $h$ \;/\; $g$ wa $h$ ni heikō de aru}{$g$ ist parallel zu $h$}
\lies{$g \perp h$}{$g$ 垂直 $h$ \;/\; $g$ は $h$ に垂直である}{$g$ suichoku $h$ \;/\; $g$ wa $h$ ni suichoku de aru}{$g$ ist senkrecht zu $h$ \;/\; $g$ steht senkrecht auf $h$}
\lies{$\sphericalangle ASB$ \newline \footnotesize jap. $\angle ASB$}{角 $ASB$（頂点を真ん中に書く）}{kaku $ASB$ (chōten o mannaka ni kaku)}{Winkel $A$ $S$ $B$ (der Scheitel steht in der Mitte)}
\lies{$\alpha = 40\degree$}{アルファは40度：鋭角}{arufa wa yonjū do: eikaku}{Alpha gleich vierzig Grad: ein spitzer Winkel}
\lies{$\alpha + \beta = 90\degree$}{アルファとベータは互いに余角：和が90度}{arufa to bēta wa tagai ni yokaku: wa ga kyūjū do}{Alpha und Beta sind Komplementwinkel: Sie ergänzen sich zu neunzig Grad}
\lies{$\alpha + \beta = 180\degree$}{アルファとベータは互いに補角：和が180度}{arufa to bēta wa tagai ni hokaku: wa ga hyakuhachijū do}{Alpha und Beta sind Supplementwinkel: Sie ergänzen sich zu hundertachtzig Grad}
\end{lesen}

\platz\zwischen{Winkel an geschnittenen Parallelen}{平行線と角}{heikōsen to kaku}

\begin{center}
\begin{tikzpicture}[font=\small, scale=1.1]
	\coordinate (S1) at (0.6,2.2); \coordinate (S2) at (-0.67,0);
	\draw[red, thin] (S1) circle (0.42); \draw[red, thin] (S2) circle (0.42);
	\draw[blue, very thick] (-4,2.2) -- (4,2.2) node[right] {$g_1$};
	\draw[blue, very thick] (-4,0) -- (4,0) node[right] {$g_2$};
	\draw[blue, very thick] ($(S2)+(240:1.5)$) -- ($(S1)+(60:1.5)$) node[right] {$t$};
	\foreach \w/\n in {30/\alpha, 120/\beta, 210/\gamma, 300/\delta} {
		\node at ($(S1)+(\w:0.78)$) {$\n_1$}; \node at ($(S2)+(\w:0.78)$) {$\n_2$};}
	\node[align=left, anchor=west] at (5,1.1) {$g_1 \parallel g_2$\\ $t$: 2直線に交わる直線\\\color{black!60}Transversale (schneidende Gerade)};
\end{tikzpicture}
\end{center}

\begin{lesen}
\lies{$\alpha_1$ und $\alpha_2$}{同位角：2直線が平行ならば等しい}{dōikaku: nichokusen ga heikō naraba hitoshii}{Stufenwinkel -- an Parallelen gleich groß}
\lies{$\gamma_1$ und $\alpha_2$}{錯角：2直線が平行ならば等しい}{sakkaku: nichokusen ga heikō naraba hitoshii}{Wechselwinkel -- an Parallelen gleich groß}
\lies{$\delta_1$ und $\alpha_2$}{同側内角：2直線が平行ならば、和は $180\degree$}{dōsoku naikaku: nichokusen ga heikō naraba, wa wa $180\degree$}{entgegengesetzt liegende Winkel -- ergänzen sich an Parallelen zu $180\degree$}
\lies{$\alpha_1$ und $\gamma_1$}{対頂角：常に等しい}{taichōkaku: tsune ni hitoshii}{Scheitelwinkel -- immer gleich groß}
\lies{$\alpha_1$ und $\beta_1$}{隣接補角：和は $180\degree$}{rinsetsu hokaku: wa wa $180\degree$}{Nebenwinkel -- ergänzen sich zu $180\degree$}
\end{lesen}

\platz\zwischen{Beispiele}{例}{rei}

\begin{lesen}
\lies{$\alpha_1 = 65\degree$ \newline $\beta_1 = 180\degree - 65\degree = 115\degree$}{アルファ1は65度。ベータ1はその隣接補角なので、115度。}{arufa ichi wa rokujūgo do. bēta ichi wa sono rinsetsu hokaku na no de, hyakujūgo do.}{Alpha eins beträgt fünfundsechzig Grad. Beta eins ist der Nebenwinkel, also hundertfünfzehn Grad.}
\lies{$\gamma_1 = \alpha_1 = 65\degree$}{ガンマ1とアルファ1は対頂角なので等しい。}{ganma ichi to arufa ichi wa taichōkaku na no de hitoshii.}{Gamma eins und Alpha eins sind Scheitelwinkel, also gleich groß.}
\lies{$\alpha_2 = \alpha_1 = 65\degree$}{アルファ2とアルファ1は平行線の同位角なので等しい。}{arufa ni to arufa ichi wa heikōsen no dōikaku na no de hitoshii.}{Alpha zwei und Alpha eins sind Stufenwinkel an Parallelen, also gleich groß.}
\lies{$\alpha = 90\degree$}{直線 $g$ と $h$ は垂直に交わる。}{chokusen $g$ to $h$ wa suichoku ni majiwaru.}{Die Geraden $g$ und $h$ schneiden einander rechtwinklig.}
\end{lesen}

%--------------------------------------------------------------------------------------------------
\newpage
\kapitel{Dreiecke, Vierecke und Vielecke}{三角形、四角形、多角形}{sankakukei, shikakukei, takakukei}\renewcommand{\lesenabstand}{1.12}\renewcommand{\arraystretch}{0.93}

\begin{center}
\begin{tikzpicture}[font=\small, scale=1.1]
	\draw[mass] (0.7,0) arc[start angle=0, end angle=59.04, radius=0.7]; \node at (28:0.45) {$\alpha$};
	\draw[mass] ($(5,0)+(136.85:0.7)$) arc[start angle=136.85, end angle=180, radius=0.7]; \node at ($(5,0)+(160:1.0)$) {$\beta$};
	\draw[mass] ($(1.8,3)+(239.04:0.6)$) arc[start angle=239.04, end angle=316.85, radius=0.6]; \node at ($(1.8,3)+(280:0.9)$) {$\gamma$};
	\draw[mass] (5.9,0) arc[start angle=0, end angle=136.85, radius=0.9]; \node at ($(5,0)+(60:1.2)$) {$\delta$};
	\draw[black!60, thick, dashed] (1.8,3) -- node[right, pos=0.6] {$h_c$} (1.8,0);
	\draw[black!60, thick] (2.15,0) arc[start angle=0, end angle=90, radius=0.35]; \fill[black!60] (1.95,0.15) circle (0.8pt);
	\draw[blue, very thick] (0,0) -- node[below] {$c$} (5,0) -- node[above right] {$a$} (1.8,3) -- node[above left] {$b$} cycle;
	\draw[blue, thick, dashed] (5,0) -- (6.6,0);
	\node[below left] at (0,0) {$A$}; \node[below] at (5,0) {$B$}; \node[above] at (1.8,3) {$C$};
	\node[align=left, anchor=west] at (7.4,2.2) {$a,\ b,\ c$: 辺 -- \color{black!60}Seiten};
	\node[align=left, anchor=west] at (7.4,1.6) {$\alpha,\ \beta,\ \gamma$: 内角 -- \color{black!60}Innenwinkel};
	\node[align=left, anchor=west] at (7.4,1.0) {$\delta$: 外角 -- \color{black!60}Außenwinkel};
	\node[align=left, anchor=west] at (7.4,0.4) {$h_c$: 高さ -- \color{black!60}Höhe};
\end{tikzpicture}
\end{center}

\begin{lesen}
\lies{$\triangle ABC$}{三角形 $ABC$}{sankakukei $ABC$}{Dreieck $A$ $B$ $C$}
\lies{$\alpha + \beta + \gamma = 180\degree$}{三角形の内角の和は180度}{sankakukei no naikaku no wa wa hyakuhachijū do}{die Innenwinkel eines Dreiecks ergeben zusammen hundertachtzig Grad}
\lies{$\delta = \alpha + \gamma$}{外角は、それと隣り合わない2つの内角の和に等しい}{gaikaku wa, sore to tonariawanai futatsu no naikaku no wa ni hitoshii}{ein Außenwinkel ist so groß wie die beiden nicht anliegenden Innenwinkel zusammen}
\lies{$A = \dfrac{1}{2} \cdot c \cdot h_c$}{面積は、底辺掛ける高さ割る2}{menseki wa, teihen kakeru takasa waru ni}{der Flächeninhalt ist gleich ein Halb mal Grundseite mal Höhe}
\lies{$a^2 + b^2 = c^2$}{$a$ 2乗プラス $b$ 2乗イコール $c$ 2乗。三平方の定理：直角三角形では、斜辺の2乗は、他の2辺の2乗の和に等しい}{$a$ nijō purasu $b$ nijō ikōru $c$ nijō. sanheihō no teiri: chokkaku sankakukei de wa, shahen no nijō wa, hoka no nihen no nijō no wa ni hitoshii}{$a$ Quadrat plus $b$ Quadrat gleich $c$ Quadrat -- Satz des Pythagoras: Im rechtwinkligen Dreieck ist das Hypotenusenquadrat so groß wie die beiden Kathetenquadrate zusammen}
\lies{$\triangle ABC \cong \triangle DEF$ \newline \footnotesize jap. $\triangle ABC \equiv \triangle DEF$}{三角形 $ABC$ と三角形 $DEF$ は合同である}{sankakukei $ABC$ to sankakukei $DEF$ wa gōdō de aru}{Dreieck $A$ $B$ $C$ ist kongruent zu Dreieck $D$ $E$ $F$}
\lies{$\triangle ABC \sim \triangle PQR$ \newline \footnotesize jap. $\triangle ABC \backsim \triangle PQR$}{三角形 $ABC$ と三角形 $PQR$ は相似である}{sankakukei $ABC$ to sankakukei $PQR$ wa sōji de aru}{Dreieck $A$ $B$ $C$ ist ähnlich zu Dreieck $P$ $Q$ $R$}
\end{lesen}

\platz\zwischen{Besondere Dreiecke}{特別な三角形}{tokubetsu na sankakukei}

\begin{lesen}
\lies{\skGleichschenklig}{二等辺三角形：2辺が等しい。残りの辺が底辺で、2つの底角は等しい}{nitōhen sankakukei: nihen ga hitoshii. nokori no hen ga teihen de, futatsu no teikaku wa hitoshii}{ein gleichschenkliges Dreieck: zwei gleich lange Seiten (Schenkel), die dritte Seite ist die Basis; die Basiswinkel sind gleich groß}
\lies{\skRechtwinklig}{直角三角形：直角に向かい合う辺が斜辺}{chokkaku sankakukei: chokkaku ni mukaiau hen ga shahen}{ein rechtwinkliges Dreieck: Die Katheten schließen den rechten Winkel ein, die Hypotenuse liegt ihm gegenüber}
\end{lesen}

\begin{center}
\begin{tikzpicture}[font=\small, scale=1.0]
	\coordinate (A) at (0,0); \coordinate (B) at (3.6,0); \coordinate (C) at (1.8,3.118); \coordinate (M) at (1.8,1.039);
	\draw[red, thick] (M) circle (1.039); \draw[red, thick] (M) circle (2.078);
	\foreach \p/\q in {A/{2.7,1.559}, B/{0.9,1.559}, C/{1.8,0}} {\draw[black!60, dash dot] ($(\p)!-0.25!(\q)$) -- ($(\p)!1.45!(\q)$);}
	\draw[blue, very thick] (A) -- (B) -- (C) -- cycle; \fill (M) circle (2pt) node[right=2pt] {$M$};
	\node[align=left, anchor=west] at (5.2,3.1) {正三角形と3本の対称の軸\\\color{black!60}ein gleichseitiges Dreieck mit seinen drei Symmetrieachsen};
	\node[align=left, anchor=west] at (5.2,1.9) {\textcolor{red}{内接円}と\textcolor{red}{外接円} --
		\color{black!60}Inkreis und Umkreis};
	\node[align=left, anchor=west, text width=10cm] at (5.2,0.2) {点 $M$ は、垂線、垂直二等分線、中線、角の二等分線の交点である。\\\color{black!60}$M$ ist der Schnittpunkt der Höhen, der Mittelsenkrechten, der Seitenhalbierenden und der Winkelhalbierenden.};
\end{tikzpicture}
\end{center}

\platz\zwischen{Vierecke und regelmäßige Vielecke}{四角形と正多角形}{shikakukei to seitakakukei}

\begin{lesen}
\lies{\skTrapez}{台形：1組の向かい合う辺が平行な四角形}{daikei: hitokumi no mukaiau hen ga heikō na shikakukei}{ein Trapez: ein Viereck mit einem Paar paralleler Seiten}
\lies{\skParallelogramm}{平行四辺形：2組の向かい合う辺がそれぞれ平行な四角形}{heikō shihenkei: futakumi no mukaiau hen ga sorezore heikō na shikakukei}{ein Parallelogramm: ein Viereck, dessen gegenüberliegende Seiten parallel sind}
\lies{\skRaute}{ひし形：4つの辺がすべて等しい四角形}{hishigata: yottsu no hen ga subete hitoshii shikakukei}{eine Raute (ein Rhombus): ein Parallelogramm mit vier gleich langen Seiten}
\lies{\skRechteck}{長方形：4つの角がすべて直角の四角形}{chōhōkei: yottsu no kaku ga subete chokkaku no shikakukei}{ein Rechteck: ein Viereck mit vier rechten Winkeln}
\lies{\skQuadrat}{正方形：4つの辺がすべて等しい長方形}{seihōkei: yottsu no hen ga subete hitoshii chōhōkei}{ein Quadrat: ein Rechteck mit vier gleich langen Seiten}
\lies{\skNeck{5} \skNeck{6} \skNeck{8}}{正多角形：正五角形、正六角形、正八角形（辺も角もすべて等しい）}{seitakakukei: seigokakukei, seirokkakukei, seihakkakukei (hen mo kaku mo subete hitoshii)}{regelmäßige Vielecke: ein Fünfeck, ein Sechseck und ein Achteck (alle Seiten und alle Winkel sind gleich groß)}
\end{lesen}

\platz\zwischen{Beispiele}{例}{rei}

\begin{lesen}
\lies{$a = 3\,\mathrm{cm},\ b = 4\,\mathrm{cm}$ \newline $c = \sqrt{3^2 + 4^2}\,\mathrm{cm} = 5\,\mathrm{cm}$}{直角をはさむ2辺が3センチメートルと4センチメートルの直角三角形では、斜辺はルート3の2乗プラス4の2乗で、5センチメートル。}{chokkaku o hasamu nihen ga san senchimētoru to yon senchimētoru no chokkaku sankakukei de wa, shahen wa rūto san no nijō purasu yon no nijō de, go senchimētoru.}{In einem rechtwinkligen Dreieck sind die Katheten drei und vier Zentimeter lang. Die Hypotenuse ist die Wurzel aus drei Quadrat plus vier Quadrat, also fünf Zentimeter.}
\lies{$A = \frac12 \cdot 8\,\mathrm{cm} \cdot 5\,\mathrm{cm} = 20\,\mathrm{cm}^2$}{三角形の面積は、8センチメートル掛ける5センチメートル割る2で、20平方センチメートル。}{sankakukei no menseki wa, hachi senchimētoru kakeru go senchimētoru waru ni de, nijū heihō senchimētoru.}{Der Flächeninhalt des Dreiecks ist ein Halb mal acht Zentimeter mal fünf Zentimeter, also zwanzig Quadratzentimeter.}
\lies{$u = 2 \cdot 6\,\mathrm{cm} + 2 \cdot 4\,\mathrm{cm} = 20\,\mathrm{cm}$}{縦4センチメートル、横6センチメートルの長方形の周の長さは、20センチメートル。}{tate yon senchimētoru, yoko roku senchimētoru no chōhōkei no shū no nagasa wa, nijū senchimētoru.}{Der Umfang eines Rechtecks mit sechs Zentimetern Länge und vier Zentimetern Breite beträgt zwanzig Zentimeter.}
\lies{$(6 - 2) \cdot 180\degree = 720\degree$}{六角形の内角の和は720度。}{rokkakukei no naikaku no wa wa nanahyakunijū do.}{Die Innenwinkel eines Sechsecks ergeben zusammen siebenhundertzwanzig Grad.}
\end{lesen}

%--------------------------------------------------------------------------------------------------
\newpage
\kapitel{Der Kreis und seine Teile}{円とその部分}{en to sono bubun}\renewcommand{\lesenabstand}{1.12}\renewcommand{\arraystretch}{0.93}

\begin{center}
\begin{tikzpicture}[font=\small, scale=1.2]
	\fill[gray!40] (120:2) arc[start angle=120, end angle=200, radius=2] -- cycle;
	\fill[blue!15] (0,0) -- (300:2) arc[start angle=300, end angle=345, radius=2] -- cycle;
	\draw[blue, very thick] (0,0) circle (2);
	\draw[blue, thick] (120:2) -- node[right, black, inner sep=2pt, pos=0.2] {$s$} (200:2);
	\draw[blue, thick] (300:2) -- (0,0) -- (345:2);
	\draw[black!70, thick] (90:2) -- node[left, pos=0.75] {$d$} (270:2);
	\draw[black!70, thick] (0,0) -- node[above] {$r$} (15:2);
	\draw[red, line width=2pt] (35:2) arc[start angle=35, end angle=75, radius=2]; \node[red] at (55:2.3) {$b$};
	\draw[red, thick] (-2.7,2) -- (2.7,2) node[right] {$t$};
	\fill[red] (0,2) circle (2pt) node[above] {$B$};
	\draw[red, thick] ($(215:2)!-0.5!(255:2)$) -- ($(215:2)!1.6!(255:2)$) node[right] {$g$};
	\fill (0,0) circle (2pt) node[left] {$M$};
	\node at (160:1.75) {$A_1$}; \node at (322:1.3) {$A_2$};
\end{tikzpicture}
\end{center}

\begin{lesen}
\lies{$M$}{中心}{chūshin}{der Mittelpunkt}
\lies{$r$}{半径}{hankei}{der Radius -- Plural: Radien}
\lies{$d = 2r$}{直径：直径は半径の2倍}{chokkei: chokkei wa hankei no nibai}{der Durchmesser: Er ist doppelt so groß wie der Radius}
\lies{$s$}{弦：円周上の2点を結ぶ線分}{gen: enshūjō no niten o musubu senbun}{eine Sehne: eine Strecke zwischen zwei Punkten des Kreises}
\lies{$g$}{割線：円と2点で交わる直線}{kassen: en to niten de majiwaru chokusen}{eine Sekante: eine Gerade, die den Kreis in zwei Punkten schneidet}
\lies{$t$, $B$}{接線：円と1点（接点 $B$）だけを共有する直線。接線は、接点を通る半径に垂直}{sessen: en to itten (setten $B$) dake o kyōyū suru chokusen. sessen wa, setten o tōru hankei ni suichoku}{eine Tangente: eine Gerade, die den Kreis in einem Punkt berührt, dem Berührungspunkt $B$; sie steht senkrecht auf dem Berührungsradius}
\lies{$b$}{弧}{ko}{ein Kreisbogen}
\lies{$A_1$}{弓形：弦と弧で囲まれた部分}{kyūkei: gen to ko de kakomareta bubun}{ein Kreisabschnitt (Kreissegment): die Fläche zwischen Sehne und Bogen}
\lies{$A_2$}{扇形：2本の半径と弧で囲まれた部分}{ōgigata: nihon no hankei to ko de kakomareta bubun}{ein Kreisausschnitt (Kreissektor): die Fläche zwischen zwei Radien und einem Bogen}
\lies{$u = 2\pi r$}{円周は2パイ $r$}{enshū wa ni pai $r$}{der Umfang ist gleich zwei pi $r$}
\lies{$A = \pi r^2$}{面積はパイ $r$ 2乗}{menseki wa pai $r$ nijō}{der Flächeninhalt ist gleich pi $r$ Quadrat}
\end{lesen}

\platz\zwischen{Beispiele}{例}{rei}

\begin{lesen}
\lies{$r = 5\,\mathrm{cm}$ \newline $u = 2\pi \cdot 5\,\mathrm{cm} \approx 31{,}4\,\mathrm{cm}$}{半径が5センチメートルなので、円周は2パイ掛ける5センチメートルで、約31.4センチメートル。}{hankei ga go senchimētoru na no de, enshū wa ni pai kakeru go senchimētoru de, yaku sanjūitten yon senchimētoru.}{Der Radius beträgt fünf Zentimeter, also ist der Umfang zwei pi mal fünf Zentimeter, ungefähr einunddreißig Komma vier Zentimeter.}
\lies{$A = \pi \cdot (5\,\mathrm{cm})^2 \approx 78{,}5\,\mathrm{cm}^2$}{面積はパイ掛ける5センチメートルの2乗で、約78.5平方センチメートル。}{menseki wa pai kakeru go senchimētoru no nijō de, yaku nanajūhatten go heihō senchimētoru.}{Der Flächeninhalt ist pi mal fünf Zentimeter zum Quadrat, ungefähr achtundsiebzig Komma fünf Quadratzentimeter.}
\lies{$b = \dfrac{60\degree}{360\degree} \cdot 2\pi \cdot 5\,\mathrm{cm} \approx 5{,}24\,\mathrm{cm}$}{中心角60度に対する弧は円周の6分の1で、約5.24センチメートル。}{chūshinkaku rokujū do ni taisuru ko wa enshū no roku bun no ichi de, yaku go ten ni yon senchimētoru.}{Der Bogen zu einem Winkel von sechzig Grad ist ein Sechstel des Umfangs, etwa fünf Komma zwei vier Zentimeter.}
\end{lesen}

%--------------------------------------------------------------------------------------------------
\newpage
\kapitel{Körper berechnen}{立体の体積と表面積}{rittai no taiseki to hyōmenseki}\renewcommand{\lesenabstand}{1.0}\renewcommand{\arraystretch}{0.93}

\begin{center}
\begin{tikzpicture}[font=\small]
	\def\dx{1.2}\def\dy{0.9}
	\draw[blue, thick, dashed] (0,0) -- (\dx,\dy) -- (4+\dx,\dy) (\dx,\dy) -- (\dx,2.5+\dy);
	\draw[red, thick, dashed] (0,0) -- (4+\dx,\dy);
	\draw[red, very thick] (0,0) -- (4+\dx,2.5+\dy);
	\draw[blue, very thick] (0,0) rectangle (4,2.5) (4,0) -- (4+\dx,\dy) -- (4+\dx,2.5+\dy) -- (\dx,2.5+\dy) -- (0,2.5) (4,2.5) -- (4+\dx,2.5+\dy);
	\node[below] at (2,0) {$a$}; \node[below right] at (4.55,0.4) {$b$}; \node[left] at (0,1.25) {$c$};
	\node[red, above left] at (2.4,1.6) {$d$}; \node[red, below] at (3.1,0.5) {$e$};
	\node[align=center] at (2.6,-0.9) {直方体 -- \color{black!60}Quader};
	\begin{scope}[shift={(9,0.2)}]
		\draw[blue, thick, dashed] (1.2,0) arc[start angle=0, end angle=180, x radius=1.2, y radius=0.35];
		\draw[blue, very thick] (-1.2,0) arc[start angle=180, end angle=360, x radius=1.2, y radius=0.35];
		\draw[blue, very thick] (0,2.8) ellipse[x radius=1.2, y radius=0.35];
		\draw[blue, very thick] (-1.2,0) -- (-1.2,2.8) (1.2,0) -- (1.2,2.8);
		\draw[red, thick] (0,2.8) -- node[above, inner sep=1.5pt] {$r$} (1.2,2.8); \fill (0,2.8) circle (1.5pt);
		\draw[red, thick, |<->|] (1.6,0) -- node[right] {$h$} (1.6,2.8);
		\node[align=center] at (0,-1.1) {円柱 -- \color{black!60}Zylinder};
	\end{scope}
\end{tikzpicture}
\end{center}

\begin{lesen}
\lies{Quader}{直方体には、頂点が8つ、辺が12本、面が6つある。}{chokuhōtai ni wa, chōten ga yattsu, hen ga jūnihon, men ga muttsu aru.}{Ein Quader hat acht Ecken, zwölf Kanten und sechs Flächen.}
\lies{$e$, $d$}{$e$ は面の対角線、$d$ は直方体の対角線}{$e$ wa men no taikakusen, $d$ wa chokuhōtai no taikakusen}{$e$ ist eine Flächendiagonale, $d$ ist eine Raumdiagonale}
\lies{$V = a \cdot b \cdot c$}{体積は、縦掛ける横掛ける高さ}{taiseki wa, tate kakeru yoko kakeru takasa}{das Volumen ist gleich Länge mal Breite mal Höhe}
\lies{$O = 2(ab + ac + bc)$}{表面積は、2掛けるかっこ $ab$ プラス $ac$ プラス $bc$ かっことじ}{hyōmenseki wa, ni kakeru kakko $ab$ purasu $ac$ purasu $bc$ kakkotoji}{der Oberflächeninhalt ist gleich zwei mal Klammer auf $a$ $b$ plus $a$ $c$ plus $b$ $c$ Klammer zu}
\lies{$V = a^3$}{立方体の体積は $a$ の3乗}{rippōtai no taiseki wa $a$ no sanjō}{das Volumen eines Würfels ist gleich $a$ hoch drei}
\lies{$V = \pi r^2 h$}{円柱の体積はパイ $r$ 2乗 $h$（底面積掛ける高さ）}{enchū no taiseki wa pai $r$ nijō $h$ (teimenseki kakeru takasa)}{das Volumen eines Zylinders ist gleich pi $r$ Quadrat $h$ (Grundfläche mal Höhe)}
\lies{$M = 2\pi r h$}{円柱の側面積は2パイ $r$ $h$}{enchū no sokumenseki wa ni pai $r$ $h$}{der Mantelflächeninhalt eines Zylinders ist gleich zwei pi $r$ $h$}
\lies{$V = \dfrac13 \cdot G \cdot h$}{角錐や円錐の体積は、3分の1掛ける底面積掛ける高さ}{kakusui ya ensui no taiseki wa, san bun no ichi kakeru teimenseki kakeru takasa}{das Volumen einer Pyramide oder eines Kegels ist gleich ein Drittel mal Grundfläche mal Höhe}
\lies{$V = \dfrac43 \pi r^3$}{球の体積は3分の4パイ $r$ 3乗}{kyū no taiseki wa san bun no yon pai $r$ sanjō}{das Volumen einer Kugel ist gleich vier Drittel pi $r$ hoch drei}
\lies{$O = 4\pi r^2$}{球の表面積は4パイ $r$ 2乗}{kyū no hyōmenseki wa yon pai $r$ nijō}{der Oberflächeninhalt einer Kugel ist gleich vier pi $r$ Quadrat}
\end{lesen}

\platz\zwischen{Besondere Körper}{特別な立体}{tokubetsu na rittai}

\begin{lesen}
\lies{\skPyramidenstumpf \; \skKegelstumpf}{角錐台と円錐台：底面に平行な平面で、頂点の側を切り取った立体}{kakusuidai to ensuidai: teimen ni heikō na heimen de, chōten no gawa o kiritotta rittai}{ein Pyramidenstumpf und ein Kegelstumpf: Die Spitze wurde parallel zur Grundfläche abgeschnitten}
\lies{\skTetraeder}{正四面体：4つの面がすべて正三角形の多面体}{seishimentai: yottsu no men ga subete seisankakukei no tamentai}{ein (regelmäßiges) Tetraeder: ein Polyeder, dessen vier Flächen gleichseitige Dreiecke sind}
\lies{\skRotation}{回転体：平面図形を軸のまわりに回転させてできる立体。ここでは半円から球ができる}{kaitentai: heimen zukei o jiku no mawari ni kaiten sasete dekiru rittai. koko de wa han'en kara kyū ga dekiru}{ein Rotationskörper: Er entsteht, wenn eine ebene Fläche um eine Achse rotiert -- hier erzeugt ein Halbkreis eine Kugel}
\end{lesen}

\platz\zwischen{Beispiele}{例}{rei}

\begin{lesen}
\lies{$a = 4\,\mathrm{cm},\ b = 3\,\mathrm{cm},\ c = 2\,\mathrm{cm}$ \newline $V = 24\,\mathrm{cm}^3$}{直方体の体積は4掛ける3掛ける2で、24立方センチメートル。}{chokuhōtai no taiseki wa yon kakeru san kakeru ni de, nijūyon rippō senchimētoru.}{Das Volumen des Quaders ist vier mal drei mal zwei, also vierundzwanzig Kubikzentimeter.}
\lies{$O = 2 \cdot (12 + 8 + 6)\,\mathrm{cm}^2 = 52\,\mathrm{cm}^2$}{表面積は52平方センチメートル。}{hyōmenseki wa gojūni heihō senchimētoru.}{Sein Oberflächeninhalt beträgt zweiundfünfzig Quadratzentimeter.}
\lies{$d = \sqrt{4^2 + 3^2 + 2^2}\,\mathrm{cm} = \sqrt{29}\,\mathrm{cm} \approx 5{,}39\,\mathrm{cm}$}{対角線の長さはルート29で、約5.39センチメートル。}{taikakusen no nagasa wa rūto nijūkyū de, yaku go ten san kyū senchimētoru.}{Die Raumdiagonale ist die Wurzel aus neunundzwanzig, ungefähr fünf Komma drei neun Zentimeter.}
\lies{$r = 3\,\mathrm{cm},\ h = 10\,\mathrm{cm}$ \newline $V = \pi \cdot 9 \cdot 10\,\mathrm{cm}^3 \approx 282{,}7\,\mathrm{cm}^3$}{円柱の体積は90パイで、約282.7立方センチメートル。}{enchū no taiseki wa kyūjū pai de, yaku nihyakuhachijūni ten nana rippō senchimētoru.}{Das Volumen des Zylinders ist neunzig pi, ungefähr zweihundertzweiundachtzig Komma sieben Kubikzentimeter.}
\lies{$r = 3\,\mathrm{cm}$ \newline $V = \frac43 \pi \cdot 27\,\mathrm{cm}^3 \approx 113{,}1\,\mathrm{cm}^3$}{球の体積は36パイで、約113.1立方センチメートル。}{kyū no taiseki wa sanjūroku pai de, yaku hyakujūsan ten ichi rippō senchimētoru.}{Das Volumen der Kugel ist sechsunddreißig pi, ungefähr hundertdreizehn Komma eins Kubikzentimeter.}
\end{lesen}

%--------------------------------------------------------------------------------------------------
\newpage
\kapitel{Trigonometrie im Dreieck}{三角比}{sankakuhi}\renewcommand{\lesenabstand}{1.12}\renewcommand{\arraystretch}{0.93}

\begin{center}
\begin{tikzpicture}[font=\small, scale=1.1]
	\draw[mass] (1.3,0) arc[start angle=0, end angle=30.96, radius=1.3]; \node at (14:0.95) {$\alpha$};
	\draw[black!60, thick] (5,0.45) arc[start angle=90, end angle=180, radius=0.45]; \fill[black!60] (4.82,0.18) circle (1pt);
	\draw[mass] ($(5,3)+(210.96:0.8)$) arc[start angle=210.96, end angle=270, radius=0.8]; \node at ($(5,3)+(242:0.5)$) {$\beta$};
	\draw[blue, very thick] (0,0) -- (5,0) -- (5,3) -- cycle;
	\node[below left] at (0,0) {$A$}; \node[below right] at (5,0) {$C$}; \node[above right] at (5,3) {$B$};
	\node[align=center, above left] at (2.4,1.6) {$c$: 斜辺\\\color{black!60}Hypotenuse};
	\node[align=left, right] at (5.15,1.5) {$a$: $\alpha$ の対辺\\\color{black!60}Gegenkathete von $\alpha$};
	\node[align=center, below] at (2.5,-0.1) {$b$: $\alpha$ の隣辺\\\color{black!60}Ankathete von $\alpha$};
\end{tikzpicture}
\end{center}

\begin{lesen}
\lies{$\sin\alpha = \dfrac{a}{c}$}{サインアルファイコール斜辺分の対辺}{sain arufa ikōru shahen bun no taihen}{Sinus Alpha gleich Gegenkathete durch Hypotenuse}
\lies{$\cos\alpha = \dfrac{b}{c}$}{コサインアルファイコール斜辺分の隣辺}{kosain arufa ikōru shahen bun no rinpen}{Kosinus Alpha gleich Ankathete durch Hypotenuse}
\lies{$\tan\alpha = \dfrac{a}{b}$}{タンジェントアルファイコール隣辺分の対辺}{tanjento arufa ikōru rinpen bun no taihen}{Tangens Alpha gleich Gegenkathete durch Ankathete}
\lies{$\cot\alpha = \dfrac{b}{a}$}{コタンジェントアルファイコール対辺分の隣辺}{kotanjento arufa ikōru taihen bun no rinpen}{Kotangens Alpha gleich Ankathete durch Gegenkathete}
\lies{$\dfrac{a}{\sin\alpha} = \dfrac{b}{\sin\beta} = \dfrac{c}{\sin\gamma}$}{正弦定理：サインアルファ分の $a$ イコールサインベータ分の $b$ イコールサインガンマ分の $c$。どんな三角形でも成り立つ}{seigen teiri: sain arufa bun no $a$ ikōru sain bēta bun no $b$ ikōru sain ganma bun no $c$. donna sankakukei de mo naritatsu}{der Sinussatz: $a$ durch Sinus Alpha gleich $b$ durch Sinus Beta gleich $c$ durch Sinus Gamma -- er gilt in jedem Dreieck}
\lies{$c^2 = a^2 + b^2 - 2ab\cos\gamma$}{余弦定理：$c$ 2乗イコール $a$ 2乗プラス $b$ 2乗マイナス2$ab$ コサインガンマ}{yogen teiri: $c$ nijō ikōru $a$ nijō purasu $b$ nijō mainasu ni $ab$ kosain ganma}{der Kosinussatz: $c$ Quadrat gleich $a$ Quadrat plus $b$ Quadrat minus zwei $a$ $b$ Kosinus Gamma}
\lies{$180\degree \mathrel{\widehat{=}} \pi$}{180度は、弧度法でパイラジアン}{hyakuhachijū do wa, kodohō de pai rajian}{hundertachtzig Grad entsprechen pi im Bogenmaß}
\end{lesen}

\anmerkung{Japanische Schülerinnen und Schüler merken sich die Definitionen oft mit den
Schreibschrift-Buchstaben s, c und t: Man schreibt den Buchstaben in das Dreieck -- die Seite, an der
der Strich beginnt, steht im Nenner, die Seite, an der er endet, im Zähler. Die Tafelwerke mit
Winkelfunktions- und Logarithmentafeln aus älteren Büchern ersetzt heute der Taschenrechner
(\jp{電卓}{dentaku}).}

\platz\zwischen{Beispiele}{例}{rei}

\begin{lesen}
\lies{$a = 3\,\mathrm{cm},\ c = 5\,\mathrm{cm}$ \newline $\sin\alpha = \frac35 = 0{,}6$ \newline $\alpha \approx 36{,}9\degree$}{サインアルファは5分の3で、0.6。よってアルファは約36.9度。}{sain arufa wa go bun no san de, rei ten roku. yotte arufa wa yaku sanjūroku ten kyū do.}{Sinus Alpha ist drei durch fünf, also null Komma sechs. Damit ist Alpha der Arkussinus von null Komma sechs, ungefähr sechsunddreißig Komma neun Grad.}
\lies{$\sin 30{,}2\degree \approx 0{,}5030$}{サイン30.2度は約0.5030}{sain sanjū ten ni do wa yaku rei ten go rei san rei}{der Sinus von dreißig Komma zwei Grad ist ungefähr null Komma fünf null drei null}
\lies{$h = 20\,\mathrm{m} \cdot \tan 35\degree \approx 14{,}0\,\mathrm{m}$}{木から20メートル離れた地点から木の先端を見上げると、仰角は35度である。木の高さは約14メートル。}{ki kara nijū mētoru hanareta chiten kara ki no sentan o miageru to, gyōkaku wa sanjūgo do de aru. ki no takasa wa yaku jūyon mētoru.}{Aus zwanzig Metern Entfernung sieht man die Spitze eines Baumes unter einem Höhenwinkel von fünfunddreißig Grad. Der Baum ist etwa vierzehn Meter hoch.}
\lies{$a = 5,\ b = 7,\ \gamma = 60\degree$ \newline $c^2 = 25 + 49 - 70 \cdot \frac12 = 39$ \newline $c \approx 6{,}24$}{余弦定理より、$c$ 2乗は25プラス49マイナス70掛ける2分の1で、39。よって $c$ は約6.24。}{yogen teiri yori, $c$ nijō wa nijūgo purasu yonjūkyū mainasu nanajū kakeru ni bun no ichi de, sanjūkyū. yotte $c$ wa yaku roku ten ni yon.}{Nach dem Kosinussatz ist $c$ Quadrat gleich fünfundzwanzig plus neunundvierzig minus siebzig mal ein Halb, also neununddreißig. Damit ist $c$ ungefähr sechs Komma zwei vier.}
\lies{$a = 8,\ \alpha = 30\degree,\ \beta = 45\degree$ \newline $b = \dfrac{8 \cdot \sin 45\degree}{\sin 30\degree} \approx 11{,}31$}{正弦定理より、$b$ はサイン30度分の8掛けるサイン45度で、約11.31。}{seigen teiri yori, $b$ wa sain sanjū do bun no hachi kakeru sain yonjūgo do de, yaku jūitten san ichi.}{Nach dem Sinussatz ist $b$ gleich acht mal Sinus fünfundvierzig Grad durch Sinus dreißig Grad, ungefähr elf Komma drei eins.}
\end{lesen}

%--------------------------------------------------------------------------------------------------
\newpage
\kapitel{Koordinaten und Geraden}{座標と直線}{zahyō to chokusen}\renewcommand{\lesenabstand}{1.12}\renewcommand{\arraystretch}{0.93}

\begin{center}
\begin{tikzpicture}[font=\small, scale=0.95]
	\draw[achse] (-0.7,0) -- (5.6,0) node[below left] {$x$}; \draw[achse] (0,-0.7) -- (0,4.2) node[below left] {$y$};
	\draw[mass] (1.2,0) arc[start angle=0, end angle=36.87, radius=1.2]; \node at (17:1.6) {$\varphi$};
	\draw[black!60, dashed] (4,0) node[below, black] {$x_P$} -- (4,3) -- (0,3) node[left, black] {$y_P$};
	\draw[red, very thick] (0,0) -- node[above left, inner sep=2pt] {$r$} (4,3);
	\fill (4,3) circle (2pt) node[above right] {$P$}; \node[below left] at (0,0) {$O$};
	\begin{scope}[shift={(9.6,0.3)}]
		\draw[achse] (-3.2,0) -- (6.2,0) node[below left] {$x$}; \draw[achse] (0,-0.9) -- (0,4.2) node[below left] {$y$};
		\draw[mass] (-1,0) arc[start angle=0, end angle=26.57, radius=1]; \node at ($(-2,0)+(12:1.4)$) {$\alpha$};
		\draw[blue, very thick] (-3,-0.5) -- (5.6,3.8) node[right] {$g$};
		\draw[black!60, thick] (2,2) -- node[below] {$\Delta x$} (4,2) -- node[right] {$\Delta y$} (4,3);
		\draw[red, thick, |<->|] (-0.3,0) -- node[left] {$n$} (-0.3,1);
		\fill[red] (0,1) circle (2pt);
	\end{scope}
\end{tikzpicture}
\end{center}

\begin{lesen}
\lies{$P(4 \mid 3)$ \newline \footnotesize jap. $P(4,\ 3)$}{点 $P$ 4、3：$x$ 座標は4、$y$ 座標は3}{ten $P$ yon, san: $x$ zahyō wa yon, $y$ zahyō wa san}{der Punkt $P$ vier drei: Seine $x$-Koordinate (Abszisse) ist vier, seine $y$-Koordinate (Ordinate) ist drei}
\lies{$P(r;\,\varphi)$}{点 $P$ の極座標：原点（極）からの距離 $r$ と偏角ファイ}{ten $P$ no kyokuzahyō: genten (kyoku) kara no kyori $r$ to henkaku fai}{die Polarkoordinaten von $P$: der Abstand $r$ vom Ursprung (Pol) und der Polarwinkel Phi}
\lies{$x = r\cos\varphi,\ y = r\sin\varphi$}{$x$ イコール $r$ コサインファイ、$y$ イコール $r$ サインファイ}{$x$ ikōru $r$ kosain fai, $y$ ikōru $r$ sain fai}{$x$ gleich $r$ Kosinus Phi, $y$ gleich $r$ Sinus Phi}
\lies{$y = mx + n$ \newline \footnotesize jap. $y = ax + b$}{直線の方程式：$y$ イコール $m$ $x$ プラス $n$}{chokusen no hōteishiki: $y$ ikōru $m$ $x$ purasu $n$}{die Normalform der Geradengleichung: $y$ gleich $m$ $x$ plus $n$}
\lies{$m = \dfrac{\Delta y}{\Delta x} = \tan\alpha$}{傾き $m$ は、$x$ の増加量分の $y$ の増加量。傾斜角アルファのタンジェントに等しい}{katamuki $m$ wa, $x$ no zōkaryō bun no $y$ no zōkaryō. keishakaku arufa no tanjento ni hitoshii}{die Steigung $m$ ist gleich Delta $y$ durch Delta $x$; sie ist der Tangens des Steigungswinkels Alpha}
\lies{$n$}{$y$ 切片：直線は点 $(0,\,n)$ で $y$ 軸と交わる}{$y$ seppen: chokusen wa ten $(0,\,n)$ de $y$ jiku to majiwaru}{der $y$-Achsenabschnitt: Die Gerade schneidet die $y$-Achse im Punkt $(0 \mid n)$}
\lies{$m_1 = m_2$}{平行な2直線の傾きは等しい}{heikō na nichokusen no katamuki wa hitoshii}{parallele Geraden haben die gleiche Steigung}
\lies{$m_1 \cdot m_2 = -1$}{傾きの積がマイナス1ならば、2直線は垂直である}{katamuki no seki ga mainasu ichi naraba, nichokusen wa suichoku de aru}{zwei Geraden sind zueinander senkrecht, wenn das Produkt ihrer Steigungen minus eins ist}
\end{lesen}

\platz\zwischen{Beispiele}{例}{rei}

\begin{lesen}
\lies{$A(1 \mid 2),\ B(3 \mid 6)$ \newline $m = \dfrac{6 - 2}{3 - 1} = 2$}{2点 $A$、$B$ を通る直線の方程式を求めなさい。傾きは、3マイナス1分の6マイナス2で、2。}{niten $A$, $B$ o tōru chokusen no hōteishiki o motomenasai. katamuki wa, san mainasu ichi bun no roku mainasu ni de, ni.}{Bestimme die Gleichung der Geraden durch $A$ und $B$. Die Steigung ist sechs minus zwei durch drei minus eins, also zwei.}
\lies{$2 = 2 \cdot 1 + n \ \Rightarrow\ n = 0$ \newline $y = 2x$}{点 $A$ の座標を代入すると、$n$ は0。直線の方程式は $y$ イコール2$x$。}{ten $A$ no zahyō o dainyū suru to, $n$ wa zero. chokusen no hōteishiki wa $y$ ikōru ni $x$.}{Setze die Koordinaten von $A$ ein: $n$ ist null. Die Geradengleichung lautet $y$ gleich zwei $x$.}
\lies{$\overline{AB} = \sqrt{2^2 + 4^2} = \sqrt{20} \approx 4{,}47$}{2点 $A$、$B$ 間の距離はルート20で、約4.47。}{niten $A$, $B$ kan no kyori wa rūto nijū de, yaku yon ten yon nana.}{Der Abstand von $A$ und $B$ ist die Wurzel aus zwanzig, ungefähr vier Komma vier sieben.}
\lies{$2x = -x + 6 \ \Rightarrow\ x = 2$ \newline $S(2 \mid 4)$}{直線 $y$ イコール2$x$ と直線 $y$ イコールマイナス $x$ プラス6は、点 $S(2,\,4)$ で交わる。}{chokusen $y$ ikōru ni $x$ to chokusen $y$ ikōru mainasu $x$ purasu roku wa, ten $S(2,\,4)$ de majiwaru.}{Die Geraden $y$ gleich zwei $x$ und $y$ gleich minus $x$ plus sechs schneiden einander im Punkt $S$ zwei vier.}
\end{lesen}

%--------------------------------------------------------------------------------------------------
\newpage
\kapitel{Kegelschnitte}{円錐曲線}{ensui kyokusen}\renewcommand{\lesenabstand}{1.12}\renewcommand{\arraystretch}{0.93}

\begin{center}
\begin{tikzpicture}[font=\footnotesize, scale=0.95]
	\begin{scope}
		\draw[achse] (-1.7,0) -- (1.9,0) node[below] {$x$}; \draw[achse] (0,-1.7) -- (0,1.9) node[left] {$y$};
		\draw[blue, very thick] (0,0) circle (1.3);
		\draw[red, thick] (0,0) -- node[above left, inner sep=1pt] {$r$} (40:1.3); \fill (40:1.3) circle (1.5pt) node[above right] {$P$};
		\node[align=center] at (0,-2.3) {円\\\color{black!60}Kreis};
	\end{scope}
	\begin{scope}[shift={(4.3,0)}]
		\draw[achse] (-2.1,0) -- (2.3,0) node[below] {$x$}; \draw[achse] (0,-1.7) -- (0,1.9) node[left] {$y$};
		\draw[blue, very thick] (0,0) ellipse[x radius=1.7, y radius=1.1];
		\coordinate (P) at ({1.7*cos(65)},{1.1*sin(65)});
		\draw[red, thick] (-1.296,0) -- (P) -- (1.296,0);
		\fill[red] (-1.296,0) circle (1.5pt) node[below] {$F_1$} (1.296,0) circle (1.5pt) node[below] {$F_2$};
		\fill (P) circle (1.5pt) node[above right] {$P$};
		\node[align=center] at (0,-2.3) {楕円\\\color{black!60}Ellipse};
	\end{scope}
	\begin{scope}[shift={(8.9,0)}]
		\draw[achse] (-2.1,0) -- (2.3,0) node[below] {$x$}; \draw[achse] (0,-1.7) -- (0,1.9) node[left] {$y$};
		\draw[black!50, dashed] (-1.6,-1.6) -- (1.6,1.6) (-1.6,1.6) -- (1.6,-1.6);
		\draw[blue, very thick, domain=-1.5:1.5, samples=40, variable=\t] plot ({0.7*cosh(\t)}, {0.7*sinh(\t)});
		\draw[blue, very thick, domain=-1.5:1.5, samples=40, variable=\t] plot ({-0.7*cosh(\t)}, {0.7*sinh(\t)});
		\fill[red] (-0.99,0) circle (1.5pt) node[below left, inner sep=1pt] {$F_1$} (0.99,0) circle (1.5pt) node[below right, inner sep=1pt] {$F_2$};
		\node[align=center] at (0,-2.3) {双曲線\\\color{black!60}Hyperbel};
	\end{scope}
	\begin{scope}[shift={(13.2,0)}]
		\draw[achse] (-1.2,0) -- (2.3,0) node[below] {$x$}; \draw[achse] (0,-1.7) -- (0,1.9) node[left] {$y$};
		\draw[blue, very thick, domain=-1.6:1.6, samples=40, variable=\t] plot ({\t*\t/1.6}, {\t});
		\draw[red, thick, dashed] (-0.4,-1.6) -- (-0.4,1.7) node[left] {$l$};
		\fill[red] (0.4,0) circle (1.5pt) node[below right, inner sep=1pt] {$F$};
		\node[align=center] at (0.5,-2.3) {放物線\\\color{black!60}Parabel};
	\end{scope}
\end{tikzpicture}
\end{center}

\begin{lesen}
\lies{$x^2 + y^2 = r^2$}{原点を中心とする半径 $r$ の円の方程式：$x$ 2乗プラス $y$ 2乗イコール $r$ 2乗}{genten o chūshin to suru hankei $r$ no en no hōteishiki: $x$ nijō purasu $y$ nijō ikōru $r$ nijō}{die Gleichung eines Kreises um den Ursprung mit dem Radius $r$: $x$ Quadrat plus $y$ Quadrat gleich $r$ Quadrat}
\lies{$\dfrac{x^2}{a^2} + \dfrac{y^2}{b^2} = 1$}{楕円の方程式：$a$ 2乗分の $x$ 2乗プラス $b$ 2乗分の $y$ 2乗イコール1。$a$ は長半径、$b$ は短半径}{daen no hōteishiki: $a$ nijō bun no $x$ nijō purasu $b$ nijō bun no $y$ nijō ikōru ichi. $a$ wa chōhankei, $b$ wa tanhankei}{die Gleichung einer Ellipse: $x$ Quadrat durch $a$ Quadrat plus $y$ Quadrat durch $b$ Quadrat gleich eins; $a$ ist die große, $b$ die kleine Halbachse}
\lies{$\overline{PF_1} + \overline{PF_2} = 2a$}{$F_1$ と $F_2$ は焦点：楕円上のどの点 $P$ でも、2つの焦点からの距離の和は一定}{$F_1$ to $F_2$ wa shōten: daenjō no dono ten $P$ de mo, futatsu no shōten kara no kyori no wa wa ittei}{$F_1$ und $F_2$ sind die Brennpunkte: Für jeden Ellipsenpunkt $P$ ist die Summe der Abstände zu den Brennpunkten konstant}
\lies{$\dfrac{x^2}{a^2} - \dfrac{y^2}{b^2} = 1$}{双曲線の方程式：$a$ 2乗分の $x$ 2乗マイナス $b$ 2乗分の $y$ 2乗イコール1。破線は漸近線}{sōkyokusen no hōteishiki: $a$ nijō bun no $x$ nijō mainasu $b$ nijō bun no $y$ nijō ikōru ichi. hasen wa zenkinsen}{die Gleichung einer Hyperbel: $x$ Quadrat durch $a$ Quadrat minus $y$ Quadrat durch $b$ Quadrat gleich eins; die gestrichelten Geraden sind ihre Asymptoten}
\lies{$y^2 = 2px$ \newline \footnotesize jap. meist $y^2 = 4px$}{原点を頂点とする放物線の方程式：$y$ 2乗イコール2$p$ $x$。$F$ は焦点、$l$ は準線}{genten o chōten to suru hōbutsusen no hōteishiki: $y$ nijō ikōru ni $p$ $x$. $F$ wa shōten, $l$ wa junsen}{die Scheitelgleichung einer Parabel: $y$ Quadrat gleich zwei $p$ $x$; $F$ ist der Brennpunkt, $l$ die Leitlinie}
\lies{$\overline{PF} = \overline{Pl}$}{放物線上の点は、焦点と準線から等しい距離にある}{hōbutsusenjō no ten wa, shōten to junsen kara hitoshii kyori ni aru}{jeder Parabelpunkt hat von Brennpunkt und Leitlinie den gleichen Abstand}
\end{lesen}

\anmerkung{Kreis, Ellipse, Parabel und Hyperbel heißen Kegelschnitte (\jp{円錐曲線}{ensui kyokusen}), weil
sie entstehen, wenn eine Ebene einen Doppelkegel schneidet. In der Schule sagt man meist
\jp{二次曲線}{niji kyokusen}, „Kurven zweiten Grades“.}

\platz\zwischen{Beispiele}{例}{rei}

\begin{lesen}
\lies{$x^2 + y^2 = 25$ \newline $P(3 \mid 4)$: $9 + 16 = 25$ \checkmark}{点 $P(3,\,4)$ はこの円周上にある。3の2乗プラス4の2乗が25だからである。}{ten $P(3,\,4)$ wa kono enshūjō ni aru. san no nijō purasu yon no nijō ga nijūgo da kara de aru.}{Der Punkt $P$ drei vier liegt auf dem Kreis, denn drei Quadrat plus vier Quadrat ist fünfundzwanzig.}
\lies{$\dfrac{x^2}{25} + \dfrac{y^2}{9} = 1$ \newline $e = \sqrt{25 - 9} = 4$}{楕円の長半径は5、短半径は3。焦点は $(4,\,0)$ と $(-4,\,0)$。}{daen no chōhankei wa go, tanhankei wa san. shōten wa $(4,\,0)$ to $(-4,\,0)$.}{Die Halbachsen der Ellipse sind fünf und drei. Die Brennpunkte liegen bei $(4 \mid 0)$ und $(-4 \mid 0)$.}
\lies{$y^2 = 8x$ \newline $p = 4$}{放物線の焦点は $(2,\,0)$、準線は直線 $x$ イコールマイナス2。}{hōbutsusen no shōten wa $(2,\,0)$, junsen wa chokusen $x$ ikōru mainasu ni.}{Der Brennpunkt der Parabel ist $(2 \mid 0)$; die Leitlinie ist die Gerade $x$ gleich minus zwei.}
\end{lesen}

%--------------------------------------------------------------------------------------------------
\newpage
\kapitel{Folgen, Vektoren, Matrizen und Kombinatorik}{数列、ベクトル、行列、組合せ}{sūretsu, bekutoru, gyōretsu, kumiawase}\renewcommand{\lesenabstand}{1.12}\renewcommand{\arraystretch}{0.93}

\platz\zwischen{Folgen und Reihen}{数列と級数}{sūretsu to kyūsū}
\begin{lesen}
\lies{$a_n = a_1 + (n - 1) \cdot d$}{等差数列の一般項：$a_n$ イコール $a_1$ プラスかっこ $n$ マイナス1かっことじ $d$（$d$ は公差）}{tōsa sūretsu no ippankō: $a_n$ ikōru $a_1$ purasu kakko $n$ mainasu ichi kakkotoji $d$ ($d$ wa kōsa)}{das $n$-te Glied einer arithmetischen Folge: $a$ $n$ gleich $a$ eins plus $n$ minus eins mal $d$}
\lies{$a_n = a_1 \cdot q^{\,n-1}$}{等比数列の一般項：$a_n$ イコール $a_1$ 掛ける $q$ の $n$ マイナス1乗（$q$ は公比）}{tōhi sūretsu no ippankō: $a_n$ ikōru $a_1$ kakeru $q$ no $n$ mainasu ichi jō ($q$ wa kōhi)}{das $n$-te Glied einer geometrischen Folge: $a$ $n$ gleich $a$ eins mal $q$ hoch $n$ minus eins}
\lies{$\displaystyle\sum_{k=1}^{n} k = \frac{n(n + 1)}{2}$}{シグマ、$k$ が1から $n$ までの $k$ は、2分の $n$ かっこ $n$ プラス1かっことじ}{shiguma, $k$ ga ichi kara $n$ made no $k$ wa, ni bun no $n$ kakko $n$ purasu ichi kakkotoji}{die Summe über $k$ für $k$ von eins bis $n$ ist gleich $n$ mal $n$ plus eins, geteilt durch zwei}
\lies{$\displaystyle\sum_{k=0}^{\infty} q^k = \frac{1}{1 - q}$ \newline für $|q| < 1$}{無限等比級数は、$q$ の絶対値が1より小さいとき収束し、そうでないとき発散する}{mugen tōhi kyūsū wa, $q$ no zettaichi ga ichi yori chiisai toki shūsoku shi, sō de nai toki hassan suru}{die unendliche geometrische Reihe konvergiert, wenn der Betrag von $q$ kleiner als eins ist; sonst divergiert sie}
\end{lesen}

\platz\zwischen{Vektoren und Matrizen}{ベクトルと行列}{bekutoru to gyōretsu}
\begin{lesen}
\lies{$\vec{a},\ |\vec{a}\,|$}{ベクトル $a$、ベクトル $a$ の大きさ}{bekutoru $a$, bekutoru $a$ no ōkisa}{Vektor $a$; der Betrag von Vektor $a$}
\lies{$\vec{a} \cdot \vec{b}$}{ベクトル $a$ とベクトル $b$ の内積 \;/\; $a$ ドット $b$}{bekutoru $a$ to bekutoru $b$ no naiseki \;/\; $a$ dotto $b$}{$a$ mal $b$ -- das Skalarprodukt}
\lies{$\vec{a} \times \vec{b}$}{ベクトル $a$ とベクトル $b$ の外積 \;/\; $a$ クロス $b$}{bekutoru $a$ to bekutoru $b$ no gaiseki \;/\; $a$ kurosu $b$}{$a$ kreuz $b$ -- das Vektorprodukt (Kreuzprodukt)}
\lies{$\vec{a} \perp \vec{b} \ \Leftrightarrow\ \vec{a} \cdot \vec{b} = 0$}{2つのベクトルが垂直であることと、内積が0であることは同値である}{futatsu no bekutoru ga suichoku de aru koto to, naiseki ga zero de aru koto wa dōchi de aru}{zwei Vektoren sind genau dann zueinander senkrecht, wenn ihr Skalarprodukt null ist}
\lies{$A = \begin{pmatrix} 1 & 2 & 3\\ 4 & 5 & 6 \end{pmatrix}$}{$A$ は2行3列の行列}{$A$ wa ni gyō san retsu no gyōretsu}{$A$ ist eine Zwei-Kreuz-Drei-Matrix: Sie hat zwei Zeilen und drei Spalten}
\lies{$a_{ij}$}{$a$ $i$ $j$：第 $i$ 行、第 $j$ 列の成分}{$a$ $i$ $j$: dai $i$ gyō, dai $j$ retsu no seibun}{$a$ $i$ $j$: das Element in Zeile $i$ und Spalte $j$}
\lies{$\det A$}{行列 $A$ の行列式}{gyōretsu $A$ no gyōretsushiki}{die Determinante von $A$}
\end{lesen}

\platz\zwischen{Kombinatorik}{順列と組合せ}{junretsu to kumiawase}
\begin{lesen}
\lies{$n!$}{$n$ の階乗：異なる $n$ 個のものを1列に並べる順列の総数}{$n$ no kaijō: kotonaru $n$ ko no mono o ichiretsu ni naraberu junretsu no sōsū}{$n$ Fakultät: die Anzahl der Permutationen von $n$ verschiedenen Elementen}
\lies{$\dbinom{n}{k}$ \newline \footnotesize jap. ${}_n\mathrm{C}_k$}{エヌシーケー：$n$ 個から $k$ 個取る組合せの総数（二項係数）}{enu shī kē: $n$ ko kara $k$ ko toru kumiawase no sōsū (nikō keisū)}{$n$ über $k$: die Anzahl der Kombinationen von $k$ aus $n$ Elementen (ohne Wiederholung)}
\lies{$\dfrac{n!}{(n - k)!}$ \newline \footnotesize jap. ${}_n\mathrm{P}_k$}{エヌピーケー：$n$ 個から $k$ 個取って並べる順列の総数}{enu pī kē: $n$ ko kara $k$ ko totte naraberu junretsu no sōsū}{die Anzahl der Variationen von $k$ aus $n$ Elementen (mit Beachtung der Reihenfolge)}
\end{lesen}

\platz\zwischen{Beispiele}{例}{rei}

\begin{lesen}
\lies{$3;\ 7;\ 11;\ 15;\ \ldots$ \newline $a_{10} = 3 + 9 \cdot 4 = 39$}{初項3、公差4の等差数列。第10項は39}{shokō san, kōsa yon no tōsa sūretsu. dai jū kō wa sanjūkyū}{eine arithmetische Folge mit dem Anfangsglied drei und der Differenz vier; das zehnte Glied ist neununddreißig}
\lies{$\vektor{1}{2}{3} \cdot \vektor{4}{-1}{2} = 4 - 2 + 6 = 8$}{2つのベクトルの内積は8}{futatsu no bekutoru no naiseki wa hachi}{das Skalarprodukt der beiden Vektoren ist acht}
\lies{$\begin{vmatrix} 1 & 2\\ 3 & 4 \end{vmatrix} = 1 \cdot 4 - 2 \cdot 3 = -2$}{行列式は、1掛ける4引く2掛ける3で、マイナス2}{gyōretsushiki wa, ichi kakeru yon hiku ni kakeru san de, mainasu ni}{die Determinante ist eins mal vier minus zwei mal drei, also minus zwei}
\lies{$4! = 4 \cdot 3 \cdot 2 \cdot 1 = 24$}{4の階乗は24：4冊の本を棚に並べる方法は、24通りある。}{yon no kaijō wa nijūyon: yonsatsu no hon o tana ni naraberu hōhō wa, nijūyontōri aru.}{Vier Fakultät ist vierundzwanzig: Vier Bücher lassen sich auf vierundzwanzig Arten ins Regal stellen.}
\lies{$\dbinom{5}{2} = \dfrac{5 \cdot 4}{2} = 10$}{5シー2は10：5人の生徒から2人を選ぶ方法は、10通りある。}{go shī ni wa jū: gonin no seito kara futari o erabu hōhō wa, juttōri aru.}{Fünf über zwei ist zehn: Es gibt zehn Möglichkeiten, aus fünf Schülern zwei auszuwählen.}
\end{lesen}

%--------------------------------------------------------------------------------------------------
\newpage
\kapitel{Das griechische Alphabet}{ギリシャ文字}{girisha moji}\renewcommand{\lesenabstand}{1.12}\renewcommand{\arraystretch}{0.93}

Griechische Buchstaben schreibt man im Japanischen in Katakana. Ihre Namen folgen meist der englischen
Aussprache -- besonders auffällig bei $\pi$, $\varphi$, $\chi$ und $\psi$.

\begin{center}
{\small\renewcommand{\arraystretch}{1.25}
\begin{tabular}{@{}c c l l @{\hspace{0.9cm}} c c l l@{}}
	\toprule
	\multicolumn{2}{@{}l}{\textbf{Zeichen}} & \textbf{Japanisch} & \textbf{Deutsch} & \multicolumn{2}{l}{\textbf{Zeichen}} & \textbf{Japanisch} & \textbf{Deutsch} \\ \midrule
	$\mathrm{A}$ & $\alpha$ & アルファ [arufa] & Alpha & $\mathrm{N}$ & $\nu$ & ニュー [nyū] & Ny \\
	$\mathrm{B}$ & $\beta$ & ベータ [bēta] & Beta & $\Xi$ & $\xi$ & クサイ [kusai] & Xi \\
	$\Gamma$ & $\gamma$ & ガンマ [ganma] & Gamma & $\mathrm{O}$ & $o$ & オミクロン [omikuron] & Omikron \\
	$\Delta$ & $\delta$ & デルタ [deruta] & Delta & $\Pi$ & $\pi$ & パイ [pai] & Pi \\
	$\mathrm{E}$ & $\varepsilon$ & イプシロン [ipushiron] & Epsilon & $\mathrm{P}$ & $\varrho$ & ロー [rō] & Rho \\
	$\mathrm{Z}$ & $\zeta$ & ゼータ [zēta] & Zeta & $\Sigma$ & $\sigma$ & シグマ [shiguma] & Sigma \\
	$\mathrm{H}$ & $\eta$ & イータ [īta] & Eta & $\mathrm{T}$ & $\tau$ & タウ [tau] & Tau \\
	$\Theta$ & $\vartheta$ & シータ [shīta] & Theta & $\Upsilon$ & $\upsilon$ & ウプシロン [upushiron] & Ypsilon \\
	$\mathrm{I}$ & $\iota$ & イオタ [iota] & Jota & $\Phi$ & $\varphi$ & ファイ [fai] & Phi \\
	$\mathrm{K}$ & $\kappa$ & カッパ [kappa] & Kappa & $\mathrm{X}$ & $\chi$ & カイ [kai] & Chi \\
	$\Lambda$ & $\lambda$ & ラムダ [ramuda] & Lambda & $\Psi$ & $\psi$ & プサイ [pusai] & Psi \\
	$\mathrm{M}$ & $\mu$ & ミュー [myū] & My & $\Omega$ & $\omega$ & オメガ [omega] & Omega \\
	\bottomrule
\end{tabular}}
\end{center}

\anmerkung{Daneben hört man ältere, am Deutschen orientierte Namen: \jp{テータ}{tēta} für $\vartheta$,
\jp{エータ}{ēta} für $\eta$, \jp{クシー}{kushī} für $\xi$ und \jp{プシー}{pushī} für $\psi$.}

\platz\zwischen{Beispiele}{例}{rei}

\begin{lesen}
\lies{$\alpha,\ \beta,\ \gamma$}{アルファ、ベータ、ガンマは三角形の内角である}{arufa, bēta, ganma wa sankakukei no naikaku de aru}{Alpha, Beta und Gamma sind die Innenwinkel des Dreiecks}
\lies{$\Delta x$}{デルタ $x$：$x$ の増分（2つの $x$ の値の差）}{deruta $x$: $x$ no zōbun (futatsu no $x$ no atai no sa)}{Delta $x$: die Differenz zweier $x$-Werte}
\lies{$\varepsilon > 0$}{任意の正の数イプシロンに対して}{nin'i no sei no sū ipushiron ni taishite}{für jedes Epsilon größer als null}
\lies{$\vec{x} = \vec{p} + \lambda \cdot \vec{u}$}{ベクトル $x$ イコールベクトル $p$ プラスラムダ掛けるベクトル $u$}{bekutoru $x$ ikōru bekutoru $p$ purasu ramuda kakeru bekutoru $u$}{Vektor $x$ gleich Vektor $p$ plus Lambda mal Vektor $u$}
\lies{$\mu = 3{,}5;\ \sigma \approx 1{,}71$}{平均ミューは3.5、標準偏差シグマは約1.71}{heikin myū wa san ten go, hyōjun hensa shiguma wa yaku itten nana ichi}{der Erwartungswert My ist drei Komma fünf, die Standardabweichung Sigma ist ungefähr eins Komma sieben eins}
\lies{$\displaystyle\sum_{k=1}^{n} a_k$}{大文字のシグマ：$k$ が1から $n$ までの $a_k$ の和}{ōmoji no shiguma: $k$ ga ichi kara $n$ made no $a_k$ no wa}{groß Sigma: die Summe über $a$ $k$ für $k$ von eins bis $n$}
\lies{$\pi \approx 3{,}14$}{パイは約3.14}{pai wa yaku san ten ichi yon}{Pi ist ungefähr drei Komma eins vier}
\end{lesen}


\newpage
%%%%  REGISTER  %%%%
\teil{Register}{索引}{sakuin}
\section{Register Deutsch -- Japanisch}\renewcommand{\lesenabstand}{1.12}\renewcommand{\arraystretch}{0.93}
\begin{multicols}{2}\small\raggedright

\medskip{\large\textbf{\color{cornflower}A}}\par\nopagebreak
\textbf{abhängige Variable} (die) -- 従属変数\dotfill\hyperref[e:107]{\pageref*{e:107}}\par
\textbf{ableiten} -- 微分する\dotfill\hyperref[e:278]{\pageref*{e:278}}\par
\textbf{Ableitung} (die) -- 導関数\dotfill\hyperref[e:273]{\pageref*{e:273}}\par
\textbf{Abstand} (der) -- 距離\dotfill\hyperref[e:152]{\pageref*{e:152}}\par
\textbf{Abszisse} (die) -- $x$ 座標\dotfill\hyperref[e:117]{\pageref*{e:117}}\par
\textbf{achsensymmetrisch zur $y$-Achse} -- $y$ 軸に関して対称\dotfill\hyperref[e:144]{\pageref*{e:144}}\par
\textbf{Achteck} (das) -- 八角形\dotfill\hyperref[e:207]{\pageref*{e:207}}\par
\textbf{addieren} -- 足す\dotfill\hyperref[e:14]{\pageref*{e:14}}\par
\textbf{ähnlich} -- 相似\dotfill\hyperref[e:176]{\pageref*{e:176}}\par
\textbf{Anfangsglied} (das) -- 初項\dotfill\hyperref[e:261]{\pageref*{e:261}}\par
\textbf{Ankathete} (die) -- 隣辺\dotfill\hyperref[e:188]{\pageref*{e:188}}\par
\textbf{arithmetische Folge} (die) -- 等差数列\dotfill\hyperref[e:259]{\pageref*{e:259}}\par
\textbf{Arkusfunktion} (die) -- 逆三角関数\dotfill\hyperref[e:136]{\pageref*{e:136}}\par
\textbf{Asymptote} (die) -- 漸近線\dotfill\hyperref[e:271]{\pageref*{e:271}}\par
\textbf{ausklammern} -- 共通因数でくくる\dotfill\hyperref[e:92]{\pageref*{e:92}}\par
\textbf{ausmultiplizieren} -- 展開する\dotfill\hyperref[e:91]{\pageref*{e:91}}\par
\textbf{Außenwinkel} (der) -- 外角\dotfill\hyperref[e:179]{\pageref*{e:179}}\par

\medskip{\large\textbf{\color{cornflower}B}}\par\nopagebreak
\textbf{Basis} (die) -- 底\dotfill\hyperref[e:43]{\pageref*{e:43}}\par
\textbf{Basis (im Dreieck)} (die) -- 底辺\dotfill\hyperref[e:183]{\pageref*{e:183}}\par
\textbf{Basiswinkel} (der) -- 底角\dotfill\hyperref[e:184]{\pageref*{e:184}}\par
\textbf{Baumdiagramm} (das) -- 樹形図\dotfill\hyperref[e:332]{\pageref*{e:332}}\par
\textbf{Berührungspunkt} (der) -- 接点\dotfill\hyperref[e:220]{\pageref*{e:220}}\par
\textbf{bestimmte Integral} (das) -- 定積分\dotfill\hyperref[e:288]{\pageref*{e:288}}\par
\textbf{Betrag} (der) -- 絶対値\dotfill\hyperref[e:53]{\pageref*{e:53}}\par
\textbf{Betrag eines Vektors} (der) -- ベクトルの大きさ\dotfill\hyperref[e:305]{\pageref*{e:305}}\par
\textbf{Beweis} (der) -- 証明\dotfill\hyperref[e:349]{\pageref*{e:349}}\par
\textbf{Binomialkoeffizient} (der) -- 二項係数\dotfill\hyperref[e:344]{\pageref*{e:344}}\par
\textbf{Binomialverteilung} (die) -- 二項分布\dotfill\hyperref[e:345]{\pageref*{e:345}}\par
\textbf{binomische Formel} (die) -- 乗法公式\dotfill\hyperref[e:93]{\pageref*{e:93}}\par
\textbf{Bogenmaß} (das) -- 弧度法\dotfill\hyperref[e:139]{\pageref*{e:139}}\par
\textbf{Brennpunkt} (der) -- 焦点\dotfill\hyperref[e:316]{\pageref*{e:316}}\par
\textbf{Bruch} (der) -- 分数\dotfill\hyperref[e:28]{\pageref*{e:28}}\par

\medskip{\large\textbf{\color{cornflower}D}}\par\nopagebreak
\textbf{daraus folgt} -- ならば\dotfill\hyperref[e:350]{\pageref*{e:350}}\par
\textbf{Deckfläche} (die) -- 上面\dotfill\hyperref[e:249]{\pageref*{e:249}}\par
\textbf{Definition} (die) -- 定義\dotfill\hyperref[e:346]{\pageref*{e:346}}\par
\textbf{Definitionsbereich} (der) -- 定義域\dotfill\hyperref[e:109]{\pageref*{e:109}}\par
\textbf{Determinante} (die) -- 行列式\dotfill\hyperref[e:325]{\pageref*{e:325}}\par
\textbf{Dezimalzahl} (die) -- 小数\dotfill\hyperref[e:13]{\pageref*{e:13}}\par
\textbf{Diagonale} (die) -- 対角線\dotfill\hyperref[e:197]{\pageref*{e:197}}\par
\textbf{Differenz} (die) -- 差\dotfill\hyperref[e:18]{\pageref*{e:18}}\par
\textbf{Differenzenquotient} (der) -- 平均変化率\dotfill\hyperref[e:272]{\pageref*{e:272}}\par
\textbf{Differenzial} (das) -- 微分\dotfill\hyperref[e:276]{\pageref*{e:276}}\par
\textbf{Differenzialquotient} (der) -- 微分係数\dotfill\hyperref[e:275]{\pageref*{e:275}}\par
\textbf{Differenzialrechnung} (die) -- 微分法\dotfill\hyperref[e:274]{\pageref*{e:274}}\par
\textbf{Differenzmenge} (die) -- 差集合\dotfill\hyperref[e:102]{\pageref*{e:102}}\par
\textbf{divergent} -- 発散する\dotfill\hyperref[e:263]{\pageref*{e:263}}\par
\textbf{Dividend} (der) -- 被除数\dotfill\hyperref[e:26]{\pageref*{e:26}}\par
\textbf{dividieren} -- 割る\dotfill\hyperref[e:24]{\pageref*{e:24}}\par
\textbf{Divisor} (der) -- 除数\dotfill\hyperref[e:27]{\pageref*{e:27}}\par
\textbf{Drachenviereck} (das) -- たこ形\dotfill\hyperref[e:203]{\pageref*{e:203}}\par
\textbf{Drehung} (die) -- 回転移動\dotfill\hyperref[e:174]{\pageref*{e:174}}\par
\textbf{Dreieck} (das) -- 三角形\dotfill\hyperref[e:177]{\pageref*{e:177}}\par
\textbf{Durchmesser} (der) -- 直径\dotfill\hyperref[e:214]{\pageref*{e:214}}\par

\medskip{\large\textbf{\color{cornflower}E}}\par\nopagebreak
\textbf{Ebene} (die) -- 平面\dotfill\hyperref[e:308]{\pageref*{e:308}}\par
\textbf{echte Bruch} (der) -- 真分数\dotfill\hyperref[e:31]{\pageref*{e:31}}\par
\textbf{echte Teilmenge} (die) -- 真部分集合\dotfill\hyperref[e:99]{\pageref*{e:99}}\par
\textbf{Ecke} (die) -- 頂点\dotfill\hyperref[e:245]{\pageref*{e:245}}\par
\textbf{eckigen Klammern} (die) -- 角括弧\dotfill\hyperref[e:89]{\pageref*{e:89}}\par
\textbf{Einheitskreis} (der) -- 単位円\dotfill\hyperref[e:137]{\pageref*{e:137}}\par
\textbf{Einheitsvektor} (der) -- 単位ベクトル\dotfill\hyperref[e:303]{\pageref*{e:303}}\par
\textbf{Element} (das) -- 要素\dotfill\hyperref[e:97]{\pageref*{e:97}}\par
\textbf{Ellipse} (die) -- 楕円\dotfill\hyperref[e:314]{\pageref*{e:314}}\par
\textbf{entgegengesetzt liegende Winkel} -- 同側内角\dotfill\hyperref[e:169]{\pageref*{e:169}}\par
\textbf{Ereignis} (das) -- 事象\dotfill\hyperref[e:328]{\pageref*{e:328}}\par
\textbf{Ergebnis} (das) -- 結果\dotfill\hyperref[e:327]{\pageref*{e:327}}\par
\textbf{Erwartungswert} (der) -- 期待値\dotfill\hyperref[e:335]{\pageref*{e:335}}\par
\textbf{erweitern} -- 倍分する\dotfill\hyperref[e:38]{\pageref*{e:38}}\par
\textbf{es gibt} -- 存在する\dotfill\hyperref[e:353]{\pageref*{e:353}}\par
\textbf{eulersche Zahl} (die) -- ネイピア数\dotfill\hyperref[e:298]{\pageref*{e:298}}\par
\textbf{Exponent} (der) -- 指数\dotfill\hyperref[e:44]{\pageref*{e:44}}\par
\textbf{Exponentialfunktion} (die) -- 指数関数\dotfill\hyperref[e:130]{\pageref*{e:130}}\par
\textbf{Exponentialgleichung} (die) -- 指数方程式\dotfill\hyperref[e:84]{\pageref*{e:84}}\par
\textbf{Extrempunkt} (der) -- 極値をとる点\dotfill\hyperref[e:281]{\pageref*{e:281}}\par

\medskip{\large\textbf{\color{cornflower}F}}\par\nopagebreak
\textbf{Faktor} (der) -- 因数\dotfill\hyperref[e:23]{\pageref*{e:23}}\par
\textbf{Fakultät} (die) -- 階乗\dotfill\hyperref[e:340]{\pageref*{e:340}}\par
\textbf{Flächendiagonale} (die) -- 面の対角線\dotfill\hyperref[e:251]{\pageref*{e:251}}\par
\textbf{flächengleich} -- 面積が等しい\dotfill\hyperref[e:195]{\pageref*{e:195}}\par
\textbf{Flächeninhalt} (der) -- 面積\dotfill\hyperref[e:210]{\pageref*{e:210}}\par
\textbf{Folge} (die) -- 数列\dotfill\hyperref[e:257]{\pageref*{e:257}}\par
\textbf{Fünfeck} (das) -- 五角形\dotfill\hyperref[e:206]{\pageref*{e:206}}\par
\textbf{Funktion} (die) -- 関数\dotfill\hyperref[e:105]{\pageref*{e:105}}\par
\textbf{Funktionsgraph} (der) -- グラフ\dotfill\hyperref[e:111]{\pageref*{e:111}}\par
\textbf{Funktionsuntersuchung} (die) -- 関数の増減を調べる\dotfill\hyperref[e:280]{\pageref*{e:280}}\par
\textbf{Funktionswert} (der) -- 関数の値\dotfill\hyperref[e:108]{\pageref*{e:108}}\par
\textbf{für alle} -- すべての\dotfill\hyperref[e:352]{\pageref*{e:352}}\par

\medskip{\large\textbf{\color{cornflower}G}}\par\nopagebreak
\textbf{ganze Zahl} (die) -- 整数\dotfill\hyperref[e:6]{\pageref*{e:6}}\par
\textbf{ganzrationale Funktion} (die) -- 多項式関数\dotfill\hyperref[e:126]{\pageref*{e:126}}\par
\textbf{gebrochenrationale Funktion} (die) -- 分数関数\dotfill\hyperref[e:127]{\pageref*{e:127}}\par
\textbf{Gegenereignis} (das) -- 余事象\dotfill\hyperref[e:329]{\pageref*{e:329}}\par
\textbf{Gegenkathete} (die) -- 対辺\dotfill\hyperref[e:189]{\pageref*{e:189}}\par
\textbf{gemischte Zahl} (die) -- 帯分数\dotfill\hyperref[e:33]{\pageref*{e:33}}\par
\textbf{genau dann, wenn} -- 同値\dotfill\hyperref[e:351]{\pageref*{e:351}}\par
\textbf{geometrische Folge} (die) -- 等比数列\dotfill\hyperref[e:260]{\pageref*{e:260}}\par
\textbf{Gerade} (die) -- 直線\dotfill\hyperref[e:147]{\pageref*{e:147}}\par
\textbf{Geradengleichung} (die) -- 直線の方程式\dotfill\hyperref[e:309]{\pageref*{e:309}}\par
\textbf{geschweiften Klammern} (die) -- 波括弧\dotfill\hyperref[e:90]{\pageref*{e:90}}\par
\textbf{gestreckte Winkel} (der) -- 平角\dotfill\hyperref[e:160]{\pageref*{e:160}}\par
\textbf{gleich} -- 等しい\dotfill\hyperref[e:60]{\pageref*{e:60}}\par
\textbf{gleichnamige Brüche} -- 同分母分数\dotfill\hyperref[e:34]{\pageref*{e:34}}\par
\textbf{gleichschenklige Dreieck} (das) -- 二等辺三角形\dotfill\hyperref[e:182]{\pageref*{e:182}}\par
\textbf{gleichseitige Dreieck} (das) -- 正三角形\dotfill\hyperref[e:185]{\pageref*{e:185}}\par
\textbf{Gleichung} (die) -- 方程式\dotfill\hyperref[e:76]{\pageref*{e:76}}\par
\textbf{Gleichungssystem} (das) -- 連立方程式\dotfill\hyperref[e:80]{\pageref*{e:80}}\par
\textbf{Grad} (der) -- 次数\dotfill\hyperref[e:75]{\pageref*{e:75}}\par
\textbf{Gradmaß} (das) -- 度数法\dotfill\hyperref[e:138]{\pageref*{e:138}}\par
\textbf{Grenzwert} (der) -- 極限値\dotfill\hyperref[e:264]{\pageref*{e:264}}\par
\textbf{größer als} -- より大きい\dotfill\hyperref[e:64]{\pageref*{e:64}}\par
\textbf{größer oder gleich} -- 以上\dotfill\hyperref[e:66]{\pageref*{e:66}}\par
\textbf{größte gemeinsame Teiler} (der) -- 最大公約数\dotfill\hyperref[e:56]{\pageref*{e:56}}\par
\textbf{Grundfläche} (die) -- 底面\dotfill\hyperref[e:248]{\pageref*{e:248}}\par

\medskip{\large\textbf{\color{cornflower}H}}\par\nopagebreak
\textbf{Halbachse} (die) -- 半軸\dotfill\hyperref[e:317]{\pageref*{e:317}}\par
\textbf{Hauptdiagonale} (die) -- 主対角線\dotfill\hyperref[e:324]{\pageref*{e:324}}\par
\textbf{Hauptnenner} (der) -- 最小公分母\dotfill\hyperref[e:36]{\pageref*{e:36}}\par
\textbf{Hochpunkt} (der) -- 極大点\dotfill\hyperref[e:282]{\pageref*{e:282}}\par
\textbf{Höhe} (die) -- 高さ\dotfill\hyperref[e:191]{\pageref*{e:191}}\par
\textbf{Hohlzylinder} (der) -- 中空円柱\dotfill\hyperref[e:237]{\pageref*{e:237}}\par
\textbf{Hyperbel} (die) -- 双曲線\dotfill\hyperref[e:315]{\pageref*{e:315}}\par
\textbf{Hypotenuse} (die) -- 斜辺\dotfill\hyperref[e:186]{\pageref*{e:186}}\par

\medskip{\large\textbf{\color{cornflower}I}}\par\nopagebreak
\textbf{identisch gleich} -- 恒等的に等しい\dotfill\hyperref[e:62]{\pageref*{e:62}}\par
\textbf{imaginäre Einheit} (die) -- 虚数単位\dotfill\hyperref[e:11]{\pageref*{e:11}}\par
\textbf{Index} (der) -- 添え字\dotfill\hyperref[e:4]{\pageref*{e:4}}\par
\textbf{Inkreis} (der) -- 内接円\dotfill\hyperref[e:193]{\pageref*{e:193}}\par
\textbf{Innenwinkel} (der) -- 内角\dotfill\hyperref[e:178]{\pageref*{e:178}}\par
\textbf{Integral} (das) -- 積分\dotfill\hyperref[e:295]{\pageref*{e:295}}\par
\textbf{Integralrechnung} (die) -- 積分法\dotfill\hyperref[e:286]{\pageref*{e:286}}\par
\textbf{Integralzeichen} (das) -- 積分記号\dotfill\hyperref[e:289]{\pageref*{e:289}}\par
\textbf{Integrand} (der) -- 被積分関数\dotfill\hyperref[e:290]{\pageref*{e:290}}\par
\textbf{Integrationsgrenze} (die) -- 上端、下端\dotfill\hyperref[e:297]{\pageref*{e:297}}\par
\textbf{Integrationskonstante} (die) -- 積分定数\dotfill\hyperref[e:292]{\pageref*{e:292}}\par
\textbf{Integrationsvariable} (die) -- 積分変数\dotfill\hyperref[e:291]{\pageref*{e:291}}\par
\textbf{integrieren} -- 積分する\dotfill\hyperref[e:296]{\pageref*{e:296}}\par
\textbf{Intervall} (das) -- 区間\dotfill\hyperref[e:270]{\pageref*{e:270}}\par
\textbf{irrationale Zahl} (die) -- 無理数\dotfill\hyperref[e:9]{\pageref*{e:9}}\par

\medskip{\large\textbf{\color{cornflower}K}}\par\nopagebreak
\textbf{Kante} (die) -- 辺\dotfill\hyperref[e:246]{\pageref*{e:246}}\par
\textbf{kartesische Produkt} (das) -- 直積\dotfill\hyperref[e:103]{\pageref*{e:103}}\par
\textbf{Kathete} (die) -- 直角をはさむ辺\dotfill\hyperref[e:187]{\pageref*{e:187}}\par
\textbf{Kegel} (der) -- 円錐\dotfill\hyperref[e:231]{\pageref*{e:231}}\par
\textbf{Kegelschnitt} (der) -- 円錐曲線\dotfill\hyperref[e:312]{\pageref*{e:312}}\par
\textbf{Kegelstumpf} (der) -- 円錐台\dotfill\hyperref[e:236]{\pageref*{e:236}}\par
\textbf{Kehrwert} (der) -- 逆数\dotfill\hyperref[e:39]{\pageref*{e:39}}\par
\textbf{Klammer} (die) -- 括弧\dotfill\hyperref[e:87]{\pageref*{e:87}}\par
\textbf{kleiner als} -- より小さい\dotfill\hyperref[e:63]{\pageref*{e:63}}\par
\textbf{kleiner oder gleich} -- 以下\dotfill\hyperref[e:65]{\pageref*{e:65}}\par
\textbf{kleinste gemeinsame Vielfache} (das) -- 最小公倍数\dotfill\hyperref[e:57]{\pageref*{e:57}}\par
\textbf{Koeffizient} (der) -- 係数\dotfill\hyperref[e:74]{\pageref*{e:74}}\par
\textbf{Kombination} (die) -- 組合せ\dotfill\hyperref[e:343]{\pageref*{e:343}}\par
\textbf{Komplementwinkel} (der) -- 余角\dotfill\hyperref[e:164]{\pageref*{e:164}}\par
\textbf{komplexe Zahl} (die) -- 複素数\dotfill\hyperref[e:10]{\pageref*{e:10}}\par
\textbf{kongruent} -- 合同\dotfill\hyperref[e:175]{\pageref*{e:175}}\par
\textbf{konvergent} -- 収束する\dotfill\hyperref[e:262]{\pageref*{e:262}}\par
\textbf{Koordinaten} (die) -- 座標\dotfill\hyperref[e:116]{\pageref*{e:116}}\par
\textbf{Koordinatensystem} (das) -- 座標系\dotfill\hyperref[e:113]{\pageref*{e:113}}\par
\textbf{Körper} (der) -- 立体\dotfill\hyperref[e:225]{\pageref*{e:225}}\par
\textbf{Kosinus} (der) -- 余弦、コサイン\dotfill\hyperref[e:133]{\pageref*{e:133}}\par
\textbf{Kosinussatz} (der) -- 余弦定理\dotfill\hyperref[e:141]{\pageref*{e:141}}\par
\textbf{Kotangens} (der) -- 余接、コタンジェント\dotfill\hyperref[e:135]{\pageref*{e:135}}\par
\textbf{Kreis} (der) -- 円\dotfill\hyperref[e:211]{\pageref*{e:211}}\par
\textbf{Kreisabschnitt} (der) -- 弓形\dotfill\hyperref[e:222]{\pageref*{e:222}}\par
\textbf{Kreisbogen} (der) -- 弧\dotfill\hyperref[e:221]{\pageref*{e:221}}\par
\textbf{Kreisfläche} (die) -- 円の面積\dotfill\hyperref[e:223]{\pageref*{e:223}}\par
\textbf{Kreisgleichung} (die) -- 円の方程式\dotfill\hyperref[e:313]{\pageref*{e:313}}\par
\textbf{Kreissektor} (der) -- 扇形\dotfill\hyperref[e:217]{\pageref*{e:217}}\par
\textbf{Kreisumfang} (der) -- 円周\dotfill\hyperref[e:215]{\pageref*{e:215}}\par
\textbf{Kreiszahl} (die) -- 円周率\dotfill\hyperref[e:216]{\pageref*{e:216}}\par
\textbf{Kreuzprodukt} (das) -- 外積\dotfill\hyperref[e:307]{\pageref*{e:307}}\par
\textbf{Kubikwurzel} (die) -- 立方根\dotfill\hyperref[e:46]{\pageref*{e:46}}\par
\textbf{Kugel} (die) -- 球\dotfill\hyperref[e:232]{\pageref*{e:232}}\par
\textbf{Kugelsegment} (das) -- 球欠\dotfill\hyperref[e:240]{\pageref*{e:240}}\par
\textbf{Kugelsektor} (der) -- 球扇形\dotfill\hyperref[e:241]{\pageref*{e:241}}\par
\textbf{kürzen} -- 約分する\dotfill\hyperref[e:37]{\pageref*{e:37}}\par

\medskip{\large\textbf{\color{cornflower}L}}\par\nopagebreak
\textbf{leere Menge} (die) -- 空集合\dotfill\hyperref[e:104]{\pageref*{e:104}}\par
\textbf{Leitlinie} (die) -- 準線\dotfill\hyperref[e:318]{\pageref*{e:318}}\par
\textbf{linear abhängig} -- 一次従属\dotfill\hyperref[e:319]{\pageref*{e:319}}\par
\textbf{lineare Funktion} (die) -- 一次関数\dotfill\hyperref[e:119]{\pageref*{e:119}}\par
\textbf{lineare Gleichung} (die) -- 一次方程式\dotfill\hyperref[e:81]{\pageref*{e:81}}\par
\textbf{logarithmische Gleichung} (die) -- 対数方程式\dotfill\hyperref[e:85]{\pageref*{e:85}}\par
\textbf{Logarithmus} (der) -- 対数\dotfill\hyperref[e:95]{\pageref*{e:95}}\par
\textbf{Logarithmusfunktion} (die) -- 対数関数\dotfill\hyperref[e:131]{\pageref*{e:131}}\par
\textbf{Lösung} (die) -- 解\dotfill\hyperref[e:78]{\pageref*{e:78}}\par
\textbf{Lösungsmenge} (die) -- 解の集合\dotfill\hyperref[e:79]{\pageref*{e:79}}\par

\medskip{\large\textbf{\color{cornflower}M}}\par\nopagebreak
\textbf{Mantelfläche} (die) -- 側面\dotfill\hyperref[e:250]{\pageref*{e:250}}\par
\textbf{Matrix} (die) -- 行列\dotfill\hyperref[e:320]{\pageref*{e:320}}\par
\textbf{Median} (der) -- 中央値\dotfill\hyperref[e:339]{\pageref*{e:339}}\par
\textbf{Menge} (die) -- 集合\dotfill\hyperref[e:96]{\pageref*{e:96}}\par
\textbf{Minuend} (der) -- 被減数\dotfill\hyperref[e:19]{\pageref*{e:19}}\par
\textbf{Mittelpunkt} (der) -- 中心\dotfill\hyperref[e:212]{\pageref*{e:212}}\par
\textbf{Mittelsenkrechte} (die) -- 垂直二等分線\dotfill\hyperref[e:170]{\pageref*{e:170}}\par
\textbf{Mittelwert} (der) -- 平均値\dotfill\hyperref[e:338]{\pageref*{e:338}}\par
\textbf{monoton fallend} -- 単調減少\dotfill\hyperref[e:143]{\pageref*{e:143}}\par
\textbf{monoton steigend} -- 単調増加\dotfill\hyperref[e:142]{\pageref*{e:142}}\par
\textbf{multiplizieren} -- 掛ける\dotfill\hyperref[e:21]{\pageref*{e:21}}\par

\medskip{\large\textbf{\color{cornflower}N}}\par\nopagebreak
\textbf{$n$-te Wurzel} (die) -- $n$ 乗根\dotfill\hyperref[e:47]{\pageref*{e:47}}\par
\textbf{natürliche Exponentialfunktion} (die) -- 自然指数関数\dotfill\hyperref[e:299]{\pageref*{e:299}}\par
\textbf{natürliche Logarithmus} (der) -- 自然対数\dotfill\hyperref[e:300]{\pageref*{e:300}}\par
\textbf{natürliche Zahl} (die) -- 自然数\dotfill\hyperref[e:5]{\pageref*{e:5}}\par
\textbf{Nebenwinkel} (der) -- 隣接補角\dotfill\hyperref[e:163]{\pageref*{e:163}}\par
\textbf{Nenner} (der) -- 分母\dotfill\hyperref[e:30]{\pageref*{e:30}}\par
\textbf{Netz} (das) -- 展開図\dotfill\hyperref[e:253]{\pageref*{e:253}}\par
\textbf{Nullfolge} (die) -- 零に収束する数列\dotfill\hyperref[e:258]{\pageref*{e:258}}\par
\textbf{Nullstelle} (die) -- 零点\dotfill\hyperref[e:122]{\pageref*{e:122}}\par
\textbf{Nullvektor} (der) -- 零ベクトル\dotfill\hyperref[e:304]{\pageref*{e:304}}\par

\medskip{\large\textbf{\color{cornflower}O}}\par\nopagebreak
\textbf{Oberflächeninhalt} (der) -- 表面積\dotfill\hyperref[e:256]{\pageref*{e:256}}\par
\textbf{Obersumme} (die) -- 過剰和\dotfill\hyperref[e:293]{\pageref*{e:293}}\par
\textbf{Ordinate} (die) -- $y$ 座標\dotfill\hyperref[e:118]{\pageref*{e:118}}\par
\textbf{Ordnungszahl} (die) -- 序数\dotfill\hyperref[e:3]{\pageref*{e:3}}\par
\textbf{Ortsvektor} (der) -- 位置ベクトル\dotfill\hyperref[e:302]{\pageref*{e:302}}\par

\medskip{\large\textbf{\color{cornflower}P}}\par\nopagebreak
\textbf{Parabel} (die) -- 放物線\dotfill\hyperref[e:125]{\pageref*{e:125}}\par
\textbf{parallel} -- 平行\dotfill\hyperref[e:150]{\pageref*{e:150}}\par
\textbf{Parallelogramm} (das) -- 平行四辺形\dotfill\hyperref[e:200]{\pageref*{e:200}}\par
\textbf{partielle Ableitung} (die) -- 偏導関数\dotfill\hyperref[e:277]{\pageref*{e:277}}\par
\textbf{Permutation} (die) -- 順列\dotfill\hyperref[e:341]{\pageref*{e:341}}\par
\textbf{Polarkoordinaten} (die) -- 極座標\dotfill\hyperref[e:311]{\pageref*{e:311}}\par
\textbf{Polyeder} (das) -- 多面体\dotfill\hyperref[e:233]{\pageref*{e:233}}\par
\textbf{Potenz} (die) -- 累乗\dotfill\hyperref[e:42]{\pageref*{e:42}}\par
\textbf{potenzieren} -- 累乗する\dotfill\hyperref[e:50]{\pageref*{e:50}}\par
\textbf{Potenzreihe} (die) -- べき級数\dotfill\hyperref[e:267]{\pageref*{e:267}}\par
\textbf{Primzahl} (die) -- 素数\dotfill\hyperref[e:12]{\pageref*{e:12}}\par
\textbf{Prisma} (das) -- 角柱\dotfill\hyperref[e:228]{\pageref*{e:228}}\par
\textbf{Produkt} (das) -- 積\dotfill\hyperref[e:22]{\pageref*{e:22}}\par
\textbf{proportional zu} -- に比例する\dotfill\hyperref[e:68]{\pageref*{e:68}}\par
\textbf{Proportionalität} (die) -- 比例\dotfill\hyperref[e:94]{\pageref*{e:94}}\par
\textbf{Prozentsatz} (der) -- 百分率\dotfill\hyperref[e:58]{\pageref*{e:58}}\par
\textbf{Punkt} (der) -- 点\dotfill\hyperref[e:146]{\pageref*{e:146}}\par
\textbf{punktsymmetrisch zum Ursprung} -- 原点に関して対称\dotfill\hyperref[e:145]{\pageref*{e:145}}\par
\textbf{Pyramide} (die) -- 角錐\dotfill\hyperref[e:230]{\pageref*{e:230}}\par
\textbf{Pyramidenstumpf} (der) -- 角錐台\dotfill\hyperref[e:235]{\pageref*{e:235}}\par

\medskip{\large\textbf{\color{cornflower}Q}}\par\nopagebreak
\textbf{Quader} (der) -- 直方体\dotfill\hyperref[e:227]{\pageref*{e:227}}\par
\textbf{Quadrat} (das) -- 正方形\dotfill\hyperref[e:198]{\pageref*{e:198}}\par
\textbf{quadratische Funktion} (die) -- 二次関数\dotfill\hyperref[e:124]{\pageref*{e:124}}\par
\textbf{quadratische Gleichung} (die) -- 二次方程式\dotfill\hyperref[e:82]{\pageref*{e:82}}\par
\textbf{quadratische Matrix} (die) -- 正方行列\dotfill\hyperref[e:321]{\pageref*{e:321}}\par
\textbf{Quadratwurzel} (die) -- 平方根\dotfill\hyperref[e:45]{\pageref*{e:45}}\par
\textbf{quadrieren} -- 二乗する\dotfill\hyperref[e:51]{\pageref*{e:51}}\par
\textbf{Quotient} (der) -- 商\dotfill\hyperref[e:25]{\pageref*{e:25}}\par

\medskip{\large\textbf{\color{cornflower}R}}\par\nopagebreak
\textbf{Radikand} (der) -- 被開方数\dotfill\hyperref[e:48]{\pageref*{e:48}}\par
\textbf{Radius} (der) -- 半径\dotfill\hyperref[e:213]{\pageref*{e:213}}\par
\textbf{radizieren, Wurzel ziehen} -- 累乗根を求める\dotfill\hyperref[e:52]{\pageref*{e:52}}\par
\textbf{rationale Zahl} (die) -- 有理数\dotfill\hyperref[e:7]{\pageref*{e:7}}\par
\textbf{Raumdiagonale} (die) -- 立体の対角線\dotfill\hyperref[e:252]{\pageref*{e:252}}\par
\textbf{Raute} (die) -- ひし形\dotfill\hyperref[e:201]{\pageref*{e:201}}\par
\textbf{rechte Winkel} (der) -- 直角\dotfill\hyperref[e:158]{\pageref*{e:158}}\par
\textbf{Rechteck} (das) -- 長方形\dotfill\hyperref[e:199]{\pageref*{e:199}}\par
\textbf{rechtwinklige Dreieck} (das) -- 直角三角形\dotfill\hyperref[e:181]{\pageref*{e:181}}\par
\textbf{reelle Zahl} (die) -- 実数\dotfill\hyperref[e:8]{\pageref*{e:8}}\par
\textbf{regelmäßige Sechseck} (das) -- 正六角形\dotfill\hyperref[e:205]{\pageref*{e:205}}\par
\textbf{Reihe} (die) -- 級数\dotfill\hyperref[e:265]{\pageref*{e:265}}\par
\textbf{relative Häufigkeit} (die) -- 相対度数\dotfill\hyperref[e:331]{\pageref*{e:331}}\par
\textbf{Rotationsellipsoid} (das) -- 回転楕円体\dotfill\hyperref[e:242]{\pageref*{e:242}}\par
\textbf{Rotationshyperboloid} (das) -- 回転双曲面\dotfill\hyperref[e:244]{\pageref*{e:244}}\par
\textbf{Rotationskörper} (der) -- 回転体\dotfill\hyperref[e:239]{\pageref*{e:239}}\par
\textbf{Rotationsparaboloid} (das) -- 回転放物面\dotfill\hyperref[e:243]{\pageref*{e:243}}\par
\textbf{runden} -- 四捨五入する\dotfill\hyperref[e:59]{\pageref*{e:59}}\par
\textbf{runden Klammern} (die) -- 丸括弧\dotfill\hyperref[e:88]{\pageref*{e:88}}\par

\medskip{\large\textbf{\color{cornflower}S}}\par\nopagebreak
\textbf{Satz} (der) -- 定理\dotfill\hyperref[e:347]{\pageref*{e:347}}\par
\textbf{Satz des Pythagoras} (der) -- 三平方の定理\dotfill\hyperref[e:190]{\pageref*{e:190}}\par
\textbf{Scheitel} (der) -- 頂点\dotfill\hyperref[e:154]{\pageref*{e:154}}\par
\textbf{Scheitelpunkt} (der) -- 頂点\dotfill\hyperref[e:129]{\pageref*{e:129}}\par
\textbf{Scheitelwinkel} (der) -- 対頂角\dotfill\hyperref[e:166]{\pageref*{e:166}}\par
\textbf{Schenkel} (der) -- 辺\dotfill\hyperref[e:155]{\pageref*{e:155}}\par
\textbf{schiefe Zylinder} (der) -- 斜円柱\dotfill\hyperref[e:238]{\pageref*{e:238}}\par
\textbf{Schnittmenge} (die) -- 共通部分\dotfill\hyperref[e:100]{\pageref*{e:100}}\par
\textbf{Schnittpunkt} (der) -- 交点\dotfill\hyperref[e:123]{\pageref*{e:123}}\par
\textbf{Schrägbild} (das) -- 見取図\dotfill\hyperref[e:254]{\pageref*{e:254}}\par
\textbf{Sehne} (die) -- 弦\dotfill\hyperref[e:218]{\pageref*{e:218}}\par
\textbf{Seitenfläche} (die) -- 面\dotfill\hyperref[e:247]{\pageref*{e:247}}\par
\textbf{Seitenhalbierende} (die) -- 中線\dotfill\hyperref[e:192]{\pageref*{e:192}}\par
\textbf{Sekante} (die) -- 割線\dotfill\hyperref[e:219]{\pageref*{e:219}}\par
\textbf{senkrecht} -- 垂直\dotfill\hyperref[e:151]{\pageref*{e:151}}\par
\textbf{Sinus} (der) -- 正弦、サイン\dotfill\hyperref[e:132]{\pageref*{e:132}}\par
\textbf{Sinussatz} (der) -- 正弦定理\dotfill\hyperref[e:140]{\pageref*{e:140}}\par
\textbf{Skalarprodukt} (das) -- 内積\dotfill\hyperref[e:306]{\pageref*{e:306}}\par
\textbf{Spalte} (die) -- 列\dotfill\hyperref[e:323]{\pageref*{e:323}}\par
\textbf{Spiegelung} (die) -- 対称移動\dotfill\hyperref[e:172]{\pageref*{e:172}}\par
\textbf{spitze Winkel} (der) -- 鋭角\dotfill\hyperref[e:157]{\pageref*{e:157}}\par
\textbf{Stammfunktion} (die) -- 原始関数\dotfill\hyperref[e:285]{\pageref*{e:285}}\par
\textbf{Standardabweichung} (die) -- 標準偏差\dotfill\hyperref[e:337]{\pageref*{e:337}}\par
\textbf{Steigung} (die) -- 傾き\dotfill\hyperref[e:120]{\pageref*{e:120}}\par
\textbf{Steigungswinkel} (der) -- 傾斜角\dotfill\hyperref[e:310]{\pageref*{e:310}}\par
\textbf{stetig} -- 連続\dotfill\hyperref[e:268]{\pageref*{e:268}}\par
\textbf{Strahl} (der) -- 半直線\dotfill\hyperref[e:149]{\pageref*{e:149}}\par
\textbf{Strecke} (die) -- 線分\dotfill\hyperref[e:148]{\pageref*{e:148}}\par
\textbf{Stufenwinkel} (der) -- 同位角\dotfill\hyperref[e:167]{\pageref*{e:167}}\par
\textbf{stumpfe Winkel} (der) -- 鈍角\dotfill\hyperref[e:159]{\pageref*{e:159}}\par
\textbf{Subtrahend} (der) -- 減数\dotfill\hyperref[e:20]{\pageref*{e:20}}\par
\textbf{subtrahieren} -- 引く\dotfill\hyperref[e:17]{\pageref*{e:17}}\par
\textbf{Summand} (der) -- 加数\dotfill\hyperref[e:16]{\pageref*{e:16}}\par
\textbf{Summe} (die) -- 和\dotfill\hyperref[e:15]{\pageref*{e:15}}\par
\textbf{Supplementwinkel} (der) -- 補角\dotfill\hyperref[e:165]{\pageref*{e:165}}\par
\textbf{Symmetrieachse} (die) -- 対称の軸\dotfill\hyperref[e:171]{\pageref*{e:171}}\par

\medskip{\large\textbf{\color{cornflower}T}}\par\nopagebreak
\textbf{Tangens} (der) -- 正接、タンジェント\dotfill\hyperref[e:134]{\pageref*{e:134}}\par
\textbf{Tangente} (die) -- 接線\dotfill\hyperref[e:224]{\pageref*{e:224}}\par
\textbf{Tangente} (die) -- 接線\dotfill\hyperref[e:279]{\pageref*{e:279}}\par
\textbf{Teiler} (der) -- 約数\dotfill\hyperref[e:54]{\pageref*{e:54}}\par
\textbf{Teilmenge} (die) -- 部分集合\dotfill\hyperref[e:98]{\pageref*{e:98}}\par
\textbf{Term} (der) -- 式\dotfill\hyperref[e:72]{\pageref*{e:72}}\par
\textbf{Tetraeder} (das) -- 四面体\dotfill\hyperref[e:234]{\pageref*{e:234}}\par
\textbf{Tiefpunkt} (der) -- 極小点\dotfill\hyperref[e:283]{\pageref*{e:283}}\par
\textbf{Trapez} (das) -- 台形\dotfill\hyperref[e:202]{\pageref*{e:202}}\par
\textbf{trigonometrische Gleichung} (die) -- 三角方程式\dotfill\hyperref[e:86]{\pageref*{e:86}}\par

\medskip{\large\textbf{\color{cornflower}U}}\par\nopagebreak
\textbf{überstumpfe Winkel} (der) -- 優角\dotfill\hyperref[e:161]{\pageref*{e:161}}\par
\textbf{Umfang} (der) -- 周の長さ\dotfill\hyperref[e:209]{\pageref*{e:209}}\par
\textbf{Umkreis} (der) -- 外接円\dotfill\hyperref[e:194]{\pageref*{e:194}}\par
\textbf{unabhängige Variable} (die) -- 独立変数\dotfill\hyperref[e:106]{\pageref*{e:106}}\par
\textbf{Unbekannte} (die) -- 未知数\dotfill\hyperref[e:73]{\pageref*{e:73}}\par
\textbf{unbestimmte Integral} (das) -- 不定積分\dotfill\hyperref[e:287]{\pageref*{e:287}}\par
\textbf{unechte Bruch} (der) -- 仮分数\dotfill\hyperref[e:32]{\pageref*{e:32}}\par
\textbf{unendlich} -- 無限大\dotfill\hyperref[e:70]{\pageref*{e:70}}\par
\textbf{unendliche Reihe} (die) -- 無限級数\dotfill\hyperref[e:266]{\pageref*{e:266}}\par
\textbf{ungefähr} -- ほぼ等しい\dotfill\hyperref[e:69]{\pageref*{e:69}}\par
\textbf{ungleich} -- 等しくない\dotfill\hyperref[e:61]{\pageref*{e:61}}\par
\textbf{ungleichnamige Brüche} -- 異分母分数\dotfill\hyperref[e:35]{\pageref*{e:35}}\par
\textbf{Ungleichung} (die) -- 不等式\dotfill\hyperref[e:77]{\pageref*{e:77}}\par
\textbf{unregelmäßige Dreieck} (das) -- 不等辺三角形\dotfill\hyperref[e:180]{\pageref*{e:180}}\par
\textbf{unstetig} -- 不連続\dotfill\hyperref[e:269]{\pageref*{e:269}}\par
\textbf{Untersumme} (die) -- 不足和\dotfill\hyperref[e:294]{\pageref*{e:294}}\par
\textbf{Ursprung} (der) -- 原点\dotfill\hyperref[e:115]{\pageref*{e:115}}\par

\medskip{\large\textbf{\color{cornflower}V}}\par\nopagebreak
\textbf{Variable} (die) -- 変数\dotfill\hyperref[e:71]{\pageref*{e:71}}\par
\textbf{Varianz} (die) -- 分散\dotfill\hyperref[e:336]{\pageref*{e:336}}\par
\textbf{Variation} (die) -- $k$ 個を並べる順列\dotfill\hyperref[e:342]{\pageref*{e:342}}\par
\textbf{Vektor} (der) -- ベクトル\dotfill\hyperref[e:301]{\pageref*{e:301}}\par
\textbf{Vereinigungsmenge} (die) -- 和集合\dotfill\hyperref[e:101]{\pageref*{e:101}}\par
\textbf{Verhältnis} (das) -- 比\dotfill\hyperref[e:40]{\pageref*{e:40}}\par
\textbf{Verhältnisgleichung} (die) -- 比例式\dotfill\hyperref[e:41]{\pageref*{e:41}}\par
\textbf{Verschiebung} (die) -- 平行移動\dotfill\hyperref[e:173]{\pageref*{e:173}}\par
\textbf{viel kleiner als} -- はるかに小さい\dotfill\hyperref[e:67]{\pageref*{e:67}}\par
\textbf{Vieleck} (das) -- 多角形\dotfill\hyperref[e:204]{\pageref*{e:204}}\par
\textbf{Vielfache} (das) -- 倍数\dotfill\hyperref[e:55]{\pageref*{e:55}}\par
\textbf{Viereck} (das) -- 四角形\dotfill\hyperref[e:196]{\pageref*{e:196}}\par
\textbf{Vierfeldertafel} (die) -- 分割表\dotfill\hyperref[e:333]{\pageref*{e:333}}\par
\textbf{Vollwinkel} (der) -- 周角\dotfill\hyperref[e:162]{\pageref*{e:162}}\par
\textbf{Volumen} (das) -- 体積\dotfill\hyperref[e:255]{\pageref*{e:255}}\par
\textbf{Voraussetzung} (die) -- 仮定\dotfill\hyperref[e:348]{\pageref*{e:348}}\par

\medskip{\large\textbf{\color{cornflower}W}}\par\nopagebreak
\textbf{Wahrscheinlichkeit} (die) -- 確率\dotfill\hyperref[e:330]{\pageref*{e:330}}\par
\textbf{Wechselwinkel} (der) -- 錯角\dotfill\hyperref[e:168]{\pageref*{e:168}}\par
\textbf{Wendepunkt} (der) -- 変曲点\dotfill\hyperref[e:284]{\pageref*{e:284}}\par
\textbf{Wertebereich} (der) -- 値域\dotfill\hyperref[e:110]{\pageref*{e:110}}\par
\textbf{Wertetabelle} (die) -- 対応表\dotfill\hyperref[e:112]{\pageref*{e:112}}\par
\textbf{Winkel} (der) -- 角\dotfill\hyperref[e:153]{\pageref*{e:153}}\par
\textbf{Winkelhalbierende} (die) -- 角の二等分線\dotfill\hyperref[e:156]{\pageref*{e:156}}\par
\textbf{Würfel} (der) -- 立方体\dotfill\hyperref[e:226]{\pageref*{e:226}}\par
\textbf{Wurzelexponent} (der) -- 根指数\dotfill\hyperref[e:49]{\pageref*{e:49}}\par
\textbf{Wurzelfunktion} (die) -- 無理関数\dotfill\hyperref[e:128]{\pageref*{e:128}}\par
\textbf{Wurzelgleichung} (die) -- 無理方程式\dotfill\hyperref[e:83]{\pageref*{e:83}}\par

\medskip{\large\textbf{\color{cornflower}X}}\par\nopagebreak
\textbf{$x$-Achse} (die) -- $x$ 軸\dotfill\hyperref[e:114]{\pageref*{e:114}}\par

\medskip{\large\textbf{\color{cornflower}Y}}\par\nopagebreak
\textbf{$y$-Achsenabschnitt} (der) -- $y$ 切片\dotfill\hyperref[e:121]{\pageref*{e:121}}\par

\medskip{\large\textbf{\color{cornflower}Z}}\par\nopagebreak
\textbf{Zahl} (die) -- 数\dotfill\hyperref[e:1]{\pageref*{e:1}}\par
\textbf{Zähler} (der) -- 分子\dotfill\hyperref[e:29]{\pageref*{e:29}}\par
\textbf{Zeile} (die) -- 行\dotfill\hyperref[e:322]{\pageref*{e:322}}\par
\textbf{Ziffer} (die) -- 数字\dotfill\hyperref[e:2]{\pageref*{e:2}}\par
\textbf{Zufallsexperiment} (das) -- 試行\dotfill\hyperref[e:326]{\pageref*{e:326}}\par
\textbf{Zufallsgröße} (die) -- 確率変数\dotfill\hyperref[e:334]{\pageref*{e:334}}\par
\textbf{Zwölfeck} (das) -- 十二角形\dotfill\hyperref[e:208]{\pageref*{e:208}}\par
\textbf{Zylinder} (der) -- 円柱\dotfill\hyperref[e:229]{\pageref*{e:229}}\par
\end{multicols}

\newpage
\section{Register Japanisch -- Deutsch}\renewcommand{\lesenabstand}{1.12}\renewcommand{\arraystretch}{0.93}

Die Begriffe sind nach ihrer Lesung geordnet, in der Reihenfolge der japanischen Silbentafel
(a, ka, sa, ta, na, ha, ma, ya, ra, wa). Begriffe, die mit einer Variablen beginnen, stehen am Anfang.

\begin{multicols}{2}\small\raggedright\raggedcolumns

\smallskip{\large\textbf{\color{cornflower}A -- Z}}\par\nopagebreak
\jfett{$k$ 個を並べる順列} {\footnotesize\color{black!70}$k$ こをならべるじゅんれつ} -- die Variation\dotfill\hyperref[e:342]{\pageref*{e:342}}\par
\jfett{$x$ 座標} {\footnotesize\color{black!70}$x$ ざひょう} -- die Abszisse\dotfill\hyperref[e:117]{\pageref*{e:117}}\par
\jfett{$y$ 座標} {\footnotesize\color{black!70}$y$ ざひょう} -- die Ordinate\dotfill\hyperref[e:118]{\pageref*{e:118}}\par
\jfett{$x$ 軸} {\footnotesize\color{black!70}$x$ じく} -- die $x$-Achse\dotfill\hyperref[e:114]{\pageref*{e:114}}\par
\jfett{$y$ 軸に関して対称} {\footnotesize\color{black!70}$y$ じくにかんしてたいしょう} -- achsensymmetrisch zur $y$-Achse\dotfill\hyperref[e:144]{\pageref*{e:144}}\par
\jfett{$n$ 乗根} {\footnotesize\color{black!70}$n$ じょうこん} -- die $n$-te Wurzel\dotfill\hyperref[e:47]{\pageref*{e:47}}\par
\jfett{$y$ 切片} {\footnotesize\color{black!70}$y$ せっぺん} -- der $y$-Achsenabschnitt\dotfill\hyperref[e:121]{\pageref*{e:121}}\par

\smallskip{\large\textbf{\color{cornflower}あ \,·\, a}}\par\nopagebreak
\jfett{以下} {\footnotesize\color{black!70}いか} -- kleiner oder gleich\dotfill\hyperref[e:65]{\pageref*{e:65}}\par
\jfett{以上} {\footnotesize\color{black!70}いじょう} -- größer oder gleich\dotfill\hyperref[e:66]{\pageref*{e:66}}\par
\jfett{一次関数} {\footnotesize\color{black!70}いちじかんすう} -- die lineare Funktion\dotfill\hyperref[e:119]{\pageref*{e:119}}\par
\jfett{一次従属} {\footnotesize\color{black!70}いちじじゅうぞく} -- linear abhängig\dotfill\hyperref[e:319]{\pageref*{e:319}}\par
\jfett{一次方程式} {\footnotesize\color{black!70}いちじほうていしき} -- die lineare Gleichung\dotfill\hyperref[e:81]{\pageref*{e:81}}\par
\jfett{位置ベクトル} {\footnotesize\color{black!70}いちベクトル} -- der Ortsvektor\dotfill\hyperref[e:302]{\pageref*{e:302}}\par
\jfett{異分母分数} {\footnotesize\color{black!70}いぶんぼぶんすう} -- ungleichnamige Brüche\dotfill\hyperref[e:35]{\pageref*{e:35}}\par
\jfett{因数} {\footnotesize\color{black!70}いんすう} -- der Faktor\dotfill\hyperref[e:23]{\pageref*{e:23}}\par
\jfett{鋭角} {\footnotesize\color{black!70}えいかく} -- der spitze Winkel\dotfill\hyperref[e:157]{\pageref*{e:157}}\par
\jfett{円} {\footnotesize\color{black!70}えん} -- der Kreis\dotfill\hyperref[e:211]{\pageref*{e:211}}\par
\jfett{円周} {\footnotesize\color{black!70}えんしゅう} -- der Kreisumfang\dotfill\hyperref[e:215]{\pageref*{e:215}}\par
\jfett{円周率} {\footnotesize\color{black!70}えんしゅうりつ} -- die Kreiszahl\dotfill\hyperref[e:216]{\pageref*{e:216}}\par
\jfett{円錐} {\footnotesize\color{black!70}えんすい} -- der Kegel\dotfill\hyperref[e:231]{\pageref*{e:231}}\par
\jfett{円錐曲線} {\footnotesize\color{black!70}えんすいきょくせん} -- der Kegelschnitt\dotfill\hyperref[e:312]{\pageref*{e:312}}\par
\jfett{円錐台} {\footnotesize\color{black!70}えんすいだい} -- der Kegelstumpf\dotfill\hyperref[e:236]{\pageref*{e:236}}\par
\jfett{円柱} {\footnotesize\color{black!70}えんちゅう} -- der Zylinder\dotfill\hyperref[e:229]{\pageref*{e:229}}\par
\jfett{円の方程式} {\footnotesize\color{black!70}えんのほうていしき} -- die Kreisgleichung\dotfill\hyperref[e:313]{\pageref*{e:313}}\par
\jfett{円の面積} {\footnotesize\color{black!70}えんのめんせき} -- die Kreisfläche\dotfill\hyperref[e:223]{\pageref*{e:223}}\par
\jfett{扇形} {\footnotesize\color{black!70}おうぎがた} -- der Kreissektor\dotfill\hyperref[e:217]{\pageref*{e:217}}\par

\smallskip{\large\textbf{\color{cornflower}か \,·\, ka}}\par\nopagebreak
\jfett{解} {\footnotesize\color{black!70}かい} -- die Lösung\dotfill\hyperref[e:78]{\pageref*{e:78}}\par
\jfett{外角} {\footnotesize\color{black!70}がいかく} -- der Außenwinkel\dotfill\hyperref[e:179]{\pageref*{e:179}}\par
\jfett{階乗} {\footnotesize\color{black!70}かいじょう} -- die Fakultät\dotfill\hyperref[e:340]{\pageref*{e:340}}\par
\jfett{外積} {\footnotesize\color{black!70}がいせき} -- das Kreuzprodukt\dotfill\hyperref[e:307]{\pageref*{e:307}}\par
\jfett{外接円} {\footnotesize\color{black!70}がいせつえん} -- der Umkreis\dotfill\hyperref[e:194]{\pageref*{e:194}}\par
\jfett{回転移動} {\footnotesize\color{black!70}かいてんいどう} -- die Drehung\dotfill\hyperref[e:174]{\pageref*{e:174}}\par
\jfett{回転双曲面} {\footnotesize\color{black!70}かいてんそうきょくめん} -- das Rotationshyperboloid\dotfill\hyperref[e:244]{\pageref*{e:244}}\par
\jfett{回転体} {\footnotesize\color{black!70}かいてんたい} -- der Rotationskörper\dotfill\hyperref[e:239]{\pageref*{e:239}}\par
\jfett{回転楕円体} {\footnotesize\color{black!70}かいてんだえんたい} -- das Rotationsellipsoid\dotfill\hyperref[e:242]{\pageref*{e:242}}\par
\jfett{回転放物面} {\footnotesize\color{black!70}かいてんほうぶつめん} -- das Rotationsparaboloid\dotfill\hyperref[e:243]{\pageref*{e:243}}\par
\jfett{解の集合} {\footnotesize\color{black!70}かいのしゅうごう} -- die Lösungsmenge\dotfill\hyperref[e:79]{\pageref*{e:79}}\par
\jfett{角} {\footnotesize\color{black!70}かく} -- der Winkel\dotfill\hyperref[e:153]{\pageref*{e:153}}\par
\jfett{角括弧} {\footnotesize\color{black!70}かくかっこ} -- die eckigen Klammern\dotfill\hyperref[e:89]{\pageref*{e:89}}\par
\jfett{角錐} {\footnotesize\color{black!70}かくすい} -- die Pyramide\dotfill\hyperref[e:230]{\pageref*{e:230}}\par
\jfett{角錐台} {\footnotesize\color{black!70}かくすいだい} -- der Pyramidenstumpf\dotfill\hyperref[e:235]{\pageref*{e:235}}\par
\jfett{角柱} {\footnotesize\color{black!70}かくちゅう} -- das Prisma\dotfill\hyperref[e:228]{\pageref*{e:228}}\par
\jfett{角の二等分線} {\footnotesize\color{black!70}かくのにとうぶんせん} -- die Winkelhalbierende\dotfill\hyperref[e:156]{\pageref*{e:156}}\par
\jfett{確率} {\footnotesize\color{black!70}かくりつ} -- die Wahrscheinlichkeit\dotfill\hyperref[e:330]{\pageref*{e:330}}\par
\jfett{確率変数} {\footnotesize\color{black!70}かくりつへんすう} -- die Zufallsgröße\dotfill\hyperref[e:334]{\pageref*{e:334}}\par
\jfett{掛ける} {\footnotesize\color{black!70}かける} -- multiplizieren\dotfill\hyperref[e:21]{\pageref*{e:21}}\par
\jfett{過剰和} {\footnotesize\color{black!70}かじょうわ} -- die Obersumme\dotfill\hyperref[e:293]{\pageref*{e:293}}\par
\jfett{数} {\footnotesize\color{black!70}かず} -- die Zahl\dotfill\hyperref[e:1]{\pageref*{e:1}}\par
\jfett{加数} {\footnotesize\color{black!70}かすう} -- der Summand\dotfill\hyperref[e:16]{\pageref*{e:16}}\par
\jfett{傾き} {\footnotesize\color{black!70}かたむき} -- die Steigung\dotfill\hyperref[e:120]{\pageref*{e:120}}\par
\jfett{括弧} {\footnotesize\color{black!70}かっこ} -- die Klammer\dotfill\hyperref[e:87]{\pageref*{e:87}}\par
\jfett{割線} {\footnotesize\color{black!70}かっせん} -- die Sekante\dotfill\hyperref[e:219]{\pageref*{e:219}}\par
\jfett{仮定} {\footnotesize\color{black!70}かてい} -- die Voraussetzung\dotfill\hyperref[e:348]{\pageref*{e:348}}\par
\jfett{仮分数} {\footnotesize\color{black!70}かぶんすう} -- der unechte Bruch\dotfill\hyperref[e:32]{\pageref*{e:32}}\par
\jfett{関数} {\footnotesize\color{black!70}かんすう} -- die Funktion\dotfill\hyperref[e:105]{\pageref*{e:105}}\par
\jfett{関数の値} {\footnotesize\color{black!70}かんすうのあたい} -- der Funktionswert\dotfill\hyperref[e:108]{\pageref*{e:108}}\par
\jfett{関数の増減を調べる} {\footnotesize\color{black!70}かんすうのぞうげんをしらべる} -- die Funktionsuntersuchung\dotfill\hyperref[e:280]{\pageref*{e:280}}\par
\jfett{期待値} {\footnotesize\color{black!70}きたいち} -- der Erwartungswert\dotfill\hyperref[e:335]{\pageref*{e:335}}\par
\jfett{逆三角関数} {\footnotesize\color{black!70}ぎゃくさんかくかんすう} -- die Arkusfunktion\dotfill\hyperref[e:136]{\pageref*{e:136}}\par
\jfett{逆数} {\footnotesize\color{black!70}ぎゃくすう} -- der Kehrwert\dotfill\hyperref[e:39]{\pageref*{e:39}}\par
\jfett{球} {\footnotesize\color{black!70}きゅう} -- die Kugel\dotfill\hyperref[e:232]{\pageref*{e:232}}\par
\jfett{弓形} {\footnotesize\color{black!70}きゅうけい} -- der Kreisabschnitt\dotfill\hyperref[e:222]{\pageref*{e:222}}\par
\jfett{球欠} {\footnotesize\color{black!70}きゅうけつ} -- das Kugelsegment\dotfill\hyperref[e:240]{\pageref*{e:240}}\par
\jfett{級数} {\footnotesize\color{black!70}きゅうすう} -- die Reihe\dotfill\hyperref[e:265]{\pageref*{e:265}}\par
\jfett{球扇形} {\footnotesize\color{black!70}きゅうせんけい} -- der Kugelsektor\dotfill\hyperref[e:241]{\pageref*{e:241}}\par
\jfett{行} {\footnotesize\color{black!70}ぎょう} -- die Zeile\dotfill\hyperref[e:322]{\pageref*{e:322}}\par
\jfett{共通因数でくくる} {\footnotesize\color{black!70}きょうつういんすうでくくる} -- ausklammern\dotfill\hyperref[e:92]{\pageref*{e:92}}\par
\jfett{共通部分} {\footnotesize\color{black!70}きょうつうぶぶん} -- die Schnittmenge\dotfill\hyperref[e:100]{\pageref*{e:100}}\par
\jfett{行列} {\footnotesize\color{black!70}ぎょうれつ} -- die Matrix\dotfill\hyperref[e:320]{\pageref*{e:320}}\par
\jfett{行列式} {\footnotesize\color{black!70}ぎょうれつしき} -- die Determinante\dotfill\hyperref[e:325]{\pageref*{e:325}}\par
\jfett{極限値} {\footnotesize\color{black!70}きょくげんち} -- der Grenzwert\dotfill\hyperref[e:264]{\pageref*{e:264}}\par
\jfett{極座標} {\footnotesize\color{black!70}きょくざひょう} -- die Polarkoordinaten\dotfill\hyperref[e:311]{\pageref*{e:311}}\par
\jfett{極小点} {\footnotesize\color{black!70}きょくしょうてん} -- der Tiefpunkt\dotfill\hyperref[e:283]{\pageref*{e:283}}\par
\jfett{極大点} {\footnotesize\color{black!70}きょくだいてん} -- der Hochpunkt\dotfill\hyperref[e:282]{\pageref*{e:282}}\par
\jfett{極値をとる点} {\footnotesize\color{black!70}きょくちをとるてん} -- der Extrempunkt\dotfill\hyperref[e:281]{\pageref*{e:281}}\par
\jfett{虚数単位} {\footnotesize\color{black!70}きょすうたんい} -- die imaginäre Einheit\dotfill\hyperref[e:11]{\pageref*{e:11}}\par
\jfett{距離} {\footnotesize\color{black!70}きょり} -- der Abstand\dotfill\hyperref[e:152]{\pageref*{e:152}}\par
\jfett{空集合} {\footnotesize\color{black!70}くうしゅうごう} -- die leere Menge\dotfill\hyperref[e:104]{\pageref*{e:104}}\par
\jfett{区間} {\footnotesize\color{black!70}くかん} -- das Intervall\dotfill\hyperref[e:270]{\pageref*{e:270}}\par
\jfett{組合せ} {\footnotesize\color{black!70}くみあわせ} -- die Kombination\dotfill\hyperref[e:343]{\pageref*{e:343}}\par
\jfett{グラフ} -- der Funktionsgraph\dotfill\hyperref[e:111]{\pageref*{e:111}}\par
\jfett{傾斜角} {\footnotesize\color{black!70}けいしゃかく} -- der Steigungswinkel\dotfill\hyperref[e:310]{\pageref*{e:310}}\par
\jfett{係数} {\footnotesize\color{black!70}けいすう} -- der Koeffizient\dotfill\hyperref[e:74]{\pageref*{e:74}}\par
\jfett{結果} {\footnotesize\color{black!70}けっか} -- das Ergebnis\dotfill\hyperref[e:327]{\pageref*{e:327}}\par
\jfett{弦} {\footnotesize\color{black!70}げん} -- die Sehne\dotfill\hyperref[e:218]{\pageref*{e:218}}\par
\jfett{原始関数} {\footnotesize\color{black!70}げんしかんすう} -- die Stammfunktion\dotfill\hyperref[e:285]{\pageref*{e:285}}\par
\jfett{減数} {\footnotesize\color{black!70}げんすう} -- der Subtrahend\dotfill\hyperref[e:20]{\pageref*{e:20}}\par
\jfett{原点} {\footnotesize\color{black!70}げんてん} -- der Ursprung\dotfill\hyperref[e:115]{\pageref*{e:115}}\par
\jfett{原点に関して対称} {\footnotesize\color{black!70}げんてんにかんしてたいしょう} -- punktsymmetrisch zum Ursprung\dotfill\hyperref[e:145]{\pageref*{e:145}}\par
\jfett{弧} {\footnotesize\color{black!70}こ} -- der Kreisbogen\dotfill\hyperref[e:221]{\pageref*{e:221}}\par
\jfett{交点} {\footnotesize\color{black!70}こうてん} -- der Schnittpunkt\dotfill\hyperref[e:123]{\pageref*{e:123}}\par
\jfett{合同} {\footnotesize\color{black!70}ごうどう} -- kongruent\dotfill\hyperref[e:175]{\pageref*{e:175}}\par
\jfett{恒等的に等しい} {\footnotesize\color{black!70}こうとうてきにひとしい} -- identisch gleich\dotfill\hyperref[e:62]{\pageref*{e:62}}\par
\jfett{五角形} {\footnotesize\color{black!70}ごかくけい} -- das Fünfeck\dotfill\hyperref[e:206]{\pageref*{e:206}}\par
\jfett{弧度法} {\footnotesize\color{black!70}こどほう} -- das Bogenmaß\dotfill\hyperref[e:139]{\pageref*{e:139}}\par
\jfett{根指数} {\footnotesize\color{black!70}こんしすう} -- der Wurzelexponent\dotfill\hyperref[e:49]{\pageref*{e:49}}\par

\smallskip{\large\textbf{\color{cornflower}さ \,·\, sa}}\par\nopagebreak
\jfett{差} {\footnotesize\color{black!70}さ} -- die Differenz\dotfill\hyperref[e:18]{\pageref*{e:18}}\par
\jfett{最小公倍数} {\footnotesize\color{black!70}さいしょうこうばいすう} -- das kleinste gemeinsame Vielfache\dotfill\hyperref[e:57]{\pageref*{e:57}}\par
\jfett{最小公分母} {\footnotesize\color{black!70}さいしょうこうぶんぼ} -- der Hauptnenner\dotfill\hyperref[e:36]{\pageref*{e:36}}\par
\jfett{最大公約数} {\footnotesize\color{black!70}さいだいこうやくすう} -- der größte gemeinsame Teiler\dotfill\hyperref[e:56]{\pageref*{e:56}}\par
\jfett{差集合} {\footnotesize\color{black!70}さしゅうごう} -- die Differenzmenge\dotfill\hyperref[e:102]{\pageref*{e:102}}\par
\jfett{錯角} {\footnotesize\color{black!70}さっかく} -- der Wechselwinkel\dotfill\hyperref[e:168]{\pageref*{e:168}}\par
\jfett{座標} {\footnotesize\color{black!70}ざひょう} -- die Koordinaten\dotfill\hyperref[e:116]{\pageref*{e:116}}\par
\jfett{座標系} {\footnotesize\color{black!70}ざひょうけい} -- das Koordinatensystem\dotfill\hyperref[e:113]{\pageref*{e:113}}\par
\jfett{三角形} {\footnotesize\color{black!70}さんかくけい} -- das Dreieck\dotfill\hyperref[e:177]{\pageref*{e:177}}\par
\jfett{三角方程式} {\footnotesize\color{black!70}さんかくほうていしき} -- die trigonometrische Gleichung\dotfill\hyperref[e:86]{\pageref*{e:86}}\par
\jfett{三平方の定理} {\footnotesize\color{black!70}さんへいほうのていり} -- der Satz des Pythagoras\dotfill\hyperref[e:190]{\pageref*{e:190}}\par
\jfett{四角形} {\footnotesize\color{black!70}しかくけい} -- das Viereck\dotfill\hyperref[e:196]{\pageref*{e:196}}\par
\jfett{式} {\footnotesize\color{black!70}しき} -- der Term\dotfill\hyperref[e:72]{\pageref*{e:72}}\par
\jfett{試行} {\footnotesize\color{black!70}しこう} -- das Zufallsexperiment\dotfill\hyperref[e:326]{\pageref*{e:326}}\par
\jfett{四捨五入する} {\footnotesize\color{black!70}ししゃごにゅうする} -- runden\dotfill\hyperref[e:59]{\pageref*{e:59}}\par
\jfett{事象} {\footnotesize\color{black!70}じしょう} -- das Ereignis\dotfill\hyperref[e:328]{\pageref*{e:328}}\par
\jfett{指数} {\footnotesize\color{black!70}しすう} -- der Exponent\dotfill\hyperref[e:44]{\pageref*{e:44}}\par
\jfett{次数} {\footnotesize\color{black!70}じすう} -- der Grad\dotfill\hyperref[e:75]{\pageref*{e:75}}\par
\jfett{指数関数} {\footnotesize\color{black!70}しすうかんすう} -- die Exponentialfunktion\dotfill\hyperref[e:130]{\pageref*{e:130}}\par
\jfett{指数方程式} {\footnotesize\color{black!70}しすうほうていしき} -- die Exponentialgleichung\dotfill\hyperref[e:84]{\pageref*{e:84}}\par
\jfett{自然指数関数} {\footnotesize\color{black!70}しぜんしすうかんすう} -- die natürliche Exponentialfunktion\dotfill\hyperref[e:299]{\pageref*{e:299}}\par
\jfett{自然数} {\footnotesize\color{black!70}しぜんすう} -- die natürliche Zahl\dotfill\hyperref[e:5]{\pageref*{e:5}}\par
\jfett{自然対数} {\footnotesize\color{black!70}しぜんたいすう} -- der natürliche Logarithmus\dotfill\hyperref[e:300]{\pageref*{e:300}}\par
\jfett{実数} {\footnotesize\color{black!70}じっすう} -- die reelle Zahl\dotfill\hyperref[e:8]{\pageref*{e:8}}\par
\jfett{四面体} {\footnotesize\color{black!70}しめんたい} -- das Tetraeder\dotfill\hyperref[e:234]{\pageref*{e:234}}\par
\jfett{斜円柱} {\footnotesize\color{black!70}しゃえんちゅう} -- der schiefe Zylinder\dotfill\hyperref[e:238]{\pageref*{e:238}}\par
\jfett{斜辺} {\footnotesize\color{black!70}しゃへん} -- die Hypotenuse\dotfill\hyperref[e:186]{\pageref*{e:186}}\par
\jfett{周角} {\footnotesize\color{black!70}しゅうかく} -- der Vollwinkel\dotfill\hyperref[e:162]{\pageref*{e:162}}\par
\jfett{集合} {\footnotesize\color{black!70}しゅうごう} -- die Menge\dotfill\hyperref[e:96]{\pageref*{e:96}}\par
\jfett{収束する} {\footnotesize\color{black!70}しゅうそくする} -- konvergent\dotfill\hyperref[e:262]{\pageref*{e:262}}\par
\jfett{従属変数} {\footnotesize\color{black!70}じゅうぞくへんすう} -- die abhängige Variable\dotfill\hyperref[e:107]{\pageref*{e:107}}\par
\jfett{十二角形} {\footnotesize\color{black!70}じゅうにかくけい} -- das Zwölfeck\dotfill\hyperref[e:208]{\pageref*{e:208}}\par
\jfett{周の長さ} {\footnotesize\color{black!70}しゅうのながさ} -- der Umfang\dotfill\hyperref[e:209]{\pageref*{e:209}}\par
\jfett{樹形図} {\footnotesize\color{black!70}じゅけいず} -- das Baumdiagramm\dotfill\hyperref[e:332]{\pageref*{e:332}}\par
\jfett{主対角線} {\footnotesize\color{black!70}しゅたいかくせん} -- die Hauptdiagonale\dotfill\hyperref[e:324]{\pageref*{e:324}}\par
\jfett{準線} {\footnotesize\color{black!70}じゅんせん} -- die Leitlinie\dotfill\hyperref[e:318]{\pageref*{e:318}}\par
\jfett{順列} {\footnotesize\color{black!70}じゅんれつ} -- die Permutation\dotfill\hyperref[e:341]{\pageref*{e:341}}\par
\jfett{商} {\footnotesize\color{black!70}しょう} -- der Quotient\dotfill\hyperref[e:25]{\pageref*{e:25}}\par
\jfett{小数} {\footnotesize\color{black!70}しょうすう} -- die Dezimalzahl\dotfill\hyperref[e:13]{\pageref*{e:13}}\par
\jfett{上端、下端} {\footnotesize\color{black!70}じょうたん、かたん} -- die Integrationsgrenze\dotfill\hyperref[e:297]{\pageref*{e:297}}\par
\jfett{焦点} {\footnotesize\color{black!70}しょうてん} -- der Brennpunkt\dotfill\hyperref[e:316]{\pageref*{e:316}}\par
\jfett{乗法公式} {\footnotesize\color{black!70}じょうほうこうしき} -- die binomische Formel\dotfill\hyperref[e:93]{\pageref*{e:93}}\par
\jfett{証明} {\footnotesize\color{black!70}しょうめい} -- der Beweis\dotfill\hyperref[e:349]{\pageref*{e:349}}\par
\jfett{上面} {\footnotesize\color{black!70}じょうめん} -- die Deckfläche\dotfill\hyperref[e:249]{\pageref*{e:249}}\par
\jfett{初項} {\footnotesize\color{black!70}しょこう} -- das Anfangsglied\dotfill\hyperref[e:261]{\pageref*{e:261}}\par
\jfett{除数} {\footnotesize\color{black!70}じょすう} -- der Divisor\dotfill\hyperref[e:27]{\pageref*{e:27}}\par
\jfett{序数} {\footnotesize\color{black!70}じょすう} -- die Ordnungszahl\dotfill\hyperref[e:3]{\pageref*{e:3}}\par
\jfett{真部分集合} {\footnotesize\color{black!70}しんぶぶんしゅうごう} -- die echte Teilmenge\dotfill\hyperref[e:99]{\pageref*{e:99}}\par
\jfett{真分数} {\footnotesize\color{black!70}しんぶんすう} -- der echte Bruch\dotfill\hyperref[e:31]{\pageref*{e:31}}\par
\jfett{垂直} {\footnotesize\color{black!70}すいちょく} -- senkrecht\dotfill\hyperref[e:151]{\pageref*{e:151}}\par
\jfett{垂直二等分線} {\footnotesize\color{black!70}すいちょくにとうぶんせん} -- die Mittelsenkrechte\dotfill\hyperref[e:170]{\pageref*{e:170}}\par
\jfett{数字} {\footnotesize\color{black!70}すうじ} -- die Ziffer\dotfill\hyperref[e:2]{\pageref*{e:2}}\par
\jfett{数列} {\footnotesize\color{black!70}すうれつ} -- die Folge\dotfill\hyperref[e:257]{\pageref*{e:257}}\par
\jfett{すべての} -- für alle\dotfill\hyperref[e:352]{\pageref*{e:352}}\par
\jfett{正弦、サイン} {\footnotesize\color{black!70}せいげん、サイン} -- der Sinus\dotfill\hyperref[e:132]{\pageref*{e:132}}\par
\jfett{正弦定理} {\footnotesize\color{black!70}せいげんていり} -- der Sinussatz\dotfill\hyperref[e:140]{\pageref*{e:140}}\par
\jfett{正三角形} {\footnotesize\color{black!70}せいさんかくけい} -- das gleichseitige Dreieck\dotfill\hyperref[e:185]{\pageref*{e:185}}\par
\jfett{整数} {\footnotesize\color{black!70}せいすう} -- die ganze Zahl\dotfill\hyperref[e:6]{\pageref*{e:6}}\par
\jfett{正接、タンジェント} {\footnotesize\color{black!70}せいせつ、タンジェント} -- der Tangens\dotfill\hyperref[e:134]{\pageref*{e:134}}\par
\jfett{正方行列} {\footnotesize\color{black!70}せいほうぎょうれつ} -- die quadratische Matrix\dotfill\hyperref[e:321]{\pageref*{e:321}}\par
\jfett{正方形} {\footnotesize\color{black!70}せいほうけい} -- das Quadrat\dotfill\hyperref[e:198]{\pageref*{e:198}}\par
\jfett{正六角形} {\footnotesize\color{black!70}せいろっかくけい} -- das regelmäßige Sechseck\dotfill\hyperref[e:205]{\pageref*{e:205}}\par
\jfett{積} {\footnotesize\color{black!70}せき} -- das Produkt\dotfill\hyperref[e:22]{\pageref*{e:22}}\par
\jfett{積分} {\footnotesize\color{black!70}せきぶん} -- das Integral\dotfill\hyperref[e:295]{\pageref*{e:295}}\par
\jfett{積分記号} {\footnotesize\color{black!70}せきぶんきごう} -- das Integralzeichen\dotfill\hyperref[e:289]{\pageref*{e:289}}\par
\jfett{積分する} {\footnotesize\color{black!70}せきぶんする} -- integrieren\dotfill\hyperref[e:296]{\pageref*{e:296}}\par
\jfett{積分定数} {\footnotesize\color{black!70}せきぶんていすう} -- die Integrationskonstante\dotfill\hyperref[e:292]{\pageref*{e:292}}\par
\jfett{積分変数} {\footnotesize\color{black!70}せきぶんへんすう} -- die Integrationsvariable\dotfill\hyperref[e:291]{\pageref*{e:291}}\par
\jfett{積分法} {\footnotesize\color{black!70}せきぶんほう} -- die Integralrechnung\dotfill\hyperref[e:286]{\pageref*{e:286}}\par
\jfett{接線} {\footnotesize\color{black!70}せっせん} -- die Tangente\dotfill\hyperref[e:224]{\pageref*{e:224}}\par
\jfett{接線} {\footnotesize\color{black!70}せっせん} -- die Tangente\dotfill\hyperref[e:279]{\pageref*{e:279}}\par
\jfett{絶対値} {\footnotesize\color{black!70}ぜったいち} -- der Betrag\dotfill\hyperref[e:53]{\pageref*{e:53}}\par
\jfett{接点} {\footnotesize\color{black!70}せってん} -- der Berührungspunkt\dotfill\hyperref[e:220]{\pageref*{e:220}}\par
\jfett{漸近線} {\footnotesize\color{black!70}ぜんきんせん} -- die Asymptote\dotfill\hyperref[e:271]{\pageref*{e:271}}\par
\jfett{線分} {\footnotesize\color{black!70}せんぶん} -- die Strecke\dotfill\hyperref[e:148]{\pageref*{e:148}}\par
\jfett{双曲線} {\footnotesize\color{black!70}そうきょくせん} -- die Hyperbel\dotfill\hyperref[e:315]{\pageref*{e:315}}\par
\jfett{相似} {\footnotesize\color{black!70}そうじ} -- ähnlich\dotfill\hyperref[e:176]{\pageref*{e:176}}\par
\jfett{相対度数} {\footnotesize\color{black!70}そうたいどすう} -- die relative Häufigkeit\dotfill\hyperref[e:331]{\pageref*{e:331}}\par
\jfett{添え字} {\footnotesize\color{black!70}そえじ} -- der Index\dotfill\hyperref[e:4]{\pageref*{e:4}}\par
\jfett{側面} {\footnotesize\color{black!70}そくめん} -- die Mantelfläche\dotfill\hyperref[e:250]{\pageref*{e:250}}\par
\jfett{素数} {\footnotesize\color{black!70}そすう} -- die Primzahl\dotfill\hyperref[e:12]{\pageref*{e:12}}\par
\jfett{存在する} {\footnotesize\color{black!70}そんざいする} -- es gibt\dotfill\hyperref[e:353]{\pageref*{e:353}}\par

\smallskip{\large\textbf{\color{cornflower}た \,·\, ta}}\par\nopagebreak
\jfett{対応表} {\footnotesize\color{black!70}たいおうひょう} -- die Wertetabelle\dotfill\hyperref[e:112]{\pageref*{e:112}}\par
\jfett{対角線} {\footnotesize\color{black!70}たいかくせん} -- die Diagonale\dotfill\hyperref[e:197]{\pageref*{e:197}}\par
\jfett{台形} {\footnotesize\color{black!70}だいけい} -- das Trapez\dotfill\hyperref[e:202]{\pageref*{e:202}}\par
\jfett{対称移動} {\footnotesize\color{black!70}たいしょういどう} -- die Spiegelung\dotfill\hyperref[e:172]{\pageref*{e:172}}\par
\jfett{対称の軸} {\footnotesize\color{black!70}たいしょうのじく} -- die Symmetrieachse\dotfill\hyperref[e:171]{\pageref*{e:171}}\par
\jfett{対数} {\footnotesize\color{black!70}たいすう} -- der Logarithmus\dotfill\hyperref[e:95]{\pageref*{e:95}}\par
\jfett{対数関数} {\footnotesize\color{black!70}たいすうかんすう} -- die Logarithmusfunktion\dotfill\hyperref[e:131]{\pageref*{e:131}}\par
\jfett{対数方程式} {\footnotesize\color{black!70}たいすうほうていしき} -- die logarithmische Gleichung\dotfill\hyperref[e:85]{\pageref*{e:85}}\par
\jfett{体積} {\footnotesize\color{black!70}たいせき} -- das Volumen\dotfill\hyperref[e:255]{\pageref*{e:255}}\par
\jfett{対頂角} {\footnotesize\color{black!70}たいちょうかく} -- der Scheitelwinkel\dotfill\hyperref[e:166]{\pageref*{e:166}}\par
\jfett{帯分数} {\footnotesize\color{black!70}たいぶんすう} -- die gemischte Zahl\dotfill\hyperref[e:33]{\pageref*{e:33}}\par
\jfett{対辺} {\footnotesize\color{black!70}たいへん} -- die Gegenkathete\dotfill\hyperref[e:189]{\pageref*{e:189}}\par
\jfett{楕円} {\footnotesize\color{black!70}だえん} -- die Ellipse\dotfill\hyperref[e:314]{\pageref*{e:314}}\par
\jfett{多角形} {\footnotesize\color{black!70}たかくけい} -- das Vieleck\dotfill\hyperref[e:204]{\pageref*{e:204}}\par
\jfett{高さ} {\footnotesize\color{black!70}たかさ} -- die Höhe\dotfill\hyperref[e:191]{\pageref*{e:191}}\par
\jfett{多項式関数} {\footnotesize\color{black!70}たこうしきかんすう} -- die ganzrationale Funktion\dotfill\hyperref[e:126]{\pageref*{e:126}}\par
\jfett{たこ形} {\footnotesize\color{black!70}たこがた} -- das Drachenviereck\dotfill\hyperref[e:203]{\pageref*{e:203}}\par
\jfett{足す} {\footnotesize\color{black!70}たす} -- addieren\dotfill\hyperref[e:14]{\pageref*{e:14}}\par
\jfett{多面体} {\footnotesize\color{black!70}ためんたい} -- das Polyeder\dotfill\hyperref[e:233]{\pageref*{e:233}}\par
\jfett{単位円} {\footnotesize\color{black!70}たんいえん} -- der Einheitskreis\dotfill\hyperref[e:137]{\pageref*{e:137}}\par
\jfett{単位ベクトル} {\footnotesize\color{black!70}たんいベクトル} -- der Einheitsvektor\dotfill\hyperref[e:303]{\pageref*{e:303}}\par
\jfett{単調減少} {\footnotesize\color{black!70}たんちょうげんしょう} -- monoton fallend\dotfill\hyperref[e:143]{\pageref*{e:143}}\par
\jfett{単調増加} {\footnotesize\color{black!70}たんちょうぞうか} -- monoton steigend\dotfill\hyperref[e:142]{\pageref*{e:142}}\par
\jfett{値域} {\footnotesize\color{black!70}ちいき} -- der Wertebereich\dotfill\hyperref[e:110]{\pageref*{e:110}}\par
\jfett{中央値} {\footnotesize\color{black!70}ちゅうおうち} -- der Median\dotfill\hyperref[e:339]{\pageref*{e:339}}\par
\jfett{中空円柱} {\footnotesize\color{black!70}ちゅうくうえんちゅう} -- der Hohlzylinder\dotfill\hyperref[e:237]{\pageref*{e:237}}\par
\jfett{中心} {\footnotesize\color{black!70}ちゅうしん} -- der Mittelpunkt\dotfill\hyperref[e:212]{\pageref*{e:212}}\par
\jfett{中線} {\footnotesize\color{black!70}ちゅうせん} -- die Seitenhalbierende\dotfill\hyperref[e:192]{\pageref*{e:192}}\par
\jfett{頂点} {\footnotesize\color{black!70}ちょうてん} -- die Ecke\dotfill\hyperref[e:245]{\pageref*{e:245}}\par
\jfett{頂点} {\footnotesize\color{black!70}ちょうてん} -- der Scheitel\dotfill\hyperref[e:154]{\pageref*{e:154}}\par
\jfett{頂点} {\footnotesize\color{black!70}ちょうてん} -- der Scheitelpunkt\dotfill\hyperref[e:129]{\pageref*{e:129}}\par
\jfett{長方形} {\footnotesize\color{black!70}ちょうほうけい} -- das Rechteck\dotfill\hyperref[e:199]{\pageref*{e:199}}\par
\jfett{直積} {\footnotesize\color{black!70}ちょくせき} -- das kartesische Produkt\dotfill\hyperref[e:103]{\pageref*{e:103}}\par
\jfett{直線} {\footnotesize\color{black!70}ちょくせん} -- die Gerade\dotfill\hyperref[e:147]{\pageref*{e:147}}\par
\jfett{直線の方程式} {\footnotesize\color{black!70}ちょくせんのほうていしき} -- die Geradengleichung\dotfill\hyperref[e:309]{\pageref*{e:309}}\par
\jfett{直方体} {\footnotesize\color{black!70}ちょくほうたい} -- der Quader\dotfill\hyperref[e:227]{\pageref*{e:227}}\par
\jfett{直角} {\footnotesize\color{black!70}ちょっかく} -- der rechte Winkel\dotfill\hyperref[e:158]{\pageref*{e:158}}\par
\jfett{直角三角形} {\footnotesize\color{black!70}ちょっかくさんかくけい} -- das rechtwinklige Dreieck\dotfill\hyperref[e:181]{\pageref*{e:181}}\par
\jfett{直角をはさむ辺} {\footnotesize\color{black!70}ちょっかくをはさむへん} -- die Kathete\dotfill\hyperref[e:187]{\pageref*{e:187}}\par
\jfett{直径} {\footnotesize\color{black!70}ちょっけい} -- der Durchmesser\dotfill\hyperref[e:214]{\pageref*{e:214}}\par
\jfett{底} {\footnotesize\color{black!70}てい} -- die Basis\dotfill\hyperref[e:43]{\pageref*{e:43}}\par
\jfett{底角} {\footnotesize\color{black!70}ていかく} -- der Basiswinkel\dotfill\hyperref[e:184]{\pageref*{e:184}}\par
\jfett{定義} {\footnotesize\color{black!70}ていぎ} -- die Definition\dotfill\hyperref[e:346]{\pageref*{e:346}}\par
\jfett{定義域} {\footnotesize\color{black!70}ていぎいき} -- der Definitionsbereich\dotfill\hyperref[e:109]{\pageref*{e:109}}\par
\jfett{定積分} {\footnotesize\color{black!70}ていせきぶん} -- das bestimmte Integral\dotfill\hyperref[e:288]{\pageref*{e:288}}\par
\jfett{底辺} {\footnotesize\color{black!70}ていへん} -- die Basis (im Dreieck)\dotfill\hyperref[e:183]{\pageref*{e:183}}\par
\jfett{底面} {\footnotesize\color{black!70}ていめん} -- die Grundfläche\dotfill\hyperref[e:248]{\pageref*{e:248}}\par
\jfett{定理} {\footnotesize\color{black!70}ていり} -- der Satz\dotfill\hyperref[e:347]{\pageref*{e:347}}\par
\jfett{点} {\footnotesize\color{black!70}てん} -- der Punkt\dotfill\hyperref[e:146]{\pageref*{e:146}}\par
\jfett{展開図} {\footnotesize\color{black!70}てんかいず} -- das Netz\dotfill\hyperref[e:253]{\pageref*{e:253}}\par
\jfett{展開する} {\footnotesize\color{black!70}てんかいする} -- ausmultiplizieren\dotfill\hyperref[e:91]{\pageref*{e:91}}\par
\jfett{同位角} {\footnotesize\color{black!70}どういかく} -- der Stufenwinkel\dotfill\hyperref[e:167]{\pageref*{e:167}}\par
\jfett{導関数} {\footnotesize\color{black!70}どうかんすう} -- die Ableitung\dotfill\hyperref[e:273]{\pageref*{e:273}}\par
\jfett{等差数列} {\footnotesize\color{black!70}とうさすうれつ} -- die arithmetische Folge\dotfill\hyperref[e:259]{\pageref*{e:259}}\par
\jfett{同側内角} {\footnotesize\color{black!70}どうそくないかく} -- entgegengesetzt liegende Winkel\dotfill\hyperref[e:169]{\pageref*{e:169}}\par
\jfett{同値} {\footnotesize\color{black!70}どうち} -- genau dann, wenn\dotfill\hyperref[e:351]{\pageref*{e:351}}\par
\jfett{等比数列} {\footnotesize\color{black!70}とうひすうれつ} -- die geometrische Folge\dotfill\hyperref[e:260]{\pageref*{e:260}}\par
\jfett{同分母分数} {\footnotesize\color{black!70}どうぶんぼぶんすう} -- gleichnamige Brüche\dotfill\hyperref[e:34]{\pageref*{e:34}}\par
\jfett{独立変数} {\footnotesize\color{black!70}どくりつへんすう} -- die unabhängige Variable\dotfill\hyperref[e:106]{\pageref*{e:106}}\par
\jfett{度数法} {\footnotesize\color{black!70}どすうほう} -- das Gradmaß\dotfill\hyperref[e:138]{\pageref*{e:138}}\par
\jfett{鈍角} {\footnotesize\color{black!70}どんかく} -- der stumpfe Winkel\dotfill\hyperref[e:159]{\pageref*{e:159}}\par

\smallskip{\large\textbf{\color{cornflower}な \,·\, na}}\par\nopagebreak
\jfett{内角} {\footnotesize\color{black!70}ないかく} -- der Innenwinkel\dotfill\hyperref[e:178]{\pageref*{e:178}}\par
\jfett{内積} {\footnotesize\color{black!70}ないせき} -- das Skalarprodukt\dotfill\hyperref[e:306]{\pageref*{e:306}}\par
\jfett{内接円} {\footnotesize\color{black!70}ないせつえん} -- der Inkreis\dotfill\hyperref[e:193]{\pageref*{e:193}}\par
\jfett{波括弧} {\footnotesize\color{black!70}なみかっこ} -- die geschweiften Klammern\dotfill\hyperref[e:90]{\pageref*{e:90}}\par
\jfett{ならば} -- daraus folgt\dotfill\hyperref[e:350]{\pageref*{e:350}}\par
\jfett{二項係数} {\footnotesize\color{black!70}にこうけいすう} -- der Binomialkoeffizient\dotfill\hyperref[e:344]{\pageref*{e:344}}\par
\jfett{二項分布} {\footnotesize\color{black!70}にこうぶんぷ} -- die Binomialverteilung\dotfill\hyperref[e:345]{\pageref*{e:345}}\par
\jfett{二次関数} {\footnotesize\color{black!70}にじかんすう} -- die quadratische Funktion\dotfill\hyperref[e:124]{\pageref*{e:124}}\par
\jfett{二次方程式} {\footnotesize\color{black!70}にじほうていしき} -- die quadratische Gleichung\dotfill\hyperref[e:82]{\pageref*{e:82}}\par
\jfett{二乗する} {\footnotesize\color{black!70}にじょうする} -- quadrieren\dotfill\hyperref[e:51]{\pageref*{e:51}}\par
\jfett{二等辺三角形} {\footnotesize\color{black!70}にとうへんさんかくけい} -- das gleichschenklige Dreieck\dotfill\hyperref[e:182]{\pageref*{e:182}}\par
\jfett{に比例する} {\footnotesize\color{black!70}にひれいする} -- proportional zu\dotfill\hyperref[e:68]{\pageref*{e:68}}\par
\jfett{ネイピア数} {\footnotesize\color{black!70}ネイピアすう} -- die eulersche Zahl\dotfill\hyperref[e:298]{\pageref*{e:298}}\par

\smallskip{\large\textbf{\color{cornflower}は \,·\, ha}}\par\nopagebreak
\jfett{倍数} {\footnotesize\color{black!70}ばいすう} -- das Vielfache\dotfill\hyperref[e:55]{\pageref*{e:55}}\par
\jfett{倍分する} {\footnotesize\color{black!70}ばいぶんする} -- erweitern\dotfill\hyperref[e:38]{\pageref*{e:38}}\par
\jfett{八角形} {\footnotesize\color{black!70}はっかくけい} -- das Achteck\dotfill\hyperref[e:207]{\pageref*{e:207}}\par
\jfett{発散する} {\footnotesize\color{black!70}はっさんする} -- divergent\dotfill\hyperref[e:263]{\pageref*{e:263}}\par
\jfett{はるかに小さい} {\footnotesize\color{black!70}はるかにちいさい} -- viel kleiner als\dotfill\hyperref[e:67]{\pageref*{e:67}}\par
\jfett{半径} {\footnotesize\color{black!70}はんけい} -- der Radius\dotfill\hyperref[e:213]{\pageref*{e:213}}\par
\jfett{半軸} {\footnotesize\color{black!70}はんじく} -- die Halbachse\dotfill\hyperref[e:317]{\pageref*{e:317}}\par
\jfett{半直線} {\footnotesize\color{black!70}はんちょくせん} -- der Strahl\dotfill\hyperref[e:149]{\pageref*{e:149}}\par
\jfett{比} {\footnotesize\color{black!70}ひ} -- das Verhältnis\dotfill\hyperref[e:40]{\pageref*{e:40}}\par
\jfett{被開方数} {\footnotesize\color{black!70}ひかいほうすう} -- der Radikand\dotfill\hyperref[e:48]{\pageref*{e:48}}\par
\jfett{引く} {\footnotesize\color{black!70}ひく} -- subtrahieren\dotfill\hyperref[e:17]{\pageref*{e:17}}\par
\jfett{被減数} {\footnotesize\color{black!70}ひげんすう} -- der Minuend\dotfill\hyperref[e:19]{\pageref*{e:19}}\par
\jfett{ひし形} {\footnotesize\color{black!70}ひしがた} -- die Raute\dotfill\hyperref[e:201]{\pageref*{e:201}}\par
\jfett{被除数} {\footnotesize\color{black!70}ひじょすう} -- der Dividend\dotfill\hyperref[e:26]{\pageref*{e:26}}\par
\jfett{被積分関数} {\footnotesize\color{black!70}ひせきぶんかんすう} -- der Integrand\dotfill\hyperref[e:290]{\pageref*{e:290}}\par
\jfett{等しい} {\footnotesize\color{black!70}ひとしい} -- gleich\dotfill\hyperref[e:60]{\pageref*{e:60}}\par
\jfett{等しくない} {\footnotesize\color{black!70}ひとしくない} -- ungleich\dotfill\hyperref[e:61]{\pageref*{e:61}}\par
\jfett{微分} {\footnotesize\color{black!70}びぶん} -- das Differenzial\dotfill\hyperref[e:276]{\pageref*{e:276}}\par
\jfett{微分係数} {\footnotesize\color{black!70}びぶんけいすう} -- der Differenzialquotient\dotfill\hyperref[e:275]{\pageref*{e:275}}\par
\jfett{微分する} {\footnotesize\color{black!70}びぶんする} -- ableiten\dotfill\hyperref[e:278]{\pageref*{e:278}}\par
\jfett{微分法} {\footnotesize\color{black!70}びぶんほう} -- die Differenzialrechnung\dotfill\hyperref[e:274]{\pageref*{e:274}}\par
\jfett{百分率} {\footnotesize\color{black!70}ひゃくぶんりつ} -- der Prozentsatz\dotfill\hyperref[e:58]{\pageref*{e:58}}\par
\jfett{標準偏差} {\footnotesize\color{black!70}ひょうじゅんへんさ} -- die Standardabweichung\dotfill\hyperref[e:337]{\pageref*{e:337}}\par
\jfett{表面積} {\footnotesize\color{black!70}ひょうめんせき} -- der Oberflächeninhalt\dotfill\hyperref[e:256]{\pageref*{e:256}}\par
\jfett{比例} {\footnotesize\color{black!70}ひれい} -- die Proportionalität\dotfill\hyperref[e:94]{\pageref*{e:94}}\par
\jfett{比例式} {\footnotesize\color{black!70}ひれいしき} -- die Verhältnisgleichung\dotfill\hyperref[e:41]{\pageref*{e:41}}\par
\jfett{複素数} {\footnotesize\color{black!70}ふくそすう} -- die komplexe Zahl\dotfill\hyperref[e:10]{\pageref*{e:10}}\par
\jfett{不足和} {\footnotesize\color{black!70}ふそくわ} -- die Untersumme\dotfill\hyperref[e:294]{\pageref*{e:294}}\par
\jfett{不定積分} {\footnotesize\color{black!70}ふていせきぶん} -- das unbestimmte Integral\dotfill\hyperref[e:287]{\pageref*{e:287}}\par
\jfett{不等式} {\footnotesize\color{black!70}ふとうしき} -- die Ungleichung\dotfill\hyperref[e:77]{\pageref*{e:77}}\par
\jfett{不等辺三角形} {\footnotesize\color{black!70}ふとうへんさんかくけい} -- das unregelmäßige Dreieck\dotfill\hyperref[e:180]{\pageref*{e:180}}\par
\jfett{部分集合} {\footnotesize\color{black!70}ぶぶんしゅうごう} -- die Teilmenge\dotfill\hyperref[e:98]{\pageref*{e:98}}\par
\jfett{不連続} {\footnotesize\color{black!70}ふれんぞく} -- unstetig\dotfill\hyperref[e:269]{\pageref*{e:269}}\par
\jfett{分割表} {\footnotesize\color{black!70}ぶんかつひょう} -- die Vierfeldertafel\dotfill\hyperref[e:333]{\pageref*{e:333}}\par
\jfett{分散} {\footnotesize\color{black!70}ぶんさん} -- die Varianz\dotfill\hyperref[e:336]{\pageref*{e:336}}\par
\jfett{分子} {\footnotesize\color{black!70}ぶんし} -- der Zähler\dotfill\hyperref[e:29]{\pageref*{e:29}}\par
\jfett{分数} {\footnotesize\color{black!70}ぶんすう} -- der Bruch\dotfill\hyperref[e:28]{\pageref*{e:28}}\par
\jfett{分数関数} {\footnotesize\color{black!70}ぶんすうかんすう} -- die gebrochenrationale Funktion\dotfill\hyperref[e:127]{\pageref*{e:127}}\par
\jfett{分母} {\footnotesize\color{black!70}ぶんぼ} -- der Nenner\dotfill\hyperref[e:30]{\pageref*{e:30}}\par
\jfett{平角} {\footnotesize\color{black!70}へいかく} -- der gestreckte Winkel\dotfill\hyperref[e:160]{\pageref*{e:160}}\par
\jfett{平均値} {\footnotesize\color{black!70}へいきんち} -- der Mittelwert\dotfill\hyperref[e:338]{\pageref*{e:338}}\par
\jfett{平均変化率} {\footnotesize\color{black!70}へいきんへんかりつ} -- der Differenzenquotient\dotfill\hyperref[e:272]{\pageref*{e:272}}\par
\jfett{平行} {\footnotesize\color{black!70}へいこう} -- parallel\dotfill\hyperref[e:150]{\pageref*{e:150}}\par
\jfett{平行移動} {\footnotesize\color{black!70}へいこういどう} -- die Verschiebung\dotfill\hyperref[e:173]{\pageref*{e:173}}\par
\jfett{平行四辺形} {\footnotesize\color{black!70}へいこうしへんけい} -- das Parallelogramm\dotfill\hyperref[e:200]{\pageref*{e:200}}\par
\jfett{平方根} {\footnotesize\color{black!70}へいほうこん} -- die Quadratwurzel\dotfill\hyperref[e:45]{\pageref*{e:45}}\par
\jfett{平面} {\footnotesize\color{black!70}へいめん} -- die Ebene\dotfill\hyperref[e:308]{\pageref*{e:308}}\par
\jfett{べき級数} {\footnotesize\color{black!70}べききゅうすう} -- die Potenzreihe\dotfill\hyperref[e:267]{\pageref*{e:267}}\par
\jfett{ベクトル} -- der Vektor\dotfill\hyperref[e:301]{\pageref*{e:301}}\par
\jfett{ベクトルの大きさ} {\footnotesize\color{black!70}ベクトルのおおきさ} -- der Betrag eines Vektors\dotfill\hyperref[e:305]{\pageref*{e:305}}\par
\jfett{辺} {\footnotesize\color{black!70}へん} -- die Kante\dotfill\hyperref[e:246]{\pageref*{e:246}}\par
\jfett{辺} {\footnotesize\color{black!70}へん} -- der Schenkel\dotfill\hyperref[e:155]{\pageref*{e:155}}\par
\jfett{変曲点} {\footnotesize\color{black!70}へんきょくてん} -- der Wendepunkt\dotfill\hyperref[e:284]{\pageref*{e:284}}\par
\jfett{変数} {\footnotesize\color{black!70}へんすう} -- die Variable\dotfill\hyperref[e:71]{\pageref*{e:71}}\par
\jfett{偏導関数} {\footnotesize\color{black!70}へんどうかんすう} -- die partielle Ableitung\dotfill\hyperref[e:277]{\pageref*{e:277}}\par
\jfett{方程式} {\footnotesize\color{black!70}ほうていしき} -- die Gleichung\dotfill\hyperref[e:76]{\pageref*{e:76}}\par
\jfett{放物線} {\footnotesize\color{black!70}ほうぶつせん} -- die Parabel\dotfill\hyperref[e:125]{\pageref*{e:125}}\par
\jfett{補角} {\footnotesize\color{black!70}ほかく} -- der Supplementwinkel\dotfill\hyperref[e:165]{\pageref*{e:165}}\par
\jfett{ほぼ等しい} {\footnotesize\color{black!70}ほぼひとしい} -- ungefähr\dotfill\hyperref[e:69]{\pageref*{e:69}}\par

\smallskip{\large\textbf{\color{cornflower}ま \,·\, ma}}\par\nopagebreak
\jfett{丸括弧} {\footnotesize\color{black!70}まるかっこ} -- die runden Klammern\dotfill\hyperref[e:88]{\pageref*{e:88}}\par
\jfett{未知数} {\footnotesize\color{black!70}みちすう} -- die Unbekannte\dotfill\hyperref[e:73]{\pageref*{e:73}}\par
\jfett{見取図} {\footnotesize\color{black!70}みとりず} -- das Schrägbild\dotfill\hyperref[e:254]{\pageref*{e:254}}\par
\jfett{無限級数} {\footnotesize\color{black!70}むげんきゅうすう} -- die unendliche Reihe\dotfill\hyperref[e:266]{\pageref*{e:266}}\par
\jfett{無限大} {\footnotesize\color{black!70}むげんだい} -- unendlich\dotfill\hyperref[e:70]{\pageref*{e:70}}\par
\jfett{無理関数} {\footnotesize\color{black!70}むりかんすう} -- die Wurzelfunktion\dotfill\hyperref[e:128]{\pageref*{e:128}}\par
\jfett{無理数} {\footnotesize\color{black!70}むりすう} -- die irrationale Zahl\dotfill\hyperref[e:9]{\pageref*{e:9}}\par
\jfett{無理方程式} {\footnotesize\color{black!70}むりほうていしき} -- die Wurzelgleichung\dotfill\hyperref[e:83]{\pageref*{e:83}}\par
\jfett{面} {\footnotesize\color{black!70}めん} -- die Seitenfläche\dotfill\hyperref[e:247]{\pageref*{e:247}}\par
\jfett{面積} {\footnotesize\color{black!70}めんせき} -- der Flächeninhalt\dotfill\hyperref[e:210]{\pageref*{e:210}}\par
\jfett{面積が等しい} {\footnotesize\color{black!70}めんせきがひとしい} -- flächengleich\dotfill\hyperref[e:195]{\pageref*{e:195}}\par
\jfett{面の対角線} {\footnotesize\color{black!70}めんのたいかくせん} -- die Flächendiagonale\dotfill\hyperref[e:251]{\pageref*{e:251}}\par

\smallskip{\large\textbf{\color{cornflower}や \,·\, ya}}\par\nopagebreak
\jfett{約数} {\footnotesize\color{black!70}やくすう} -- der Teiler\dotfill\hyperref[e:54]{\pageref*{e:54}}\par
\jfett{約分する} {\footnotesize\color{black!70}やくぶんする} -- kürzen\dotfill\hyperref[e:37]{\pageref*{e:37}}\par
\jfett{優角} {\footnotesize\color{black!70}ゆうかく} -- der überstumpfe Winkel\dotfill\hyperref[e:161]{\pageref*{e:161}}\par
\jfett{有理数} {\footnotesize\color{black!70}ゆうりすう} -- die rationale Zahl\dotfill\hyperref[e:7]{\pageref*{e:7}}\par
\jfett{要素} {\footnotesize\color{black!70}ようそ} -- das Element\dotfill\hyperref[e:97]{\pageref*{e:97}}\par
\jfett{余角} {\footnotesize\color{black!70}よかく} -- der Komplementwinkel\dotfill\hyperref[e:164]{\pageref*{e:164}}\par
\jfett{余弦、コサイン} {\footnotesize\color{black!70}よげん、コサイン} -- der Kosinus\dotfill\hyperref[e:133]{\pageref*{e:133}}\par
\jfett{余弦定理} {\footnotesize\color{black!70}よげんていり} -- der Kosinussatz\dotfill\hyperref[e:141]{\pageref*{e:141}}\par
\jfett{余事象} {\footnotesize\color{black!70}よじしょう} -- das Gegenereignis\dotfill\hyperref[e:329]{\pageref*{e:329}}\par
\jfett{余接、コタンジェント} {\footnotesize\color{black!70}よせつ、コタンジェント} -- der Kotangens\dotfill\hyperref[e:135]{\pageref*{e:135}}\par
\jfett{より大きい} {\footnotesize\color{black!70}よりおおきい} -- größer als\dotfill\hyperref[e:64]{\pageref*{e:64}}\par
\jfett{より小さい} {\footnotesize\color{black!70}よりちいさい} -- kleiner als\dotfill\hyperref[e:63]{\pageref*{e:63}}\par

\smallskip{\large\textbf{\color{cornflower}ら \,·\, ra}}\par\nopagebreak
\jfett{立体} {\footnotesize\color{black!70}りったい} -- der Körper\dotfill\hyperref[e:225]{\pageref*{e:225}}\par
\jfett{立体の対角線} {\footnotesize\color{black!70}りったいのたいかくせん} -- die Raumdiagonale\dotfill\hyperref[e:252]{\pageref*{e:252}}\par
\jfett{立方根} {\footnotesize\color{black!70}りっぽうこん} -- die Kubikwurzel\dotfill\hyperref[e:46]{\pageref*{e:46}}\par
\jfett{立方体} {\footnotesize\color{black!70}りっぽうたい} -- der Würfel\dotfill\hyperref[e:226]{\pageref*{e:226}}\par
\jfett{隣接補角} {\footnotesize\color{black!70}りんせつほかく} -- der Nebenwinkel\dotfill\hyperref[e:163]{\pageref*{e:163}}\par
\jfett{隣辺} {\footnotesize\color{black!70}りんぺん} -- die Ankathete\dotfill\hyperref[e:188]{\pageref*{e:188}}\par
\jfett{累乗} {\footnotesize\color{black!70}るいじょう} -- die Potenz\dotfill\hyperref[e:42]{\pageref*{e:42}}\par
\jfett{累乗根を求める} {\footnotesize\color{black!70}るいじょうこんをもとめる} -- radizieren, Wurzel ziehen\dotfill\hyperref[e:52]{\pageref*{e:52}}\par
\jfett{累乗する} {\footnotesize\color{black!70}るいじょうする} -- potenzieren\dotfill\hyperref[e:50]{\pageref*{e:50}}\par
\jfett{零点} {\footnotesize\color{black!70}れいてん} -- die Nullstelle\dotfill\hyperref[e:122]{\pageref*{e:122}}\par
\jfett{零に収束する数列} {\footnotesize\color{black!70}れいにしゅうそくするすうれつ} -- die Nullfolge\dotfill\hyperref[e:258]{\pageref*{e:258}}\par
\jfett{零ベクトル} {\footnotesize\color{black!70}れいベクトル} -- der Nullvektor\dotfill\hyperref[e:304]{\pageref*{e:304}}\par
\jfett{列} {\footnotesize\color{black!70}れつ} -- die Spalte\dotfill\hyperref[e:323]{\pageref*{e:323}}\par
\jfett{連続} {\footnotesize\color{black!70}れんぞく} -- stetig\dotfill\hyperref[e:268]{\pageref*{e:268}}\par
\jfett{連立方程式} {\footnotesize\color{black!70}れんりつほうていしき} -- das Gleichungssystem\dotfill\hyperref[e:80]{\pageref*{e:80}}\par

\smallskip{\large\textbf{\color{cornflower}わ \,·\, wa}}\par\nopagebreak
\jfett{和} {\footnotesize\color{black!70}わ} -- die Summe\dotfill\hyperref[e:15]{\pageref*{e:15}}\par
\jfett{和集合} {\footnotesize\color{black!70}わしゅうごう} -- die Vereinigungsmenge\dotfill\hyperref[e:101]{\pageref*{e:101}}\par
\jfett{割る} {\footnotesize\color{black!70}わる} -- dividieren\dotfill\hyperref[e:24]{\pageref*{e:24}}\par
\end{multicols}

\clearpage
\end{CJK}
\end{document}
